\documentclass[11pt,oneside,openany]{book}
\usepackage[T1]{fontenc}
\usepackage[utf8]{inputenc}
\usepackage{lmodern}
\usepackage[a4paper,inner=28mm,outer=25mm,top=26mm,bottom=27mm,headheight=15pt]{geometry}
\usepackage{amsmath,amssymb,amsthm,mathtools,bm,mathrsfs}
\usepackage{microtype,booktabs,longtable,array,tabularx,enumitem,xcolor}
\usepackage{fancyhdr}
\usepackage[breaklinks=true,colorlinks=true,linkcolor=black,citecolor=black,urlcolor=black,pdfusetitle]{hyperref}
\usepackage{aliascnt}
\usepackage[nameinlink,noabbrev]{cleveref}
\setlength{\emergencystretch}{2em}
\setlength{\parskip}{3pt}
\setlength{\parindent}{1.2em}
\setcounter{tocdepth}{1}
\setcounter{secnumdepth}{2}
\raggedbottom
\allowdisplaybreaks[2]
\pagestyle{fancy}
\fancyhf{}
\fancyhead[L]{\small\nouppercase{\leftmark}}
\fancyfoot[C]{\thepage}
\renewcommand{\headrulewidth}{0.3pt}
\renewcommand{\chaptermark}[1]{\markboth{\thechapter.\ #1}{}}
\fancypagestyle{plain}{\fancyhf{}\fancyfoot[C]{\thepage}\renewcommand{\headrulewidth}{0pt}}
\newtheorem{theorem}{Theorem}[chapter]
\newaliascnt{proposition}{theorem}
\newtheorem{proposition}[proposition]{Proposition}
\aliascntresetthe{proposition}
\newaliascnt{lemma}{theorem}
\newtheorem{lemma}[lemma]{Lemma}
\aliascntresetthe{lemma}
\newaliascnt{corollary}{theorem}
\newtheorem{corollary}[corollary]{Corollary}
\aliascntresetthe{corollary}
\newaliascnt{criterion}{theorem}
\newtheorem{criterion}[criterion]{Conditional criterion}
\aliascntresetthe{criterion}
\theoremstyle{definition}
\newaliascnt{definition}{theorem}
\newtheorem{definition}[definition]{Definition}
\aliascntresetthe{definition}
\newaliascnt{certificate}{theorem}
\newtheorem{certificate}[certificate]{Reported finite certificate}
\aliascntresetthe{certificate}
\newaliascnt{hypothesis}{theorem}
\newtheorem{hypothesis}[hypothesis]{Hypothesis}
\aliascntresetthe{hypothesis}
\newaliascnt{problem}{theorem}
\newtheorem{problem}[problem]{Remaining problem}
\aliascntresetthe{problem}
\theoremstyle{remark}
\newaliascnt{remark}{theorem}
\newtheorem{remark}[remark]{Remark}
\aliascntresetthe{remark}
\newaliascnt{scope}{theorem}
\newtheorem{scope}[scope]{Scope}
\aliascntresetthe{scope}
\crefname{theorem}{theorem}{theorems}
\crefname{proposition}{proposition}{propositions}
\crefname{lemma}{lemma}{lemmas}
\crefname{corollary}{corollary}{corollaries}
\crefname{criterion}{criterion}{criteria}
\crefname{definition}{definition}{definitions}
\crefname{hypothesis}{hypothesis}{hypotheses}
\crefname{certificate}{finite certificate}{finite certificates}
\Crefname{certificate}{Finite certificate}{Finite certificates}
\crefname{problem}{remaining problem}{remaining problems}
\crefname{remark}{remark}{remarks}
\crefname{scope}{scope statement}{scope statements}
\newcommand{\R}{\mathbb R}
\newcommand{\C}{\mathbb C}
\newcommand{\N}{\mathbb N}
\newcommand{\Z}{\mathbb Z}
\newcommand{\T}{\mathbb T}
\newcommand{\E}{\mathbb E}
\newcommand{\cH}{\mathcal H}
\newcommand{\cA}{\mathcal A}
\newcommand{\cG}{\mathcal G}
\newcommand{\cK}{\mathcal K}
\newcommand{\one}{\mathbf1}
\newcommand{\dd}{\mathrm d}
\newcommand{\ii}{\mathrm i}
\newcommand{\Tr}{\operatorname{Tr}}
\newcommand{\tr}{\operatorname{tr}}
\newcommand{\diag}{\operatorname{diag}}
\newcommand{\Ran}{\operatorname{Ran}}
\newcommand{\Ker}{\operatorname{Ker}}
\newcommand{\Var}{\operatorname{Var}}
\newcommand{\Cov}{\operatorname{Cov}}
\newcommand{\supp}{\operatorname{supp}}
\newcommand{\dist}{\operatorname{dist}}
\newcommand{\osc}{\operatorname{osc}}
\newcommand{\gap}{\operatorname{gap}}
\newcommand{\SU}{\mathrm{SU}}
\newcommand{\OS}{\mathrm{OS}}
\newcommand{\op}{\mathrm{op}}
\newcommand{\rank}{\operatorname{rank}}
\newcommand{\norm}[1]{\left\lVert#1\right\rVert}
\newcommand{\abs}[1]{\left\lvert#1\right\rvert}
\newcommand{\ip}[2]{\left\langle#1,#2\right\rangle}
\newcommand{\avg}[1]{\left\langle#1\right\rangle_{\T^4}}
\newcommand{\src}[2]{\cite[#2]{#1}}
\newcommand{\sourcebasis}[1]{\par\smallskip\noindent #1\par}
\newcolumntype{L}[1]{>{\raggedright\arraybackslash}p{#1}}
\newcolumntype{Y}{>{\raggedright\arraybackslash}X}
\title{\textbf{Renewal Yang--Mills Theory}\\[0.55em]\Large Finite Construction, Renormalization,\\and Conditional Mass-Gap Closure}
\author{A.~P\'elissier}
\date{Specialized monograph, Version 8\\21 September 2026}
\hypersetup{
  pdftitle={Renewal Yang--Mills Theory: Finite Construction, Renormalization, and Conditional Mass-Gap Closure},
  pdfauthor={A. Pélissier},
  pdfsubject={Finite Yang--Mills construction, gauge-compatible Gaussian shells, and conditional continuum mass-gap closure}
}
\begin{document}
\frontmatter
\begin{titlepage}
\centering
\vspace*{25mm}
{\small THE RENEWAL PROGRAMME\par}
\vspace{15mm}
{\Huge\bfseries Renewal\\[3mm]Yang--Mills Theory\par}
\vspace{12mm}
{\Large Finite Construction, Renormalization,\\[2mm]and Conditional Mass-Gap Closure\par}
\vspace{20mm}
{\large A.~P\'elissier\par}
\vfill
{\normalsize Specialized monograph\quad\textbar\quad Version 8\par}
\vspace{3mm}
{\normalsize 21 September 2026\par}
\end{titlepage}
\chapter*{Abstract}
\addcontentsline{toc}{chapter}{Abstract}
The Yang--Mills sector of the Renewal Grand Tensor admits a finite,
source-resolved formulation in which gauge reduction, conditional
elimination, response, and reflected transfer arise from a common
occurrence law. This monograph develops the conditional route from
that construction to a nontrivial continuum theory with a physical
mass gap. Exact Schur identities, normalized connected-source calculus,
measured field redefinitions, degree-faithful analytic forests, and
calibrated spectral criteria keep the finite law and its physical
normalization visible throughout. The nested block-average Gaussian
tower has positive conditional shells and quantitative transverse
control. Its complete response includes mixed-shell contractions;
Schur ownership assigns them once without assuming that retained
endpoint terms are marginally invisible.

The nonlinear root-path block has a smooth compact completion and an
exact inverse-Haar density. A depth-independent critical-norm domain,
weak continuity of the limiting source, and exact-source variational
recovery identify its finite-amplitude local classical continuum
minimum. On the controlled local branch, the measured two-loop
expansion has a volume-uniform $O(\beta^{-2})$ remainder through four
retained-source derivatives under explicit higher-order analytic
hypotheses. Finite measured composition closes the field-degree packet
$(8,6,4;7)$ for the classical action, one-loop density, two-loop term,
and block map. The all-level one-cell determinant retains the normal,
constrained, gauge, and Haar factors. Schatten-ideal convergence controls the subtracted
tail, while low-order traces require separate measured cancellations.
An exact flat-holonomy formula has a dyadically summable $O(L^{-2})$
cofinal error. For the balanced-root $\SU(2)$ convention, the complete
noncommuting quartic coefficient is
$-11E_4\log L/(3\pi^2)$, with a convergent finite quartic part.
A separate fourth-source subtraction has a second-chaos logarithmic
counterterm and a cofinal finite part with summable local tails.
The source counterterm retains the full matrix filter rather than its
measured scalar trace. Both coefficient results are distinguished from
a joint nonlinear remainder bound and physical coupling identification.

For the unrestricted bare Wilson law, full-pressure and reflection
estimates yield a selected near-critical source window, sharp
cell-energy covariance control, total-action Gaussian fluctuations,
and finite-entropy bounds for changes of law. Exact generated
likelihoods propagate those bounds to the actual blocked measures;
their entropy decomposition retains the output normalization, and
conditional regression tests which source information survives.
These estimates have no small-field cutoff, but their selected coupling
regime is not identified with a prescribed interacting trajectory. For the actual
compact parent, stationary-orbit certificates establish both local
tube coercivity and hidden-diffusion metastability. Same-law resolvent
comparison nevertheless gives original-clock visible repair floors.
Inverse-paired noncommuting sources extend those floors through the
complete radial cutoff transition on a certified family, without
claiming an all-exterior cover.

The native physical transfer is treated separately from auxiliary
diffusions. Exact magnetic normalization produces a conditional-affinity
kernel, not by itself a kinetic semigroup. Its Gaussian magnetic image
has an explicit occupation law: a subextensive cutoff in the original
total oscillator degree loses its weight. The corresponding genuine
compact projections have a separate fixed-volume limit.
Resistance-normalized gauge geometry and an all-representation
subordinated comparison retain the
complete magnetic action and each operator's actual Perron value.
The ordered critical spectrum is governed by an invariant quartic
Hamiltonian, both at fixed regulator and on a growing-regulator branch
with $L^2\gamma^{-1/3}+L^4\gamma^{-2/3}+L^3\gamma^{-1}\to0$.
Its neutral logarithmic gap is asymptotic to
$(\gamma L^3)^{-1/3}(e_1-e_0)$; center-sector tops occur on a
different scale. This joint domain is not a calibrated fixed physical
box. Exact kinetic centering, cubic-chaos accounting, and Perron
transport identify requirements for continuation beyond that domain.

At the native Hamiltonian level, a separate electric-dominated regime
admits an actual finite-lattice ground state with a volume-uniform gap,
and its first physical plaquette band is analytic, exponentially local,
and has a volume-independent cubic spectral remainder. Without assuming
small interaction, one-link high-degree elimination has a
volume-independent Schur floor and factorial tails for every fixed
low-energy spectral window. Exact equal-electric-energy resonances
prevent a whole-head homological inverse; retaining those resonances
instead yields an exact local unitary normal form with a spatially
summable connected remainder when $13\pi b/\kappa<1$. These are
strong-electric finite-lattice statements and are not identified with
the weak-coupling continuum trajectory.

Physically bounded contractible Wilson loops become insensitive to the
constant soft coordinate. A reducing-band scalar-compression criterion
converts that fact into a vanishing local Gram contribution without
inverting a possibly singular entrance matrix. In contrast, a
canonically smeared magnetic-energy probe in a selected actual-vacuum
mesoscopic window has nonzero limiting norm and a continuous spectral
measure reaching zero. These complementary source tests distinguish
invisible finite-volume soft modes from locally observable soft
spectrum; neither is substituted for the interacting physical theory.

The terminal theorem remains a quantitative conditional statement.
A compatible nontrivial continuum law, noncollapsed local probes,
and contraction on the complete centered reflected algebra at one
fixed physical time imply a positive Hamiltonian mass gap. Uniform
nonlinear transport, complete compact control, continuum construction,
and the calibrated physical contraction remain explicit obligations;
no local coefficient, selected microscopic or mesoscopic limit, or
critical spectral asymptotic is substituted for that packet.

\chapter*{Scope and logical conventions}
\addcontentsline{toc}{chapter}{Scope and logical conventions}
The exposition is ordered by mathematical dependency. Finite algebraic
identities, estimates on declared analytic domains, and continuum closure
criteria are distinguished throughout. Alternative sufficient spectral
mechanisms are not treated as simultaneous requirements. Repeated uses of
an identity refer to the same law, source space, quotient, and physical
normalization unless an explicit transport is given.

An \emph{exact identity} is an algebraic, probabilistic, or operator
identity on its stated domain. An \emph{estimate} retains its gauge group,
regulator, analytic norm, and uniformity hypotheses. A \emph{reported
finite certificate} is a computational identity or inequality documented
in a cited reference; it is not presented as an independently rerun
calculation. A \emph{conditional criterion} is a proved implication whose
input estimates must still be verified for the proposed four-dimensional
Yang--Mills family.

Three distinctions govern the Gaussian-to-interacting interface. The
gauge quotient does not determine the representative in which individual
vertex legs are evaluated. The fluctuation representative does not choose
the background with respect to which a response is differentiated. A
fixed-background calculation does not identify the physical anomaly
without the corresponding determinant and background comparison. These
distinctions are implemented in \cref{ch:tower} and carried into the
closure hypotheses of \cref{ch:closure}.

Two further distinctions govern the nonlinear continuation. Equality of
local block germs does not identify complete compact pushforwards, and
a one-loop remainder on a restricted box is not a derivative bound on
its compact complement. The conditional construction in
\cref{ch:compact-response} retains both distinctions. In
\cref{ch:compact-repair}, a local Neumann gap, a killed exit floor, an
unrestricted auxiliary gap, and a physical Hamiltonian gap are separate
objects. The counterexample for the hidden diffusion is kept alongside
the local repair theorems rather than superseded by them.

The full-law source results of \cref{ch:full-compact-sources} concern
the unrestricted bare Wilson law at a selected microscopic sequence.
The native spectral results of \cref{ch:native-transfer} concern the
physical Wilson transfer at fixed regulator and on a stated joint
critical domain. The infinite-time vacuum probe in
\cref{sec:smeared-vacuum-probe} has its own selected mesoscopic
window and normalization. None is silently transferred to a
renewal-loaded continuum family. A measured quartic logarithm, a
renormalized random-source coefficient, a critical toron level, and
a fixed-physical-time continuum gap are different conclusions.

The electric-Hamiltonian results later in \cref{ch:native-transfer}
have a different parameter regime again. Their small parameter is the
magnetic-to-electric ratio $b/\kappa$, and their uniformity is in the
spatial torus size at fixed lattice normalization. A volume-uniform
lattice gap in that regime, an analytic low band, a high-harmonic
Schur floor, or a resonance-preserving normal form is therefore not
relabelled as a continuum mass. Conversely, the exact resonances found
there are retained as an obstruction to one particular elimination
scheme, not as an obstruction to Yang--Mills mass generation.

Local invisibility is likewise source-dependent. Excision of a soft
reducing band is justified only by vanishing overlap with the actual
normalized local source family. A low-energy band with surviving
probe weight cannot be removed by this argument. In Gaussian
comparisons, a homogeneous higher Wiener chaos is not treated as a
Gaussian random variable, and low-order moments are not promoted to
an all-order exponential or operator-norm estimate.

The mathematical dependencies and precise result locations are collected
in \cref{app:concordance}. The finite construction inherited from the
Grand Tensor is recalled only to the extent needed to fix the common
Yang--Mills law, source algebra, and physical clock. No independent
verification of the large computational certificates, or external
literature audit, is asserted.

\tableofcontents
\mainmatter
\chapter{The target and the conditional architecture}
\label{ch:architecture}
\section{What has to be constructed}
The target is a nontrivial four-dimensional Yang--Mills quantum theory in a specified gauge group and vacuum sector, obtained as a compatible continuum limit of the finite renewal loading. Its local gauge-invariant Euclidean correlations must support Osterwalder--Schrader reconstruction, and the reconstructed Hamiltonian must satisfy
\begin{equation}\label{eq:target}
H\Omega=0,\qquad \Ker H=\C\Omega,\qquad
H|_{\Omega^\perp}\succeq mI,\qquad m>0.
\end{equation}
The adjective \emph{physical} is essential: the lower bound is measured in a declared time unit, not in an adjustable update counter or an auxiliary heat-bath time. The algebra must also be declared. A gap on bounded-support local observables is not automatically a gap in every wrapping sector of every finite torus.

Most finite constructions are formulated for a compact gauge group. The earlier colour-traced response banks use the declared $\SU(3)$ normalization, whereas the actual one-loop checks and compact parent, orbit, and shell certificates in $\SU(2)$ have their own Hilbert--Schmidt and action conventions. The smooth pivot construction is formulated for the stated compact matrix-group setting. None of these group-specific quantitative results is silently promoted to every compact simple group; the closure inputs must be supplied for the group under consideration.

\section{The dependency chain}
The principal construction has the following order:
\begin{equation}\label{eq:main-chain}
\begin{aligned}
&\text{common finite history law and gauge-invariant sources}\\
&\quad\longrightarrow\text{exact compact reduction and blocking}\\
&\quad\longrightarrow\text{Gaussian principal tower and controlled interactions}\\
&\quad\longrightarrow\text{calibrated marginal trajectory and local continuum law}\\
&\quad\longrightarrow\text{nontrivial reflected local Hilbert space}\\
&\quad\longrightarrow\text{uniform fixed-physical-time contraction}\\
&\quad\longrightarrow\text{positive Hamiltonian mass gap}.
\end{aligned}
\end{equation}
Each arrow means a theorem with a specific input interface. In particular, the final arrow is a short spectral argument; the preceding contraction estimate is not obtained merely by renaming a finite-dimensional Hessian gap.

The finite steps in this chain have explicit interfaces. The exact
blocking differential retains a unit pullback sector and a conditionally
centered nonlinear remainder. Its measured redundant quotient is
preserved through the pushed conditional current. The third current
jet is tied to four-entry cumulants by the normalized score hierarchy;
the degree-faithful forest and corrected contour estimate provide a
sufficient marked analytic route without replacing its uniformity
hypotheses.

There are alternative ways to supply the spectral input: a complete
effective-action Dobrushin bound, a source-resolved physical
Poincar\'e estimate with a proved same-generator incidence, a two-level
physical form comparison, an orbit-coercive Feshbach argument, or a direct
reflected Gram contraction. These alternatives are not extra conjunctive
requirements. The exact two-level identity identifies the retained
infrared estimate that a detail-fibre gap cannot supply. A successful
route need not solve every historical sufficient condition.

\begin{definition}[Three logically different closure tasks]
\label{def:three-tasks}
\emph{Construction closure} supplies compatible interacting measures, local limits, the required Euclidean axioms, and nontriviality. \emph{Renormalization closure} controls the stable coordinates and identifies the physical marginal trajectory, including its anomaly coefficient and clock. \emph{Spectral closure} gives a uniform contraction or equivalent physical coercivity on the complete centered local sector. None of the three is substituted for either of the others.
\end{definition}

\section{The Gaussian endpoint and its response interface}
The Gaussian construction is an actual nested block-average covariance
tower, rather than a repeatedly reset endpoint chart. Its positive shells,
fixed point, convergence rate, and transverse inverse bounds are explicit.
The full fine-edge Wilson vertices can be coupled to those shells only
after a matching fluctuation representative is chosen. The tree-gauge map
in \cref{sec:tree-gauge} supplies that finite algebraic interface.

The next response requirement is not simply a scalar evaluation with
unspecified vertices. It is a complete full-torus calculation with common
normalization, momentum-derivative control, and a declared background.
The fixed fine-background prescription uses the existing vertex banks;
the harmonic prescription follows the minimizing blocked action and
contains nonzero background aliases. A calculation in the former supplies
the physical Gaussian anomaly only after the averaged comparison in
\cref{crit:background-ADMC}, or an equivalent direct identification, has
been proved. The finite map and corner identities do not establish that
comparison.

The nonlinear compact continuation has an explicit one-step carrier:
\cref{cor:compact-product} gives solved compact pivots and their
positive density. Under the higher-order local hypotheses,
\cref{thm:measured-two-loop} strengthens the one-loop estimate to a
measured two-loop expansion through four source derivatives.
The all-level local determinant and its evaluated quartic logarithm
are treated in \cref{sec:cofinal-measured-response,sec:measured-quartic-log}.
The fourth-source finite part of
\cref{thm:fourth-source-finite-part} supplies a distinct local source
renormalization with the full matrix filter retained. These results
do not identify the full compact complement or establish a compatible
interacting trajectory. The degree-faithful forests, owner decomposition,
and reciprocal-coupling shielding remain the interfaces for that
continuation. Likewise, the actual visible repair floors in
\cref{ch:compact-repair} need either complete coverage and assembly or a
populated conditional-patch criterion, followed by physical form
transport. Even a correctly matched ultraviolet trajectory does not
replace that infrared requirement.

The later native Hamiltonian analysis supplies an independent finite-
lattice control block. In an explicit electric-dominated domain the
actual ground state and first physical band are controlled uniformly in
volume. At arbitrary fixed coupling, sufficiently high one-link
harmonics can be eliminated with a volume-independent Schur estimate,
and every bounded low-energy spectral window has factorial one-link
tails. These facts sharpen the finite spectral side of the programme,
but exact equal-electric-energy matrix elements show that they cannot
be promoted by a global Schrieffer--Wolff or homological inversion of an
entire link head. The resonance-preserving averaging construction in
\cref{sec:native-electric-control} is the compatible replacement on its
stated strong-electric domain.

\begin{scope}[Meaning of the word closure]
Throughout this monograph, a \emph{conditional closure theorem} means that the conclusion follows from the displayed packet. It does not mean that the complete packet has already been populated. The results established here do not yet imply \cref{eq:target} without the remaining continuum and physical-contraction hypotheses.
\end{scope}
\sourcebasis{Local closure and physical source criteria are established in \src{F1}{V1--V4}; the nested Gaussian construction and representative-dependent response in \src{F11}{V31--V41}. The common finite source and physical-calibration framework is developed in \cite{GT}.}

\chapter{A short reconstruction of the finite Yang--Mills loading}
\label{ch:finite}
\section{From histories to the readable finite state}
The Grand Tensor begins with a typed event grammar and actual occurrence records. Two past histories are predictively equivalent when every admitted record-conditioned future has the same law. The quotient is the coarsest persistent predictive carrier. A finite ledger is a compiled realization of this quotient, not an independently postulated hidden state space.

Prediction alone does not supply spatial routing, orientation, or physical time. The readable enrichment retains the same-cut realization, protected orientation, root and scheduler marks, calibration, determining-field data, and action records needed by the downstream construction. Forgetting those added readable marks recovers the original Store cylinders. Across cutoffs, old--detail polarization tables carry the refinement information. In the Yang--Mills specialization, this common interface fixes the finite gauge variables, their action, the admitted gauge-invariant writers, and the reflected time transfer. It is not necessary to repeat the Standard Model or spacetime classifications to formulate the Yang--Mills closure problem. \cite{GT}

\section{One law, its reductions, and its sources}
At a finite regulator let $X$ denote the field configuration, $\mu$ its probability law, and $\mathfrak A$ the admitted gauge-invariant source algebra. The essential discipline is that a source, a conditional expectation, a Schur complement, and a response derivative refer to this same law or to an explicitly specified pushforward of it.

A subgroup restriction, conditioning on a retained field, integrating out a detail field, and pushing a measure to a quotient are different operations. Gauge invariance alone does not make their resulting laws identical. In particular, a reduced Wilson fibre at fixed coarse background is not automatically the physical Perron ground law of a transfer operator.

\begin{proposition}[Exact Gaussian elimination]\label{prop:schur}
Let a quadratic form on retained and detail variables have matrix
\[
L=\begin{pmatrix}A&B\\B^*&C\end{pmatrix},\qquad C\succ0.
\]
The minimizing detail is $y=-C^{-1}B^*x$, and
\begin{equation}\label{eq:schur}
\ip{(x,y)}{L(x,y)}
=\ip{x}{(A-BC^{-1}B^*)x}
 +\ip{y+C^{-1}B^*x}{C(y+C^{-1}B^*x)}.
\end{equation}
Gaussian integration retains both the Schur form $A-BC^{-1}B^*$ and the determinant contribution $\frac12\log\det C$ to the effective action, up to background-independent normalization.
\end{proposition}
\begin{proof}
Expand the last quadratic term in \cref{eq:schur}. Its cross terms recover those of $L$, and its retained square cancels the subtracted Schur term. Integration of the centered detail Gaussian gives $(\det C)^{-1/2}$ times the fixed Gaussian normalization.
\end{proof}

This identity fixes two independent pieces of data: the retained response and the normalization. Keeping the first while discarding a background-dependent determinant changes the effective action. Repeating an exact elimination also requires the retained law from the preceding step, rather than an unrelated reference with the same quadratic Taylor jet.

\section{Half-slab amplitudes and the physical Doob law}
Let $F(y,x)$ be the amplitude for a half-slab, and define
\[
(Af)(y)=\int F(y,x)f(x)\,\dd\mu_0(x),\qquad T=A^*A.
\]
For $F\in L^2$, $A$ is Hilbert--Schmidt and
\begin{equation}\label{eq:HS}
\Tr T=\norm{A}_{\rm HS}^2=\int|F|^2,\qquad
\Tr(\wedge^rT)=\norm{\wedge^rA}_{\rm HS}^2\ge0.
\end{equation}
Thus the exterior-power trace coefficients are genuine positive finite data. They do not alone control the ratio of the first two singular values uniformly in the regulator.

Assume now that $F$ is strictly positive and continuous on compact spaces with full-support reference measures. Write its largest singular value as $s_1$, with positive normalized singular functions $u_0,v_0$. The joint probability law
\begin{equation}\label{eq:half-law}
\dd\Lambda(x,y)=s_1^{-1}u_0(x)F(y,x)v_0(y)
\,\dd\mu_0(x)\dd\mu_{1/2}(y)
\end{equation}
has marginals $u_0^2\mu_0$ and $v_0^2\mu_{1/2}$. Conditional expectation between the two marginals is unitarily equivalent to $s_1^{-1}A$.

\begin{theorem}[Half-slab correlation and full transfer]\label{thm:half}
Under these finite-regulator hypotheses, the centered half-slab maximal correlation is
\[
\rho_{1/2}=\frac{s_2}{s_1}.
\]
The centered normalized full transfer has norm $q=\rho_{1/2}^2$. For a full slab of physical duration $t$, its spectral mass is
\begin{equation}\label{eq:half-mass}
m_t=-\frac{2}{t}\log\rho_{1/2}.
\end{equation}
\end{theorem}
\begin{proof}
Multiplication by the positive Perron functions identifies the two marginal $L^2$ spaces with the original reference spaces. In these coordinates the conditional channel is $A/s_1$, the constant functions correspond to the top singular vectors, and its centered singular norm is $s_2/s_1$. Composing the half channels gives $A^*A/s_1^2$. Functional calculus gives \cref{eq:half-mass}.
\end{proof}

If a common measurable factor survives on both boundaries, the relevant constant projection is replaced by the projection $P_{\rm com}$ onto that common factor. With boundary projections $P_0,P_1$,
\[
\rho_{\rm rel}=\norm{(P_1-P_{\rm com})(P_0-P_{\rm com})}.
\]
The two-boundary averaging chain has relative gap $(1-\rho_{\rm rel})/2$. This is a gap relative to retained common labels. It does not justify deleting protected sectors and calling the remainder the full theory.

\subsection{The exact conditional-variance form}
The half-slab description has a useful variational formulation which
retains the physical Perron factors. Let $\mathsf Kf(y)=\E_\Lambda[f(X)\mid
Y=y]$ and let $P_{\rm slab}=\mathsf K^*\mathsf K$ on the initial marginal
$\pi=u_0^2\mu_0$. Write $\pi_{1/2}=v_0^2\mu_{1/2}$ for the
other marginal. The ground-transform identification above gives the normalized
full transfer, not an auxiliary conditional sampler.

\begin{proposition}[Half-slab variance identity]
\label{prop:half-variance}
For every $f\in L^2(\pi)$,
\begin{equation}\label{eq:half-variance}
\begin{aligned}
\ip{f}{(I-P_{\rm slab})f}_{L^2(\pi)}
&=\E_\Lambda\!\left[\Var_\Lambda(f(X)\mid Y)\right]\\
&=\inf_{h\in L^2(\pi_{1/2})}
  \E_\Lambda|f(X)-h(Y)|^2.
\end{aligned}
\end{equation}
In particular, a bound by $\gamma\Var_\pi(f)$ from below is equivalent
to $\norm{P_{\rm slab}|_{\one^\perp}}\le1-\gamma$. For a full physical
slab of duration $t$, $0<\gamma<1$ therefore gives
\begin{equation}\label{eq:half-variance-mass}
 m_t\ge-\frac1t\log(1-\gamma).
\end{equation}
\end{proposition}
\begin{proof}
Conditional expectation is the orthogonal projection of $f(X)$ onto
functions of $Y$. Its squared projection norm is $\norm{\mathsf Kf}^2$;
subtracting this from $\norm f^2$ proves both expressions in
\cref{eq:half-variance}. Positivity and self-adjointness of
$P_{\rm slab}$ identify the best centered form bound with one minus
its centered operator norm. Functional calculus proves
\cref{eq:half-variance-mass}. \src{F6}{V58}
\end{proof}

\begin{lemma}[Same-carrier density comparison]
\label{lem:half-density-comparison}
Let $\Lambda_0$ be a reference probability law on the same boundary
spaces and suppose
\[
0<m_0\le\frac{\dd\Lambda}{\dd\Lambda_0}\le M_0<\infty.
\]
If the reference conditional-variance form is at least
$\gamma_0\Var_{\pi_0}(f)$, then the physical form in
\cref{eq:half-variance} is at least
\begin{equation}\label{eq:half-density-comparison}
\frac{m_0}{M_0}\gamma_0\Var_\pi(f).
\end{equation}
\end{lemma}
\begin{proof}
The lower density bound applies before the infimum over $h$ in
\cref{eq:half-variance}. The upper marginal density bound gives
$\Var_\pi(f)\le M_0\Var_{\pi_0}(f)$ by minimizing over constants.
Combining the two estimates proves the assertion. \src{F6}{V58}
\end{proof}

The comparison must include the singular ground profiles in
\cref{eq:half-law}. For a Gibbs density ratio, $M_0/m_0$ is controlled
by the oscillation of the complete density difference, which can grow
with volume even when its local summands are small. The lemma gives an
explicit law-transport cost; it does not make that cost uniform. Local
conditional comparison or a signed form estimate is needed whenever a
global ratio loses the required physical margin.

\section{Physical update versus auxiliary dynamics}
For a normalized positive transfer $T\Omega=\Omega$, the Doob operator
\[
Pf=\Omega^{-1}T(\Omega f),\qquad \dd\pi=|\Omega|^2\dd\nu,
\]
is a self-adjoint Markov contraction. In its stationary two-sided path realization, let $U$ be the shift and $J$ the time-zero embedding. Then
\begin{equation}\label{eq:koopman}
J^*U^nJ=P^n,\qquad
J^*(U-I)^*(U-I)J=2(I-P).
\end{equation}
These are identities for the actual transfer. A spatial update chain, a conditional Gibbs sampler, or a comparison Witten operator can be useful upstream, but is not made into $-\log T/t$ by notation.

\begin{remark}[Why the finite recap is enough]
The rest of the monograph uses only the admitted gauge-invariant source algebra, the actual compact action and its exact pushforwards, the same-law reflected transfer, and the calibrated time. The upstream predictive and readable construction explains why these objects belong to one interface; it does not remove the analytic obligations of continuum Yang--Mills theory.
\end{remark}
\sourcebasis{Grand Tensor Yang--Mills loading, half-slab and Doob results \cite{GT}; Perron-law and same-generator analysis \src{F1}{V7--V11, V21}; exact half-slab variance and same-carrier density comparison \src{F6}{V58}.}

\chapter{Source-resolved coercivity and the law-matching problem}
\label{ch:source}
\section{The physical ground law cannot be guessed}
Let a symmetric transfer have the form
\[
T=e^{-G/2}\,\mathcal C\,e^{-G/2},\qquad Th=\lambda h,
\]
where $h>0$. Its one-slice law is proportional to $h^2\dd\mathfrak m$. A candidate density $Z^{-1}e^{-F}\dd\mathfrak m$ agrees with that law precisely when
\begin{equation}\label{eq:law-incidence}
\mathcal C(e^{-(G+F)/2})=\lambda e^{(G-F)/2}.
\end{equation}

\begin{proposition}[No naive same-exponent identification]\label{prop:law-no-go}
If $\mathcal C1=1$ and $\mathcal C$ is injective, then the candidate with $F=G$ is the Perron law if and only if $G$ is constant almost everywhere.
\end{proposition}
\begin{proof}
With $F=G$, \cref{eq:law-incidence} gives $\mathcal C(e^{-G})=\lambda1$. Hence $\mathcal C(e^{-G}-\lambda)=0$. Injectivity forces $e^{-G}=\lambda$. The converse is immediate.
\end{proof}

The temporal Wilson convolution in the stated finite setting has strictly positive character multipliers, so this obstruction applies to its full configuration-space law. A conditional identity at fixed coarse field is a different statement and must be checked by disintegration. The valid options are to prove that conditional incidence, transport through an explicit density defect with a controlled loss, or rebuild the source calculus directly on the Perron law. \src{F1}{V21}

\section{Scores, curvature, and the complete signed source}
For a normalized conditional density proportional to $e^{-S(a,x)}$, parameter differentiation gives
\begin{equation}\label{eq:score}
\partial_a\E_a f=\E_a(\partial_af)-\Cov_a(f,\partial_aS),
\end{equation}
and the effective action $\mathcal F(a)=-\log\int e^{-S(a,x)}\dd x$ satisfies
\begin{equation}\label{eq:free-hess}
\partial_b\partial_a\mathcal F
=\E_a(\partial_b\partial_aS)-\Cov_a(\partial_aS,\partial_bS).
\end{equation}
The covariance is a signed part of the complete curvature source. Bounding a direct term and later charging the same covariance again changes the operator. The Perron and sequential-marginal analogues retain the terminal Perron weight and its derivatives; the unweighted conditional fibre is not a replacement for them.

For the moving conditional projection $D_i=I-\E_i$, with normalized conditional density $\rho_i$, the source identity is
\begin{equation}\label{eq:moving-score}
X_aD_if=D_iX_af-\Cov_i(f,X_a\log\rho_i).
\end{equation}
Different sign conventions for the density require the corresponding score sign. The identity, not a memorized sign, determines the response.

\subsection{Higher normalized responses}
The first-order score identity extends exactly to the higher responses
needed by blocking. Work on a fixed finite conditional chart with
normalized law
\[
\dd\nu_a(x)=Z(a)^{-1}e^{-S(a,x)}\dd m(x).
\]
Assume the indicated derivatives may pass under the integral. A moving
reference density is included in $S$ after transport to this chart; it
cannot be omitted from the score. For a finite set $I$ of labelled,
commuting parameter directions, write $\partial_I=\prod_{i\in I}\partial_i$,
$F_J=\partial_JF$, and $S_B=\partial_BS$. Denote joint cumulants under
$\nu_a$ by $\kappa_a$, with $\kappa_a(G)=\E_aG$, and let $\Pi(A)$ be
the set of partitions of $A$.

\begin{theorem}[Complete connected score hierarchy]
\label{thm:connected-score-hierarchy}
With the empty-partition convention,
\begin{equation}\label{eq:connected-score-hierarchy}
\partial_I\E_aF
=\sum_{J\subseteq I}\ \sum_{\pi\in\Pi(I\setminus J)}
 (-1)^{|\pi|}\kappa_a\bigl(F_J,(S_B)_{B\in\pi}\bigr).
\end{equation}
In particular, a third derivative of a normalized current can contain a
four-entry connected cumulant. All denominator derivatives are already
included in this identity.
\end{theorem}
\begin{proof}
Differentiate the logarithm of the joint moment-generating function to
obtain
\[
\partial_i\kappa_a(G_1,\ldots,G_k)
=\sum_{\ell=1}^k\kappa_a(G_1,\ldots,\partial_iG_\ell,\ldots,G_k)
 -\kappa_a(G_1,\ldots,G_k,S_{\{i\}}).
\]
For $k=1$ this is \cref{eq:score}. Induct on $|I|$: a new direction
joins $J$, joins an existing score block, or creates a singleton score
block. These possibilities enumerate each term in the new partition
sum exactly once, with the sign displayed in
\cref{eq:connected-score-hierarchy}. \src{F8}{V60}
\end{proof}

For example, writing subscripts for parameter derivatives, the second
response is
\begin{equation}\label{eq:score-second-explicit}
\begin{aligned}
\partial_b\partial_a\E F
={}&\E F_{ab}-\kappa(F_a,S_b)-\kappa(F_b,S_a)\\
 &-\kappa(F,S_{ab})+\kappa(F,S_a,S_b).
\end{aligned}
\end{equation}
At third order the term with three singleton score blocks is
$-\kappa(F,S_a,S_b,S_c)$. Pair-covariance control cannot pay for this
term. The four-mark criterion in \cref{thm:four-mark-clustering} supplies
a sufficient connected estimate once its complete source-dependent
activity hypothesis is verified.

\begin{scope}[An exact hierarchy is not an all-order norm estimate]
Formula \cref{eq:connected-score-hierarchy} holds at every fixed order
under its differentiation hypotheses. Bounding its partitions separately
does not prove a uniform exponential-grade analytic estimate. The
factorial restrictions in \cref{ch:trees} and the radius budgets in
\cref{ch:forests} continue to apply.
\end{scope}

\section{A physical response-matrix compiler}
Suppose the physical ground-state representation has Dirichlet form
\begin{equation}\label{eq:groundform}
\ip{\Omega f}{(H-E_0)\Omega f}
=\kappa_H\int|\nabla f|^2\dd\pi.
\end{equation}
For the sequential filtration $\mathcal F_i$, let $M_i f=\E[f\mid\mathcal F_i]$ and $D_i=M_i-M_{i-1}$. The exact variance ledger is
\[
\Var_\pi(f)=\sum_i\norm{D_if}_{L^2(\pi)}^2.
\]

\begin{criterion}[Sequential source response]\label{thm:response-matrix}
Assume every sequential single-coordinate conditional has Poincar\'e floor at least $\rho_{\rm site}>0$, and a nonnegative matrix $\mathsf A$ satisfies
\[
\norm{\nabla_iM_if}_{L^2(\pi)}
\le\sum_bA_{ib}\norm{\nabla_bf}_{L^2(\pi)}
\]
for every admitted smooth gauge-invariant $f$. Then
\begin{equation}\label{eq:response-gap}
\Var_\pi(f)\le\frac{\norm{\mathsf A}_{2\to2}^2}{\rho_{\rm site}}
\int|\nabla f|^2\dd\pi,\qquad
\gap(H)\ge\frac{\kappa_H\rho_{\rm site}}{\norm{\mathsf A}_{2\to2}^2}.
\end{equation}
\end{criterion}
\begin{proof}
Conditional Poincar\'e gives $\norm{D_if}^2\le\rho_{\rm site}^{-1}\norm{\nabla_iM_if}^2$. Apply the matrix norm to the vector of gradient norms, sum the martingale differences, and use \cref{eq:groundform}.
\end{proof}

A simpler sufficient packet is $\norm{D_if}^2\le\sum_bw_{ib}\int|\nabla_bf|^2\dd\pi$ with bounded column congestion $C_{\rm col}=\sup_b\sum_iw_{ib}$. It gives $\gap(H)\ge\kappa_H/C_{\rm col}$. Conditional site gaps alone are insufficient: their responses to the previously exposed field must also be controlled. The source-resolved Green-column construction is designed to populate these response matrices on the actual signed source range, rather than by the norm of a full inverse. \src{F1}{V7, V14--V16}

\section{Comparison-minus-defect and physical landing}
On a declared one-form source carrier, write
\[
L=\mathcal A-\mathcal V\mathcal V^*,\qquad \mathcal A\succ0,
\]
and define
\begin{equation}\label{eq:woodbury-objects}
\mathcal K=\mathcal V^*\mathcal A^{-1}\mathcal V,
\quad \mathcal B=\mathcal A^{-1}\mathcal V,
\quad \mathcal M=\mathcal B^*\mathcal B.
\end{equation}
When the inverse is defined and $\norm{\mathcal K}<1$, Woodbury gives
\[
L^{-1}=\mathcal A^{-1}
+\mathcal A^{-1}\mathcal V(I-\mathcal K)^{-1}\mathcal V^*\mathcal A^{-1}.
\]
The exact identities are more informative than this sufficient Neumann condition.

\begin{theorem}[Channel-to-physical landing]\label{thm:landing}
For $\omega_z=\mathcal Bz$,
\begin{equation}\label{eq:landing}
\begin{aligned}
\norm{\omega_z}^2&=\ip{z}{\mathcal Mz},\\
\ip{\omega_z}{\mathcal A\omega_z}&=\ip{z}{\mathcal Kz},\\
\ip{\omega_z}{L\omega_z}&=\ip{z}{\mathcal K(I-\mathcal K)z}.
\end{aligned}
\end{equation}
On the closure of $\Ran\mathcal B$ in the comparison-form norm,
\begin{equation}\label{eq:relative-edge}
1-\norm{\mathcal K}
=\inf_{\omega\ne0}
\frac{\ip{\omega}{L\omega}}{\ip{\omega}{\mathcal A\omega}}.
\end{equation}
\end{theorem}
\begin{proof}
The first two equalities are direct substitutions. Also $\mathcal V^*\omega_z=\mathcal Kz$, giving the third. The supremum of $\ip{z}{\mathcal K^2z}/\ip{z}{\mathcal Kz}$ is $\norm{\mathcal K}$ by spectral calculus, which proves \cref{eq:relative-edge}.
\end{proof}

A near-unit channel sequence therefore need not be a physical infrared sequence. Its landing norm can degenerate while its comparison energy escapes to high frequencies. A uniform landing frame or a bounded comparison-energy window is needed before a channel survivor gives a normalized physical soft vector. Conversely, a uniform bound on the complete defect synthesis can force the required landing window. The source-carrier Combes--Thomas estimate then uses a genuine spectral edge and a controlled conjugation modulus to obtain spatially summable response columns. It does not manufacture the edge. \src{F1}{V16--V20}

\section{Bad-set capacity follows source incidence}
Let $B_i$ be a predictable bad event for martingale increment $D_i$. Define $Z_i(f)=\norm{\one_{B_i}D_if}$ and $X_b(f)=\norm{\nabla_bf}$. Suppose an assembled source vector $R(f)$ obeys
\[
Z(f)\le\mathsf I R(f)+\mathsf E X(f),\qquad R(f)\le\mathsf QX(f)
\]
entrywise, with nonnegative matrices. Then
\begin{equation}\label{eq:bad-capacity}
\sum_i Z_i(f)^2\le
\norm{\mathsf I\mathsf Q+\mathsf E}_{2\to2}^2\int|\nabla f|^2\dd\pi.
\end{equation}
With good-site floor $\rho_g$ and good response matrix $\mathsf A$, the hybrid physical gap is bounded below by
\begin{equation}\label{eq:hybrid-gap}
\gap(H)\ge
\frac{\kappa_H}{\rho_g^{-1}\norm{\mathsf A}_{2\to2}^2+
\norm{\mathsf I\mathsf Q+\mathsf E}_{2\to2}^2}.
\end{equation}
The first matrix is source incidence; the second is capacity routing. A small probability of $B_i$ is neither one. This distinction is the reason that a Fisher-information average or a Chebyshev estimate cannot replace a restricted operator estimate. \src{F1}{V19--V20}

\section{What these reductions actually close}
The source calculus closes the algebraic relation between physical response, signed curvature, comparison return, and the martingale gap. It also identifies exact failure witnesses. It does not yet supply a regulator-uniform response matrix or bad-capacity stock for the weak-coupling continuum trajectory.

The finite temporal Wilson-to-Hamiltonian construction is another retained result: calibrated character multipliers converge to their Casimir symbols on fixed representation banks, heat-kernel domination controls high representations, and a diagonal growing schedule gives the fixed-spatial-volume differential-Hamiltonian limit. Strong semigroup convergence at fixed volume is not a uniform thermodynamic mass theorem. The order of limits and the physical coefficient in \cref{eq:groundform} remain part of the application. \src{F1}{V10--V11}
\sourcebasis{The actual-law, sequential-response, source-carrier, and capacity results are developed in \src{F1}{V6--V22}; the moving conditional calculus is continued in \src{F6}{terminal developments} and \src{F7}{opening conditional analysis}.}

\chapter{Harmonic normalization and connected-source estimates}
\label{ch:trees}
\section{Peter--Weyl bookkeeping before inequalities}
The earliest response programme studies products of representation-valued sources and the cancellation of disconnected partitions. A recurring obstruction is incorrect normalization: dimensions, raw Fourier coefficients, representation probabilities, and marked-coordinate probabilities cannot be interchanged.

Let
\[
\bigotimes_{i=1}^m H_{R_i}
=\bigoplus_U W_U(H_U\otimes M_U),\qquad
X_U=W_U^*\left(\bigotimes_i A_i\right)W_U,
\]
where $W_U$ denotes the isometric block embedding. With the Fourier and kernel conventions of the source, the exact coefficient is
\begin{equation}\label{eq:PW}
\widehat h(U)=\frac{d_{\rm in}}{d_U}
\Tr_{M_U}\bigl[X_U(I\otimes K_U)\bigr],\qquad
 d_{\rm in}=\prod_i d_{R_i}.
\end{equation}
The trace-class pinching inequality is
\begin{equation}\label{eq:pinching}
\sum_U\norm{X_U}_1\le\prod_i\norm{A_i}_1.
\end{equation}
The output factor $d_U$ cancels only against its prescribed Fourier weight. A dual-singlet test retains raw mass $1/d_R$; replacing it by a probability inferred from representation dimensions changes the estimate. These elementary tests are retained because later tree estimates depend on their exact grade. \src{F1}{harmonic source programme}; \src{F2}{opening normalization audit}

For nonnegative coordinate stocks $a_{i,e}$ define
\begin{equation}\label{eq:marks}
\Xi_i=\sum_e a_{i,e},\quad
H_{ij}=\sum_e a_{i,e}a_{j,e},\quad
\theta_{ij}=\frac{H_{ij}}{\Xi_i\Xi_j}\le1.
\end{equation}
Zero-stock sources are removed before normalization. The $\theta_{ij}$ are genuine probability-mark overlaps. A tree numerator in the unnormalized $H_{ij}$ is not yet the normalized marked-tree statement in the $\theta_{ij}$.

\section{First connectivity and the quartic closure}
A connected coordinate word has a unique first coordinate at which the current row partition becomes connected. Before that coordinate, the word factors over a proper partition. At that coordinate, the complete joining cumulant acts. After it, the suffix is a complete moment. The final resolvent is paid once after this assembly, rather than separately in every predecessor block.

The quartic problem is resolved only after the normalization and suffix corrections. Its fourteen proper predecessor partitions consist of four of type $3+1$, three of type $2+2$, six of type $2+1+1$, and one all-singleton partition. Certain literal local weighted estimates fail because a suffix can rescue a channel that is small in its prefix. The successful bounds sum the rescued square globally over coordinates before taking the final norm.

\begin{theorem}[Completed fixed-arity quartic marked-tree estimate]\label{thm:quartic}
For the declared normalized quartic source problem, with its single final resolvent and probability-mark convention, the complete connected coefficient satisfies
\begin{equation}\label{eq:quartic}
\kappa_{\otimes}^{\rm abs}(\mathcal C^{\rm pay}_4)
\le C_{G,s}\sum_{T\in\mathcal T_4}\prod_{\{i,j\}\in T}\theta_{ij}.
\end{equation}
Here $\kappa_{\otimes}^{\rm abs}$ is the source-defined absolute product-state coefficient capacity, and $C_{G,s}$ has the group and regulator dependence specified in that construction.
\end{theorem}
\begin{proof}[Proof structure]
Split at first connectivity. Use the corrected normalized $3+1$ atlas and the globally coordinate-summable rescued squares for the $3+1$ and $2+1+1$ cases; treat $2+2$ and the singleton case by their completed collision bounds. Summing the fourteen proper partitions and retaining one final payer yields \cref{eq:quartic}. The conclusion is the normalized coefficient-capacity estimate, not an unnormalized operator-tree numerator; the globally summed suffix is essential. \src{F3}{V89--V96}; \src{F4}{inherited quartic ledger}
\end{proof}

At arity five there are fifty-one proper predecessor partitions and $5^3=125$ labeled trees. The local cumulant and projected-tail development closes substantial fixed-arity estimates, including the Gaussian tangent degree filtration. It also produces counterexamples: one Casimir payment per input is not sufficient in the stated coherent $\SU(3)$ packet; a whole $\SU(2)$ heat slot retains singlet capacity of order $(1+C_2)^{-1/2}$; and endpoint leakage cannot be removed by separately bounding merge differences. These results constrain the all-order extension rather than establishing it. \src{F3}{V97--V100}; \src{F4}{quintic and projected-tail developments}

\section{The factorial-grade correction}
The all-order question is not settled by proving each fixed arity with a large constant.

\begin{definition}[Strong and factorial-graded tree bounds]\label{def:factorial}
A strong marked-tree bound has, for each fixed labeled tree,
\[
\norm{\mathcal K_{m,T}}\le K_0^{m-1}\mathcal W_T.
\]
A factorial-graded bound instead has
\[
\norm{\mathcal K_{m,T}}\le m!K^m\mathcal W_T.
\]
The distinction is uniform in $m$, not a change in the numerical constant at one arity.
\end{definition}

\begin{proposition}[The extra factorial cannot be hidden]\label{prop:factorial}
The all-order $\SU(2)$ estimates of \src{F5}{V1--V11} have the factorial-graded form. They do not imply the strong marked-tree estimate or its uniform field-cumulant consequence.
\end{proposition}
\begin{proof}
The exact predecessor-partition sum is
\begin{equation}\label{eq:partition-factorial}
\sum_{\rho\in\Pi([m])}|\rho|!\prod_{B\in\rho}|B|!
=m!2^{m-1}.
\end{equation}
Rooted labeled-tree enumeration already uses the allowed factorial,
\[
m^{m-1}\le e^{m-1}m!.
\]
A factorial per fixed tree therefore produces order $(m!)^2K_1^m$ after spatial compilation. No fixed exponential base absorbs the extra $m!$. The coefficient $1/m!$ in an exponential generating series cancels only one factorial; the remaining ordinary series is still of factorial grade. This is the factorial-grade obstruction; it leaves the finite inequalities valid but rules out the stronger uniform analytic inference. \src{F5}{V12}
\end{proof}

The correction leaves the displayed all-order inequalities valid. It removes only their overstrong analytic interpretation. It is therefore inappropriate either to discard the finite estimates or to describe the strong all-order source gate as closed.

\section{Exact Haar assembly recovers one exponential-grade component}
The $\SU(2)$ exponential-chart Jacobian, in the source convention, is
\begin{equation}\label{eq:Haar-series}
J_\epsilon(y)=\left(\frac{\sin(\sqrt\epsilon|y|)}{\sqrt\epsilon|y|}\right)^2
=\sum_{n\ge0}\frac{(-1)^n2^{2n+1}}{(2n+2)!}\epsilon^n|y|^{2n}.
\end{equation}
Its logarithmic vertices have large isolated high derivatives; bounding them before recombination loses the factorial denominators of the full Jacobian.

\begin{theorem}[Integrated exact-Haar derivative bound]\label{thm:Haar}
There are $C<\infty$ and $\epsilon_0>0$ such that for every centered Gaussian $Y\in\R^3$ with covariance bounded by $I$, every $0<\epsilon\le\epsilon_0$, and every $v\ge1$,
\begin{equation}\label{eq:Haar-bound}
\E\norm{\nabla^vJ_\epsilon(Y)}_{\rm row}
\le C^v\epsilon^{v/2}.
\end{equation}
The zeroth-order expectation is bounded by a fixed constant.
\end{theorem}
\begin{proof}[Proof outline]
Differentiate \cref{eq:Haar-series} term by term. The radial derivative row is at most
\[
C^v\frac{(2n)!}{(2n-v)!}(1+|y|)^{2n-v}.
\]
Gaussian moments give $\E(1+|Y|)^q\le C_0^q q^{q/2}$. The denominator $(2n+2)!$ cancels the dangerous derivative factorial, while $q^{q/2}/q!\le e^q$. The sum begins at $n\ge\lceil v/2\rceil$ and is bounded by a convergent geometric series in $C_1\epsilon$, giving \cref{eq:Haar-bound}. \src{F5}{V12}
\end{proof}

This closes the exact-Haar component, not the whole Wilson-density transfer. The exact action, normalization, assembled interpolation, and tree assignment must still have the correct fixed-tree grade.

\section{Forest projection versus physical source occurrences}
The later compact heat and L\'evy formulation replaces several coefficientwise arguments by an exact probabilistic interpolation. Optional support labels are summed before norms using convex contractions of the form $(1-\lambda)I+\lambda D^R(U)$. At the actual hierarchical interpolation matrices, the resulting weights are subprobabilities rather than Bell-number counts. Weighted Cayley disintegration, random-root orientation, and marginalization of unused completion branches give the strong vertex-accreting forest-projection estimate with total base $K^m$.

A physical source occurrence is not an unused completion branch. It owns a disjoint nonempty subforest and can carry a forced component excess. Marginalizing that ownership would remove a real derivative payment. Thus the complete forest-projection result and the source-occurrence closure are different achievements. \src{F6}{V1--V7}

One useful conditional reduction is as follows. On a carrier graph of degree $\Delta$, suppose source activities $z_A$ satisfy, for each component $\alpha$ and connected carrier $\Gamma$,
\begin{equation}\label{eq:carrier-profile}
\sum_{\substack{A\ni\alpha\\\Gamma(A)=\Gamma}}|A|z_A
\le C_0|\Gamma|^r a^{|\Gamma|}.
\end{equation}
Then the all-rank incidence row is bounded by
\begin{equation}\label{eq:carrier-row}
\kappa_{\rm hyp}\le C_0a\sum_{n\ge1}n^r(\Delta^2a)^{n-1},
\qquad \Delta^2a<1.
\end{equation}
The graph enumeration and the implication are closed. The physical carrier-activity profile \cref{eq:carrier-profile} is an input, not a consequence of the abstract forest projection. \src{F6}{V16}
\sourcebasis{The harmonic, collision, arity, and source-occurrence developments are distributed across \cite{F1,F2,F3,F4,F5,F6}. The factorial-grade restriction of \src{F5}{V12} governs every use of the all-order bound.}

\chapter{Conditional geometry, collars, and bad-field localization}
\label{ch:geometry}
\section{Local control must include the conditioned exterior}
An isolated tree gauge can simplify a finite block, but changing that gauge may also change its exterior data. Such a transformation relates conditional fibres; it is not necessarily a symmetry within one fixed conditional fibre. Weak-coupling exterior staples can create competing central wells, so an unconditional good-field probability cannot justify a uniform conditional convexity assertion.

The source analysis keeps the fixed exterior explicit. For six $\SU(3)$ staples $S_j$, set
\[
F=\min_V\sum_{j=1}^{6}\bigl(3-\Re\Tr(VS_j)\bigr).
\]
After the declared alignment, the staple sum differs from $6I$ by at most $\sqrt{12F}$ in Frobenius norm. A sufficiently small value of $F$ gives the single-well local quadratic control used by the good-field analysis. The direct star defect
\[
D_e=\sum_{j=1}^{6}(3-\Re\Tr S_j)
=\frac12\sum_{j=1}^{6}\norm{S_j-I}_{\rm F}^2
\]
satisfies $F_e\le D_e$. These are local, conditioned statements. \src{F6}{terminal staple analysis}; \src{F7}{opening geometric audit}

\section{The finite collar reconstruction}
The one-layer collar used in the source consists of $72$ plaquettes, $123$ edges, and $64$ vertices. Its scalar curvature map has rank $60$ and kernel equal to the $63$-dimensional gauge image. Fixing the full tree leaves $60$ chord variables. The reported integer certificate contains a $60\times60$ minor with determinant $-1$ and $102$ nonzero entries, giving the explicit bound
\begin{equation}\label{eq:collar-sigma}
\sigma_{\min}^2\ge102^{-59}.
\end{equation}
A $60$-generator, $72$-relator elimination certificate identifies the relevant complex as simply connected. Consequently the zero-curvature completion is unique after all prescribed tree frames are fixed, and a small-defect estimate extends to
\begin{equation}\label{eq:collar-distance}
\dist(U,\mathcal Z_{\rm flat})^2\le C_{\rm glob}D(U)
\end{equation}
on the declared finite collar with its boundary-frame convention. These are strong finite reconstruction results. The finite certificate is reported in the cited construction and is not independently recomputed here.

The geometric reconstruction does not close three further estimates: restricted smoothing for the full conditional action, stability of the physical response and ground law, and bounded-congestion coverage of the translated collars with the required physical form incidence. Keeping these separate prevents a topological statement from being used as a global Poincar\'e theorem. \src{F6}{final collar ledger}; \src{F7}{initial continuation}

\begin{proposition}[No finite Wilson-action-closed update block]\label{prop:no-finite-block}
On the infinite lattice there is no nonempty finite edge set closed under adding every edge of every plaquette touching an updated edge.
\end{proposition}
\begin{proof}[Geometric argument]
Take an updated edge at an extremal coordinate of the finite set. A plaquette containing it and extending in an available transverse direction introduces an edge beyond that extremum. Repeated plaquette closure therefore cannot terminate in a nonempty finite set. The conditional formulation must retain the boundary interactions or control their influence; it cannot assume a literal finite action-closed block. \src{F7}{V1}
\end{proof}

\section{Smoothing and rarity are different estimates}
Let $\E$ be the conditional projection, $D=I-\E$, and $P$ a conditional Markov smoothing with $\E P=P\E=\E$. Suppose its restricted bad-set multiplier has the bound encoded by $\eta_B$. The proved projected smoothing inequality is
\begin{equation}\label{eq:smoothbad}
\norm{M_BDF}^2
\le(1+\theta)\eta_B\norm{DF}^2
 +(1+\theta^{-1})\ip{F}{(I-P)F},\qquad\theta>0.
\end{equation}
Its useful datum is the restricted operator stock $\eta_B$, not merely $\mu(B)$. Ground transformation also changes the conditional projections and their multipliers by explicit Perron factors. These factors must be retained before comparing bad capacity between a reference law and the physical one.

\subsection{A quantitative core--bad-set compiler}
The restricted multiplier in \cref{eq:smoothbad} can be combined with a
conditional core gap without treating rarity as a functional inequality.
Let $\mathsf E_e$ be conditional expectation on a declared enlarged
block, $\mathsf D_e=I-\mathsf E_e$, and $\mathsf P_e$ a positive
self-adjoint conditional Markov transfer with
$\mathsf E_e\mathsf P_e=\mathsf P_e\mathsf E_e=\mathsf E_e$.
All operators act on the same physical law $\mu$. Let $B_e$ be the bad
set, $K_e=B_e^c$, and let $w_e\ge0$ be the incidence weights. Norms of
conditional operators below are uniform over the admitted exteriors.

\begin{lemma}[Gluing the conditional core mean]
\label{lem:core-mean-gluing}
For a probability law $\nu$, $p=\nu(B)<1$, $q=1-p$, and $\nu g=0$,
\begin{equation}\label{eq:core-mean-gluing}
\int_{B^c}|g|^2\dd\nu
\le q\Var_{\nu(\cdot\mid B^c)}(g)
  +\frac pq\int_B|g|^2\dd\nu.
\end{equation}
Moreover, if $P$ is Markov, then
$\nu(B)\le\norm{M_BP}^2$.
\end{lemma}
\begin{proof}
Writing the two conditional means as $m_K,m_B$, the centering identity
$q m_K+p m_B=0$ gives
$q|m_K|^2\le(p/q)\int_B|g|^2\dd\nu$. Add the conditional variance.
For the second assertion apply $M_BP$ to the unit constant, using
$P\one=\one$. \src{F6}{V97}
\end{proof}

\begin{criterion}[Core gap and restricted smoothing]
\label{crit:core-smoothing}
Let $\mathcal E_\mu$ be the declared physical energy form. Assume
approximate tensorization and core incidence in the form
\begin{align}
\Var_\mu(f)&\le C_*\mathcal S(f),
 &\mathcal S(f)&=\sum_e w_e\norm{\mathsf D_ef}^2,
 \label{eq:core-tensorization}\\
\sum_e w_e\E\!\left[q_e\Var_{\mu_e(\cdot\mid K_e)}
                    (\mathsf D_ef)\right]
 &\le C_{\rm core}\mathcal E_\mu(f),
 &q_e&=\mu_e(K_e). \label{eq:core-incidence}
\end{align}
Assume also
\begin{equation}\label{eq:core-transfer-incidence}
\sup_e\norm{M_{B_e}\mathsf P_e}^2\le\eta<1,
\qquad
\sum_e w_e\ip{f}{(I-\mathsf P_e)f}
\le m_P\delta\,\mathcal E_\mu(f).
\end{equation}
Here $\delta$ is the transfer duration in the declared incidence
normalization. For any $\theta>0$ with $(2+\theta)\eta<1$,
\begin{equation}\label{eq:core-smoothing-gap}
\Var_\mu(f)
\le
\frac{C_*\{(1-\eta)C_{\rm core}
 +(1+\theta^{-1})m_P\delta\}}
 {1-(2+\theta)\eta}\,\mathcal E_\mu(f).
\end{equation}
\end{criterion}
\begin{proof}
Set $\mathcal A(f)=\sum_e w_e\norm{M_{B_e}\mathsf D_ef}^2$.
On each conditional fibre, \cref{lem:core-mean-gluing} and
$p_e\le\eta$ give
\[
\mathcal S(f)\le C_{\rm core}\mathcal E_\mu(f)
 +\frac{\mathcal A(f)}{1-\eta}.
\]
The projected smoothing estimate and
\cref{eq:core-transfer-incidence} imply
\[
\mathcal A(f)\le(1+\theta)\eta\mathcal S(f)
 +(1+\theta^{-1})m_P\delta\,\mathcal E_\mu(f).
\]
Substitute, multiply by $1-\eta$, and absorb the coefficient
$(1+\theta)\eta$. The positive remaining coefficient is
$1-(2+\theta)\eta$. Apply \cref{eq:core-tensorization}.
\src{F6}{V97}
\end{proof}

If the normalized physical core differs from a bare core with
Poincar\'e floor $s c_a$ by a potential of oscillation at most $K$,
and the gradient incidence is $m_{\rm core}$, the source comparison
supplies
\begin{equation}\label{eq:core-oscillation}
C_{\rm core}\le\frac{m_{\rm core}e^K}{s c_a}.
\end{equation}
This is useful only when the oscillation bound includes the crossing
plaquettes, response terms, and Perron profile of the complete
conditional action. Neither the collar rank nor the bare Wilson
minimum proves such a bound. Likewise, the approximate tensorization
and both incidences in \cref{crit:core-smoothing} are inputs, not
consequences of conditional site gaps. The criterion gives the exact
constant obtained when these distinct physical estimates are available.

\section{A curvature-to-defect bridge}
The later exact-block analysis constructs tree--chord fillings that express local displacement in terms of plaquette curvature. In the stated saturation convention,
\begin{equation}\label{eq:filling}
\sum_{y\in Y}D_y(z;W)
\le20\sum_{p\in\mathcal P(Y)}\norm{U_p-I}_{\rm F}^2.
\end{equation}
For a coarse-flat background, including the declared commuting-toron family, the zero-defect set is the unique moving section. For a nonflat coarse background it is empty. The section is therefore background dependent and cannot be replaced by a fixed coordinate origin without paying the induced response.

If a connected set $Y$ is $r$-bad, the filling estimate gives the Wilson-action lower bound
\begin{equation}\label{eq:Peierls}
S_{\rm W}(Y)\ge\frac{r^2}{120}|Y|.
\end{equation}
With $r_\beta=\beta^{-2/5}$ and the source small-ball normalization, the primitive far-field stock is
\begin{equation}\label{eq:qfar}
q_{\rm far}\le C_-^{-45}(\beta/120)^{180}
\exp(-\beta^{1/5}/120),
\end{equation}
while the regular Wilson remainder satisfies
\begin{equation}\label{eq:etaW}
\eta_{\rm W}\le7680\pi^3\beta^{-1/5}.
\end{equation}
The power, volume cost, and exponential belong to the same normalized conditional integral. A bare action penalty alone is not the probability estimate. \src{F8}{curvature localization and final good--bad development}; \src{F9}{initial compact rebase}

\section{What is transported, and what is not}
The collar and curvature estimates provide a finite geometric scaffold for good-field localization, rare exits, and analytic normalization. They do not bound the complete retained slow field after conditional elimination. The identity
\begin{equation}\label{eq:totalcov}
\Cov(F,G)=\E\bigl[\Cov(F,G\mid W)\bigr]
+\Cov\bigl(\E[F\mid W],\E[G\mid W]\bigr)
\end{equation}
exhibits the missing slow contribution directly. Strong mixing inside every detail fibre leaves the second term untouched. The purpose of the renormalization construction is to keep, rather than silently erase, that retained law.
\sourcebasis{Conditional score and capacity distinctions: \cite{F6,F7}. Curvature localization, moving sections, and bad-field normalization: \cite{F8,F9}.}

\chapter{Exact compact blocking and the Gaussian principal part}
\label{ch:blocking}
\section{The action is the logarithm of the exact pushforward}
For a declared compact block map $P:X\to Y$, disintegrate the reference Haar measure and define the coarse action by
\begin{equation}\label{eq:exact-block}
e^{-S_+(y)}\dd\nu_+(y)
=P_*\bigl(e^{-S(x)}\dd\nu(x)\bigr),
\end{equation}
with its normalization fixed. This is the definition of an exact blocking step. It carries all generated interactions, the appropriate Jacobian, and any retained sectors. The quadratic principal part is extracted from this object; it is not used to redefine it.

For the nonlinear root-path block, \cref{thm:global-pivot,cor:compact-product}
construct such a pushforward explicitly on a smooth compact carrier.
The Haar product is retained as an inverse-pivot amplitude; its bounds
are local per factor, not a uniform bound on the total product. Agreement
with the original block near the identity is an equality of germs, with
the separate global-law comparison specified in
\cref{scope:completion-germ}.

When a tree--chord change of variables yields product Haar measure, its normalization is exact. Overlapping path blocks do not automatically have a product law: shared edges require their common conditional kernel. Under absolute convergence, regrouping the interaction by its declared coarse footprint is exact. For the selected-tree construction in the source, the resulting norm has the schematic but explicit form
\begin{equation}\label{eq:selected-tree}
\mathfrak O_\mu(\overline\Phi)
\le D_{\rm dec}\left(\rho+\frac{rs}{1-\epsilon}\right),
\end{equation}
with $\rho,r,s,\epsilon$ the inherited primitive, marked-response, and propagation stocks. This implication requires the full source-dependent conditional polymer hypothesis. A pair-covariance estimate alone is not sufficient for all higher cumulants. \src{F7}{V92--V94}

\subsection{The exact retained sector and conditional remainder}
Let $\mathcal R_P(S)=S_+$ denote the fixed-reference blocking map in
\cref{eq:exact-block}, and let $\nu_{S,y}$ be its normalized conditional
law. Define retained pullback and conditional expectation by
\[
(\mathsf J_P q)(x)=q(Px),\qquad
(\mathsf K_S h)(y)=\E_{\nu_{S,y}}h.
\]
The reference measures and block map are held fixed in the following
variation. Any later calibration or reference reset has its own
measured contribution.

\begin{theorem}[Exact differential and retained-pullback sector]
\label{thm:exact-block-differential}
For bounded perturbations $h_1,\ldots,h_n$,
\begin{equation}\label{eq:exact-block-cumulants}
D^n\mathcal R_{P,S}[h_1,\ldots,h_n](y)
=(-1)^{n-1}\kappa_{S,y}(h_1,\ldots,h_n).
\end{equation}
In particular, $D\mathcal R_{P,S}=\mathsf K_S$ and
\begin{equation}\label{eq:retained-pullback}
\mathcal R_P(S+\mathsf J_Pq)=\mathcal R_P(S)+q,
\qquad \mathsf K_S\mathsf J_P=I.
\end{equation}
Every bounded real perturbation has the exact decomposition
$h=\mathsf J_Pm+v$, where $m=\mathsf K_Sh$ and $\mathsf K_Sv=0$, and
\begin{equation}\label{eq:exact-centered-remainder}
\mathcal R_P(S+h)-\mathcal R_P(S)
=m-\log\E_{\nu_{S,y}}e^{-v}.
\end{equation}
If $\delta_v(y)$ is the essential range length of $v$ on the fibre,
then
\begin{equation}\label{eq:centered-range-bound}
-\frac{\delta_v(y)^2}{8}
\le\mathcal R_P(S+h)(y)-\mathcal R_P(S)(y)-m(y)\le0.
\end{equation}
\end{theorem}
\begin{proof}
The difference of blocked actions is the negative logarithm of
$\E_{\nu_{S,y}}\exp(-\sum_i t_i h_i)$. Its mixed derivatives give
\cref{eq:exact-block-cumulants}. A retained pullback is constant on each
fibre and factors out of this expectation, proving
\cref{eq:retained-pullback,eq:exact-centered-remainder}. Jensen's
inequality gives the upper bound in \cref{eq:centered-range-bound}.
For the lower bound let $\psi(t)=\log\E e^{-tv}$. Its tilted variance
obeys $\psi''(t)\le\delta_v(y)^2/4$, while
$\psi(0)=\psi'(0)=0$. Integrating twice up to $t=1$ gives
$\psi(1)\le\delta_v(y)^2/8$. \src{F7}{V99}
\end{proof}

If the chosen fine and coarse norms make $\mathsf J_P$ isometric,
\cref{eq:retained-pullback} forces
$\norm{D\mathcal R_{P,S}}\ge1$. Thus strict contraction on the complete
interaction space is excluded by an exact identity, not merely by a
weak estimate. Conditional centering removes the first-order detail
response, but the pointwise bound \cref{eq:centered-range-bound} is not
a volume-uniform polymer norm: its range may be extensive. Calibrated
marginal transport, local marked estimates, and stable-complement
contraction remain separate parts of the renormalization argument.

\section{Gaussian conditioning and exact harmonic extension}
Let $L$ be positive on the declared gauge quotient and let $Q$ be a full-rank retained linear map. Put $G=L^{-1}$ and
\begin{equation}\label{eq:gauss-conditional}
G_+=QGQ^*,\qquad J=GQ^*G_+^{-1},\qquad
\Gamma=G-GQ^*G_+^{-1}QG.
\end{equation}
Then
\begin{equation}\label{eq:harmonic}
QJ=I,\qquad J^*LJ=G_+^{-1},\qquad Q\Gamma=0,\qquad\Gamma\succeq0.
\end{equation}
The retained harmonic field and the detail innovation are energy orthogonal. In particular, the exactly Gaussian detail has no residual retained feedback after the harmonic recentering. Gauge gradients and any retained flat cohomology are removed or retained according to the declared quotient; they are not inverted as if massive.

\begin{proof}[Proof of the conditioning identities]
The first two identities follow by substitution. For the last, write
\[
\Gamma=G^{1/2}(I-\mathsf P)G^{1/2},\quad
\mathsf P=G^{1/2}Q^*(QGQ^*)^{-1}QG^{1/2}.
\]
The operator $\mathsf P$ is an orthogonal projection. This also proves positivity and the annihilation by $Q$.
\end{proof}

\section{Determinant ownership and complete derivatives}
A normalized stationary Gaussian reference carries its full background-dependent determinant. It is not a small polymer vertex. For its eliminated Hessian $L(a)$,
\begin{align}
D\log\det L&=\Tr(L^{-1}DL),\label{eq:detfirst}\\
D_2D_1\log\det L
&=\Tr\bigl(L^{-1}D_2D_1L-L^{-1}D_2L\,L^{-1}D_1L\bigr).
\label{eq:detsecond}
\end{align}
The second formula includes the full return term. Taking norms before the trace words and their Ward partners are assembled loses cancellations needed by the marginal response.

These determinant derivatives are taken in the declared background of the
exact pushed action. A later change of fluctuation representative must
transport the covariance and vertex contractions coherently; a change of
background is a separate operation. In particular, the fixed-background
banks of \cref{sec:background-prescription} are not substituted for the
harmonic background of the blocked action without the comparison in
\cref{crit:background-ADMC}.

Uniformly gapped finite-range eliminated Hessians admit exponential locality bounds. The corresponding absolute four-dimensional massless response can have logarithmic growth. A locality theorem with a gap hypothesis cannot be applied to a massless retained mode by omitting that hypothesis. For a normalized conditional current, three background derivatives can require a four-entry connected cumulant. The Hessian of its pushed Haar-Jacobian term requires the third current jet, as specified in \cref{thm:connected-score-hierarchy,prop:Jacobian-third-jet}. \src{F7}{V95--V100}; \src{F8}{V60--V61}

\section{Mixed-shell response and exact determinant ownership}
\label{sec:mixed-shells}
Independent Gaussian innovations do not make a background-dependent
quadratic response diagonal in scale. The distinction matters before
any anomaly coefficient is extracted. Let $z_j$ be independent standard
Gaussian vectors, $U=(U_0,\ldots,U_{n-1})$, and
\[
 C=UU^*=\sum_j\Gamma_j,\qquad \Gamma_j=U_jU_j^*.
\]
All operators below act on the declared finite quotient. For a
self-adjoint analytic perturbation $V(t)$ with $V(0)=0$, write
\[
 R=V'(0),\quad Q=V''(0),\quad
 K(t)=I+U^*V(t)U,\quad W(t)=\tfrac12\log\det K(t),
\]
and assume $K(t)>0$ near zero. Put $A_{ij}=U_i^*RU_j$ and
$B_{ij}=U_i^*QU_j$.

\begin{proposition}[Complete quadratic shell response]
\label{prop:mixed-response}
The exact second response is
\begin{align}
 W''(0)&=\tfrac12\Tr(CQ)-\tfrac12\Tr(CRCR)\notag\\
 &=\sum_j\left\{\tfrac12\Tr B_{jj}
                    -\tfrac12\Tr A_{jj}^2\right\}
                -\sum_{i<j}\norm{A_{ij}}_{\mathrm{HS}}^2.
 \label{eq:mixed-response}
\end{align}
Consequently the sum of the same-shell responses is exact if and only
if all off-diagonal blocks $A_{ij}$ vanish.
\end{proposition}
\begin{proof}
Differentiate $\log\det K$: at zero its second derivative is
$\Tr(K''-K'^2)$. The diagonal trace of $K'^2$ is
$\sum_{i,j}\Tr(A_{ij}A_{ji})$. Self-adjointness gives
$A_{ji}=A_{ij}^*$; the two orientations of every off-diagonal pair
therefore give the last term of \cref{eq:mixed-response}.
\end{proof}

The missing term is not necessarily an endpoint correction. For
example, let the shells be the coordinate lines of $\R^{n+1}$, let
$C=I$, and let $R$ be the adjacency matrix of the path with $n$ edges.
For $V(t)=tR$ and $|t|<1/2$, all same-shell responses vanish while
$W''(0)=-n$. Positivity of the perturbed precision and independence of
the innovations coexist with a linear-in-depth discrepancy.

\begin{proposition}[Schur ownership of mixed returns]
\label{prop:Schur-owned-response}
Eliminate blocks in increasing order and let $E_j=\{0,\ldots,j-1\}$.
The successive Schur pivots
\[
 S_j=K_{jj}-K_{jE_j}K_{E_jE_j}^{-1}K_{E_jj}
\]
satisfy
\begin{align}
 S_j'(0)&=A_{jj},&
 S_j''(0)&=B_{jj}-2\sum_{i<j}A_{ji}A_{ij},\notag\\
 \left(\tfrac12\log\det S_j\right)''(0)
 &=\tfrac12\Tr B_{jj}-\tfrac12\Tr A_{jj}^2
                         -\sum_{i<j}\norm{A_{ji}}_{\mathrm{HS}}^2.
 \label{eq:Schur-owned-response}
\end{align}
Moreover $W=\sum_j\tfrac12\log\det S_j$. Thus every mixed return is
owned once by its later endpoint. If a final block is retained, its
second vertex contains the corresponding return
$B_{FF}-2A_{FE}A_{EF}$ rather than just $B_{FF}$.
\end{proposition}
\begin{proof}
At zero all diagonal blocks of $K$ are identities and all off-diagonal
blocks vanish. Two derivatives of the Schur product can therefore
survive only when one falls on each off-diagonal factor. The
block-determinant identity then gives the claimed decomposition.
\end{proof}

For an orthogonal innovation decomposition after whitening, write
$\Gamma_j=C^{1/2}P_jC^{1/2}$, where $\sum_jP_j=I$ on the Gaussian
support and the $P_j$ are orthogonal projections. With
$\widehat A=C^{1/2}RC^{1/2}$, the error in discarding cross-scale terms
has the equivalent form
\begin{equation}
 W''(0)-W''_{\rm same}(0)
 =-\sum_{i<j}\norm{P_i\widehat A P_j}_{\mathrm{HS}}^2
 =-\tfrac14\sum_j\norm{[\widehat A,P_j]}_{\mathrm{HS}}^2.
 \label{eq:shell-commutator-defect}
\end{equation}
Orthogonality in the Gaussian energy space is therefore sufficient
only when the response operator preserves those subspaces. It is not
an orthogonality statement about arbitrary fine-field vertex legs.

\begin{scope}[Amplitude response, momentum extraction, and endpoints]
The sign in \cref{eq:mixed-response} concerns the scalar amplitude
parameter $t$. Differentiation in external momentum, followed by the
calibrated transverse extractor $\mathfrak b$, need not preserve that
sign. Neither positivity of the mixed norm nor Schur ownership proves
that the retained endpoint has zero marginal response. Such an endpoint
must be retained explicitly or estimated in the same normalization.
\end{scope}
\sourcebasis{The mixed response, Schur allocation, and whitened
commutator identity are developed in \src{SM}{V1--V4}.}


\section{Measured coordinate changes and the stable complement}
\label{sec:measured-redundancy}
An analytic change of variables acts on the measure, not only on the
displayed action. Infinitesimally its response combines action transport
with the divergence of the coordinate vector field. The exact
pushforward identifies how such a redundant tangent is represented on
the retained field, even when the fine vector field is not vertical for
the block map.

Let $\mathcal B:X\to Y$ be a smooth submersion between compact oriented
finite-dimensional manifolds without boundary, with smooth positive
reference measures $\nu_X,\nu_Y$. Set
\[
\rho\,\dd\nu_Y=\mathcal B_*(e^{-S}\dd\nu_X),
\qquad A=-\log\rho.
\]
For a smooth fine vector field $V$, define its vector-density current
$\mathcal J_V$ weakly by
\begin{equation}\label{eq:pushed-vector-current}
\int_Y\alpha(\mathcal J_V)\dd\nu_Y
=\int_Xe^{-S(x)}\alpha_{\mathcal B(x)}
       (D\mathcal B(x)V(x))\dd\nu_X(x)
\end{equation}
for every smooth one-form $\alpha$ on $Y$.

\begin{theorem}[Measured redundancy commutes with exact blocking]
\label{thm:measured-redundancy}
For the measured tangent
$\mathcal O_V=DS[V]-\operatorname{div}_{\nu_X}V$,
\begin{equation}\label{eq:redundant-pushforward}
\mathcal B_*(\mathcal O_Ve^{-S}\dd\nu_X)
=-\operatorname{div}_{\nu_Y}\mathcal J_V\,\dd\nu_Y.
\end{equation}
If $\rho>0$ and $\overline V=\mathcal J_V/\rho$, the first variation of
$A$ under $S+\varepsilon\mathcal O_V$ is
\begin{equation}\label{eq:coarse-redundant-tangent}
\dot A=DA[\overline V]-\operatorname{div}_{\nu_Y}\overline V,
\qquad
\overline V(y)=\E[D\mathcal B\,V\mid\mathcal B=y].
\end{equation}
Consequently the exact RG differential descends to the quotient by
measured field-redefinition tangents.
\end{theorem}
\begin{proof}
Apply the divergence theorem to $e^{-S}(f\circ\mathcal B)V$:
\[
\int_X(f\circ\mathcal B)\mathcal O_Ve^{-S}\dd\nu_X
=\int_Xe^{-S}Df_{\mathcal B(x)}(D\mathcal B V)\dd\nu_X
=-\int_Y f\operatorname{div}_{\nu_Y}\mathcal J_V\dd\nu_Y.
\]
This proves \cref{eq:redundant-pushforward}. Differentiating the pushed
density gives $\dot\rho=\operatorname{div}_{\nu_Y}(\rho\overline V)$.
Since $D\rho=-\rho\,DA$, division by $-\rho$ proves
\cref{eq:coarse-redundant-tangent}. The resulting tangent has the same
measured form on $Y$, which is the required quotient invariance.
\src{F8}{V58}
\end{proof}

Strict verticality $D\mathcal B\,V=0$ is sufficient to make the pushed
current vanish, but is not necessary for redundancy. In particular,
the natural equation-of-motion vector is generally not vertical for
straight dyadic blocking. The theorem replaces that false requirement
by the conditional current of the actual pushed measure. For compact
link groups it is formulated directly in product Haar measure and does
not require a global logarithmic chart.

\subsection{The Gaussian equation-of-motion tangent}
Use the Gaussian data $L,Q,J,G_+$ of
\cref{eq:gauss-conditional}, and put $L_+=G_+^{-1}$. For the linear
fine vector field $V(x)=\tfrac12Lx$, its coarse current is
\begin{equation}\label{eq:Gaussian-coarse-velocity}
\overline V(y)=R_Vy,
\qquad R_V=\frac12QLJ=\frac12QQ^*L_+.
\end{equation}
The integration-by-parts argument remains valid for this Gaussian
law and polynomial vector field by its integrable tails.

\begin{corollary}[No Gaussian marginal feed from the measured EOM tangent]
\label{cor:Gaussian-EOM-no-feed}
The pushed tangent is
\begin{equation}\label{eq:Gaussian-EOM-owner}
\mathcal O_{\overline V}(y)
=\ip{L_+y}{R_Vy}-\Tr R_V.
\end{equation}
If $L_+(k)=O(|k|^2)$ and $QQ^*(k)=O(1)$ near zero in the declared
transverse sector, its quadratic symbol is $O(|k|^4)$. It therefore
has no transverse $F^2$ moment; the displayed constant is a vacuum
contact.
\end{corollary}
\begin{proof}
The conditional mean is $Jy$, so the conditional fine velocity is
$\tfrac12LJy$. Applying $Q$ gives
\cref{eq:Gaussian-coarse-velocity}. Insert this linear field into
\cref{eq:coarse-redundant-tangent}; its divergence is $\Tr R_V$.
The symmetric quadratic kernel is
$\tfrac12(L_+R_V+R_V^*L_+)=O(|k|^4)$ under the stated symbol bounds.
\src{F8}{V58}
\end{proof}

This is an identity for the complete measured Gaussian tangent. It does
not say that the raw Jacobian term of an arbitrary nonlinear current
has zero marginal coefficient. Its subtraction is a coordinate choice
only after the complete pushed tangent and calibrated quotient have
been identified.

\subsection{The derivative order required by the Haar Jacobian}
\label{sec:Jacobian-third-jet}
In retained product-Haar exponential coordinates, write
$\dd\nu_Y=h(y)\dd y$ and $\gamma_i=\partial_i\log h$. At the identity,
$\gamma_i(0)=0$. For $U=U^i\partial_i$,
\[
\operatorname{div}_{\rm H}U=\partial_iU^i+\gamma_iU^i,
\]
where repeated component indices are summed.

\begin{proposition}[Third-current-jet requirement]
\label{prop:Jacobian-third-jet}
If $U(0)=0$, then
\begin{equation}\label{eq:Jacobian-third-jet}
\begin{aligned}
\partial_a\partial_b\operatorname{div}_{\rm H}U\big|_0
={}&\partial_a\partial_b\partial_iU^i\big|_0\\
 &+(\partial_a\gamma_i)(0)\,\partial_bU^i(0)
  +(\partial_b\gamma_i)(0)\,\partial_aU^i(0).
\end{aligned}
\end{equation}
Thus the Hessian of the pushed Haar-Jacobian term requires the third
background jet of $\overline V$.
\end{proposition}
\begin{proof}
Differentiate the divergence twice. In the product rule for
$\gamma_iU^i$, the terms with $U(0)$ and with $\gamma_i(0)$ vanish;
the two mixed first-derivative terms remain. The divergence of $U$
contributes its third derivative. \src{F8}{V60}
\end{proof}

Apply \cref{thm:connected-score-hierarchy} to the local current
$j=D\mathcal B\,V$. Local score derivatives through order three and
connected decay for up to four marked observables provide the input
for $D^3\overline V$. With the local Haar coefficients and summability
of the support incidence, \cref{eq:Jacobian-third-jet} then gives a
local Jacobian Hessian. The explicit sufficient activity criterion is
\cref{crit:nonlinear-current-gate}. A determinant locality estimate or
a two-point covariance bound alone does not supply these hypotheses.

In four dimensions, the $F^2$ and $F\widetilde F$ quadratic structures
are marginal. The properly identified dimension-six and higher local
jets contract at tree level by at most $b^{-2}$ under a block factor
$b>1$. The physical marginal and measured redundant directions must
be projected out before asserting contraction of the stable
complement. The retained unit sector is the exact one in
\cref{thm:exact-block-differential}; invariance of the redundant quotient
is supplied by \cref{thm:measured-redundancy}. Neither statement proves
the nonlinear stable estimate by itself. \src{F7}{V99--V100};
\src{F8}{V58--V61}; \src{F9}{V9--V11}

\begin{scope}[Linear versus compact blocking]
The explicit Gaussian block-average map in \cref{ch:tower} is a linear map on the gauge-field tangent problem. A complete interacting construction must supply a compatible compact nonabelian blocking, its measure disintegration, and the uniform estimates connecting its transported vertices to the Gaussian tower. These are not consequences of the linear formulas alone.
\end{scope}
\sourcebasis{Exact coarse action, normalized Gaussian extraction, retained pullback, and conditional expansions: \src{F7}{V92--V100}. Measured redundant pushforward and the higher current/Jacobian hierarchy: \src{F8}{V58--V61}. Measured calibration and sector transport: \src{F9}{V9--V11}.}

\chapter{Compact analytic forests and nonduplicated ownership}
\label{ch:forests}
\section{Soft reference factors and compact heat kernels}
A nontrivial compactly supported analytic cutoff does not exist. The source therefore replaces differentiated hard selectors by analytic soft factors and exact add--subtract identities. The regular soft-good factor is absorbed into the stronger normalized reference. It is not reintroduced as a forest vertex. The rebased bad factor is nonnegative and is estimated in that same law.

A global logarithmic chart with a subtracted Euclidean Gaussian is invalid across compact-group sheets. The compact heat-kernel reference retains all sheets, and its replica interpolation factorizes at the endpoints required by the forest formula. Right-invariant derivatives are arranged to commute with the left-weakened generators. This avoids the derivative commutator loss of a mismatched interpolation. \src{F9}{V1--V4}

\section{The correct all-tree estimate}
Let $z$ be the primitive analytic activity on radius $R$, let $r<R$ be the output radius, let $h_*$ be the fixed local incidence stock, and let $G_\mu$ be the weighted covariance stock. Define
\begin{equation}\label{eq:KHK}
K_{\rm HK}=\frac{8h_*G_\mu}{(R-r)^2}.
\end{equation}
In a local real coordinate chart, the Taylor--Wiener norm is
\[
\norm{F}_R=\sup_{\varphi\in\R^q}\sum_{\alpha\in\N^q}
\frac{R^{|\alpha|}}{\alpha!}|\partial^\alpha F(\varphi)|.
\]
Its compact heat-kernel version uses the declared invariant derivative bank and source radii. For a multi-index $\nu$,
\[
\norm{\partial^\nu F}_r\le
\frac{\nu!}{(R-r)^{|\nu|}}\norm{F}_R.
\]
The Taylor--Wiener derivative estimate spends the available radius once for all derivatives incident at a vertex. The fixed-tree bound retains the product of vertex-degree factorials. Removing that product is neither necessary nor justified.

\begin{lemma}[Degree-faithful labeled-tree census]\label{lem:degree-census}
For $n\ge2$,
\begin{equation}\label{eq:degree-census}
\sum_{T\in\mathcal T_n}\prod_{i=1}^n d_i(T)!
=(n-2)!\binom{3n-3}{n-2}.
\end{equation}
\end{lemma}
\begin{proof}
A degree sequence $d_i\ge1$ with $\sum_i d_i=2n-2$ occurs $(n-2)!/\prod_i(d_i-1)!$ times by the Pr\"ufer code. Multiplying by $\prod_i d_i!$ gives $(n-2)!\prod_i d_i$. With $e_i=d_i-1$, the remaining sum is the coefficient of $x^{n-2}$ in
\[
\left(\sum_{e\ge0}(e+1)x^e\right)^n=(1-x)^{-2n},
\]
which is the binomial coefficient in \cref{eq:degree-census}.
\end{proof}

\begin{theorem}[Connected compact forest bound]\label{thm:forest}
Under the compact analytic reference, covariance, incidence, and endpoint-factorization hypotheses above, if $8zK_{\rm HK}<1$, then
\begin{equation}\label{eq:forest}
u_{\rm good}\le C_{\rm own}
\left(z+\frac{4z^2K_{\rm HK}}{1-8zK_{\rm HK}}\right).
\end{equation}
The degree factorials in \cref{eq:degree-census} are included in this estimate.
\end{theorem}
\begin{proof}[Proof structure]
Apply the compact interpolation forest formula with the one-radius derivative bound. Retain $\prod_i d_i!$ at fixed tree, sum it by \cref{lem:degree-census}, and include the connected-series normalization. The resulting geometric majorant has ratio $8zK_{\rm HK}$, yielding \cref{eq:forest}. The owner embedding contributes $C_{\rm own}$. \src{F9}{V1--V4}
\end{proof}

A degree-free strengthening is not an additional requirement for this argument. The degree-faithful theorem already closes the one-scale combinatorics; what remains is compatibility of its constants along the chosen regulator family. \src{F10}{V100}

\section{An explicit regulator-growth criterion}
Assume
\[
z_\beta\le C_{\rm an}\bigl(e^{c_*\beta^{-1/5}}-1\bigr),
\qquad c_*=7680\pi^3,
\qquad K_{\rm HK}\le C_K\beta^p,
\]
with constants uniform in lattice spacing and volume. Put $A=eC_{\rm an}c_*$. For $p<1/5$ and
\begin{equation}\label{eq:beta-threshold}
\beta\ge\max\left\{1,c_*^5,(16AC_K)^{1/(1/5-p)}\right\},
\end{equation}
we have $8z_\beta K_{\rm HK}\le1/2$ and
\begin{equation}\label{eq:forest-beta}
u_{\rm good}\le C_{\rm own}(A+8A^2C_K)\beta^{-1/5}.
\end{equation}
At $p=1/5$, the sufficient critical condition is $Q=8AC_K<1$, with coefficient $C_{\rm own}[A+4A^2C_K/(1-Q)]$. These follow by $e^x-1\le ex$ for $0\le x\le1$ and substitution into \cref{eq:forest}.

Thus the surviving forest compatibility question is a concrete growth estimate for $8h_*G_\mu/(R-r)^2$, not a new enumeration problem. Fixed positive radius loss and fixed covariance/incidence stocks give $p=0$. The theorem does not establish that those stocks remain fixed in a full interacting continuum construction. \src{F10}{V100}

\section{Morse gluing and all-colour scheduling}
The parameter-dependent Morse construction exactly Gaussianizes the local action around its moving stationary center. The remaining regular interaction is the transformed Jacobian and the explicitly extracted higher-order density. Exact add--subtract gluing keeps the Gaussian principal part whole and avoids derivatives of a cutoff indicator. Constant, linear, and quadratic jets are absorbed into the calibrated reference; interior third-order terms and exterior exits have different owners.

The all-colour extension requires a triangular elimination schedule.
Every primitive has one historical birth support, one saturation, and
one owner. A later colour update may use an already exposed contribution
but cannot charge it again as a new primitive. Under the common
analytic-domain and source hypotheses, this gives the complete one-scale
owner product. The marked criterion below states how its activity must
be controlled to obtain the higher normalized current derivatives; its
activity is not the physical inverse-coupling coordinate.
\src{F8}{V61--V65}; \src{F9}{V5--V8}

\section{Four-mark clustering and the nonlinear current interface}
\label{sec:four-mark-interface}
The scalar owner bound becomes a statement about differentiated
conditional currents only after local complex sources have been
included. Work in a finite volume of the actual contour graph
$\mathcal G$, with maximum degree $\Delta$, and let polymers be its nonempty finite connected
vertex sets. Two polymers are incompatible when they intersect. Assume
an exact source-dependent representation
\begin{equation}\label{eq:marked-polymer-partition}
\mathcal Z(t)=\sum_{\Gamma\ {\rm compatible}}
                 \prod_{Y\in\Gamma}z_t(Y),
\qquad |z_t(Y)|\le u^{|Y|}
\end{equation}
uniformly on the declared complex source polydisc. Source-independent
normalizing factors are kept outside $\mathcal Z$; every
source-dependent local factor must be retained in this representation
or in its explicitly bounded local contribution.

\begin{lemma}[Explicit contour-entropy criterion]
\label{lem:polymer-entropy-gate}
For $a,\mu>0$ put
\begin{equation}\label{eq:polymer-Q}
Q_{a,\mu}(u)=\frac{u e^{a+\mu}}
                 {1-\Delta^2u e^{a+\mu}}.
\end{equation}
If $\Delta^2u e^{a+\mu}<1$ and $Q_{a,\mu}(u)\le a$, then
\begin{equation}\label{eq:polymer-KP}
\sum_{Y':\,Y'\cap Y\ne\varnothing}
 |z_t(Y')|e^{(a+\mu)|Y'|}\le a|Y|.
\end{equation}
For the filling contours of \cref{eq:filling}, the graph is
$y\sim_* y'$ when $0<\norm{y-y'}_\infty\le1$, hence
$\Delta=3^4-1=80$. A sufficient choice with $a=1$ is
\begin{equation}\label{eq:four-mark-6401}
6401e^{1+\mu_*}u\le1.
\end{equation}
\end{lemma}
\begin{proof}
A deterministic depth-first traversal encodes a connected set of size
$n$ containing a fixed vertex by $2(n-1)$ steps. Its number is therefore
at most $\Delta^{2(n-1)}$. Summing these sets with the weight in
\cref{eq:marked-polymer-partition} gives $Q_{a,\mu}(u)$ at each vertex.
Sum over the vertices of $Y$ to obtain \cref{eq:polymer-KP}. For $a=1$
the condition is $(\Delta^2+1)u e^{1+\mu}\le1$, which gives
\cref{eq:four-mark-6401}. The smaller constant $65$ applies only to a
degree-eight primitive graph, not to these filling contours.
\src{F8}{V61, V63}
\end{proof}

For local supports $X_0,\ldots,X_r$, $r\le3$, assume that $z_t(Y)$
depends on $t_i$ only if $Y\cap X_i\ne\varnothing$. Let $\rho_i>0$ be
the source radii. Assume the activities are holomorphic on the source
polydisc and obey the displayed bound on every smaller closed
polydisc. Let $\tau_\mathcal G(X_0,\ldots,X_r)$ be the length
of a shortest graph tree meeting all the supports.

\begin{theorem}[Volume-uniform four-mark clustering]
\label{thm:four-mark-clustering}
Under \cref{eq:marked-polymer-partition,eq:polymer-KP}, the partition
function has a nonvanishing branch connected to $\mathcal Z=1$ at zero
activities, its connected expansion is normally convergent, and
\begin{equation}\label{eq:four-mark-clustering}
\left|\partial_{t_0}\cdots\partial_{t_r}\log\mathcal Z(0)\right|
\le 2^{r+1}|X_0|Q_{a,\mu}(u)
 \left(\prod_{i=0}^r\rho_i^{-1}\right)
 e^{-\mu\tau_\mathcal G(X_0,\ldots,X_r)}.
\end{equation}
The constant is independent of total volume.
\end{theorem}
\begin{proof}[Connected expansion and marked estimate]
Expand the hard-core logarithm in its connected Ursell coefficients.
The connected-graph deletion bound majorizes each coefficient by the
sum of incompatibility spanning trees. Rooting at a polymer $Y_0$,
\cref{eq:polymer-KP} sums the terminal branches with a total rooted
majorant at most $e^{a|Y_0|}$. This proves absolute normal convergence
and the nonvanishing logarithmic branch.

Inclusion--exclusion with $t_i=0$ selects clusters depending on every
marked source and costs at most $2^{r+1}$. Their connected union meets
every $X_i$, so their total polymer size is at least
$\tau_\mathcal G(X_0,\ldots,X_r)$. Extract the resulting exponential
weight, root at a polymer meeting $X_0$, and sum the root activities
using \cref{eq:polymer-Q}. The anchor sum contributes $|X_0|$.
Multivariate Cauchy differentiation contributes
$\prod_i\rho_i^{-1}$, proving \cref{eq:four-mark-clustering}.
\src{F8}{V61}
\end{proof}

\begin{criterion}[Complete nonlinear-current locality]
\label{crit:nonlinear-current-gate}
Suppose the exact good--bad expansion of the normalized reduced Wilson
fibre, including the current source and up to three local score
sources, satisfies \cref{eq:marked-polymer-partition} on the filling
contour graph, uniformly in the retained background chart. Assume its
nonduplicated activity obeys
\begin{equation}\label{eq:current-activity-sum}
u(\beta)\le q_{\rm far}(\beta)+u_{\rm good}(\beta),
\qquad
6401e^{1+\mu_*}\{q_{\rm far}(\beta)+u_{\rm good}(\beta)\}\le1,
\end{equation}
where the two stocks are in the same rebased law, owner convention,
and source polydisc. If the required local current and score jets
have bounded incidence and uniform local norms, then the conditional
current has three exponentially summable background jets in every
smaller decay weight. Its pushed Haar-Jacobian Hessian is therefore a
well-defined local kernel.
\end{criterion}
\begin{proof}
Apply \cref{thm:four-mark-clustering} to the mixed derivatives defining
the connected cumulants of the current and score insertions. Every
term of \cref{eq:connected-score-hierarchy} at order at most three has
at most four entries. The finite partition sum and weighted
support-tree sums give the stated locality, with the source radii and
local jet constants retained. Then use
\cref{prop:Jacobian-third-jet}. \src{F8}{V59--V63}
\end{proof}

The forest bound \cref{thm:forest} and far-field stock \cref{eq:qfar}
supply activities only through a common owner embedding in the exact
marked representation. With shared source domains and covariance
stocks, $p<1/5$ in \cref{eq:forest-beta} sends both activities to zero
and yields \cref{eq:current-activity-sum} for sufficiently large $\beta$;
the critical case retains its separate strict condition. This implication
neither permits a hard or merely smooth selector in place of analytic
rebasing nor restores the degree-free tree hypothesis.


\begin{remark}[Two distinct factorial issues]
The factorial-grade obstruction of \cref{prop:factorial} concerns the strong all-order source coefficient assigned to each labeled tree. The degree factorials of \cref{lem:degree-census} belong to a different compact analytic forest expansion and are correctly summed by its own generating coefficients. Closing the latter does not retroactively prove the former, and the former does not invalidate the latter.
\end{remark}
\sourcebasis{Analytic forests and ownership: \src{F9}{V1--V8}; corrected four-mark and current estimates: \src{F8}{V60--V65}; degree-faithful regulator growth: \src{F10}{V100}.}

\chapter{Calibrated renormalization and the physical scale}
\label{ch:rg}
\section{Coordinates and the direction of the iteration}
Write $u=\beta^{-1}$ for the declared inverse-coupling coordinate. A calibrated renormalization state consists of $u$, the retained marginal anisotropies and sector labels, and stable coordinates $z_i$ carrying specified powers $p_i$. Measured redundant directions have already been removed. The flow is not a contraction on the entire effective action: it is a controlled marginal recursion coupled to a contracting stable complement.

We use increasing $j$ for the ultraviolet orientation, so $a_j\downarrow0$ and an asymptotically free branch has
\[
\beta_{j+1}-\beta_j\longrightarrow b>0,
\qquad u_{j+1}=u_j-bu_j^2+o(u_j^2).
\]
A coarse blocking computation naturally gives the opposite orientation. Its shell coefficient must be converted before it is compared with $b$. The current response programme keeps this sign and its field normalization explicit.

\section{Invariant power cones}
Suppose the stable recursion has the nonnegative majorant
\begin{equation}\label{eq:cone-rec}
z_i^+\le\sum_jQ_{ij}z_j+d_i u^{p_i},
\qquad \frac{u^+}{u}\ge\sigma>0,
\qquad 0<u\le1.
\end{equation}
Assume $Q_{ij}\ne0$ only when $p_j\ge p_i$. Put
\[
(Q_p)_{ij}=\sigma^{-p_i}Q_{ij},\qquad
(d_p)_i=\sigma^{-p_i}d_i.
\]

\begin{criterion}[Power-preserving stable cone]\label{thm:cone}
If the spectral radius of $Q_p$ is less than one, then
\[
H=(I-Q_p)^{-1}d_p\ge0
\]
defines an invariant cone $0\le z_i\le H_i u^{p_i}$ for \cref{eq:cone-rec}. For a triangular $Q_p$, it is enough that every diagonal entry be less than one.
\end{criterion}
\begin{proof}
If $z_j\le H_ju^{p_j}$, then $u^{p_j}\le u^{p_i}$ wherever $Q_{ij}\ne0$. Consequently
\[
z_i^+\le\left(\sum_jQ_{ij}H_j+d_i\right)u^{p_i}
\le H_i(\sigma u)^{p_i}\le H_i(u^+)^{p_i}.
\]
The middle inequality is the resolvent equation for $H$. Positivity follows from the convergent Neumann series for the nonnegative matrix $Q_p$. \src{F9}{V12}
\end{proof}

\section{Reciprocal coupling shields stable errors}
The exact physical coupling reset has the form
\begin{equation}\label{eq:reciprocal}
F_r(u,x)=\frac1{u^{-1}+c_r(u^{-1},x)}.
\end{equation}
If $|c_r|\le M$, $\norm{D_xc_r}\le L$, and $u\le(2M)^{-1}$, then
\begin{equation}\label{eq:shield}
\norm{D_xF_r}\le4Lu^2.
\end{equation}
Therefore a stable displacement $x=O(u^{6/5})$ affects the physical coupling only at order $u^{16/5}$. The factor $u^2$ follows from differentiating the actual reciprocal map; it is not an independently chosen smallness assignment.

If $c_r(\beta,0)=b+O(\beta^{-\delta})+\varepsilon_r$, then expansion of \cref{eq:reciprocal} gives
\[
F_r(u,0)=u-bu^2+O(u^{2+\delta})-\varepsilon_r u^2+O(u^3)
\]
under the displayed boundedness assumptions. The exact sign of the $\varepsilon_r$ term follows the convention in $c_r$; its absolute size is the relevant error stock. \src{F9}{V12--V14}

\begin{lemma}[Ces\`aro marginal recursion]\label{lem:cesaro-beta}
Suppose a positive sequence satisfies
\[
\beta_{j+1}=\beta_j+b+d_j,\qquad b>0,
\qquad \frac1N\sum_{j<N}d_j\longrightarrow0.
\]
Then
\begin{equation}\label{eq:beta-linear}
\frac{\beta_N}{N}\longrightarrow b,\qquad
u_N:=\beta_N^{-1}\sim\frac1{bN}.
\end{equation}
\end{lemma}
\begin{proof}
Sum the recursion and divide by $N$. The fixed initial value contributes $\beta_0/N\to0$, and the assumed average error vanishes. Invert the resulting positive asymptotics.
\end{proof}

\section{Logarithmic ultraviolet control is not a power of the cutoff}
For geometric spacing $a_j=a_0B^{-j}$ and \cref{eq:beta-linear},
\begin{equation}\label{eq:clock}
\frac{\beta_j}{\log(a_0/a_j)}\longrightarrow\frac b{\log B},
\qquad
u_j\log(a_0/a_j)\longrightarrow\frac{\log B}{b}.
\end{equation}
Thus $\beta_j^{-p}$ corresponds to $(\log(a_0/a_j))^{-p}$, not to $a_j^q$ for a positive $q$. The ratio $j^{-p}B^{qj}$ diverges. Any regulator estimate that requires a power of $a_j$ needs an additional argument; the marginal recursion alone does not provide it.

\begin{proposition}[Stable memory preserves logarithmic powers]\label{prop:stable-memory}
If $x_{j+1}\le\kappa x_j+C\beta_j^{-p}$ with $0\le\kappa<1$, $p>0$, and $\beta_j/j\to b>0$, then
\begin{equation}\label{eq:stable-memory}
\limsup_{j\to\infty}j^p x_j\le\frac{Cb^{-p}}{1-\kappa}.
\end{equation}
Its Ces\`aro average is $O(N^{-p})$ for $p<1$, $O(\log N/N)$ for $p=1$, and $O(N^{-1})$ for $p>1$, apart from the exponentially decaying initial contribution.
\end{proposition}
\begin{proof}[Proof outline]
Iterate the recursion into a geometric convolution. For each fixed backward offset, $(j/\beta_{j-\ell})^p\to b^{-p}$. The recent part is bounded by a geometric series and the remote history has an exponentially small weight. This proves \cref{eq:stable-memory}; summing $j^{-p}$ gives the stated averages. \src{F11}{V1}
\end{proof}

The contraction constant in this proposition must be smaller than one. The finite Wilson covariance-derivative ceiling smaller than $4/3$ in \cref{ch:oracle} is a different object and cannot be used as $\kappa$.
\sourcebasis{Calibrated power cones, reciprocal shielding, and the marginal recursion: \src{F9}{V12--V14}. Cofinal clock and stable-tail estimates: \src{F11}{V1}.}

\chapter{Anomaly matching and determinant endpoints}
\label{ch:anomaly}
\section{Determinant accumulation without duplicate resets}
Let $\mathfrak b$ extract the calibrated transverse $F^2$ coefficient
from the complete background germ. Exact nested elimination factors
the declared determinant into its successive Schur pivots and, when
present, a retained endpoint determinant. Thus the correct ledger is
\begin{equation}\label{eq:det-telescope}
 \Gamma_G^{[0,N)}=\sum_{j<N}\Gamma_{G,j}^{\rm owned}
                   +\mathcal E_N+C_N,\qquad \mathfrak b(C_N)=0.
\end{equation}
Here $C_N$ contains only genuinely background-independent
normalizations. The owned terms include the mixed returns in
\cref{prop:Schur-owned-response}; $\mathcal E_N$ retains the final
Schur response and any endpoint not already assigned to a pivot.
Neither $\mathfrak b(\mathcal E_N)=0$ nor
$\mathfrak b(\mathcal E_N)=o(N)$ follows from independence of the
Gaussian innovations. A reference reset contributes an additional
term only when it is not already in this ledger. All such assignments
are made once, in the same background and normalization.

\begin{definition}[Averaged determinant matching card]\label{def:ADMC}
The averaged determinant matching card, abbreviated ADMC, requires
\begin{equation}\label{eq:ADMC}
\mathfrak b\bigl(\Gamma_G^{[0,N)}\bigr)=Nb_{\rm YM}+o(N),
\qquad R_N^{\rm reset}+R_N^{\rm boundary}=o(N),
\end{equation}
with one common compact determinant, Haar, quotient, sector, and normalization convention throughout the block family.
\end{definition}

The endpoint $\mathcal E_N$ is included in the complete determinant
in \cref{eq:ADMC}; a shellwise proof must estimate its contribution
explicitly. The evaluated one-cell quartic coefficient in
\cref{thm:measured-quartic-log} supplies a concrete measured
coefficient and finite-part limit, but physical calibration and the
nonlinear trajectory remain additional interfaces.

Together with the calibrated analytic domain and the marked-response bounds, ADMC supplies the Ces\`aro error required by \cref{lem:cesaro-beta}. It is stronger than a correct coefficient at one scale and weaker than demanding a pointwise error estimate at every scale. The nonlinear inverse-coupling remainder is controlled separately; it is not inserted into the principal determinant a second time. \src{F9}{V13--V15}

\section{Compact--continuum intertwiners}
\begin{definition}[Sublinear compact--continuum intertwiner card]\label{def:SCIC}
The sublinear compact--continuum intertwiner card, abbreviated SCIC, is a decomposition
\begin{equation}\label{eq:SCIC}
\Gamma^{\rm cmp}_j-\Gamma^{\rm cov}_j
=\mathcal A_{j+1}-\mathcal A_j+\mathcal E_j+C_j,
\end{equation}
with
\[
\mathfrak b(C_j)=0,\qquad
\mathfrak b(\mathcal A_N)=o(N),\qquad
\frac1N\sum_{j<N}\mathfrak b(\mathcal E_j)\longrightarrow0,
\]
uniformly over the retained sectors needed by the application.
\end{definition}

If the covariant comparison has averaged coefficient $b_{\rm YM}$, telescoping \cref{eq:SCIC} supplies the Gaussian part of ADMC. An exactly measured change of variables can contribute a coboundary, and an exactly compensated pure-gauge quotient can contribute no anomaly. These facts do not show that every compact--continuum difference is one of those terms.

The Ward-recombined homotopy formulation makes the residual explicit. For the complete scalar response $W_{j,t}(k)$,
\begin{equation}\label{eq:Wardhomotopy}
\partial_tW_{j,t}
=\operatorname{div}_kV_{j,t}
+D_{j+1,t}-D_{j,t}+R_{j,t}+C_{j,t},
\qquad \avg{C_{j,t}}=0.
\end{equation}
Periodic divergences integrate to zero. Artificial low/edge partition derivatives cancel only after both pieces have been recombined. With sublinear endpoint and Ces\`aro-null residual integrals, \cref{eq:Wardhomotopy} implies SCIC. \src{F9}{V15--V16}

\section{Analytic alias cohomology}
Use normalized Haar average on $\T^4$ and define the dyadic inverse-branch operator
\begin{equation}\label{eq:alias-P}
(\mathcal Pf)(q)=\frac1{16}\sum_{\sigma\in\{0,1\}^4}
f(q/2+\pi\sigma).
\end{equation}
For the analytic Wiener norm
\[
\norm{f}_\rho=\sum_{m\in\Z^4}e^{\rho|m|_1}|\widehat f(m)|,
\]
one has $\avg{\mathcal Pf}=\avg f$ and
\begin{equation}\label{eq:Fourier-decimation}
\widehat{\mathcal P^nf}(m)=\widehat f(2^nm).
\end{equation}

\begin{theorem}[The constant mode is the only analytic obstruction]\label{thm:alias}
If $\avg f=0$ and $0<\rho'\le\rho$, then
\begin{equation}\label{eq:alias-decay}
\norm{\mathcal P^nf}_{\rho'}
\le e^{-(\rho2^n-\rho')}\norm f_\rho,\qquad n\ge1.
\end{equation}
Consequently every analytic real scalar response has the unique decomposition
\begin{equation}\label{eq:alias-decomp}
f=\avg f+(I-\mathcal P)a,\qquad \avg a=0,
\end{equation}
where $a=\sum_{n\ge0}\mathcal P^n(f-\avg f)$ converges in each stated smaller analytic norm. The quotient by analytic alias coboundaries is one dimensional.
\end{theorem}
\begin{proof}
For $m\ne0$, the ratio between the output and input analytic weights is bounded by $e^{-(\rho2^n-\rho')}$ because $|m|_1\ge1$. Apply \cref{eq:Fourier-decimation} and sum to prove \cref{eq:alias-decay}. The Poisson series converges, and its partial sums telescope under $I-\mathcal P$. If $\mathcal Pa=a$, then $\widehat a(m)=\widehat a(2m)$; absolute summability forces all nonzero Fourier coefficients to vanish. Haar invariance forces the constant in \cref{eq:alias-decomp}. \src{F9}{V17}
\end{proof}

For a uniformly analytic, fully Ward-recombined edge family $f_{j,t}$ in the common rescaled variable, put $m_{j,t}=\avg{f_{j,t}}$. Under the source's compatible edge-transport identification, the nonconstant part gives a bounded endpoint cochain. The exact remaining edge criterion is
\begin{equation}\label{eq:edge-scalar}
\frac1N\sum_{j<N}\int_0^1m_{j,t}\,\dd t\longrightarrow0.
\end{equation}
If $m_{j,t}\to m_{\infty,t}$ in $L^1(\dd t)$, it reduces to $\int_0^1m_{\infty,t}\dd t=0$. The analytic cohomology theorem does not evaluate that integral. Subtracting an unknown mean in order to solve the Poisson equation is not a proof that the mean vanishes.

\section{What a finite quadrature can certify}
For the $M^4$ equally spaced grid on $\T^4$, let $\langle f\rangle_M$ denote its normalized mean. The source analytic stock gives the explicit certificate
\begin{equation}\label{eq:quadrature}
\abs{\langle f\rangle_M-\avg f}
\le\norm f_\rho
\left[\left(\frac{1+e^{-\rho M}}{1-e^{-\rho M}}\right)^4-1\right].
\end{equation}
The proof sums precisely the nonzero Fourier modes divisible by $M$ in each coordinate. Thus a finite grid can certify a mean when paired with a uniform analytic norm. A grid value without that stock is not a Haar integral, and a pointwise sign at finitely many momenta is not an all-torus sign theorem.

\section{The obstruction to a fixed-reference comparison}
The finite response calculations produce exact local algebra. They do not license a global Neumann approximation outside its certified domain. The finite test in \src{F11}{V31} finds that the relevant frozen-chart residual test holds at only one of the $256$ quarter-turn points and fails at the other $255$. A constant Laurent coefficient of a truncated polynomial remains exact as polynomial algebra, but is not an approximation to the rational Haar mean without a global remainder bound.

There is also a more structural obstruction. If every rescaled step uses the same frozen endpoint pair and its scalar mean difference is a constant $m_G$, then the average in \cref{eq:edge-scalar} is $m_G$. A certified nonzero value would refute SCIC for that reset comparison, not prove anomaly matching. The correct nested relation has the form
\begin{equation}\label{eq:nested-shell-cob}
b_{G,j}=b_{\rm YM}+c_j-c_{j+1}+e_j,
\qquad c_N=o(N),\qquad \frac1N\sum_{j<N}e_j\to0.
\end{equation}
The nested Gaussian tower of \cref{ch:tower} avoids this repeatedly reset comparison; it is not inferred from a frozen-chart sign calculation. \src{F11}{V31--V32}
\sourcebasis{ADMC, SCIC, Ward recombination, and alias cohomology: \src{F9}{V13--V17}. Correct next marginal row: \src{F10}{V100}. Frozen-chart limitations and nested replacement: \src{F11}{V31--V34}.}

\chapter{The finite response oracle and its actual scope}
\label{ch:oracle}
\section{The reduced Hessian and its full edge origin}
The parity-tree calculation starts from a $64$-edge fibre and removes the declared gauge directions to obtain a $45$-coordinate reduced Hessian. The response oracle retains the noncommuting colour algebra, harmonic routers, cubic and quartic vertices, and the retree transformations. A coordinate restriction is valid only for the vertical space for which it was constructed. A different block map can require full-edge vertex banks before projection.

The source gives all-torus floors
\begin{equation}\label{eq:finite-H-floor}
H_{\rm W}(k)\succeq\frac{49}{1250}I_{45},\qquad
H_*(k)\succeq\frac{343}{1250}I_{45}
\end{equation}
for its declared endpoint Hessians. The proof uses the exact alias-energy and reciprocal estimates, not interpolation between a finite set of samples. These are finite eliminated-fibre statements, not floors for a physical continuum Hamiltonian. \src{F9}{V63 and endpoint development}; \src{F10}{inherited oracle}

\section{Moving-frame covariance}
Let $V$ be the normalized alias isometry. In the source convention, the frame connection is
\[
\Gamma_\mu=-\frac{\ii}{2}\diag(s_\mu),\qquad
\boldsymbol\nabla_\mu X=\partial_\mu X-[X,\Gamma_\mu].
\]
The full compressed symbol satisfies the exact moving-frame derivative identity, while commutators telescope in complete typed traces. A derivative of an individual matrix entry is not the corresponding covariant response. Likewise, an isolated fine $4\times4$ curl block has a moving gauge kernel; dropping the $45$-dimensional constrained range can destroy the positivity needed by a proposed estimate.

The tree-gauge map used for the tower is not merely a unitary change of
frame inside this compressed symbol. It selects a representative of the
fine gauge quotient before contraction with fixed fluctuation-leg banks.
The finite discrepancy in \cref{prop:representative-dependence} therefore
cannot be removed by invoking moving-frame covariance alone. The
covariance must be pushed through the specified representative map.

\section{A global relative covariance-derivative margin}
Write the exact finite Laurent representation
\[
H_{\rm W}(k)=\sum_{\alpha\in\mathcal S}G_\alpha e^{\ii\alpha\cdot k},
\qquad \overline M_{\rm W}=\sum_{\alpha\in\mathcal S}\norm{G_\alpha}_{\rm F}.
\]
The completed finite atlas reports the strictly positive rational margin
\begin{equation}\label{eq:theta97}
\theta_{97}=
\frac{23784808996629854813771747874389}
{241558005344776833024000000000000000000}
\end{equation}
for
\begin{equation}\label{eq:relative-floor}
\frac43H_{\rm W}\pm\boldsymbol\nabla_1H_{\rm W}\succeq\theta_{97}I_{45}.
\end{equation}

\begin{theorem}[Normalized derivative consequence]\label{thm:covariance-margin}
For this fixed-cutoff Wilson fibre,
\begin{align}
2H_{\rm W}\pm\boldsymbol\nabla_1H_{\rm W}
&\succeq\left(\frac{49}{1875}+\theta_{97}\right)I_{45},\label{eq:H2margin}\\
\norm{H_{\rm W}^{-1/2}(\boldsymbol\nabla_1H_{\rm W})H_{\rm W}^{-1/2}}_{\op}
&\le\kappa_{\rm W}:=\frac43-\frac{\theta_{97}}{\overline M_{\rm W}}<\frac43.
\label{eq:kappaW}
\end{align}
The same bound holds for
$H_{\rm W}^{1/2}(\boldsymbol\nabla_1H_{\rm W}^{-1})H_{\rm W}^{1/2}$.
\end{theorem}
\begin{proof}
Add $(2/3)H_{\rm W}$ to \cref{eq:relative-floor} and use \cref{eq:finite-H-floor}. Since $H_{\rm W}\preceq\overline M_{\rm W}I$, congruence of \cref{eq:relative-floor} by $H_{\rm W}^{-1/2}$ gives eigenvalue bounds $\pm(4/3-\theta_{97}/\overline M_{\rm W})$ for the normalized derivative. Finally,
\[
\boldsymbol\nabla_1H^{-1}=-H^{-1}(\boldsymbol\nabla_1H)H^{-1}
\]
proves the inverse-covariance assertion. \src{F10}{V97, V100}
\end{proof}

The rational atlas is a reported finite certificate. The consequence just proved is an elementary deduction from that certificate and the stated Hessian floor. Neither inequality evaluates the complete Ward scalar. In particular, $\kappa_{\rm W}<4/3$ is not a contraction smaller than one, not a regulator-uniform Schwinger-function derivative bound, and not a mass gap.

\section{The retained role of the oracle}
The oracle supplies exact vertices, normalized derivative control, and
finite tests of complete trace-word cancellation. For the block-average
tower, its extension to the full fine-edge space is necessary but not
sufficient: the shell covariance must also be put in the tree-gauge
representative used by the banks. The resulting corrected finite route
agrees with the slice route at the reported nonzero corners. The fixed
background and the block-rule harmonic background then give different
second-order corner means, recorded separately in
\cref{sec:fixed-background-cards}. This preserves the finite oracle as a
usable component without confusing its cards with a physical anomaly.
\sourcebasis{The finite oracle and its global relative margin are
developed in \cite{F9,F10}; full-edge reconstruction and the representative
correction are in \src{F11}{V39, V41}.}

\chapter{The explicit nested block-average Gaussian tower}
\label{ch:tower}
\section{Why this is a new principal object}
The earlier base-point path rules solve the axis problem and, after symmetrization, yield an off-axis fixed point with contraction rate $17/36$. Their global nondegeneracy requires an additional inequality. The block-average rule changes the linear map upstream: its alias blocks are diagonal, so the covariance tower and the limiting nondegeneracy can be written as positive lattice sums. The unresolved inequality for the older rule is not an extra hypothesis of the new one. \src{F11}{V32--V37}

Let
\begin{equation}\label{eq:qPi}
q_\mu(k)=e^{\ii k_\mu}-1,\qquad
\lambda(k)=\sum_{\mu=1}^{4}|q_\mu(k)|^2,
\qquad \Pi(k)=I-\frac{q(k)q(k)^*}{\lambda(k)}
\end{equation}
for $k\in\T^4\setminus\{0\}$. The free transverse covariance is
\[
G^{(0)}(k)=\frac{\Pi(k)}{\lambda(k)}.
\]
All inverses below are transverse inverses, or equivalently the indicated Moore--Penrose inverses on the gauge quotient.

\section{The block rule and its Ward identity}
For a cell $\{0,1\}^4$ and direction $\mu$, define
\begin{equation}\label{eq:block-average}
c_\mu=\frac1{16}\sum_{x\in\{0,1\}^4}
\bigl(a_{(x,\mu)}+a_{(x+e_\mu,\mu)}\bigr).
\end{equation}
Links leaving the cell carry their neighbor-cell phases. For the aliases $p_\epsilon=k/2+\pi\epsilon$, $\epsilon\in\{0,1\}^4$, put
\begin{equation}\label{eq:alias-B}
s_\epsilon=\frac1{16}\prod_{\nu=1}^{4}(1+e^{\ii p_{\epsilon,\nu}}),
\qquad
\Delta_\epsilon=\diag_\mu(1+e^{\ii p_{\epsilon,\mu}}),
\qquad B_\epsilon=s_\epsilon\Delta_\epsilon.
\end{equation}
The basic Ward identity is
\begin{equation}\label{eq:block-Ward}
B_\epsilon(k)q(p_\epsilon)=s_\epsilon q(k).
\end{equation}
It follows from $(1+w_\mu)(w_\mu-1)=w_\mu^2-1$. Thus the average of a fine pure gauge is a coarse pure gauge, and fine gauge changes with zero averaged potential are annihilated.

In the slice normalization of the source, the covariance recursion is
\begin{equation}\label{eq:tower-rec}
G^{(j+1)}(k)=\frac1{16}\Pi(k)
\sum_{\epsilon}B_\epsilon(k)G^{(j)}(p_\epsilon)B_\epsilon(k)^*\Pi(k).
\end{equation}
The segment family in \cref{eq:block-average} is covariant under coordinate permutations and reflections up to the stated cell-translation phase. This gives the correct hypercubic symmetry convention; reflection does not mean merely changing a momentum sign while leaving the edge frame fixed.

\section{A closed formula at every level}
For $x=k+2\pi M$ set
\begin{equation}\label{eq:Fj}
F_\nu^{(j)}(x)=\prod_{\ell=1}^{j}|1+e^{\ii x_\nu/2^\ell}|^2
=\frac{\sin^2(k_\nu/2)}{\sin^2(x_\nu/2^{j+1})}.
\end{equation}
At $k_\nu=0$, the convention is $F_\nu^{(j)}=4^j$ when $M_\nu\equiv0\pmod{2^j}$ and zero otherwise. Empty products at $j=0$ equal one.

\begin{theorem}[Explicit finite-level propagator]\label{thm:tower-explicit}
For every $j\ge0$ and $k\ne0$,
\begin{equation}\label{eq:Gj}
G^{(j)}(k)=\Pi(k)\diag_\mu\bigl(\gamma_\mu^{(j)}(k)\bigr)\Pi(k),
\end{equation}
where
\begin{equation}\label{eq:gammaj}
\gamma_\mu^{(j)}(k)=16^{-3j}
\sum_{M\bmod2^j}
\frac{\bigl(\prod_\nu F_\nu^{(j)}(k+2\pi M)\bigr)
F_\mu^{(j)}(k+2\pi M)}
{\lambda((k+2\pi M)/2^j)}.
\end{equation}
\end{theorem}
\begin{proof}
Iterate \cref{eq:tower-rec}. The chain block $W_M$ is the product of the diagonal $B$ matrices along the alias path. The inner fine transverse projection can be removed after the outer compression because the iterated Ward identity sends its gauge direction to $q(k)$. The squared scalar part of $W_M$ is $16^{-2j}\prod_\nu F_\nu^{(j)}$, its directional squared part is $F_\mu^{(j)}$, and the slice covariance contributes $16^{-j}$. This gives \cref{eq:gammaj}. \src{F11}{V37}
\end{proof}

\section{The positive lattice-sum fixed point}
Define
\[
\sigma_\nu(k,M)=\frac{\sin^2(k_\nu/2)}{(k_\nu+2\pi M_\nu)^2},
\]
with $\sigma_\nu=(1/4)\delta_{M_\nu,0}$ at $k_\nu=0$.

\begin{theorem}[Explicit fixed point and uniform difference control]\label{thm:tower-fixed}
The limiting propagator is
\begin{align}
G^*(k)&=\Pi(k)\diag_\mu\bigl(\Gamma_\mu^*(k)\bigr)\Pi(k),\label{eq:Gstar}\\
\Gamma_\mu^*(k)&=1024\sum_{M\in\Z^4}
\frac{\bigl(\prod_\nu\sigma_\nu(k,M)\bigr)\sigma_\mu(k,M)}
{|k+2\pi M|^2}.
\label{eq:positive-lattice-sum}
\end{align}
The series is absolutely convergent with nonnegative terms for $k\ne0$. If
\[
d_1=\sup_{k\ne0}\norm{G^{(1)}(k)-G^{(0)}(k)},
\]
then $d_1<\infty$ and
\begin{equation}\label{eq:tower-rate}
\sup_{k\ne0}\norm{G^{(j)}(k)-G^*(k)}
\le\frac43d_1\,4^{-j}.
\end{equation}
The uniform statement concerns the bounded difference of covariances sharing an infrared singularity; it does not assert that $G^*$ itself is bounded near zero.
\end{theorem}
\begin{proof}[Proof structure]
Choose symmetric representatives of the aliases in \cref{eq:gammaj}. The sine inequalities on $[-\pi/2,\pi/2]$ compare each summand to a fixed multiple of the positive summand in \cref{eq:positive-lattice-sum}; the coordinatewise product gives the summable lattice majorant. Taking the limit yields the displayed series.

For the linear covariance map $\mathcal L$ in \cref{eq:tower-rec}, direct alias summation gives
\begin{equation}\label{eq:Lcontraction}
\mathcal L[\Pi](k)=\frac14\Pi(k)
\diag_\mu\left(1-\frac{|q_\mu(k)|^2}{8}\right)\Pi(k)
\preceq\frac14\Pi(k).
\end{equation}
By positivity, this controls any bounded transverse difference by a factor $1/4$. The common infrared leading term makes the first difference bounded. Therefore consecutive differences are bounded by $d_1 4^{-j}$; summing their geometric tail gives \cref{eq:tower-rate}. \src{F11}{V35, V37}
\end{proof}

\begin{theorem}[Global transverse nondegeneracy]\label{thm:tower-floor}
Let $c_0=256/\pi^{12}$. Then
\begin{equation}\label{eq:Gfloor}
G^*(k)\succeq c_0\Pi(k),\qquad
\norm{S^*(k)}\le\frac{\pi^{12}}{256},\qquad S^*=(G^*)^+.
\end{equation}
For
\[
c_j=c_0-\frac43d_1 4^{-j}>0,
\]
the kernels $S^{(j)}=(G^{(j)})^+$ satisfy
\begin{equation}\label{eq:Srate}
\sup_{k\ne0}\norm{S^{(j)}(k)-S^*(k)}
\le\frac{d_1 4^{-j}}{(3/4)c_j^2}.
\end{equation}
\end{theorem}
\begin{proof}
Take the $M=0$ summand in \cref{eq:positive-lattice-sum} with $k\in[-\pi,\pi]^4$. Each of the five sine ratios is at least $\pi^{-2}$, and $|k|^2\le4\pi^2$, giving $1024\pi^{-10}/(4\pi^2)=c_0$. Compression by $\Pi$ proves \cref{eq:Gfloor}. The difference estimate gives the finite-level floor $c_j\Pi$. The inverse identity and \cref{eq:tower-rate} bound the inverse difference by $(4/3)d_1 4^{-j}/c_j^2$. \src{F11}{V37}
\end{proof}

\begin{scope}[Nondegeneracy is not a mass]
The bound is a lower bound on the covariance and an upper bound on the kernel. On a coordinate axis, for a transverse component,
\begin{equation}\label{eq:axis}
G^*_{22}(t,0,0,0)=\frac{2+\cos t}{12\sin^2(t/2)},
\qquad
S^*_{22}(t,0,0,0)=\frac{12\sin^2(t/2)}{2+\cos t}
=t^2+O(t^4).
\end{equation}
The Gaussian fixed point is infrared massless. Its uniform invertibility away from the gauge direction, in the precise sense of \cref{eq:Gfloor}, is useful analytic control and not the conclusion of the Yang--Mills mass-gap problem.
\end{scope}

\section{Harmonic extensions and positive shells}
Let $p_M=(k+2\pi M)/2^j$ and
\[
W_M=\prod_{\ell=1}^{j}B(e^{\ii(k+2\pi M)/2^\ell}).
\]
The minimizing fine field generated by a level-$j$ transverse coarse field $c$ is
\begin{equation}\label{eq:tower-minimizer}
a(p_M)=\mathcal M_M^{(j)}(k)c,
\qquad
\mathcal M_M^{(j)}=16^{-j}G^{(0)}(p_M)W_M^*S^{(j)}(k).
\end{equation}
The exact constraint and energy identities are
\begin{equation}\label{eq:tower-min-identities}
\Pi\sum_MW_M\mathcal M_M^{(j)}=\Pi,
\qquad
16^j\sum_M(\mathcal M_M^{(j)})^*S^{(0)}(p_M)\mathcal M_M^{(j)}=S^{(j)}.
\end{equation}
Successive minimizer maps compose. Intermediate transverse projectors cancel in the appropriate places by the Ward identity, not by treating them as identity matrices on the full edge space.

On a fixed finest alias class, let $\mathcal C_\ell$ denote the part of the fine covariance explained by the level-$\ell$ retained field. At the top level $j$ its blocks are
\begin{equation}\label{eq:explained-cov}
\mathcal C_j(M,M')=16^{-2j}G^{(0)}(p_M)W_M^*S^{(j)}W_{M'}G^{(0)}(p_{M'}).
\end{equation}
The unrestricted fine covariance in these units is $16^{-j}G^{(0)}(p_M)\delta_{MM'}$.

\begin{theorem}[Positive Gaussian shell decomposition]\label{thm:shells}
The differences
\[
\Gamma_\ell=\mathcal C_\ell-\mathcal C_{\ell+1}
\]
are positive semidefinite and
\begin{equation}\label{eq:shell-telescope}
\mathcal C_0=\sum_{\ell=0}^{j-1}\Gamma_\ell+\mathcal C_j.
\end{equation}
They are the one-step conditional vertical covariances lifted to the fine field. In the fixed-fine-momentum sense stated in the source, the explained tail tends to zero, giving the corresponding infinite shell expansion of the free covariance.
\end{theorem}
\begin{proof}
For nested Gaussian retained sigma fields, the explained covariances are obtained by orthogonal projection in the Gaussian Hilbert space. Projection to a coarser retained field can only reduce the explained covariance. Its difference is the covariance of the conditional innovation, hence positive. Telescope the differences and apply the source's explained-tail limit. The explicit formulas follow from \cref{eq:tower-minimizer}. \src{F11}{V38}
\end{proof}

The one-step vertical block formula makes its normalization explicit:
\begin{equation}\label{eq:vertical}
\begin{aligned}
\operatorname{Cov}^{\rm vert}_{\epsilon\epsilon'}
={}&\frac{\delta_{\epsilon\epsilon'}}{16}G^{(j)}(p_\epsilon)\\
&-\frac{\overline s_\epsilon s_{\epsilon'}}{256}
G^{(j)}(p_\epsilon)\Delta_\epsilon^*S^{(j+1)}(k)
\Delta_{\epsilon'}G^{(j)}(p_{\epsilon'}).
\end{aligned}
\end{equation}
It is positive and annihilated by the retained block map. Shell positivity gives an exact Gaussian multiscale decomposition; it does not give a nonlinear measure or a physical transfer contraction.

\section{Full-edge vertices and the fluctuation representative}
\label{sec:shell-functional}
The positive shells of \cref{thm:shells} are defined on transverse alias
spaces. A background-response calculation requires, in addition, a
representative of each fluctuation and a background at which the vertices
are differentiated. These are distinct choices. Positivity of a covariance
on the gauge quotient does not make a partially assembled vertex
contraction independent of its fluctuation representative.

Let $Q_0(b)$ and $R_0(b)$ denote the full fine-edge quadratic and linear
background vertices, with the colour trace and the paired-vertex signs
fixed as in the finite response construction. The full-edge construction
extends the cubic and quartic banks from the $45$ retained coordinates to
all $64$ edges. Their restrictions reproduce the original banks. The
additional columns comprise fifteen tree links and four second links; the
cubic extension has $378$ nonzero entries across the three
external-momentum orders. These are reported exact finite identities of
\src{F11}{V39}; they make the vertex forms available on the $49$-coordinate
tree slice before the coarse constraint is imposed.

\begin{proposition}[Representative dependence of the finite response]
\label{prop:representative-dependence}
The Wilson response formed by contracting these vertex banks with a
transverse alias covariance need not equal the response in the tree-gauge
slice, even when the two fluctuation fields represent the same gauge
class. In the stated normalization, the reported order-zero, order-one,
and order-two cards for the straight rule at $k=(\pi,0,0,0)$ are
\begin{align}
\mathfrak c_{\perp}(k)
 &=\left(-\frac{118777}{51450},\ 0,
          -\frac{3401839261}{1512630000}\right),
 \label{eq:transverse-corner}\\
\mathfrak c_{\rm tree}(k)
 &=\left(\frac{636}{245},\ 0,
          -\frac{1110886837}{43218000}\right).
 \label{eq:tree-corner}
\end{align}
Consequently, direct substitution of the transverse covariance into the
fixed tree-gauge vertex banks is not a valid representative-independent
prescription.
\end{proposition}
\begin{proof}[Finite-certificate implication]
The unequal zeroth-order entries already disprove representative
independence. The exact cards are reported in \src{F11}{V41}; their
full computational derivation is not repeated here. The distinction
concerns the fluctuation legs of the specified response banks. It is not a
failure of gauge invariance of the complete Wilson action or of a fully
matched change of variables in the measure.
\end{proof}

\section{The tree-gauge map and covariance transport}
\label{sec:tree-gauge}
Fix a nonzero coarse Bloch momentum $k$, a rooted spanning tree of the
sixteen sites of a cell, and the coarse-link map $L$ of the chosen rule.
Let $a$ be a cell edge field. Integrate $a$ along the fifteen tree links to
obtain a site potential $\phi_a$ with $\phi_a(0)=0$, and extend that
potential to neighbouring cells with the Bloch phases. Then
$a'=a-d\phi_a$ vanishes on the tree. Define the residual gauge field
\[
v_{(s,\mu)}(k)=q_\mu(k)\one_{\{s_\mu=1\}},
\]
which is generated by a constant potential on each cell, and set
\begin{equation}\label{eq:tree-gauge-map}
c_a(k)=\frac{q(k)^*L(a')}{\lambda(k)},\qquad
T_{\rm g}(k)a=a-d\phi_a-c_a(k)v(k).
\end{equation}
This formula specifies the map on $k\ne0$; it does not prescribe an inverse
of $\lambda$ on the toron sector.

\begin{proposition}[Exact representative constraints]
\label{prop:tree-gauge-constraints}
The map $T_{\rm g}$ is linear. Its image has zero tree links and zero
longitudinal coarse link, and
\begin{equation}\label{eq:tree-gauge-constraints}
q(k)^*L(T_{\rm g}a)=0,\qquad
\Pi(k)L(T_{\rm g}a)=\Pi(k)L(a).
\end{equation}
In particular, a transversely vertical field is sent to the full
coarse-link constraint $L(T_{\rm g}a)=0$. On the specified nonzero momentum
sector this is the unique representative in its fine gauge class with
these constraints, and $T_{\rm g}^2=T_{\rm g}$.
\end{proposition}
\begin{proof}
Tree integration is linear and fixes the potential relative to the root.
The residual field $v$ vanishes on the tree and obeys $L(v)=q(k)$. The
Ward identity of the block map sends every fine gradient to a coarse
gradient. Thus subtracting $d\phi_a$ and $c_av$ preserves the transverse
coarse link, while the definition of $c_a$ removes its longitudinal
component. If two representatives in the same fine gauge class have zero
tree links, their potential difference is constant within each cell. Their
remaining difference is therefore a multiple of $v$. The longitudinal
condition forces that multiple to vanish because $\lambda(k)>0$.
Applying the map to a field already satisfying the constraints makes both
subtractions zero. These are the representative identities underlying the
construction of \src{F11}{V41}.
\end{proof}

The normalization of the frame change is part of the covariance transport.
At level zero write
\[
D_0(k)=\diag_{\epsilon}G^{(0)}(p_\epsilon),\qquad
\mathsf B(k)=\bigl(B_\epsilon(k)\bigr)_\epsilon,
\]
and let
$E_{(s,\mu),(\epsilon,\nu)}=e^{\ii p_\epsilon\cdot s}\delta_{\mu\nu}$.
In the colour and Fourier convention of the response banks,
\begin{align}
\widehat C_0(k)
 &=3\left[D_0(k)-\frac1{16}D_0(k)\mathsf B(k)^*
 S^{(1)}(k)\mathsf B(k)D_0(k)\right],
 \label{eq:response-alias-cov}\\
C^{\rm edge}_0(k)&=\frac1{16}E(k)\widehat C_0(k)E(k)^*,
 \label{eq:response-edge-cov}\\
\Xi_0(k)&=T_{\rm g}(k)C^{\rm edge}_0(k)T_{\rm g}(k)^*.
 \label{eq:response-tree-cov}
\end{align}
The factor $3$ is the inverse of the one-colour Hessian normalization; the
factor $1/16$ in the edge conversion is a Fourier-frame factor. They are
not interchangeable. In particular, the bracket in
\cref{eq:response-alias-cov} is sixteen times the level-zero covariance in
\cref{eq:vertical}; this reconciles the two displayed covariance
conventions without inserting a further colour factor into the vertices.

\begin{proposition}[Covariance transport on a fixed quotient]
\label{prop:tree-covariance}
For a centered edge field with covariance $C^{\rm edge}\succeq0$, its
representative $T_{\rm g}a$ has covariance
\[
\Xi=T_{\rm g}C^{\rm edge}T_{\rm g}^*\succeq0.
\]
If the original covariance is supported on the transversely vertical
space, then $L\Xi=0$. When an alias calculation and a slice calculation
realize the same quotient Gaussian law in this unique representative,
their covariances and their contractions with the same vertex banks agree.
\end{proposition}
\begin{proof}
Covariance is transported by linear congruence. For every test vector $h$,
$\ip{h}{\Xi h}=\ip{T_{\rm g}^*h}{C^{\rm edge}T_{\rm g}^*h}\ge0$.
The constraint follows from \cref{prop:tree-gauge-constraints}. Equality
for two realizations of the same quotient law follows from uniqueness of
the representative; equality of their matrix contractions is then
algebraic. The same-quotient-law premise is essential.
\end{proof}

For the finite level-zero response, the exact-arithmetic comparison in
\src{F11}{V41} reports agreement of the corrected alias route with the
straight-rule reference cards and with the block-average slice route at
all fifteen nonzero corners and all three external-momentum orders.
The zero-momentum toron corner is treated by the regular slice route:
the transverse alias representation there has a singular constant-alias
kernel, so the ordinary Taylor inversion used away from zero is not
applicable. The finite agreement does not certify a full-torus quadrature.
The precise punctured-torus derivative bounds, rather than a uniform
pointwise bound at the toron, are given in
\cref{prop:toron-regularity}.

For a higher shell, let $C_j^{\rm edge}$ denote the lifted covariance
$\Gamma_j$ after the corresponding edge-frame and colour-normalization
conversion. The required covariance in the fixed fine-edge banks is
\begin{equation}\label{eq:nested-tree-cov}
\Xi_j=T_{{\rm g},j}C_j^{\rm edge}T_{{\rm g},j}^*.
\end{equation}
The nested representative is fixed scale by scale, first on the retained
lattice and then on each finer cell. This gives an algebraic prescription;
regulator-uniform estimates for that prescription and its momentum jets
remain separate analytic requirements. Positivity of each congruence does
not furnish those derivative bounds.

\subsection{Toron regularity of the representative}
\label{sec:toron-regularity}
The projector in \cref{eq:tree-gauge-map} cannot be treated as a finite
Laurent polynomial. Write $T_0$ for tree zeroing before the final
longitudinal subtraction, and let $R_{\partial}:\C^4\to\C^{64}$ place
one value on the eight boundary edges of its family. Thus
$T_0R_{\partial}=R_{\partial}$ and $LR_{\partial}=I$. At nonzero
momentum,
\[
 T_{\rm g}=(I-R_{\partial}P_{\rm long}L)T_0,\qquad
 P_{\rm long}=qq^*/\lambda,\qquad
 T_{\rm g}R_{\partial}=R_{\partial}(I-P_{\rm long}).
\]

\begin{proposition}[Punctured-torus regularity and integrable jets]
\label{prop:toron-regularity}
The implemented representative is bounded and smooth away from zero,
but has no continuous extension at zero on the full edge space. For
$j=0,1,2$ it satisfies
\[
 \norm{D^jT_{\rm g}(k)}\le C_j|k|^{-j}
\]
near zero. With any value assigned at zero it belongs to
$W^{1,p}(\T^4)$ for $1\le p<4$ and to $W^{2,p}(\T^4)$ for
$1\le p<2$. In particular $DT_{\rm g}\in L^2$ and
$D^2T_{\rm g}\in L^1$.
\end{proposition}
\begin{proof}
Along $k=te_1$ the map sends $R_{\partial}e_1$ to zero; along
$k=te_2$ it fixes that nonzero vector. The limits disagree. Boundedness
follows from $\norm{P_{\rm long}}=1$ and the finite Laurent matrices
$R_{\partial},L,T_0$. Since $\lambda\asymp|k|^2$, differentiating
$qq^*/\lambda$ gives the displayed bounds. Radial integration in four
dimensions gives $jp<4$. Integration by parts outside a radius-$r$
ball has boundary errors $O(r^3)$ for the first derivative and
$O(r^2)$ for the second; hence these are also the weak derivatives.
\src{F11}{V42}
\end{proof}
Finite products of translated such representatives and compatible
jointly $C^2$ bounded matrix factors are twice differentiable into
$L^1$. A second derivative contains either one $L^1$ second derivative
or two $L^2$ first derivatives; all other factors are bounded.
Consequently their second derivatives may be integrated. If a product
trace has bound $|\partial_t^2F(k,0)|\le B_2|k|^{-2}$ near zero, its
omitted normalized integral over $|k|<\delta$ is at most
$\pi^2B_2\delta^2/(2\pi)^4$. These finite-product estimates do not by
themselves bound a growing all-level alias trace. They refine, rather
than eliminate, the toron requirement in \cref{crit:background-ADMC}.

\section{Fixed and harmonic background prescriptions}
\label{sec:background-prescription}
The natural background of an exactly blocked tree-level action is the
harmonic extension
\begin{equation}\label{eq:harmonic-background}
b_j^{\rm harm}=\mathcal M^{(j+1)}\beta
\end{equation}
of the level-$(j+1)$ coarse polarization. This remains the background
associated with the minimizing construction. It cannot be inserted into
zero-background-momentum vertex banks as though it were a single fine
plane wave. Its nonconstant aliases have fine momenta
$2^{-j-1}(q+2\pi M)$ with $M\ne0$ and amplitudes of order $q$; they can
therefore contribute to the second external-momentum derivative.

A second, explicitly fixed prescription uses the fine tree-gauge router
extension of the reference coarse polarization plane wave in direction
$2$, with the same external-momentum jets at zero at every level. Denote
this background by $b^{\rm fix}$. Its vertex banks are fixed, while the
scale dependence enters through \cref{eq:nested-tree-cov}. The two
prescriptions must be kept distinct. Their limiting $F^2$ coefficients
are not identified by the gauge-representative map, which acts on
fluctuation legs rather than replacing a background variation.
\src{F11}{V41}

With this fixed prescription the tree-level Wilson shell density is
\begin{equation}\label{eq:shell-functional}
\begin{aligned}
\Omega_j^{\rm fix}(k,q)
={}&\frac12\Tr\bigl[\Xi_j(k)Q_0(b^{\rm fix};q)\bigr]\\
&-\frac13\Tr\bigl[\Xi_j(k+qe_1)R_0(b^{\rm fix};q)
\Xi_j(k)R_0(b^{\rm fix};q)^*\bigr].
\end{aligned}
\end{equation}
The trace runs over the full corresponding fine alias class, in the same
normalization as the banks. The downward coefficient is
\begin{equation}\label{eq:b-shell}
b_{G,j}^{\mathrm{fix},\downarrow}
=\frac32\int_{\T^4}
\left.\partial_q^2\Omega_j^{\rm fix}(k,q)\right|_{q=0}\dd k.
\end{equation}
Both the tadpole and bubble are differentiated completely, including the
momentum dependence of the covariance, its representative map, and the
vertices. They are assembled before absolute estimates. A coefficient for
the harmonic prescription instead requires the full background jets of
\cref{eq:harmonic-background}, including its nonzero aliases. No equality
between these two responses is imposed by notation.

\begin{scope}[What the fixed prescription supplies]
The fixed background makes the existing fine-edge banks usable in an
explicit shell-response prescription at every level. It does not itself
prove that this prescription equals the derivative of the exact nested
blocked action in its harmonic background. Such an identification requires
a background-compatibility estimate, or a direct calculation in the
harmonic prescription. It is also distinct from the complete Haar,
quotient, and determinant normalization required by ADMC.
\end{scope}

\section{Exact finite cards in the declared convention}
\label{sec:fixed-background-cards}
The three entries below are the corner means of the order-zero,
order-one, and order-two true external-momentum derivatives of the
level-zero Wilson response. For the block-average rule with the fixed
background, the reported exact values are
\begin{equation}\label{eq:corners-fixed}
\overline{\mathfrak c}^{\,\rm fix}_{\rm block}
=\left(\frac{5587}{1568},\ 0,
-\frac{407872449589}{47416320000}\right).
\end{equation}
For the same block rule with its own harmonic background, they are
\begin{equation}\label{eq:corners}
\overline{\mathfrak c}^{\,\rm harm}_{\rm block}
=\left(\frac{5587}{1568},\ 0,
-\frac{254367096889}{47416320000}\right),
\end{equation}
whereas the straight-rule reference gives
\begin{equation}\label{eq:corners-straight}
\overline{\mathfrak c}_{\rm straight}
=\left(\frac{176943}{627200},\ 0,
-\frac{2822949104411}{265531392000}\right).
\end{equation}
The first two block-rule entries agree, while their order-two difference
is exactly
\begin{equation}\label{eq:background-corner-difference}
\overline{\mathfrak c}^{\,\rm fix}_{{\rm block},2}
-\overline{\mathfrak c}^{\,\rm harm}_{{\rm block},2}
=-\frac{3480847}{1075200}.
\end{equation}
These identities are reported finite certificates from
\src{F11}{V39, V41}. Their difference shows why the background convention
must accompany any quoted second-order card. It does not prove that the
two full-torus integrals or their limiting tower coefficients are unequal.
Conversely, calling the difference a lattice-scale effect is not a proof
that its accumulated contribution vanishes. The cards are level-zero
quantities, not higher-shell coefficients, not exact torus integrals, and
not an evaluation of $b^*$.

\section{The fixed-point response and background-compatibility criteria}
\label{sec:background-matching}
\begin{criterion}[Gaussian matching from a nested fixed point]
\label{thm:fixedpoint-anomaly}
Suppose the nested kernels converge at rate $C\rho^j$, $\rho<1$, and the
complete shell response in a declared prescription has a common Lipschitz
control along the tower. That control includes the vertex banks,
representative maps, external derivatives, infrared treatment, trace
normalization, and all Gaussian-principal contributions needed by the
chosen determinant. If the consistently oriented coefficients satisfy
\[
b_{G,j}=b^*+O(\rho^j),
\]
then
\begin{equation}\label{eq:fixedpoint-ADMC}
\sum_{j<N}b_{G,j}=Nb^*+O(1).
\end{equation}
If these are the coefficients of the actual nested determinant and
$b^*=b_{\rm YM}$ in the same convention, its Gaussian ADMC follows.
Exact nesting creates no reference-reset term.
\end{criterion}
\begin{proof}
Sum the geometric coefficient error. For coefficients of the actual
nested determinant, the normalization telescopes as in
\cref{eq:det-telescope}. The physical coefficient identification is an
independent input. This is the nested fixed-point implication of
\src{F11}{V34, V38}, with the representative and background requirements
made explicit.
\end{proof}

The $4^{-j}$ covariance and inverse-kernel estimates establish the kernel
part of this criterion. Multilinearity of \cref{eq:shell-functional}
gives Lipschitz bounds on any fixed finite alias class. It does not remove
the need for a common bound as the class grows, or justify differentiating
through the singular toron representation without an appropriate
alternative chart or integrable control.

For a fixed-background evaluation, the missing connection to the actual
nested determinant can be stated at precisely the averaged strength
required by ADMC. Let $b_{G,j}^{\rm fix}$ and $b_{G,j}^{\rm harm}$ use the
same upward orientation, physical normalization, and complete principal
contributions, and define
\begin{equation}\label{eq:background-defect}
\delta_j^{\rm bg}=b_{G,j}^{\rm harm}-b_{G,j}^{\rm fix}.
\end{equation}

\begin{criterion}[Background compatibility at the ADMC interface]
\label{crit:background-ADMC}
Assume the fixed-background coefficients obey
$b_{G,j}^{\rm fix}=b_{\rm fix}^*+O(\rho^j)$ and
\begin{equation}\label{eq:background-Cesaro}
\frac1N\sum_{j<N}\delta_j^{\rm bg}\longrightarrow0.
\end{equation}
Then
\[
\sum_{j<N}b_{G,j}^{\rm harm}=Nb_{\rm fix}^*+o(N).
\]
In particular, $b_{\rm fix}^*=b_{\rm YM}$ supplies the Gaussian
coefficient part of ADMC for the harmonic determinant, provided its other
normalization and endpoint requirements hold. A sufficient form of
\cref{eq:background-Cesaro} is
\begin{equation}\label{eq:background-coboundary}
\delta_j^{\rm bg}=d_{j+1}-d_j+r_j,\qquad
d_N=o(N),\qquad \frac1N\sum_{j<N}r_j\longrightarrow0.
\end{equation}
\end{criterion}
\begin{proof}
Sum \cref{eq:background-defect} and apply
\cref{eq:fixedpoint-ADMC,eq:background-Cesaro}. For the sufficient
condition, the $d_j$ terms telescope to $d_N-d_0$, which is sublinear.
This is the SCIC telescoping argument of \cref{def:SCIC} applied to the
specific background discrepancy; it does not establish the discrepancy
estimate itself.
\end{proof}

The criterion expresses, rather than resolves, the background-independence
assumption. It permits either direct computation with the full harmonic
background or a fixed-background computation accompanied by a tested
comparison. The level-zero difference in
\cref{eq:background-corner-difference} alone decides neither route.

\section{Measured coefficients and physical normalization}
The colour and action convention is
\[
-\tr_{\rm ad}(\operatorname{ad}X\operatorname{ad}Y)=C_A\ip{X}{Y},
\qquad S_{\rm bg}=\frac1{4u}\int|\overline F|^2.
\]
Its upward dyadic comparison coefficient is
\begin{equation}\label{eq:bYM}
b_{\rm YM}=\frac{22C_A}{3(4\pi)^2}\log2.
\end{equation}
The one-colour incidence Hessian carries a factor $1/3$ relative to the
colour-traced convention, hence the propagator factor $3$ in
\cref{eq:response-alias-cov}. The reference flat Schur form has
$s_0(q)=q^2/3+O(q^4)$, so a response
$\gamma(q)=\kappa_{\rm tree}q^2+O(q^4)$ shifts the downward inverse
coupling by $3\kappa_{\rm tree}$. This accounts for the factor $3/2$
multiplying the second true derivative in \cref{eq:b-shell}. The direction
of the iteration must then be converted to the upward convention of
\cref{eq:bYM}; a downward card is not compared to that coefficient without
this conversion.

The one-step nonlinear response of \cref{sec:actual-theta} supplies an
additional actual-fibre integral, with its minimizing lift, fourth jets,
and Haar terms. Its $\SU(2)$ coefficient $\theta$ is not by definition
the complete all-level colour-normalized coefficient in
\cref{eq:bYM}. The flat cancellation of \cref{thm:flat-bloch} validates
one normalization interface, not that physical identification.

The finite representative correction and frame factors are explicit.
There is now also an evaluated complete noncommuting one-cell
coefficient for the balanced-root $\SU(2)$ branch:
\cref{thm:measured-quartic-log} identifies its quartic logarithm and
finite part, rather than leaving its leading measured scalar
unevaluated. This is a different statement from identifying the
colour convention and physical background of \cref{eq:bYM}.
The remaining comparison must control the toron region, growing
alias traces, the background change where used, and physical
normalization on one compatible trajectory. The mixed-shell and
endpoint terms of \cref{sec:mixed-shells} belong to that comparison.
No corner mean or principal-margin protection theorem supplies it.
Thus measured coefficient evaluation, interacting continuation, and
physical infrared contraction remain distinct mathematical tasks.
\sourcebasis{The explicit propagator, inverse bounds, extensions, and
positive shells are established in \src{F11}{V37--V38}; the full-edge
banks and finite response cards in \src{F11}{V39}; the representative map
and fixed-background prescription in \src{F11}{V41}. The averaged
background criterion is the elementary specialization of
\cref{def:SCIC} displayed above.}

\chapter{Nonlinear compact fibres and the measured local response}
\label{ch:compact-response}
The Gaussian tower of \cref{ch:tower} fixes a principal quadratic
object, but not a nonlinear nonabelian fibre or its measured response.
This chapter keeps the root-path block, its solved pivots, the moving
constrained minimum, and the Haar density together. The local measured
expansion is controlled through two loops, and a closed finite
field-degree packet specifies its composition. A common critical-norm
domain then supports the actual nested source and the finite-amplitude
classical continuum minimum, providing the branch on which the cofinal
measured determinant is evaluated. These results have distinct roles:
classical convergence does not prove convergence of a quantum law,
finite coefficient composition does not supply all-scale remainder
bounds, and local estimates do not control the entire compact fibre.

\section{A nonlinear block with the prescribed linear germ}
\label{sec:nonlinear-root-block}
Let $G$ be a compact matrix gauge group. Divide the fine lattice into
sixteen-site parity cells and use the rooted tree which clears the
highest active parity bit. Tree retraction leaves fifteen tree links
fixed to identity in each cell. Of the remaining forty-nine group
coordinates, seventeen are internal non-tree links and thirty-two are
boundary links. In each of the four coarse directions choose one root
boundary link as a pivot. The forty-five other groups are independent
fibre coordinates.

For a coarse edge $e$, the sixteen chronological two-link paths have two
identical root products $B_e$ on this tree slice. Write
$P_{e,1},\ldots,P_{e,14}$ for the other products. None contains a pivot;
each depends only on the nonpivot boundary groups of this edge and the
internal groups of its two endpoint cells. On the common logarithm
domain the block is
\begin{equation}\label{eq:root-block-local}
K_e(Y,B)=B_e\exp\left\{\frac1{16}\sum_{j=1}^{14}
                  \log(B_e^{-1}P_{e,j}(Y))\right\}.
\end{equation}
Its differential at the identity is the sixteen-path linear average
$L$. A right inverse $R$ places the same coarse Lie-algebra value in all
eight boundary links of a family and zero in its internal links; then
$LR=I$. A free-coordinate map $E$ parametrizes $\Ker L$. In particular,
the four pivot rows of $E$ are determined by $LE=0$, not selected as
additional independent directions. \src{F11}{V57--V58, V64--V66, V76}

\begin{proposition}[Why the linear average is not a nonlinear gauge map]
\label{prop:nonabelian-average}
For a nonabelian Lie algebra $\mathfrak g$, a scalar-weight map
$A:\mathfrak g^{\oplus m}\to\mathfrak g$ of the form
$A(X_1,\ldots,X_m)=\sum_i w_iX_i$ can be a Lie algebra homomorphism
only if $w_iw_j=0$ for $i\ne j$ and $w_i^2=w_i$ for every $i$.
Consequently a genuine multi-entry average is not the differential of
a Lie-group homomorphism with these domain and range groups.
\end{proposition}
\begin{proof}
Elements supported in different summands commute in the direct-sum
algebra. Their images have bracket $w_iw_j[X,Y]$, which must vanish
for every $X,Y$. Within one summand, bracket preservation instead
requires $w_i^2[X,Y]=w_i[X,Y]$. Nonabelianity gives both assertions.
The derivative of a Lie-group homomorphism preserves the bracket.
\end{proof}
This obstruction concerns a proposed homomorphism, not gauge covariance
of \cref{eq:root-block-local}. Endpoint path dressing makes that block
covariant under the residual cell-constant gauge transformations. The
field dependence of this dressing is precisely what a frozen linear
average omits. \src{F11}{V56--V57}

\subsection{The actual constrained Hessian and density}
Near the identity write the local fibre as
\[
x=X(z,c)=Ez+Rv(z,c),\qquad F(X(z,c))=c,\qquad F=\log K.
\]
The implicit construction has a common local maximum-norm domain, and
its fixed-order local derivatives are uniform in volume. Put
$f(z,c)=S_0(\exp X(z,c))$. At a constrained stationary point, with
$\dd S_0=\lambda\,\dd F$, the pullback Hessian is
\begin{equation}\label{eq:actual-constrained-Hessian}
D_z^2 f=(D_zX)^*\bigl(D^2S_0-\lambda\cdot D^2F\bigr)D_zX.
\end{equation}
Indeed the ordinary chain rule contains the additional term
$\dd S_0[D_z^2X]$, while twice differentiating $F(X)=c$ gives
$DF[D_z^2X]=-D^2F[D_zX,D_zX]$. Thus neither the moving tangent nor its
acceleration can be dropped from an off-shell response.

With a fixed ordering of the free coordinates, the local measure
relative to $\dd z\,\dd c$ is, up to a constant independent of $z,c$,
\begin{equation}\label{eq:root-local-density}
 e^{-\beta f(z,c)+\ell(z,c)}\dd z\,\dd c,\qquad
 \ell=\sum_{\text{non-tree }l}\log j_G(X_l)
             -\log\det\bigl(DF(X)R\bigr).
\end{equation}
The determinant is positive on the chosen local branch. The pivot
structure makes its relevant coarse blocks finite and local. Its
second derivative involves the third jet of $F$; the quadratic action
alone does not determine the measured order-one response. This is the
concrete fibre realization of the Jacobian-jet requirement in
\cref{prop:Jacobian-third-jet}. \src{F11}{V55, V58--V60}

\section{A smooth global compact completion}
\label{sec:global-pivots}
Fix a smooth nonincreasing cutoff $\chi:[0,\infty)\to[0,1]$, equal to
one on $[0,1/2]$ and zero on $[1,\infty)$. For a sufficiently small
fixed $\rho>0$ inside an injective logarithm neighbourhood, define
\[
 g_\rho(U)=\chi(\norm{\log U}/\rho)\log U
\]
there and extend it by zero. This is a smooth conjugation-equivariant
function on $G$. Define the completed block $\widetilde K$ by replacing
every logarithm in \cref{eq:root-block-local} with $g_\rho$.

\begin{lemma}[An upper derivative bound for the truncated logarithm]
\label{lem:pivot-cutoff}
For a fixed cutoff and sufficiently small $\rho$, constants independent
of the group input satisfy
\[
 \norm{g_\rho}\le\rho,\qquad
 \norm{D_Lg_\rho}\le C_\chi,\qquad
 \operatorname{Sym}D_Lg_\rho\preceq(1+C\rho)I.
\]
Here $D_Lg(U)v=\left.\frac{\dd}{\dd t}g(e^{tv}U)\right|_{t=0}$.
\end{lemma}
\begin{proof}
In logarithm coordinates $h(z)=\chi(\norm z/\rho)z$ has derivative
\[
 Dh(z)=\chi(s)I+s\chi'(s)\frac{z\otimes z}{\norm z^2},
 \qquad s=\norm z/\rho.
\]
Its eigenvalues are at most one, and their absolute values are bounded
by a constant depending only on $\chi$. The left logarithmic derivative
is $I+O(\norm z)$ uniformly on the support. Composing these two
derivatives proves the estimates; the zero extension is smooth.
No positive lower bound for the cutoff derivative is asserted.
\end{proof}

\begin{theorem}[Globally invertible compact pivots]
\label{thm:global-pivot}
For sufficiently small fixed $\rho$, arbitrary compact inputs
$P_1,\ldots,P_{14}$, and $0\le t\le1$, the maps
\[
 F_t(B)=B\exp(tX(B)),\qquad
 X(B)=\frac1{16}\sum_{j=1}^{14}g_\rho(B^{-1}P_j)
\]
are smooth diffeomorphisms of $G$. All differential singular values are
bounded above and bounded below by $1/16$, after decreasing $\rho$.
Their inverses have uniform derivatives of every fixed order in the
compact source parameters.
\end{theorem}
\begin{proof}
For the input variation $Be^{sv}$, put
$A_B=16^{-1}\sum_jD_Lg_\rho(B^{-1}P_j)$. Up to an orthogonal adjoint
factor the derivative of $F_t$ is
\[
 M_t=I-t\mathcal J(tX)A_B,\qquad
 \mathcal J(Y)=\int_0^1e^{s\operatorname{ad}Y}\dd s.
\]
The preceding lemma gives $\norm X\le14\rho/16$,
$\norm{A_B}\le14C_\chi/16$, and
$\operatorname{Sym}A_B\preceq(14/16)(1+C\rho)I$.
Since $\mathcal J(tX)=I+O(\rho)$,
\[
 \operatorname{Sym}M_t\succeq(1-14t/16-C'\rho)I
                         \succeq\tfrac1{16}I
\]
for sufficiently small $\rho$. This bounds the inverse differential.
The two duplicate root paths are essential to the strict margin.

To pass from local to global invertibility, solve the time-dependent
vector field $V_t=-(DF_t)^{-1}\partial_tF_t$ on the compact group.
Its complete flow $\Psi_t$ satisfies
$\frac{\dd}{\dd t}F_t(\Psi_t(B))=0$, hence
$F_t\circ\Psi_t=I$. The flow is a diffeomorphism, so $F_t=\Psi_t^{-1}$.
Smooth parameter dependence and repeated implicit differentiation on
the fixed compact parameter carrier give all stated derivative bounds.
They may depend on $\rho$, which is not sent to zero in this theorem.
\end{proof}

\begin{corollary}[Exact global product carrier and smooth density]
\label{cor:compact-product}
For $n$ coarse cells, the map
\[
 (Y,B)\longmapsto(Y,\widetilde K(Y,B)),\qquad
 Y\in G^{45n},\quad B\in G^{4n},
\]
is a global diffeomorphism. Its inverse $B=B(Y,C)$ has
\begin{equation}\label{eq:compact-haar-density}
 \dd Y\,\dd B=j(Y,C)\dd Y\,\dd C,\qquad
 j=\prod_e j_e,\qquad 0<c\le j_e\le C_1<\infty.
\end{equation}
Each factor and each solved pivot has a fixed finite stencil and
uniform fixed-order derivative bounds. For every finite volume and
finite $\beta$, the completed blocked law has a positive smooth density
relative to coarse Haar,
\begin{equation}\label{eq:completed-Z}
 \widetilde Z_\beta(C)=\int_{G^{45n}}
       e^{-\beta S_0(Y,B(Y,C))}j(Y,C)\dd Y.
\end{equation}
It is gauge invariant. If
$V_\beta=\beta S_0(Y,B(Y,C))-\log j(Y,C)$ and
$\widetilde W=-\log\widetilde Z_\beta$, then in any fixed coarse chart
\begin{equation}\label{eq:compact-score}
 \partial_a\widetilde W=\E V_{\beta,a},\qquad
 \partial_a\partial_b\widetilde W
   =\E V_{\beta,ab}-\Cov(V_{\beta,a},V_{\beta,b}).
\end{equation}
\end{corollary}
\begin{proof}
At fixed $Y$, each output depends only on its own pivot. Apply
\cref{thm:global-pivot} to each edge. The total Jacobian is triangular,
so its inverse Haar factor is the product in
\cref{eq:compact-haar-density}. Compactness permits differentiation of
\cref{eq:completed-Z}. Endpoint gauge transformations conjugate each
relative source, preserve Haar measure and the Wilson action, and
transport the inverse pivots; this proves invariance. Differentiating
the normalized integral gives \cref{eq:compact-score}.
\end{proof}
The per-factor bounds in \cref{eq:compact-haar-density} are not a
volume-independent bound on the complete product $j$. The local score
coefficients have size $O(\beta+1)$, but this alone gives no uniform
spatial decay of their covariance. Also, \cref{eq:completed-Z} is
already a density relative to coarse Haar: no second coarse-Haar
factor is inserted in \cref{eq:compact-score}. \src{F12}{V18}

\begin{scope}[Agreement of germs and identification of laws]
\label{scope:completion-germ}
The completed block agrees with the original logarithmic block wherever
all relative logarithms lie inside the inner cutoff region. Their local
action and measured jets therefore agree on a possibly smaller common
box. A piecewise compact fallback and this smooth completion need not
have the same global pushforward. Equality of their germs does not
identify their normalized compact conditionals, their reflected laws,
or their iterates. A change of completion requires its own comparison
in \cref{hyp:compact}; it is not absorbed into Gaussian notation.
\end{scope}

\subsection{The metric and completion equation in the explicit SU(2) fibre}
\label{sec:su2-pivot-equation}
The explicit compact spectral application uses $G=\SU(2)$ with its
invariant Hilbert--Schmidt Riemannian metric. Write $d$ for its geodesic
distance and $d_{\rm ch}(U,V)=\norm{U-V}_{\rm HS}$ for chordal distance.
On the logarithm support $d(I,U)=\norm{\log U}_{\rm HS}$, and
$d_{\rm ch}\le d$. The working cutoff is $\rho=1/64$ in the geodesic
convention. Chordal lower margins used below are sufficient inactivity
conditions; they do not redefine the cutoff.

There is also a direct characterization at this specified radius. Set
\[
 \Phi_\rho(r)=\int_0^r t\chi(t/\rho)\dd t,\qquad
 R_\rho=7\rho/8,\qquad
 \kappa_\rho=\tfrac18-\tfrac{49\rho^2}{256}.
\]
For arbitrary fourteen source products define, on the geodesic ball
$\overline B_{R_\rho}(C)$,
\[
 \mathcal V_C(B)=\tfrac12d(B,C)^2
             -\tfrac1{16}\sum_j\Phi_\rho(d(B,P_j)).
\]

\begin{proposition}[Strongly convex characterization of the actual pivot]
\label{prop:pivot-potential}
For $0<\rho\le1/4$,
$\operatorname{Hess}\mathcal V_C\succeq\kappa_\rho I>0$ on this
ball. Its unique critical point $B_*$ is the unique solution of the
completed coarse equation. In right-trivialized coordinates its
residual and a certified error estimate are
\begin{align}
 \mathcal R_C(B)&=\log(C^{-1}B)
             +\tfrac1{16}\sum_j g_\rho(B^{-1}P_j),
             \label{eq:pivot-residual}\\
 d(B,B_*)&\le\kappa_\rho^{-1}\norm{\mathcal R_C(B)}.
             \label{eq:pivot-residual-bound}
\end{align}
In particular, $B_*=C$ if and only if
$\sum_jg_\rho(C^{-1}P_j)=0$, and
$\kappa_{1/64}=131023/1048576$.
\end{proposition}
\begin{proof}
The group is the round sphere of radius $\sqrt2$. A radial potential
$f(r)$ has radial Hessian $f''(r)$ and transverse Hessian
$f'(r)\cot(r/\sqrt2)/\sqrt2$. For $r^2/2$ the latter is at least
$1-r^2/4$ on this ball. For $\Phi_\rho$, the radial eigenvalue is
$\chi(s)+s\chi'(s)\le1$ and the transverse eigenvalue is at most one;
outside its support both vanish. Subtracting fourteen such upper
bounds gives $1-R_\rho^2/4-14/16=\kappa_\rho$.
The radial derivative at the boundary is strictly positive because
$\norm{g_\rho}<\rho$. Hence the minimum is interior, and strict
geodesic convexity gives uniqueness. Its gradient is
\cref{eq:pivot-residual}. Every solution of the completed equation
lies within the same ball, so the characterization is global for the
pivot equation. Integrating the Hessian lower bound along the geodesic
from $B_*$ to $B$ and applying Cauchy--Schwarz proves
\cref{eq:pivot-residual-bound}.
\end{proof}
This convexity belongs to an auxiliary pivot-solving potential, not to
the Wilson action. Small independent pivot caps are not replacements
for \cref{eq:pivot-residual}: actual source balance, including possible
cancellation of active sources, is required. \src{F12}{V98}

\section{The moving minimum and its measured coefficient}
\label{sec:actual-one-loop}
Return to the local germ of \cref{sec:nonlinear-root-block} and to
orthonormal Lie coordinates. The free fibre Hessian obeys the
volume-uniform eliminated-space bound
\[
 f_{zz}(0,0)=E^*d^*dE\succeq hI,\qquad h=24/2327.
\]
The finite incidence construction and its coercivity estimate concern
this fibre, not the retained gauge field. On a common smaller
maximum-norm domain there is a unique nearby stationary section
$z_*(c)$, with
\begin{equation}\label{eq:moving-minimum}
 f_z(z_*,c)=0,\qquad H_c=f_{zz}(z_*,c)\succeq h_1I,\qquad
 D z_*=-H_c^{-1}f_{zc},
\end{equation}
where $h_1>0$ is uniform in volume. The inverse Hessian and the
fixed-order local derivatives of this section have exponential spatial
locality. Defining $G(c)=f(z_*(c),c)$ gives
\begin{equation}\label{eq:minimized-Schur}
 D^2G=f_{cc}-f_{cz}H_c^{-1}f_{zc}.
\end{equation}
This is the nonlinear constrained counterpart of the Schur identity
in \cref{prop:schur}; its retained gauge null directions at zero have
not been removed by the fibre bound. \src{F11}{V66, V80--V81}

Let $\mathcal D$ be the fixed local integration box and
$N=45(\dim G)n$. With the actual density \cref{eq:root-local-density},
put
\begin{align}
 Z_\beta^{\rm loc}(c)&=\int_{\mathcal D}
                e^{-\beta f(z,c)+\ell(z,c)}\dd z,
     &W_\beta^{\rm loc}&=-\log Z_\beta^{\rm loc},\notag\\
 \widehat W_\beta^{\rm loc}&=W_\beta^{\rm loc}+\log j_c(c),
     &\widehat Q(c)&=\tfrac12\log\det H_c-\ell(z_*(c),c)
                                     +\log j_c(c).
 \label{eq:normalized-local-Q}
\end{align}
Here $j_c(c)=\prod_ej_G(c_e)$ converts the coordinate density to a
density relative to coarse Haar. This conversion is needed in
\cref{eq:normalized-local-Q}, unlike \cref{eq:completed-Z}.

\begin{theorem}[Volume-uniform measured one-loop approximation on the box]
\label{thm:uniform-local-loop}
On the declared common local domain, let
\[
 \mathcal R_\beta=\widehat W_\beta^{\rm loc}-\beta G
                  -\tfrac N2\log(\beta/2\pi)-\widehat Q.
\]
For sufficiently large $\beta$, constants independent of finite volume
and of the admitted small coarse field satisfy
\begin{align}
 N^{-1}|\mathcal R_\beta|+|\partial_a\mathcal R_\beta|
      &\le C\beta^{-1/2},\label{eq:loop-error-zero}\\
 |\partial_a\partial_b\mathcal R_\beta|
      &\le C\beta^{-1/2}e^{-\gamma d(a,b)},\label{eq:loop-error-local}\\
 \norm{D_c^2\mathcal R_\beta}_{\rm row}
       +\norm{D_c^2\mathcal R_\beta}_{\op}
      &\le C\beta^{-1/2}.\label{eq:loop-error-row}
\end{align}
The indices $a,b$ are fixed local coarse directions, and the row norm
sums the norms of their spatial Hessian blocks.
\end{theorem}
\begin{proof}
Use the centered convex extension on the common local domain. Its
localized expansion, with the actual minimum, Hessian, and density,
has an $O(\beta^{-1/2})$ error per degree of freedom and the localized
first and second derivative bounds of \cref{eq:loop-error-local}.
The comparison with the original potential on the same box is
$O(e^{-\kappa\beta})$ in those norms. These are the local extension
and same-box estimates of \src{F12}{V10--V11}.

For completeness, the nontrivial boundary comparison is a normalized
wall ratio, not the removal of a rare global union. Introduce a
nonnegative coarse-independent wall $\Psi$ and the potential
$E_\beta+\lambda\Psi$, $0\le\lambda<\infty$. For its free energy
$F_\lambda$, normalized differentiation gives
\begin{align*}
 \partial_\lambda F_\lambda&=\E_\lambda\Psi,\\
 \partial_\lambda\partial_aF_\lambda
              &=-\Cov_\lambda(E_{\beta,a},\Psi),\\
 \partial_\lambda\partial_a\partial_bF_\lambda
              &=-\Cov_\lambda(E_{\beta,ab},\Psi)
                +\kappa_{3,\lambda}(E_{\beta,a},E_{\beta,b},\Psi).
\end{align*}
The localized wall estimates give, after summing the local wall pieces,
\[
 |\partial_\lambda\partial_a\partial_bF_\lambda|
 \le C\beta^2e^{-a_0\beta}(1+\lambda)^{-9/8}
                                      e^{-\gamma d(a,b)},
\]
with the analogous first-derivative and per-site zeroth-order bounds.
Since the strength integral is $8$, the limiting hard-wall ratio and
its first two derivatives are exponentially small. At finite volume,
dominated convergence identifies the limit with the actual box integral;
the wall has no omitted coarse derivative. Adding the extension error,
same-box comparison, and wall ratio proves
\cref{eq:loop-error-zero,eq:loop-error-local}. Summability of the
spatial exponential gives \cref{eq:loop-error-row}.
\src{F12}{V12}
\end{proof}
This theorem strengthens a fixed-volume Laplace expansion. It is not
an all-orders volume-uniform expansion, and it gives no derivative
estimate for the logarithm of the full compact complement.

\subsection{The measured two-loop coefficient and four source derivatives}
\label{sec:measured-two-loop}
The preceding Hessian estimate admits a sharper expansion when the
higher local derivative and moment bounds are retained. These are
additional analytic inputs, not a consequence of a positive Hessian
alone. On a common smaller source domain assume bounded-support
analytic stencils through order twelve, an interior stationary section,
exponential locality of $H_c^{-1}$, and the centered convex extension
and normalized wall estimates through four source derivatives. The
connected estimates use rooted spatial tree sums and local moments
through the orders required by four noise and four source derivatives.
These hypotheses are realized for the controlled local root chart in
\src{SM}{V7}; they do not impose convexity on its compact complement.

For an $r$-fold local derivative, let $\tau(a_1,\ldots,a_r)$ be the
minimal spatial tree length connecting its marks. One convenient norm is
\begin{equation}
 \norm{F}_{a,r}=
 \sup_{a_1}\sum_{a_2,\ldots,a_r}
 e^{a\tau(a_1,\ldots,a_r)}\sup_c
 \abs{\partial_{a_1}\cdots\partial_{a_r}F(c)},\qquad r\ge1,
 \label{eq:two-loop-mark-norm}
\end{equation}
with component norms understood for vector-valued marks. At order zero
use $N^{-1}\sup_c|F(c)|$. A possible reduction of $a>0$ is fixed once
and is independent of volume.

Let $\eta$ be centered Gaussian with covariance $C_c=H_c^{-1}$ and
set, with all derivatives evaluated at $(z_*(c),c)$,
\begin{align}
 P_c(\eta)&=\tfrac16 f_{zzz}[\eta,\eta,\eta]-\ell_z[\eta],\notag\\
 Q_c(\eta)&=\tfrac1{24}f_{zzzz}[\eta,\eta,\eta,\eta]
                                  -\tfrac12\ell_{zz}[\eta,\eta],\notag\\
 \mathcal B_1(c)&=\E Q_c-\tfrac12\E P_c^2.
 \label{eq:measured-two-loop-coefficient}
\end{align}
The density $\ell$ here is the complete constrained density, including
the solved-pivot factor. Replacing it by fine Haar alone changes
$\mathcal B_1$.

\begin{proposition}[Complete Wick contraction]
\label{prop:two-loop-Wick}
Writing $T_{abc}=f_{abc}$, $V_{abcd}=f_{abcd}$,
$p_a=\ell_a$, $L_{ab}=\ell_{ab}$, and
$\vartheta_i=T_{abi}(C_c)_{ab}$, one has
\begin{align}
 \mathcal B_1={}&\tfrac18V_{abcd}(C_c)_{ab}(C_c)_{cd}
             -\tfrac12L_{ab}(C_c)_{ab}\notag\\
 &-\tfrac1{12}T_{abc}T_{def}
                  (C_c)_{ad}(C_c)_{be}(C_c)_{cf}\notag\\
 &-\tfrac12(p-\vartheta/2)^*C_c(p-\vartheta/2).
 \label{eq:complete-two-loop-Wick}
\end{align}
Repeated fibre indices are summed. All terms have uniform per-site
bounds and bounded marked norms through order four under the stated
hypotheses.
\end{proposition}
\begin{proof}
The quartic Gaussian moment has three pairings. In the square of the
cubic term, six pairings connect all three legs across the two cubic
vertices, while nine pairings contract two legs within each vertex.
The latter, together with the cubic--linear and linear--linear terms,
form the last square in \cref{eq:complete-two-loop-Wick}. Exponential
covariance locality and finite stencils permit one vertex to be rooted
before summing the remaining positions. Source derivatives preserve
this connected summation after a fixed reduction of the exponent.
\end{proof}

\begin{theorem}[Local measured two-loop expansion in four marked norms]
\label{thm:measured-two-loop}
Under the higher-order local hypotheses above, on the original fixed
box $\mathcal D$ one has
\begin{equation}
 \widehat W_\beta^{\rm loc}
 =\beta G+\tfrac N2\log(\beta/2\pi)+\widehat Q
                         +\beta^{-1}\mathcal B_1+\mathcal R_\beta^{(2)},
 \label{eq:two-loop-expansion}
\end{equation}
where, for $0\le r\le4$,
\begin{equation}
 \norm{\mathcal R_\beta^{(2)}}_{a,r}\le C_r\beta^{-2}.
 \label{eq:two-loop-C4}
\end{equation}
The constants are independent of the finite volume. At $r=0$ the norm
has the per-degree-of-freedom normalization stated above. Any fixed
number of external momentum derivatives is controlled by the
corresponding exponentially weighted spatial moments.
\end{theorem}
\begin{proof}
Put $t=\beta^{-1/2}$ and center the fibre variable as $z=z_*+t\eta$.
After removing the Gaussian normalization and the displayed minimum
and density terms, the extended exponent is
\[
 \tfrac12\eta^*H_c\eta+tP_c(\eta)+t^2Q_c(\eta)+O(t^3).
\]
Use the centered extension in its signed-noise form. Its full exponent
satisfies $U_{-t}(\eta,c)=U_t(-\eta,c)$; hence its normalized free
energy is exactly even in $t$, not merely even at the formal level.
The second noise derivative at zero is $2\mathcal B_1$.

The fourth noise derivative with at most four source marks is a finite
sum of connected expectations with at most eight entries. The stated
local moment and rooted tree estimates bound these derivatives in
\cref{eq:two-loop-mark-norm}, without an extensive unrooted sum. Taylor's
formula with integral remainder therefore gives $O(t^4)$ in all five
norms. Odd powers vanish by the exact parity identity.

It remains to recover the original domain. For a coarse-independent
wall with strength $\lambda$, normalized score differentiation expresses
each $\partial_\lambda\partial_I F_\lambda$ as a sum of connected
expectations containing one local wall insertion. For $|I|\le4$ the
higher-order wall estimate is
\[
 \norm{\partial_\lambda F_\lambda}_{a,|I|}
 \le C_{|I|}\beta^{|I|}e^{-b\beta}
                                    (1+\lambda)^{-21/20}.
\]
Its strength integral is $20$. The hard-wall limit and the same-box
extension comparison are thus exponentially small in every required
norm. They can be absorbed in $C_r\beta^{-2}$, proving
\cref{eq:two-loop-C4} for the actual local integral.
\end{proof}

\begin{corollary}[Summable remainder, retained two-loop term]
\label{cor:two-loop-summability}
Suppose the hypotheses and source transports of
\cref{thm:measured-two-loop} remain uniform along a blocking trajectory,
with uniformly bounded amplification of these norms and
$\beta_j\ge b(j+j_0)$ for some $b,j_0>0$. Then
$\sum_j\norm{\mathcal R_{\beta_j}^{(2)}}_{a,r}<\infty$ for $0\le r\le4$.
The coefficient $\beta_j^{-1}\mathcal B_{1,j}$ is not thereby summable
and must remain in the running measured action unless a separate
cancellation or improved bound is proved.
\end{corollary}
\begin{proof}
The remainder is dominated by a constant multiple of
$\sum_j(j+j_0)^{-2}$. No analogous argument applies to
$\sum_j(j+j_0)^{-1}$.
\end{proof}
The theorem sharpens \cref{thm:uniform-local-loop} on its strengthened
analytic domain. It is still a local-box theorem. Its four field marks
must not be confused with a fourth momentum derivative, nor can the
exponentially small wall ratio be replaced by a claim about the entire
compact exterior.
\sourcebasis{The higher-order centered expansion, complete measured
coefficient, and normalized wall return are established in
\src{SM}{V7--V8}.}


\subsection{Measured composition and the closed field--loop packet}
\label{sec:measured-loop-composition}
A uniform one-step remainder does not by itself specify which
coefficients must enter the next step. The incoming one-loop density
and two-loop action are part of the measured pushforward. To state the
finite coefficient rule, let $h=\beta^{-1}$ and consider
\[
 e^{-S(x)/h-Q(x)-hB(x)}\dd x.
\]
The parameter $h$ is neither a mesh size nor a physical time. Let
$F$ be an $h$-independent local submersion, with $F(0)=0$ and
$DS(0)=0$, and choose a regular fibre chart $x=X(y,c)$ satisfying
$F(X(y,c))=c$. Define
\[
 f=S\circ X,\qquad q=Q\circ X-\log|\det DX|,\qquad b=B\circ X.
\]
The complete chart or inverse-pivot density is included once in $q$.
Let $y_*(c)$ be the local classical minimum, $H=f_{yy}(y_*,c)>0$,
$C=H^{-1}$, and $G(c)=f(y_*,c)$.

\begin{theorem}[Finite measured two-loop composition]
\label{thm:measured-loop-composition}
After separating the source-independent Gaussian normalization, the
pushforward has negative logarithm
\[
 h^{-1}G(c)+Q_+(c)+hB_+(c)+O(h^2),
\]
with the finite-regulator coefficient identities
\begin{align}
 Q_+&=q(y_*,c)+\tfrac12\log\det H,\notag\\
 B_+&=b(y_*,c)+
 \E_C\!\left[\tfrac1{24}f_{yyyy}[\eta^4]
                      +\tfrac12q_{yy}[\eta^2]\right]
 -\tfrac12\E_C\!\left[\left(\tfrac16f_{yyy}[\eta^3]+q_y[\eta]\right)^2\right].
 \label{eq:measured-loop-step}
\end{align}
All fibre derivatives are evaluated at $(y_*(c),c)$. For two compatible
submersions with positive eliminated Hessians and with all intermediate
measured terms retained, the coefficient map satisfies
\begin{equation}
 \mathfrak R^{[2]}_{F_2}\circ\mathfrak R^{[2]}_{F_1}
                  =\mathfrak R^{[2]}_{F_2\circ F_1}.
 \label{eq:measured-loop-composition}
\end{equation}
The equality concerns the classical, one-loop, and two-loop germs on
the same stationary branch.
\end{theorem}
\begin{proof}
Set $y=y_*+\sqrt h\,\eta$. After removing the stationary factors, the
order-$\sqrt h$ perturbation is
$f_{yyy}[\eta^3]/6+q_y[\eta]$, which is odd. The order-$h$ term is
$f_{yyyy}[\eta^4]/24+q_{yy}[\eta^2]/2+b(y_*,c)$. Expanding the
normalized exponential and then its negative logarithm yields
\cref{eq:measured-loop-step}. It reduces to
\cref{eq:measured-two-loop-coefficient} when the incoming $B$ vanishes
and $q$ is the negative logarithm of the full density.

For composition, use the exact local pushforward identity
\[
 (F_2)_*\bigl((F_1)_*\nu_h\bigr)=(F_2\circ F_1)_*\nu_h
\]
with a common localization around the compatible stationary branch.
Classical minimization is associative there, and positivity of the
second eliminated Hessian is positivity of the relevant Schur
complement. At finite regulator, the inner Laplace expansion and its
required parameter derivatives are uniform on a smaller neighbourhood
of the outer saddle. Its $hB_+$ is evaluated at that saddle at the
next order, while its $Q_+$ contributes both first and second fibre
derivatives. An $O(h^2)$ term cannot alter the coefficient of $h$.
Uniqueness of the stationary expansion gives
\cref{eq:measured-loop-composition}. Equivalently this is an identity
of formal Gaussian coefficient germs. Neither interpretation supplies
a growing-volume remainder without additional estimates.
\src{SM}{V8}
\end{proof}

An $h$-independent retained reference $w(c)\dd c$ adds $\log w$ to
$Q_+$ and does not change $B_+$. That conversion must be undone before
using a Lebesgue-coordinate formula at a later stage. An $h$-dependent
reference or a redefinition of $h$ requires its own coefficient
conversion. Dropping the incoming one-loop density would lose the
$q_y$ and $q_{yy}$ terms; shifting the saddle again after including
the last square in \cref{eq:measured-loop-step} would count the same
order-$h$ contribution twice.

\begin{theorem}[Closed field-degree budget through two loops]
\label{thm:closed-loop-jet-packet}
Let $J^d$ denote the Taylor polynomial through total field degree $d$
at the stationary origin, counting both retained and eliminated
fields. Under the local hypotheses above, the input packet
\begin{equation}
 \bigl(J^8S,\ J^6Q,\ J^4B;\ J^7F\bigr)
 \label{eq:closed-loop-jet-packet}
\end{equation}
determines $J^8G$, $J^6Q_+$, and $J^4B_+$. The packet is closed
under any fixed finite number of measured two-loop steps on compatible
branches. Its field degrees are not external-momentum orders.
\end{theorem}
\begin{proof}
Write $L=DF(0)$, choose $LR=I$, and let $E$ include $\Ker L$.
Construct the chart as $X(y,c)=Ey+Rv(y,c)$ with $F(X)=c$. The
homogeneous degree-$d$ equation for $v$ has identity linear
coefficient, so $J^7F$ determines $J^7X$. Since $DS(0)=0$, an
unknown degree-eight term in $X$ first affects $S\circ X$ in degree
nine. Hence $J^8f$ is determined. The degree-six chart logarithmic
Jacobian uses $J^7X$, giving $J^6q$; $J^4b$ is also determined.

The equation $f_y(y_*(c),c)=0$ has invertible fibre linearization and
determines $J^7y_*$ from $J^8f$. Its unknown next coefficient cannot
enter $J^8G$ because the first derivative at zero vanishes. The
sixth jet of $H$ uses only $J^8f$, so $J^6Q_+$ is fixed. Four
retained derivatives of $f_{yyyy}$ use at most $J^8f$, those of
$q_{yy}$ use at most $J^6q$, and those of $b(y_*,c)$ use at most
$J^4b$. Differentiating $HC=I$ introduces no higher primitive jets.
This gives $J^4B_+$. The output has the same degree budget and a
stationary classical origin, while seventh jets compose by the chain
rule. Apply \cref{thm:measured-loop-composition} inductively.
\src{SM}{V8}
\end{proof}

The higher density degree is needed even when the incoming density
is constant. Consider the finite model
\[
 S(c,y,z)=\tfrac12(c^2+y^2+z^2)+\lambda c^4y^2z^2,
 \qquad Q=B=0,\qquad \lambda\ge0.
\]
Eliminating $y$ first is an exact Gaussian integral and gives
\[
 G_1(c,z)=\tfrac12(c^2+z^2),\qquad
 Q_1(c,z)=\tfrac12\log(1+2\lambda c^4z^2),\qquad B_1=0.
\]
At the next step, $\tfrac12\partial_z^2Q_1(c,0)=\lambda c^4$
contributes to $B_2(c)$. Truncating the intermediate density at
field degree four would erase it, despite the target being a quartic
retained coefficient. Likewise, a degree-eight term $c^4y^4$ of
the action affects $f_{yyyy}$, and a degree-seven change of the
constraint affects the quartic two-loop output through the chart.
These tests explain why the field budget cannot be replaced by a
quartic truncation of every input. \src{SM}{V8}

The finite packet specifies the algebra to be transported.
\Cref{thm:measured-two-loop,cor:two-loop-summability} supply its
analytic remainder interface only while their common domains,
marked norms, and transport bounds remain valid. Finite-stage
coefficient closure is not an all-scale nonlinear induction.

\subsection{A complete fourth-jet formula}
At the identity set $P=f_{zz}(0)^{-1}$, $A=f_{zc}(0)$ and
$Ju=(-PAu,u)$. Let $f_3,f_4$ be the third and fourth derivatives of
the actual pullback action and define fibre matrices
\[
 B(\xi)_{ij}=f_3[e_i^z,e_j^z,\xi],\qquad
 C(\xi,\eta)_{ij}=f_4[e_i^z,e_j^z,\xi,\eta].
\]
For the semisimple gauge groups and the normalized convention under
consideration, invariance makes the relevant first derivative at zero
vanish. The complete measured Hessian is
\begin{equation}\label{eq:complete-fourth-jet}
\begin{split}
 D^2\widehat Q(0)[u,v]={}&\tfrac12\Tr\{P C(Ju,Jv)
                       -P B(Ju)P B(Jv)\}\\
 &-\ell_2[Ju,Jv]
       +\tfrac1{12}\sum_e\Tr_{\mathfrak g}
                       (\operatorname{ad}u_e\operatorname{ad}v_e).
\end{split}
\end{equation}
Indeed differentiating $\tfrac12\log\det H_c$ gives the two trace
terms; differentiating the density gives the last line. Invariance
removes the term in the second derivative of the minimizing section,
not its first derivative $J$. The formula retains the block constraint,
pivot acceleration, and Haar normalization through the actual jets.
\src{F11}{V91--V93}

Exponential locality of these kernels gives an exponentially local
infinite-volume coefficient kernel and a row-summable finite-volume
periodization error. Its Fourier symbol $\mathcal K(p)$ is analytic
and obeys the Ward identity
\[
 \mathcal K(p)g(p)=0,\qquad g_\mu(p)=1-e^{\ii p_\mu},\qquad
 \mathcal K(0)=0.
\]
For the last assertion differentiate the Ward identity once at zero;
the derivatives of $g$ span the four coarse directions. This is a
thermodynamic limit of a coefficient, not of the entire quantum law.

As an exact finite-regulator normalization test, the reported
$\SU(2)$ checkerboard calculation uses
$c_{r,1}=(-1)^{r_0}T$, $\norm T_{\rm HS}^2=2$, on sixteen coarse cells.
The contributions to $n^{-1}D^2\widehat Q(0)[c,c]$ are
\[
 \frac{19}{3}-\frac{689}{384}
          +\frac{79227}{4480}-\frac{166901}{5760}
                      =-\frac{272249}{40320}.
\]
These are respectively the normalized measure contribution, neutral
term, charged second-vertex term, and subtracted first-vertex square.
The equality is a reported finite certificate in \src{F11}{V94--V99}.
The external momentum is $\pi$; its negative value is not a mass or an
infrared coefficient.

\section{The actual small-momentum coefficient}
\label{sec:actual-theta}
The remaining one-step transverse coefficient can be formulated without
a small real-momentum difference quotient. For the $45$-component
scalar-colour fibre write
\[
 H(k)=E^\#(k)D^\#(k)D(k)E(k),\qquad P(k)=H(k)^{-1}.
\]
Here $A^\#(k)=A(-k)^T$ is the Laurent continuation of the real-torus
adjoint. It is not complex conjugation of the complex momentum.
The finite coefficient enumeration reports $33$ nonzero Laurent
coefficients $H_r$, with
\[
 R_H=\max|r|_1=2,\qquad
 W_H=\sum_r\max(\norm{H_r}_1,\norm{H_r}_\infty)=1041/2.
\]
Together with $H(k)\succeq hI$, $h=24/2327$, these data give
\begin{equation}\label{eq:actual-holomorphic-strip}
 s=\frac{h}{12W_HR_H}=\frac2{2422407},\qquad
 \norm{P(k)}\le2/h\quad(\norm{\operatorname{Im}k}_\infty\le3s).
\end{equation}
To see this, bound the perturbation off the real torus by
$W_H(e^{3sR_H}-1)\le h/2$ and use a Neumann series.
The integer coefficient data are reported in \src{F12}{V5}; the
analytic implication is independent of momentum sampling.

Fix a unit real external polarization $v$ and external momentum $pe_0$.
Let $V(k,p)$ be the actual cubic vertex along $J(p)v$, carrying fibre
momentum $k$ to $k+p$, and let $C(k,p)$ be the zero-transfer fourth
vertex for the pair of opposite external modes. Define
$\widetilde V(k,p)=V(k+p,-p)$. These are the vertices of
\cref{eq:complete-fourth-jet}, not frozen fine-action vertices.

\begin{theorem}[Differentiated determinant representation]
\label{thm:theta-integral}
In the stated unit-polarization normalization,
\begin{align}
 g(k,p)&=\tfrac12\Tr\{P(k)C(k,p)
       -P(k+p)V(k,p)P(k)\widetilde V(k,p)\},\notag\\
 d_2(k)&=[p^2]g(k,p),\qquad
 \theta=\frac{19}{48}
           +\int_{\T^4}d_2(k)\frac{\dd k}{(2\pi)^4}.
 \label{eq:actual-theta}
\end{align}
Coefficient extraction commutes with integration. If $M$ is an upper
bound for $|g|$ on the complex domain
$\norm{\operatorname{Im}k}_\infty\le s$, $|p|\le s$, then the
four-dimensional tensor-grid mean $\theta_N$ satisfies
\begin{equation}\label{eq:theta-quadrature}
 |\theta_N-\theta|\le\frac{M}{s^2}
 \left\{\left(\frac{1+e^{-sN}}{1-e^{-sN}}\right)^4-1\right\}.
\end{equation}
\end{theorem}
\begin{proof}
Apply \cref{eq:complete-fourth-jet} to opposite commensurate plane waves
on finite tori. The fourth vertex has zero transfer, while the cubic
pair shifts momentum by $p$ and back. The trace per cell is therefore
the discrete mean of $g$. Locality and the inverse bound give the
thermodynamic integral; the normalized Haar term has quadratic
coefficient $19/48$ in this convention. The finite-range jets and
\cref{eq:actual-holomorphic-strip} make the integrand jointly
holomorphic on a neighbourhood of the indicated compact complex set.
Cauchy's coefficient formula justifies interchange of extraction and
integration and bounds $|d_2|$ there by $M/s^2$.

Contour shifting bounds the Fourier coefficient at $r$ by
$(M/s^2)e^{-s|r|_1}$. The tensor grid retains exactly the Fourier
indices $r=Nl$. Summing the nonzero aliases gives
\cref{eq:theta-quadrature}. \src{F12}{V5}
\end{proof}
For direct evaluation, the coefficients of $P(k+p)$ are
\[
 P_0=P,\qquad P_1=-PH'P,\qquad
 P_2=PH'PH'P-\tfrac12PH''P,
\]
and
\[
 d_2=\tfrac12\Tr\left\{PC_2-
             \sum_{a+b+c=2}P_aV_bP\widetilde V_c\right\}.
\]
All quantities are Taylor coefficients, including their factorials.
The opposite vertex must be differentiated in both its arguments, as
must the momentum-dependent minimizing lift. A uniform absolute error
$\varepsilon$ in each $d_2$ contributes at most $\varepsilon$ to the
average, in addition to separately bounded summation errors.

The strip in \cref{eq:actual-holomorphic-strip} is very conservative.
The supplied calculation does not populate a practically useful bound
$M$ and a certified numerical enclosure of $\theta$. Thus no sign of
$\theta$, and no identification with the all-level physical coefficient
$b_{\rm YM}$ of \cref{eq:bYM}, follows from this representation.

\section{Exact flat-background cancellation and finite-volume aliases}
\label{sec:flat-bloch-cancellation}
For the $\SU(2)$ constant flat coarse field in direction $1$, with
amplitude $tT$ and $\norm T_{\rm HS}^2=2$, put
$q(t)=e^{\ii t}\sin(t)/t$. The actual constrained charged fibre is
related to the free fibre at $k'=k+2te_1$ by a $45$-dimensional
coordinate map $J_t(k)$ satisfying
\begin{equation}\label{eq:flat-fibre-congruence}
 \det_{\C}J_t=q(t)^7,\qquad
 H_t^{\rm ch}(k)=J_t(k)^*H^{\rm free}(k')J_t(k).
\end{equation}
There are eight fine boundary coordinate factors and one coarse factor;
the exponent seven is their quotient. The neutral fibre is unchanged.
Taking the determinant in \cref{eq:flat-fibre-congruence} contributes
$14\log(\sin t/t)$. The fine Haar density contributes
$16\log(\sin t/t)$ and the coarse Haar factor contributes
$2\log(\sin t/t)$, so the normalized measure cancels that term exactly.

\begin{theorem}[Flat one-loop identity at every real momentum]
\label{thm:flat-bloch}
Let $f(k)=\log\det H^{\rm free}(k)$, in the charged-coordinate
normalization of \cref{eq:flat-fibre-congruence}. The normalized
order-one integrand changes by
\[
 \delta\widehat q_t(k)=f(k+2te_1)-f(k),\qquad
 \mathcal I_0(k)=2\partial_{k_1}^2f(k).
\]
Consequently its infinite-volume flat variation vanishes exactly. On a
side-$L$ coarse torus, if $f_r$ are the Fourier coefficients of $f$,
\begin{align}
 \frac{\widehat Q_L(c(t))-\widehat Q_L(0)}{L^4}
   &=\sum_{l\in\Z^4}f_{Ll}(e^{2\ii tLl_1}-1),\label{eq:flat-alias}\\
 L^{-4}\sum_{k\in\mathcal G_L}\mathcal I_0(k)
   &=-2L^2\sum_{l\ne0}l_1^2 f_{Ll}.\label{eq:flat-alias-second}
\end{align}
\end{theorem}
\begin{proof}
The charged real determinant contributes one complex log determinant
to $\tfrac12\log\det_{\R}H$. Combining the congruence and the exact
Haar cancellation gives the first identity. Half its second amplitude
derivative is $2\partial_{k_1}^2f$, the indicated unit-polarization
normalization. Translation invariance of Haar measure on momentum
torus makes the integral of the difference zero.

The gapped fibre gives a periodic holomorphic logarithm of the positive
real-torus determinant on a smaller strip. Its Fourier coefficients
are exponentially summable, including two derivatives. A side-$L$
grid averages $e^{\ii k\cdot r}$ to zero unless $r=Ll$, proving
\cref{eq:flat-alias,eq:flat-alias-second}. \src{F12}{V9}
\end{proof}
The reported side-two constant mean $7499/15680$ is such a Fourier
alias. It is neither a nonzero bulk constant nor a correction which
may be subtracted from a nonflat quadratic card to produce $\theta$.
The cancellation is integrated, not pointwise. In particular, the
absence of a constant term does not determine the transverse $p^2$
coefficient of \cref{thm:theta-integral}.

\section{The nonlinear classical continuum branch}
\label{sec:classical-root-continuum}
The one-cell determinant requires a common nonlinear source and a common
classical minimum as the fine lattice is refined. A finite Taylor
calculation alone does not identify either object at nonzero source.
This section constructs that interface for the balanced-root $\SU(2)$
convention and the full Wilson action. All domains are small classical
domains; no statement about a compact Gibbs measure follows from them.

\subsection{A critical norm for the actual nested root observable}
Let $t_{r,s}(e^a)$ be the ordered tree transport from $2r$ to $2r+s$,
where $s\in\{0,1\}^4$. After evaluating the root block
\cref{eq:root-block-local}, define
\[
 \phi_r(a)=\frac1{16}\sum_s\log t_{r,s}(e^a),\qquad
 h_r(a)=e^{-\phi_r(a)},\qquad
 \mathfrak T_M(a)_{r,\mu}=\log\bigl(h_rK_{r,\mu}h_{r+e_\mu}^{-1}\bigr).
\]
This is an actual coarse gauge transformation of the root output.
Gauge covariance of the chronological paths makes every iterate
orbit-equivalent to the literal nested root block. It does not equate
their representative-valued densities: the balancing gauge depends on
the fine field, and its measured coordinate terms remain necessary.
The identity logarithms below use one common primitive analytic domain.

All discrete $\ell^p$ norms in this section use unnormalized counting
measure. The derivative $\mathsf B_M=D\mathfrak T_M(0)$ is
\begin{equation}
 (\mathsf B_Ma)_{r,\mu}=\frac1{16}\sum_{s\in\{0,1\}^4}
       \bigl(a_{2r+s,\mu}+a_{2r+s+e_\mu,\mu}\bigr).
 \label{eq:classical-linear-root}
\end{equation}
Repeated occurrences on a short torus are counted with multiplicity.
Write $\mathfrak T_M=\mathsf B_M+\mathsf E_M$.

\begin{lemma}[Critical linear scale and contracting nonlinear remainder]
\label{lem:classical-two-norm}
For $1\le p\le\infty$,
\begin{equation}
 \norm{\mathsf B_M}_{p\to p}=2^{1-4/p}.
 \label{eq:classical-block-norm}
\end{equation}
On a common primitive maximum-norm ball there are constants independent
of $M$ such that
\begin{align}
 \norm{\mathsf E_M(a)}_2&\le C\norm a_4^2,\notag\\
 \norm{\mathsf E_M(a)-\mathsf E_M(a')}_2
    &\le C(\norm a_4+\norm{a'}_4)\norm{a-a'}_4,\notag\\
 \norm{D\mathsf E_M(a)}_{2\to2}&\le D\norm a_\infty.
 \label{eq:classical-root-remainder}
\end{align}
In particular the linear map has norm one on $\ell^4$ and norm
$1/2$ on the space $\ell^2$ in which the remainder is estimated.
\end{lemma}
\begin{proof}
Each row of \cref{eq:classical-linear-root} has weight two and each
column has weight $1/8$. Jensen gives
$\norm{\mathsf B_Ma}_p^p\le2^{p-4}\norm a_p^p$; constant fields
attain equality. The endpoint norms follow from the same row and
column sums. The nonlinear map is an analytic bounded stencil with
vanishing zeroth and first remainder jets. Each output uses at most
54 fine links and each link occurs in at most ten output supports.
Cauchy estimates on a smaller primitive polydisc therefore give
pointwise quadratic bounds in the local input sizes. H\"older with
two $\ell^4$ factors and the bounded incidence sums gives the first
two estimates. For the derivative use one local $\ell^\infty$
factor and the bounded $\ell^2$ incidence operator. These arguments
also hold on the complexified analytic domain. \src{SM}{V10, V12--V13}
\end{proof}

Let $L=2^n$, start on side $NL$, and let $\mathcal F_{L,N}$ be the
$n$-fold balanced root map to side $N$. If $a_s$ is the actual
intermediate field, set
\[
 R_s=a_s-\mathsf B^{[s]}a,
 \qquad R_{s+1}=\mathsf B_{NL/2^s}R_s+\mathsf E_{NL/2^s}(a_s).
\]
Here $\mathsf B^{[s]}$ is the corresponding linear composition.

\begin{theorem}[Uniform nonlinear root domain]
\label{thm:classical-root-domain}
There is $r_*>0$, independent of $L$ and $N$, such that for
$\rho=\norm a_4<r_*$ every actual intermediate logarithm is defined,
and
\begin{equation}
 \norm{R_s}_2\le4C\rho^2,\qquad
 \norm{a_s}_4\le\rho+4C\rho^2\le\tfrac54\rho.
 \label{eq:classical-root-domain}
\end{equation}
The composite is holomorphic on this common ball, with uniformly
bounded fixed-order derivatives on smaller balls. For every
$0<\alpha<1$, a further radius reduction gives
\begin{equation}
 \norm{\mathcal F_{K,N}(a)-\mathcal F_{K,N}(a')}_2
             \le K^{-\alpha}\norm{a-a'}_2,
 \qquad K=2^m,
 \label{eq:classical-error-contraction}
\end{equation}
when both inputs lie in that smaller $\ell^4$ ball.
\end{theorem}
\begin{proof}
The counting-norm inclusion $\ell^2\subset\ell^4$ has norm one.
If $Cr_*\le1/16$ and the primitive domain contains the ball of
radius $5r_*/4$, the induction hypothesis gives
\[
 \norm{R_{s+1}}_2\le\tfrac12\norm{R_s}_2
                     +C(\rho+\norm{R_s}_2)^2
 \le\bigl(2+(5/4)^2\bigr)C\rho^2<4C\rho^2.
\]
Starting from $R_0=0$ proves the bounds and the validity of every
primitive logarithm. The same induction holds on a complex ball.
Cauchy's formula there bounds every fixed derivative of the
nonlinear remainder uniformly; its linear part is already controlled.
For the last assertion, \cref{eq:classical-root-remainder} gives
$\norm{D\mathfrak T_M(a)}_{2\to2}\le1/2+D\norm a_\infty$.
Choose the smaller radius so this is at most $2^{-\alpha}$ at all
intermediate fields. Integrate along each segment in the primitive
ball and compose the $m$ estimates. \src{SM}{V12--V13}
\end{proof}

The marginal $\ell^4$ component has been retained rather than
iteratively charged as a nonlinear error. This distinction avoids a
spurious depth-dependent loss of the root domain. It controls the
observable, not contraction of the measured effective action.

\subsection{Weak source continuity at critical regularity}
Fix a coarse torus $\T_N^4$ and sample a continuum one-form by
\[
 (\mathsf S_LA)_\mu(x)=L^{-1}
       \operatorname{avg}_{x/L+[0,1/L)^4}A_\mu.
\]
This is the specified cell-average sampling, not a Wilson line
integral. It is contractive from $L^4$ to counting $\ell^4$.
Put $\mathcal O_{L,N}=\mathcal F_{L,N}\mathsf S_L$.
For dyadic $K$, define the finite-rank linear averages
\[
 (\mathsf A_KA)_{r,\mu}=K^{-1}\int_0^1
   \operatorname{avg}_{r/K+[0,1/K)^4+t e_\mu/K}A_\mu\,\dd t,
 \qquad \mathcal P_{K,N}=\mathcal F_{K,N}\mathsf A_K.
\]
The rank refers to the input averages; $\mathcal P_{K,N}$ remains a
nonlinear function of them.

\begin{theorem}[Uniform root-source limit on a critical ball]
\label{thm:classical-weak-source}
On a sufficiently small complex $L^4(\T_N^4)$ ball the dyadic limit
$\mathcal O_{\infty,N}$ exists and is holomorphic. For each
$0<\alpha<1$ and a suitable smaller radius $\rho$,
\begin{align}
 \norm{\mathcal O_{\infty,N}(A)-\mathcal P_{K,N}(A)}_4
      &\le4C\norm A_{L^4}^2K^{-\alpha},\notag\\
 \sup_{\norm A_{L^4}\le\rho}
 \norm{\mathcal O_{L,N}(A)-\mathcal O_{\infty,N}(A)}_4
      &\le C_{\rho,\alpha}L^{-\alpha/(1+\alpha)}.
 \label{eq:classical-source-limit}
\end{align}
At fixed $N$, the limit is weakly sequentially continuous on smaller
closed balls. Every fixed Fr\'echet derivative has the same mesh
power on smaller complex balls. Its derivative at zero is
\begin{equation}
 (\mathcal B_{\infty,N}A)_{r,\mu}
       =\int_0^1\int_{r+[0,1)^4}A_\mu(x+t e_\mu)\dd x\dd t.
 \label{eq:classical-continuum-linear-source}
\end{equation}
\end{theorem}
\begin{proof}
Stop the first $\log_2(L/K)$ steps before the final $\log_2K$
steps. The intermediate nonlinear remainder has $\ell^2$ norm at
most $4C\norm A_4^2$ and the remaining steps contract it by
$K^{-\alpha}$. The intermediate linear average differs from
$\mathsf A_K$ by at most $4K/L$ in $L^4\to\ell^4$ operator norm.
Indeed its kernels are Riemann sums of translated box indicators;
for almost every point, the indicator has at most two jumps in the
translation parameter. The resulting absolute row and column kernel
bounds give this estimate by H\"older. The uniform $\ell^4$
Lipschitz bound for the remaining composite therefore yields
\[
 \norm{\mathcal O_{L,N}(A)-\mathcal P_{K,N}(A)}_4
 \le4C\rho^2K^{-\alpha}+C_1\rho K/L.
\]
Comparing two refinements through a fixed $K$ proves that the maps
are uniformly Cauchy on the ball: first increase $K$, then the two
refinements. Passing to the limit gives the first bound; choosing
$K$ within a factor two of $L^{1/(1+\alpha)}$ gives the second.
For fixed $K,N$, weak convergence passes through the finitely many
linear averages and then their continuous nonlinear composition.
The uniform $K^{-\alpha}$ approximation passes this continuity to
the limit. Uniform convergence on complex balls and Cauchy's formula
give holomorphy and the differentiated bounds. The linear path sums
converge to \cref{eq:classical-continuum-linear-source}.
\src{SM}{V12--V13}
\end{proof}

For a fine field $a_L$ define its physical reconstruction by
$(\mathsf I_La_L)_\mu(x)=L(a_L)_\mu(z)$ on
$z/L+[0,1/L)^4$. Then $\mathsf S_L\mathsf I_La_L=a_L$ and
$\norm{\mathsf I_La_L}_{L^4}=\norm{a_L}_4$. Consequently, for fields
in the common ball,
\begin{equation}
 \mathsf I_La_L\rightharpoonup A\text{ in }L^4
 \quad\Longrightarrow\quad
 \mathcal F_{L,N}(a_L)\longrightarrow\mathcal O_{\infty,N}(A).
 \label{eq:classical-varying-source}
\end{equation}
Thus the source passes to the limit on varying minimizing fields;
compactness of the critical Sobolev embedding is not assumed.

\subsection{The full constrained minimum, not a homogeneous trial}
For the variational conclusion set $N=1$. In the fine Coulomb space
write $a=s/L+w$, with $d^*w=0$ and mean-zero $w$, and use the
detail energy norm $\norm{w}_V=\norm{dw}_2$. The discrete Sobolev
bound is $\norm w_4\le C_S\norm{w}_V$, uniformly in $L$.
The one-cell linear source is exactly $s$.

\begin{lemma}[Uniform local stationary branches]
\label{lem:classical-branches}
On fixed small source and detail balls the exact constraint
$\mathcal F_{L,1}(s/L+w)=v$ has a holomorphic solution
$s=s_L(v,w)=v+O((|v|+\norm{w}_V)^2)$. The full Wilson chart action
has a unique stationary point $w_L(v)$ in the detail ball and
\begin{equation}
 \norm{w_L(v)}_V\le C|v|^2,\qquad
 \tfrac12\norm{\xi}_V^2\le
 D_w^2\mathcal S_L\bigl(s_L(v,w)/L+w\bigr)[\xi,\xi]
                  \le\tfrac32\norm{\xi}_V^2.
 \label{eq:classical-branch-convexity}
\end{equation}
The same construction holds on
$V=\{w\in H^1(\T^4):\operatorname{div}w=0,\ \int w=0\}$
for the continuum constraint $\mathcal O_{\infty,1}(s+w)=v$ and
\begin{equation}
 \mathcal S_\infty(A)=\frac12\sum_{\mu<\nu}\int_{\T^4}
 |\partial_\mu A_\nu-\partial_\nu A_\mu+[A_\mu,A_\nu]|^2.
 \label{eq:classical-continuum-action}
\end{equation}
The bracket uses the same $\mathfrak{su}(2)$ normalization as the
fine Wilson action. Positivity statements concern real fields.
\end{lemma}
\begin{proof}
The nonlinear source remainder and its derivative are respectively
quadratic and linear in the critical field norm. The equation for
$s$ is therefore a contraction of a fixed finite-dimensional ball.
In this chart the quadratic action is $\tfrac12\norm{dw}_2^2$.
The full Wilson nonlinear force is $O((|v|+\norm{w}_V)^2)$ in the
dual detail norm, and its differentiated force is
$O(|v|+\norm{w}_V)$; these follow from its analytic plaquette
stencils, the critical Sobolev estimate, and the uniform chart
derivatives. The detail Riesz inverse has norm one, so its
stationary equation is another contraction. A radius reduction
makes the Hessian perturbation at most $1/2$, proving the bounds.
The continuum cubic and quartic terms obey the same estimates by
H\"older with exponents $2,4,4$ and four $L^4$ factors. The
complex contractions give holomorphic branches. The noncommuting
constant-field energy remains in the nonlinear action throughout.
\src{SM}{V9, V12--V13}
\end{proof}

Set $a_L^*(v)=s_L(v,w_L(v))/L+w_L(v)$ and
$\mathcal G_L(v)=\mathcal S_L(a_L^*(v))$, and denote their
continuum counterparts by $A_\infty^*(v)$ and $\mathcal G_\infty(v)$.

\begin{theorem}[Finite-amplitude local classical continuum limit]
\label{thm:classical-minimum-limit}
For every real $v$ in a common smaller source ball,
\begin{align}
 \mathcal G_L(v)&\longrightarrow\mathcal G_\infty(v),\notag\\
 \mathsf I_La_L^*(v)&\longrightarrow A_\infty^*(v)
                                          \quad\text{in }L^4,\notag\\
 D^L\mathsf I_La_L^*(v)&\longrightarrow\nabla A_\infty^*(v)
                                          \quad\text{in }L^2.
 \label{eq:classical-minimum-limit}
\end{align}
Here $D^L$ is the physical forward difference, with mesh $1/L$.
The scalar minima converge locally uniformly on a common smaller
complex source ball, with every fixed source derivative. This
statement gives no quantitative rate for the minimizing fields or
minimum values.
\end{theorem}
\begin{proof}
Uniform discrete Coulomb energy and bounded physical means give,
after interpolation on the fixed torus, strong compactness in
$L^p$ for $p<4$, weak compactness in $L^4$, and weak compactness of
physical first differences in $L^2$. Discrete summation by parts
identifies a limiting $H^1$ Coulomb field. Its source is exactly
$v$ by \cref{eq:classical-varying-source}.

The full Wilson action has an exact square representation. Regard
$U_p$ as a unit quaternion and reconstruct
$\mathcal C_L=L^2(U_p-I)$ on each plaquette cell. Then
\[
 \mathcal S_L(a_L)=\tfrac12\norm{\mathcal C_L}_{L^2}^2.
\]
A bounded-energy subsequence has an $L^2$-weakly convergent
curvature. Expand the ordered plaquette through quadratic order
only to identify its distributional limit. After reconstruction,
the cubic remainder is $O(L^{-1}\norm{\mathsf I_La_L}_3^3)$ in
$L^1$, the shifted quadratic terms are bounded by
$CL^{-1}\norm{\mathsf I_La_L}_2\norm{D^L\mathsf I_La_L}_2$,
and the real quadratic part by $CL^{-2}\norm{da_L}_2^2$.
They vanish. Strong $L^2$ field convergence identifies the bracket
in distributions, while the $L^4$ bound bounds it in $L^2$.
Thus $\mathcal C_L\rightharpoonup dA+[A,A]$ in $L^2$ and
\[
 \mathcal S_\infty(A)\le\liminf_L\mathcal S_L(a_L).
\]
This lower bound allows, rather than assumes away, concentration.

For recovery, first approximate a continuum Coulomb detail by a
finite Fourier polynomial. Project each of its finitely many modes
onto the discrete Coulomb space; these modewise projections converge
to their continuum values. Correct only the twelve constant source
coordinates by the uniform constraint chart. The correction tends
to zero, preserves the detail energy, and makes the nonlinear source
exact. Smooth plaquette consistency gives energy convergence. A
diagonal Fourier approximation then provides exact-source recovery
for every detail strictly inside the continuum chart, strongly in
$L^4$ and in physical derivative $L^2$.

Apply recovery to $A_\infty^*(v)$. The finite minima have no larger
action than these competitors, while the preceding lower bound and
continuum constrained convexity give the reverse limit inequality
and identify every weak subsequential limit with $A_\infty^*(v)$.
For a recovery field $\widetilde a_L$ with detail $\widetilde w_L$,
finite constrained convexity gives
\begin{equation}
 \tfrac14\norm{d(\widetilde w_L-w_L(v))}_2^2
 \le\mathcal S_L(\widetilde a_L)-\mathcal G_L(v)\longrightarrow0.
 \label{eq:classical-minimum-defect}
\end{equation}
Discrete Sobolev, the chart Lipschitz estimate, and the discrete
Hodge identity now give the two strong limits. Finally the scalar
minima are uniformly holomorphic and bounded on the common complex
ball. Every locally convergent subsequence agrees with the
continuum minimum on the real ball, hence everywhere by the
holomorphic identity theorem. Cauchy's formula gives convergence of
all fixed source derivatives. The compactness and recovery argument
supplies no mesh rate for these minima. \src{SM}{V12--V13}
\end{proof}

The distinction between source and energy convergence is necessary.
For example, a small commuting signed circulation $aE$ on the four
edges of one elementary square is discrete Coulomb, has zero mean,
and has exactly zero one-cell root source. Nevertheless
\[
 \norm{a_L}_4=\sqrt2|a|,\qquad
 \mathcal S_L(a_L)=1-\cos(4a)+20(1-\cos a)>0
\]
for nonzero small $a$. Its physical reconstruction converges weakly
to zero in $L^4$ while retaining this energy. The variational defect
\cref{eq:classical-minimum-defect}, not weak source continuity,
excludes that phenomenon for the actual minimizers. These results
identify a local finite-amplitude classical functional; they do not
identify it with a homogeneous trial, establish a quantum continuum
law, or replace the measured determinant estimates below.
\src{SM}{V11--V13}

\section{Cofinal measured response on the one-cell branch}
\label{sec:cofinal-measured-response}
The all-level local coefficient is now evaluated on the one-cell
$\SU(2)$ branch of \cref{sec:classical-root-continuum}, with Wilson
plaquette action $s_p=1-\tfrac12\Tr U_p$ and the balanced root source
$v=(v_1,\ldots,v_4)$ fixed while $L=2^n$ grows.
\Cref{thm:classical-root-domain,thm:classical-weak-source,thm:classical-minimum-limit}
supply respectively the common nonlinear domain, the limiting source,
and the full local classical minimum. Only the source-map estimate has
the stated power rate. The normal source differential and constrained
precision below are evaluated at that same minimum; independently
freezing the background or block map would change the coefficient.
The resulting determinant remains a local measured response, not the
complete compact free energy.

\subsection{Normal, constrained, gauge, and Haar factors}
Let $M_L(v)$ be the normal differential of the balanced constraint and
let $I+K_L(v)$ be the relative Lagrangian precision on the constrained
detail tangent space. Let $I+G_L^{\rm gh}(a_L^*(v))$ denote the relative
based-gauge Jacobian, expressed in the same logarithmic coordinates.
The superscript labels a determinant, not an added stochastic field.
Let $\ell_{\rm H}$ be the complete fine Haar logarithm and
\[
 \ell_c(v)=\sum_{\mu=1}^4 2\log\frac{\sin|v_\mu|}{|v_\mu|}
\]
the coarse product-Haar logarithm, continued analytically through zero.
All expressions are relative to their values at $v=0$.

\begin{proposition}[Complete finite measured determinant]
\label{prop:cofinal-measured-determinant}
The local one-loop coefficient relative to coarse Haar is
\begin{align}
 \widehat Q_L(v)={}&\log\det M_L(v)
               +\tfrac12\log\det(I+K_L(v))\notag\\
 &-\log\det(I+G_L^{\rm gh}(a_L^*(v)))
                  -\ell_{\rm H}(a_L^*(v))+\ell_c(v).
 \label{eq:cofinal-measured-determinant}
\end{align}
The normal determinant and the based-gauge determinant have different
domains and cannot replace one another.
\end{proposition}
\begin{proof}
The coarea density of the normal constraint is $(\det M_L)^{-1}$.
Gaussian integration on its actual stationary tangent gives
$\det(I+K_L)^{-1/2}$ after the source-independent normalization.
The based-gauge Jacobian and fine Haar amplitude are evaluated at the
same minimum. Their derivatives contribute at the next Laplace order,
not as extra Hessian terms at this order. Taking the negative logarithm
and converting the retained coordinate density to Haar gives
\cref{eq:cofinal-measured-determinant}.
\end{proof}

A change from logarithmic to right-Lie tangent coordinates transforms
the full bordered stationary Hessian and all constraint rows. Its
Jacobian then cancels the corresponding Haar factor. Dropping Haar
while leaving the old Hessian and constraint rows fixed does not
implement this coordinate change.

\begin{theorem}[Trace-ideal convergence of the subtracted determinant]
\label{thm:cofinal-trace-tail}
On a common smaller complex source ball, the constrained relative
precisions converge in the Schatten ideal $\mathfrak S_5$, as do the
relative gauge operators, with all fixed source derivatives. The
finite-dimensional normal matrices also converge. For
\[
 \mathcal L_5(K)=\sum_{r\ge5}\frac{(-1)^{r+1}}r\Tr K^r,
\]
using the small analytic branch, the quantity
\begin{equation}
 \mathcal T_L=\log\det M_L+\tfrac12\mathcal L_5(K_L)
                  -\mathcal L_5(G_L^{\rm gh})
 \label{eq:cofinal-tail}
\end{equation}
therefore converges locally uniformly with all fixed derivatives.
\end{theorem}
\begin{proof}
The common Coulomb realization and the local constrained-minimum
estimates give the stated ideal convergence for the literal precision
and gauge operators; the source-space convergence gives the normal
matrix convergence. These are the operator estimates of
\src{SM}{V15--V16}. On a smaller ball the operator norms are bounded by
a fixed number below one. Products of five $\mathfrak S_5$ factors are
trace class, and the remaining factors are bounded. The geometric tail
then converges normally in trace norm. Differentiation is justified by
the same bound, or by Cauchy's formula on nested source balls.
\end{proof}

The complementary finite expression is not negligible:
\begin{equation}
 \mathcal H_L=
 \sum_{r=1}^4\frac{(-1)^{r+1}}r
 \left\{\tfrac12\Tr K_L^r-\Tr(G_L^{\rm gh})^r\right\}
          -\ell_{\rm H}(a_L^*)+\ell_c,
 \qquad
 \widehat Q_L=\mathcal T_L+\mathcal H_L.
 \label{eq:cofinal-Born-head}
\end{equation}
The tail $\mathcal T_L$ uses the same normalization as
$\widehat Q_L$. In particular, ideal convergence does not imply convergence of its first four traces.
For the actual Wilson kinetic fourth-order contribution, the explicit
calculation gives
\[
 \Tr E_{{\rm kin},L}^{[4]}
   =-\tfrac32(1-L^{-4})E_4(v),\qquad
 \norm{E_{{\rm kin},L}^{[4]}}_{\mathfrak S_5}
                   \le C|v|^4L^{-16/5},
 \qquad
 E_4(v)=\tfrac12\sum_{\mu<\nu}|[v_\mu,v_\nu]|^2.
\]
Thus an operator can disappear in $\mathfrak S_5$ while its first
Born contribution tends to the nonzero polynomial $-3E_4/4$.
The complete low-order measured expression, rather than an absolute
bound on each determinant tail, is the object to renormalize.

\subsection{An exact all-level flat-holonomy coefficient}
\label{sec:cofinal-flat-response}
The full expression can be evaluated on the commuting source
$v=(xE_3,0,0,0)$, where $E_3$ is a unit imaginary quaternion. The
fine minimum is $a_1=xE_3/L$, $a_2=a_3=a_4=0$. All constant source
coordinates are retained; nonzero Fourier modes alone are integrated.
Define, for centered $m\in\mathcal K_L^3$,
\begin{align}
 h_L(x)&=1-\tfrac23(1-L^{-2})\sin^2x,\notag\\
 t_L(m)&=2L\,\operatorname{arsinh}
       \sqrt{\sum_{\nu=1}^3\sin^2(\pi m_\nu/L)},&
 D_L(m)&=\cosh t_L(m)-1.
 \label{eq:flat-cofinal-symbols}
\end{align}
Here $\mathcal K_L$ is any centered complete set of residues modulo $L$.

\begin{theorem}[Complete flat coefficient and its cofinal rate]
\label{thm:cofinal-flat}
On a sufficiently small fixed complex disc about zero,
\begin{equation}
 \widehat Q_L(x)=3\log h_L(x)+
 2\sum_{m\in\mathcal K_L^3\setminus\{0\}}
       \log\left(1+\frac{1-\cos2x}{D_L(m)}\right).
 \label{eq:flat-cofinal-exact}
\end{equation}
It converges to
\begin{align}
 \widehat Q_\infty(x)={}&3\log\left(1-\tfrac23\sin^2x\right)\notag\\
 &+2\sum_{m\in\Z^3\setminus\{0\}}
 \log\left(1+\frac{1-\cos2x}{\cosh(2\pi|m|)-1}\right).
 \label{eq:flat-cofinal-limit}
\end{align}
For a sufficiently small $\rho<1/4$,
$\sup_{|x|\le\rho}|\widehat Q_L-\widehat Q_\infty|
\le C_\rho L^{-2}$. Every fixed source derivative has the same rate
on a smaller disc. The dyadic errors are summable.
\end{theorem}
\begin{proof}
In each charged mode the logarithmic tangent Jacobian, gauge
determinant, and transverse Wilson determinant cancel to a shifted
scalar dispersion. The normal root differential and retained tangent
correction contain powers of
$\sin x/(L\sin(x/L))$; the complete fine and coarse Haar conversion
cancels those powers. The surviving root factor is $h_L^3$.
Taking the exact product over the longitudinal Fourier coordinate
by the Chebyshev product identity yields
\cref{eq:flat-cofinal-exact}. This cancellation uses the complete
finite expression \cref{eq:cofinal-measured-determinant}.

For centered modes, $2|m|\le t_L(m)\le2\pi|m|$. On a growing
low-frequency region,
$|t_L(m)-2\pi|m||\le C|m|^3L^{-2}$.
The derivative of each logarithm with respect to $t_L$ has an
exponential bound in $|m|$, uniformly on the fixed small disc.
Summing the low-frequency errors gives $O(L^{-2})$; the remaining
modes have exponentially small total contribution. Finally
$h_L-h_\infty=O(L^{-2})$. Normal convergence and Cauchy's estimate
prove the derivative statement.
\end{proof}
This is an all-level one-final-cell local coefficient, not the
one-step thermodynamic Bloch identity of \cref{thm:flat-bloch} and not
a sum over all compact flat sectors. The two identities check different
parts of the same measured normalization.
\sourcebasis{The common source and minimum are treated in
\src{SM}{V12--V14}; trace-ideal control and the low-Born obstruction in
\src{SM}{V15}; the complete gauge--Haar formula and flat evaluation in
\src{SM}{V16}.}

\section{The complete noncommuting quartic logarithm}
\label{sec:measured-quartic-log}
A flat source has $E_4=0$ and cannot determine the nonabelian marginal
coefficient. Retain the balanced-root and coarse-Haar conventions of
\cref{sec:cofinal-measured-response}, and write
$\mathcal Q_L^{\rm H}=\widehat Q_L$ for the complete relative measured
coefficient. The notation $[4]$ means its homogeneous Taylor polynomial
of degree four, not an unnormalized fourth derivative. The norm in
$E_4=\tfrac12\sum_{\mu<\nu}|[v_\mu,v_\nu]|^2$ is the imaginary-quaternion
norm. For example, $v=(xE_1,yE_2,0,0)$ gives $E_4=2x^2y^2$.

The source-response matrix $\Gamma_*$ and the precision-response
matrix $h_*$ enter only through the measured combination
$\mathsf C_*=\Gamma_*-2h_*^{\mathsf T}$. The individual source matrix
need not be diagonal. A scalar measured combination is not a scalar
replacement of the microscopic source filter.

\begin{certificate}[Finite reduction of the complete measured banks]
\label{cert:measured-quartic-bank}
For the literal balanced-root mixed source and quadratic banks, the
finite Laurent identities in \src{F14}{V28--V30} cancel the
 off-diagonal entries of $\mathsf C_*$ and reduce every diagonal entry to
\begin{equation}
 (\mathsf C_*)_{\mu\mu}=
 \frac1{(2\pi)^4}\int_{\R^4}
                 \{\Phi_\mu(\xi)-16\Phi_\mu(2\xi)\}\dd\xi.
 \label{eq:quartic-flux-reduction}
\end{equation}
Writing $r=|\xi|$ and
\begin{align*}
 S(x)&=\left(\frac{\sin(x/2)}{x/2}\right)^2,&
 \chi(\xi)&=\prod_{\rho=1}^4S(\xi_\rho),\\
 u(x)&=4\frac{\sin x}{x}-\tfrac23(5\cos x+1),&
 v_0(x)&=\tfrac23(\cos x+2)-8\frac{\sin x}{x},
\end{align*}
the reduced kernel is
\[
 \Phi_\mu(\xi)=\chi(\xi)
 \left\{\frac{u(\xi_\mu)}{r^2\xi_\mu^2}
                         -\frac{v_0(\xi_\mu)}{r^4}\right\}.
\]
The singularities at $\xi_\mu=0$ are interpreted by analytic
continuation of $u(x)/x^2$. The reported rational calculation checks
64 mixed-stencil identities, 16 unprojected gradient identities,
20 primitive reductions, 20 mixed parity reductions, and 10 scaling
or Taylor identities. Its finite verification is distinct from the
real-variable integral argument below.
\end{certificate}

\begin{lemma}[Critical scale flux]
\label{lem:critical-scale-flux}
Suppose $\Phi(\xi)=a|\xi|^{-4}+O(|\xi|^{-2})$ near zero,
$\Phi$ is integrable outside the unit ball, and
$\Phi(\xi)-16\Phi(2\xi)$ is integrable. Then
\begin{equation}
 \int_{\R^4}\{\Phi(\xi)-16\Phi(2\xi)\}\dd\xi
                        =a\,|S^3|\log2.
 \label{eq:critical-scale-flux}
\end{equation}
\end{lemma}
\begin{proof}
Integrate on $\varepsilon<|\xi|<R$ and change variables in the second
term before removing either cutoff. The difference equals the integral
on $\varepsilon<|\xi|<2\varepsilon$ minus that on $R<|\xi|<2R$.
The latter tends to zero. In the former, the leading radial integral
is $a|S^3|\log2$ and the error is $O(\varepsilon^2)$.
\end{proof}

\begin{proposition}[Scalar measured matrix]
\label{prop:measured-matrix-scalar}
The finite bank reduction in \cref{cert:measured-quartic-bank}, with its
actual source and precision terms, gives
\begin{equation}
 \mathsf C_* =\frac{3\log2}{4\pi^2}I_4.
 \label{eq:measured-matrix-scalar}
\end{equation}
\end{proposition}
\begin{proof}
The elementary expansions are $u(x)=x^2+O(x^4)$,
$v_0(x)=-6+x^2+O(x^4)$, and $\chi(\xi)=1+O(r^2)$.
Hence $\Phi_\mu=6r^{-4}+O(r^{-2})$. The box filters give integrability
at infinity; the leading singularity cancels in the scale difference.
Apply \cref{lem:critical-scale-flux}, using $|S^3|=2\pi^2$ and the
normalization $(2\pi)^{-4}$. The certified off-diagonal cancellation
then proves \cref{eq:measured-matrix-scalar}.
\end{proof}
The two critical integrals in \cref{eq:quartic-flux-reduction} must not
be separated before their common cutoff is removed. Doing so would
incorrectly discard the scale flux.

\begin{theorem}[Measured noncommuting logarithm and finite quartic part]
\label{thm:measured-quartic-log}
For the specified local balanced-root branch, coefficientwise in the
finite-dimensional space of quartic source polynomials,
\begin{align}
 q_{*,4}^{\rm rel}&=-\frac{3\log2}{\pi^2}E_4,\notag\\
 \frac{(\mathcal Q_{2^n}^{\rm H})_{[4]}}n
       &\longrightarrow-\frac{11\log2}{3\pi^2}E_4,\notag\\
 (\mathcal Q_{2^n}^{\rm H})_{[4]}
         +n\frac{11\log2}{3\pi^2}E_4
       &\longrightarrow\mathcal P_{{\rm fin},4}.
 \label{eq:measured-quartic-log}
\end{align}
The finite polynomial $\mathcal P_{{\rm fin},4}$ retains both scale
endpoints. Equivalently the coefficient divided by $\log L$ tends to
$-11E_4/(3\pi^2)$ along dyadic $L$.
\end{theorem}
\begin{proof}
The complete measured expansion identifies the relative term as
$-g(v)\cdot\mathsf C_*v$, where $g=\nabla E_4$ in the stated
normalization. Euler's identity gives $g(v)\cdot v=4E_4(v)$, so
\cref{eq:measured-matrix-scalar} yields the first line. The remaining
Wilson, gauge, and Haar contribution has logarithmic coefficient
$-2\log2\,E_4/(3\pi^2)$ in the same convention. Their sum gives
$-11\log2\,E_4/(3\pi^2)$. The cofinal finite-part decomposition of the
actual measured expression, with both endpoint series retained, gives
the last line. The decomposition and endpoint convergence are the
analytic inputs of \src{F14}{V19--V22}; the matrix evaluation is
\cref{prop:measured-matrix-scalar}.
\end{proof}

\begin{scope}[An evaluated coefficient is not a completed trajectory]
At the level of these Taylor polynomials, the classical coefficient
$\beta$ of $E_4$ acquires the correction $-11\log L/(3\pi^2)$.
This resolves the scalar evaluation for this measured $\SU(2)$ branch.
It does not establish a joint $(\beta,L)$ remainder estimate, control
the full compact density, or identify $\beta$ with the physical
coupling convention of \cref{eq:bYM}. In particular, the source matrix
$\Gamma_*$ can remain non-diagonal even though
$\Gamma_*-2h_*^{\mathsf T}$ is scalar. Its conditional innovations
are not removed by the measured cancellation.
\end{scope}
\sourcebasis{The complete noncommuting finite-part decomposition is
proved in \src{F14}{V19--V22}; the measured mixed-bank reduction and
critical scale evaluation are in \src{F14}{V28--V30}.}


\subsection{Renormalized fourth source jets and their distinct counterterm}
\label{sec:fourth-source-renormalization}
The fourth coefficient of a random blocked source is not the quartic
coefficient of its measured effective action. The scalar cancellation in
\cref{prop:measured-matrix-scalar} concerns the latter. The source itself
continues to depend on the complete, generally cross-oriented matrix
$\Gamma_*$. This distinction has a quantitative consequence on the common
free Gaussian coupling of the balanced-root tower.

For dyadic $L=pM$, let $\mathcal F^{\rm root}_{L\to M}$ be the actual
iterated root map and set
\[
 \mathcal V_{L,M}=\mathcal F^{\rm root}_{L\to M,[4]}(Z_L),
 \qquad j=\log_2(L/M).
\]
Here $Z_L$ is the transverse Gaussian input, and $[d]$ denotes a homogeneous
Taylor coefficient, without a factorial. At an intermediate side $R=sM$,
write $Y_{L,R},Q_{L,R},\mathcal R_{L,R}$ for the first three incoming
coefficients. Denote the degree-$d$ primitive block by $E_R^{[d]}$ and the
exact linear transport from side $R/2$ to $M$ by $B_{s/2}$. The contracted
cubic has the finite-filter decomposition
\begin{equation}
 \mathcal R_{L,R}=\log_2(L/R)D_{L,R}+\widehat{\mathcal R}_{L,R},
 \qquad D_{L,R}=\mathcal A_{L/R,R}(\Gamma_*)Z_L,
 \label{eq:source-cubic-finite-filter}
\end{equation}
where $\mathcal A$ is the propagated linear source filter of this same
block. Its finite endpoint terms remain in
$\widehat{\mathcal R}_{L,R}$. Below, local $L^q$ bounds are
one-link moment bounds, not unnormalized sums over the output volume.
Convergence on a fixed finite output carrier follows componentwise.
Define
\begin{align}
 \mathcal H_{L,M;s}
   &=B_{s/2}D^2E_R^{[2]}[Y_{L,R},D_{L,R}],\notag\\
 \mathcal J_{L,M;s}
   &=B_{s/2}\bigl\{
       D^2E_R^{[2]}[Y_{L,R},\widehat{\mathcal R}_{L,R}]
       +E_R^{[2]}(Q_{L,R})\notag\\
   &\hspace{23mm}+\tfrac12D^3E_R^{[3]}[Y_{L,R},Y_{L,R},Q_{L,R}]
       +E_R^{[4]}(Y_{L,R})\bigr\},\notag\\
 \mathcal H^{\rm loc}_{L,M}&=\sum_{\ell=1}^{j}\mathcal H_{L,M;2^\ell}.
 \label{eq:fourth-source-shells}
\end{align}
The derivatives in this formula are the multilinear derivatives of the
homogeneous primitive polynomials. Colour invariance and the Lie bracket
make $\mathcal H_{L,M;s}$ a centered second Gaussian chaos.

\begin{proposition}[Exact fourth-source subtraction]
\label{prop:fourth-source-subtraction}
On the finite Gaussian carrier,
\begin{align}
 \mathcal V_{L,M}
   &=\sum_{\ell=1}^{j}(j-\ell)\mathcal H_{L,M;2^\ell}
      +\sum_{\ell=1}^{j}\mathcal J_{L,M;2^\ell},\notag\\
 \widehat{\mathcal V}_{L,M}
   &:=\mathcal V_{L,M}-j\mathcal H^{\rm loc}_{L,M}
     =\sum_{\ell=1}^{j}
        \bigl(\mathcal J_{L,M;2^\ell}
                         -\ell\mathcal H_{L,M;2^\ell}\bigr).
 \label{eq:fourth-source-subtraction}
\end{align}
\end{proposition}
\begin{proof}
Substitution of $\epsilon Y+\epsilon^2Q+\epsilon^3\mathcal R$
into the primitive map gives, at degree four,
$D^2E^{[2]}[Y,\mathcal R]+E^{[2]}(Q)
 +D^3E^{[3]}[Y,Y,Q]/2+E^{[4]}(Y)$.
Transport each term through the remaining linear blocks and insert
\cref{eq:source-cubic-finite-filter}. At side $R=2^\ell M$ the incoming
cubic clock is $\log_2(L/R)=j-\ell$. This gives the first identity;
subtracting $j\mathcal H^{\rm loc}_{L,M}$ gives the second.
\end{proof}

\begin{theorem}[Cofinal finite part of the fourth source]
\label{thm:fourth-source-finite-part}
For the specified balanced-root Gaussian tower, the local shell estimates
of \src{F14}{V31--V32} give, for every finite $q\ge1$,
\begin{align}
 \|\mathcal H_{L,M;2^\ell}\|_{L^q}
      &\le C_q2^{-\ell}\sqrt{1+\ell},\notag\\
 \|\mathcal J_{L,M;2^\ell}
                       -\ell\mathcal H_{L,M;2^\ell}\|_{L^q}
      &\le C_q2^{-\ell}(1+\ell)^{3/2}.
 \label{eq:fourth-source-summable-shells}
\end{align}
The constants are independent of incoming depth and output side in the
local normalization of \cref{eq:fourth-source-shells}. For fixed $M$,
jointly on each finite output carrier,
\begin{align}
 \mathcal H^{\rm loc}_{L,M}&\longrightarrow
                            \mathcal H^{\rm loc}_{\infty,M},\notag\\
 \widehat{\mathcal V}_{L,M}&\longrightarrow
 \mathcal V^{\rm fin}_{\infty,M}
   =\sum_{\ell\ge1}
      \bigl(\mathcal J_{\infty,M;2^\ell}
                  -\ell\mathcal H_{\infty,M;2^\ell}\bigr)
 \label{eq:fourth-source-finite-part}
\end{align}
in every finite $L^q$. The last series has tail
$C_q2^{-b}(1+b)^{3/2}$ beyond shell $2^b$. At each output link there
are constants $\eta,C>0$, uniform in $L,M$, such that
\begin{equation}
 \E\exp\{\eta|\widehat{\mathcal V}_{L,M}|^{1/2}\}\le C.
 \label{eq:fourth-source-stretched}
\end{equation}
\end{theorem}
\begin{proof}
The estimates in \cref{eq:fourth-source-summable-shells} use the
contracted cubic decomposition before summing its absolute kernels.
Its individual inner errors have both bounded local kernel mass and a
decaying local moment bound. Applying the exact outer averaging row
before summing those errors yields the second, summable estimate;
summing absolute kernels first would retain an unnecessary incoming-depth
factor. These are the local analytic estimates of \src{F14}{V32}.
For a fixed outer shell all incoming coefficients and their finite endpoint
terms converge on the common Gaussian coupling. Truncation to a finite
head, followed by the uniform summable tails, proves
\cref{eq:fourth-source-finite-part}. The second- and fourth-chaos
projections converge separately. Degree-four hypercontractivity and the
uniform $L^2$ estimate give
$\|\widehat{\mathcal V}_{L,M}\|_{L^q}\le Cq^2$ for $q\ge2$.
Expansion of the exponential then proves
\cref{eq:fourth-source-stretched} for sufficiently small $\eta$.
\end{proof}

At the terminal cell this gives
$\mathcal F^{\rm root}_{2^n\to1,[4]}(Z_{2^n})
 -n\mathcal H^{\rm loc}_{2^n,1}\to\mathcal V^{\rm fin}_{\infty,1}$.
The counterterm is a second-chaos source observable, not the scalar
$E_4$ coefficient of \cref{thm:measured-quartic-log}. In particular,
replacing $\Gamma_*$ by its measured scalar combination would change
\cref{eq:source-cubic-finite-filter}. The finite incoming counterterm
is retained in the subtraction; convergence of that counterterm alone
does not justify multiplying its limiting error by $\log_2L$.
The stretched-exponential estimate does not permit exponentiation of an
uncut quartic polynomial or supply a remainder bound for the full compact
nonlinear map.
\sourcebasis{The fourth-source chain rule and quadratic-chaos logarithm
are developed in \src{F14}{V31}; the locally renormalized finite part,
finite endpoint convention, and summable shell bounds in
\src{F14}{V32}.}


\section{What remains between the local box and the full compact law}
\label{sec:compact-local-interface}
The global completed integral and the controlled local integral solve
different parts of the construction. The former has a fixed compact
carrier and exact score identities; the latter has the stronger
localized derivative estimates in
\cref{thm:uniform-local-loop,thm:measured-two-loop}.
Neither positivity of the former nor concentration in the latter
bounds the logarithmic derivatives of their ratio without additional
normalized estimates.

For a retained sigma-field $\mathscr S$, the exact identity
\begin{equation}\label{eq:conditional-means-retained}
 \Cov(f,g)=\E\Cov(f,g\mid\mathscr S)
          +\Cov\bigl(\E[f\mid\mathscr S],\E[g\mid\mathscr S]\bigr)
\end{equation}
keeps the conditional-mean term visible. A decay bound for the first
term alone cannot be transferred to the full covariance. Similarly,
a small exceptional-set probability is not a bound on the Dirichlet
energy of every normalized function concentrated on that set.
These distinctions are already present in the normalized source
calculus and the core--bad-set compiler; they remain essential after
compact completion. \src{F12}{V19, V23--V26, V33}

The full compact estimates of \cref{ch:full-compact-sources} supply
additional, genuinely unrestricted source control on their stated
bare-Wilson sequence. They do not identify the ratio between this
local root integral and a complete blocked interacting law.
The nonlinear construction therefore strengthens the input to the
owner and forest analysis of \cref{ch:remainder}, without discharging
its cofinal uniformity hypotheses. In the auxiliary spectral problem,
\cref{ch:compact-repair} makes the obstruction concrete: a well with
exponentially small mass can still prevent a temperature-uniform
hidden diffusion gap. The valid replacement uses same-law visible
repair and states separately the coverage needed to turn local
estimates into an unrestricted interface inequality.

\chapter{Full compact Wilson sources and entropy transport}
\label{ch:full-compact-sources}
The local measured expansion does not exhaust the available control of
the compact law. This chapter records estimates which are instead
proved for the unrestricted Wilson measure. Their carrier is explicit:
they concern the bare $\SU(2)$ Wilson law, without additional renewal
weights. Applying them to a renewal-loaded or blocked law requires the
same-law identification or an actual comparison budget. No small-field
cutoff occurs in the statements.

\section{Thermal sources under the unrestricted law}
On a four-dimensional periodic lattice of side $M$, write $V=M^4$ and
\begin{equation}
 \dd\mu^{\rm W}_{M,\beta}
 =\frac{e^{-\beta S(U)}}{Z^{\rm W}_{M,\beta}}
                  \prod_e\dd H(U_e),\qquad
 S=\sum_p\left(1-\tfrac12\Tr U_p\right),\qquad
 X_p=\beta\left(1-\tfrac12\Tr U_p\right).
 \label{eq:full-Wilson-law}
\end{equation}
Haar is normalized. Gauge-tree integration and an explicit small-link
lower region give the full-pressure estimate
\begin{equation}
 e^{-48V}c_b^{D}\beta^{-3D/2}
 \le Z^{\rm W}_{M,\beta}
 \le C_1^{N_0}\beta^{-3N_0/2},\quad
 D=3V+1,\quad N_0=3V-4M^3+1,
 \label{eq:full-Wilson-pressure}
\end{equation}
where $c_b=8/(3\pi^3)$ and $C_1=(\pi/2)^{5/2}$, for $M\ge3$,
$\beta\ge1$. The boundary loss is explicit; this is not an expansion
of a restricted partition function. Together with reflection estimates,
it yields uniform local exponential moments of $X_p$ on dyadic
families with $(\log\beta)/M$ bounded. \src{F14}{V95--V97}

Distribute each plaquette over its four incident unit four-cells by
\begin{equation}
 T_x=\frac14\sum_{\mu<\nu}
 \sum_{\substack{\varepsilon\in\{0,1\}^4\\
                      \varepsilon_\mu=\varepsilon_\nu=0}}
 X_{(x+\varepsilon;\mu,\nu)},\qquad
                 \sum_xT_x=\beta S.
 \label{eq:full-Wilson-cell-energy}
\end{equation}
The shared faces are part of this definition. They are not removed
when reflection positivity is used. Put
$\widetilde T_x=T_x-\E T_x$ and
$\chi_{M,\beta}=\beta^2\Var(S)/V$.

\begin{proposition}[Exact covariance norm of cell-energy sources]
\label{prop:full-Wilson-covariance}
At a dyadic finite regulator, the covariance matrix of the cell energies
satisfies
\begin{equation}
 \norm{\Cov(T)}_{\ell^2\to\ell^2}=\chi_{M,\beta},\qquad
 \Var\left(\sum_xa_xT_x\right)
                       \le\chi_{M,\beta}\sum_x|a_x|^2.
 \label{eq:full-Wilson-covariance}
\end{equation}
There is no restriction on the support of $a$.
\end{proposition}
\begin{proof}
The shared-face chessboard inequality bounds the inhomogeneous
normalized pressure by the mean of uniform cell-source pressures.
Differentiating twice at zero gives the upper quadratic-form bound.
Translation invariance makes the constant vector an eigenvector of
the covariance matrix, with eigenvalue
$V^{-1}\Var(\sum_xT_x)=\chi_{M,\beta}$. This proves equality of the
operator norm. \src{F14}{V98--V99}
\end{proof}

\section{A selected near-critical source window}
The following result uses an existential coupling selection. It is
not a uniform assertion for every large bare coupling or a prescribed
renormalization trajectory. Set
\[
 J(t)=-\log(1-t)-t,\qquad t<1.
\]

\begin{theorem}[Full-law pressure and concentration at selected couplings]
\label{thm:full-Wilson-source-window}
There are dyadic couplings $\beta_M^\diamond$ with
\[
 e^{\sqrt{\log\log M}}\le\beta_M^\diamond\le\log M,
\]
for sufficiently large $M$, and numbers $\lambda_M\downarrow0$ after
replacement by a decreasing upper envelope, $R_M\to1$, such that,
with $k_M=9/2+\lambda_M$ and expectation at this selected law,
\begin{align}
 \frac{\beta_M^\diamond}{V}\E S&\longrightarrow\frac92,&
 \chi_{M,\beta_M^\diamond}&\longrightarrow\frac92,\notag\\
 0\le \log\E\exp\left(\sum_xb_x\widetilde T_x\right)
       &\le k_M\sum_xJ(b_x),& \norm{b}_\infty&\le R_M.
 \label{eq:full-Wilson-source-window}
\end{align}
In addition, for the normalized uniform pressure
\[
 \Psi_M(t)=V^{-1}\log\E e^{t\beta_M^\diamond(S-\E S)},
\]
one has $|\Psi_M(t)-(9/2)J(t)|\le\lambda_MJ(t)$ on a source
window containing $[-R_M,R_M]$.
\end{theorem}
\begin{proof}
The full-pressure bounds and thermal moment control first give an
integrated bound on the susceptibility over logarithmic coupling
windows. Energy saturation identifies its limiting averaged value
as $9/2$. The normalized energy and susceptibility obey the exact
logarithmic-derivative relation obtained by differentiating the
partition function. The window-averaged absolute susceptibility
defect therefore tends to zero. A maximal-average selection inside
the window supplies a center for which all admitted one-sided
pressure integrals have the same small defect. Integrating twice in
the source variable gives the stated uniform-pressure estimate.
The shared-face chessboard inequality then gives the upper bound
for arbitrary arrays $b$; Jensen's inequality gives the lower bound.
The coupling-window construction and its quantitative selection are
proved in \src{F14}{V99--V100}.
\end{proof}

\begin{corollary}[Bernstein bound and total-action fluctuations]
\label{cor:full-Wilson-Bernstein}
For a real array $a$, put $X(a)=\sum_xa_x\widetilde T_x$,
$v(a)=k_M\norm{a}_2^2$, and $B(a)=\norm{a}_\infty/R_M$. Then
\begin{align}
 \log\E e^{\pm tX(a)}
   &\le\frac{v(a)t^2}{2(1-B(a)t)},\qquad 0\le t<B(a)^{-1},\notag\\
 \Pr\{|X(a)|>\sqrt{2v(a)u}+B(a)u\}&\le2e^{-u},\qquad u>0.
 \label{eq:full-Wilson-Bernstein}
\end{align}
Zero arrays are interpreted by continuity. Along the same selection,
\begin{equation}
 \frac{\beta_M^\diamond(S-\E S)}{M^2}
                         \ \Longrightarrow\ N(0,9/2),
 \label{eq:full-Wilson-action-CLT}
\end{equation}
with convergence of every fixed moment.
\end{corollary}
\begin{proof}
Use $J(s)\le s^2/[2(1-|s|)]$ in
\cref{eq:full-Wilson-source-window}, and enlarge the denominator using
$R_M\le1$. Chernoff optimization gives
\cref{eq:full-Wilson-Bernstein}. For the last assertion, insert
$t/\sqrt V$ into the uniform pressure estimate. Since
$VJ(t/\sqrt V)\to t^2/2$, the moment generating functions converge
near zero to $\exp(9t^2/4)$. Their common local exponential bound
also gives convergence of fixed moments.
\end{proof}
The central limit theorem is for the total thermal action in this
microscopic lattice regime. It neither makes individual cells
independent nor constructs an interacting continuum field theory.
A bound on the pressure around zero also must not be differentiated
at an arbitrary nonzero tilt to assert an unproved tilted covariance
estimate.

\section{Finite entropy budgets for changes of law}
\label{sec:entropy-source-transport}
Fix one selected finite law $\mu=\mu^{\rm W}_{M,\beta_M^\diamond}$.
For any probability $\nu\ll\mu$, let
$H=D(\nu\Vert\mu)<\infty$. A change of law need not be an
infinitesimal tilt.

\begin{theorem}[Entropy-controlled full-law source transport]
\label{thm:entropy-source-transport}
For every real array $a$,
\begin{equation}
 \left|\E_\nu\sum_xa_xT_x-\E_\mu\sum_xa_xT_x\right|
 \le\sqrt{2k_MH}\,\norm{a}_2
                         +\frac H{R_M}\norm{a}_\infty.
 \label{eq:entropy-source-transport}
\end{equation}
Equivalently the mean-response vector admits a decomposition
$m=u+w$ with $\norm{u}_2\le\sqrt{2k_MH}$ and
$\norm{w}_1\le H/R_M$. The decomposition need not be unique.
\end{theorem}
\begin{proof}
The entropy variational inequality and
\cref{eq:full-Wilson-Bernstein} give, for either sign,
\[
 \pm\E_\nu X(a)\le\frac Ht+
                         \frac{v(a)t}{2(1-B(a)t)}.
\]
For $H,v(a)>0$, choose
$t=\sqrt{2H/v(a)}/(1+B(a)\sqrt{2H/v(a)})$; the result is
$\sqrt{2v(a)H}+B(a)H$. Limiting cases follow directly.
The right side of \cref{eq:entropy-source-transport} is the support
function of the sum of the indicated Euclidean and $\ell^1$ balls.
Finite-dimensional separation proves the decomposition.
\end{proof}
One may replace $H$ by the relative entropy of the induced cell-energy
laws, since the source is measurable with respect to $T$. Thus the
comparison can exploit information at the level of the actual
observable rather than charging unrelated hidden variables. For
$a_x=1/V$, an entropy budget $H=o(V)$ implies convergence of mean
energy densities.

There is also an explicit class of perturbations for which the budget
is controlled. Define
\[
 \frac{\dd\mu_a}{\dd\mu}
   =\exp\{\sum_xa_x\widetilde T_x-K(a)\},\qquad
 K(a)=\log\E_\mu e^{\sum_xa_x\widetilde T_x}.
\]
This is precisely a Wilson law with plaquette couplings
$\beta_p=\beta(1-\alpha_p)$, where
\[
 \alpha_{(y;\mu,\nu)}
  =\tfrac14\sum_{\substack{\varepsilon\in\{0,1\}^4\\
                 \varepsilon_\mu=\varepsilon_\nu=0}}a_{y-\varepsilon},
 \qquad \norm{\alpha}_\infty\le\norm a_\infty,
 \qquad \sum_p\alpha_p^2\le6\sum_xa_x^2.
\]
For $p>1$ with $\norm{pa}_\infty\le R_M$, convexity of $K$ gives
\begin{equation}
 D(\mu_a\Vert\mu)=a\cdot\nabla K(a)-K(a)
       \le\frac{k_M}{p-1}\sum_xJ(pa_x).
 \label{eq:Wilson-tilt-entropy}
\end{equation}
Indeed $K(pa)\ge K(a)+(p-1)a\cdot\nabla K(a)$ and $K(a)\ge0$.
This supplies an actual nonperturbative source comparison on the
selected bare law. It does not assert that a general renormalization
step has a small entropy budget, preserves this selected sequence,
or satisfies a temporal mixing estimate.
\sourcebasis{The full compact pressure, reflection-based cell response,
and selected source window are developed in \src{F14}{V95--V100}.
The entropy transport and actual plaquette-coupling interpretation are
proved in \src{F15}{V1--V2}.}


\section{Generated compact laws and retained source information}
\label{sec:generated-source-laws}
The entropy budget becomes useful for blocking only after identifying the
actual output likelihood. Let $\mathcal B:\Omega_M\to\mathcal C$ be a
fixed measurable compact block map, with $\mathcal C$ a standard Borel
space. It may be a finite composition of the completed maps of
\cref{ch:compact-response}. The map and its reference coordinates are
held fixed as the source varies. Set
\[
 \rho=\mathcal B_*\mu,\qquad \rho_a=\mathcal B_*\mu_a,\qquad
 F_a(C)=\log\E_\mu[e^{X(a)}\mid C],\qquad
 A(a)(C)=\E_\mu[X(a)\mid C].
\]
Here $X(a)$, $K(a)$, and $\mu_a$ have the normalization of
\cref{sec:entropy-source-transport}. In particular, $X(a)$ is centered
under the baseline fine law. All conditional integrations use that law,
including its solved-pivot densities and generated interactions.

\begin{theorem}[Exact generated likelihood and information split]
\label{thm:generated-likelihood}
At each finite regulator,
\begin{align}
 \frac{\dd\rho_a}{\dd\rho}(C)&=L_a(C):=e^{F_a(C)-K(a)},\notag\\
 \frac{\dd\mu_a(\cdot\mid C)}{\dd\mu(\cdot\mid C)}(U)
       &=e^{X(a)(U)-F_a(C)},\label{eq:generated-likelihood}\\
 D(\mu_a\Vert\mu)&=D(\rho_a\Vert\rho)
  +\int D\bigl(\mu_a(\cdot\mid C)\Vert\mu(\cdot\mid C)\bigr)
                                               \dd\rho_a(C).
 \label{eq:generated-entropy-split}
\end{align}
Consequently, if $p>1$ and $\norm{pa}_\infty\le R_M$, then
\begin{align}
 D(\rho_a\Vert\rho)&\le
       \frac{k_M}{p-1}\sum_xJ(pa_x),\notag\\
 \E_\rho L_a^p&\le e^{K(pa)-pK(a)}
                    \le\exp\left\{k_M\sum_xJ(pa_x)\right\}.
 \label{eq:generated-source-budgets}
\end{align}
No locality of $F_a$ in the retained variables is asserted.
\end{theorem}
\begin{proof}
For bounded $\varphi$, disintegration gives
\[
 \E_{\mu_a}\varphi(\mathcal B(U))
 =e^{-K(a)}\E_\mu[\varphi(\mathcal B(U))e^{X(a)}]
 =\int\varphi(C)e^{F_a(C)-K(a)}\dd\rho(C).
\]
This proves the first identity; normalization inside each fibre proves
the second. Split the fine logarithmic likelihood into
$X(a)-K(a)=(F_a-K(a))+(X(a)-F_a)$ and integrate under $\mu_a$.
This gives \cref{eq:generated-entropy-split}, with the conditional
entropy averaged under the changed output law, not the baseline law.
Its nonnegativity and \cref{eq:Wilson-tilt-entropy} give the first
budget. Since $L_a(\mathcal B(U))=\E_\mu[e^{X(a)-K(a)}\mid C]$,
conditional Jensen for the power $p$ gives the second. The last
inequality uses $K(a)\ge0$ and
\cref{thm:full-Wilson-source-window}. All likelihoods are positive
and finite because the fine source is bounded at fixed regulator.
\src{F15}{V2}
\end{proof}

The effective action relative to the baseline output law changes by
$-F_a+K(a)$. This is the normalized version of
\cref{thm:exact-block-differential}; it is not obtained by replacing
fine plaquettes with coarse plaquettes. For nested maps
$C^{(j+1)}=\mathcal B_{j+1}(C^{(j)})$, one has exactly
\[
 L_{j+1,a}=\E_{\rho_j}[L_{j,a}\mid C^{(j+1)}].
\]
Thus the entropy and $L^p$ budgets are nonincreasing under further
observation. Resetting $\rho_j$ to product Haar would change this
identity and require a separate comparison.

\begin{proposition}[Retained Fisher information and a noncollapse test]
\label{prop:retained-source-information}
Write $A_x(C)=\E_\mu[\widetilde T_x\mid C]$. At zero source,
\begin{align}
 \partial_{a_x}\log L_a\big|_0&=A_x,\notag\\
 \partial_{a_x}\partial_{a_y}\log L_a\big|_0
   &=\Cov_\mu(T_x,T_y\mid C)-\Cov_\mu(T_x,T_y).
 \label{eq:retained-source-derivatives}
\end{align}
Moreover,
\begin{align}
 \log\E_\rho e^{A(b)}&\le K(b)
                \le k_M\sum_xJ(b_x),\qquad \norm b_\infty\le R_M,
                  \notag\\
 0\preceq\Cov_\rho(A)&\preceq\Cov_\mu(T)
                         \preceq\chi_{M,\beta}I.
 \label{eq:retained-source-information}
\end{align}
For any square-integrable retained observable $Y=Y(C)$ with positive
variance and any fine source $X(b)$,
\begin{equation}
 \Var_\rho(A(b))\ge
       \frac{\Cov_\mu(X(b),Y(\mathcal B(U)))^2}{\Var_\rho(Y)}.
 \label{eq:retained-source-regression}
\end{equation}
In particular, a nonzero limiting mixed covariance and a finite
positive limiting retained variance certify survival of that score.
\end{proposition}
\begin{proof}
Differentiate the bounded conditional moment-generating function and
subtract the derivatives of $K$. Conditional Jensen gives the
exponential bound. The exact total-covariance identity is
\[
 \Cov_\mu(T)=\E_\rho\Cov_\mu(T\mid C)+\Cov_\rho(A).
\]
Both summands are positive semidefinite; apply
\cref{prop:full-Wilson-covariance}. Finally conditional expectation
preserves covariance with $Y(C)$, so Cauchy--Schwarz under $\rho$
gives \cref{eq:retained-source-regression}. The same argument permits a
growing fine source provided its mixed-moment errors have already
been summed over its full support. \src{F15}{V2, V4, V7}
\end{proof}

These derivatives are with respect to the original fine source
parameters, not retained link coordinates. The upper information bound
does not imply noncollapse, and a positive lower bound for one growing
score need not be a positive fraction of its full fine variance. The
regression test supplies a static retained-source criterion; the local
reflected noncollapse required in \cref{hyp:noncollapse} additionally
needs the physical source and time identification.

\begin{proposition}[Measured error of linearizing the generated source]
\label{prop:generated-source-linearization}
Define
\[
 \mathcal R_a=F_a-A(a),\qquad
 K_{\rm lin}(a)=\log\E_\rho e^{A(a)},\qquad
 \dd\widehat\rho_a=e^{A(a)-K_{\rm lin}(a)}\dd\rho.
\]
Then $\mathcal R_a\ge0$, $0\le K_{\rm lin}(a)\le K(a)$, and
\begin{align}
 \E_\rho\mathcal R_a&\le K(a),\notag\\
 \norm{\rho_a-\widehat\rho_a}_{\rm TV}
   &\le1-e^{K_{\rm lin}(a)-K(a)}
                     \le K(a)-K_{\rm lin}(a),\notag\\
 D(\widehat\rho_a\Vert\rho_a)
   &=K(a)-K_{\rm lin}(a)-\E_{\widehat\rho_a}\mathcal R_a
                     \le K(a)-K_{\rm lin}(a),\notag\\
 \E_\rho e^{\mathcal R_a/2}
   &\le e^{(K(a)+K(-a))/2}.
 \label{eq:generated-source-linearization}
\end{align}
Total variation is half the $L^1$ distance of densities. The order of
the two laws in the displayed relative entropy is part of the result.
\end{proposition}
\begin{proof}
Conditional Jensen gives $F_a\ge A(a)$. Since $\E_\rho A(a)=0$,
ordinary Jensen and $\E_\rho e^{F_a}=e^{K(a)}$ give the stated bounds
on $K_{\rm lin}$ and $\E\mathcal R_a$. For nonnegative $f\ge g$
with masses $m\ge n>0$,
$\tfrac12\int|f/m-g/n|\le1-n/m$; apply this with
$f=e^{F_a}$ and $g=e^{A(a)}$. The logarithmic ratio of the two
normalized laws is $K-K_{\rm lin}-\mathcal R_a$, giving the entropy
identity. Finally Cauchy--Schwarz gives
\[
 \E_\rho e^{(F_a-A(a))/2}
 \le(\E_\rho e^{F_a})^{1/2}(\E_\rho e^{-A(a)})^{1/2}
 \le e^{(K(a)+K(-a))/2}.
\]
\src{F15}{V2}
\end{proof}

For sources uniformly inside the admitted amplitude window,
$\norm a_2\to0$ makes this a uniform $O(\norm a_2^2)$ approximation
of the generated probability law, independently of the blocking map.
It does not make $\mathcal R_a(C)$ pointwise small or spatially local.
Nor does conditional control alone pay for the output normalizer:
\cref{eq:generated-entropy-split} retains its separate entropy cost.
These distinctions are essential when the full-law estimates are used
to populate the nonlinear transport assumptions of \cref{ch:closure}.

\chapter{Nonlinear remainder protection and sector discipline}
\label{ch:remainder}
The local compact construction of \cref{ch:compact-response} supplies
an actual action and density to which this bookkeeping can be applied.
Its measured two-loop expansion has an $O(\beta^{-2})$ remainder
through four retained-source derivatives under the strengthened local
hypotheses of \cref{thm:measured-two-loop}. This remains a restricted-box
estimate and does not replace the source-dependent owner stock or its
cofinal propagation assumptions below. For a fourth-order measured
output, the finite incoming jet packet is
$(J^8S,J^6Q,J^4\mathcal B;J^7F)$, where $F$ is the constraint map and
$Q,\mathcal B$ are the incoming measured loop coefficients. The
coarea density is counted once. This finite derivative bookkeeping
from \src{SM}{V8} states what must be transported, not a proof of
uniform transport on the entire compact family. Localized well removal in \cref{prop:well-removal}
is likewise an explicitly compared finite contribution, not a deletion
of the entire bad-field class.

\section{An exhaustive owner ledger}
The remainder analysis begins with the one-scale compact analytic and ownership estimates. At a finite regulator, fixed retained zero-theta sector, and exact triangular elimination, primitives carry birth tags for regular local, bad, reset, or genuine boundary contributions. The Gaussian principal shell and extracted finite jets are outside this remainder alphabet.

A connected word is classified by the first applicable alternative
\[
\text{boundary}\succ\text{bad}\succ\text{reset}\succ
\text{regular multi-vertex}\succ\text{regular single-vertex}.
\]
This is an exhaustive disjoint partition, not an additional payment rule. Historical support and ownership are unchanged. Under normal convergence on the common smaller source polydisc, the normalized marginal response therefore has the exact identity
\begin{equation}\label{eq:owner-ledger}
\mathcal R_{\rm exact}=\mathcal R_{\rm main}
+E_{\rm vert}+E_{\rm polymer}+E_{\rm bad}
+E_{\rm reset}+E_{\rm multiscale}.
\end{equation}
The last term includes the genuine boundary class and the endpoint of propagated stable history. Derivatives pass through the series only after the normal-convergence hypothesis is supplied. Identity \cref{eq:owner-ledger} is not by itself a smallness estimate. \src{PAR}{V1}

\section{Power-preserving nonlinear control}
The common primitive stock is built from the two-jet-subtracted regular density, the normalized kernel defect, saturated exits, and compensation terms. For $u=\beta^{-1}$, the owner expansion gives, on its declared common domain,
\begin{equation}\label{eq:owner-power}
w(u)\le C_wu^{6/5}.
\end{equation}
The stretched exponentials are converted to this common power by the elementary maximum of $\beta^Ae^{-c\beta^d}$. Radius restriction gives, for the high-grade local density,
\[
\norm{D^kV_{\rm hi}}_r\le\frac{k!C_{\rm loc}}{(R-r)^k}u^{6/5},
\qquad k=0,1,2,
\]
with the zeroth-order formula interpreted on its original radius. Source derivatives have their corresponding Cauchy-radius losses. The marked connected-polymer and bad-field bounds use the same $w$, and the stable-history propagation preserves the $6/5$ power under the imported cone hypothesis. \src{PAR}{V2--V3}

\begin{criterion}[Physical marginal remainder and sign protection]\label{thm:remainder}
Assume the common-regulator owner, analytic-domain, fixed-sector, and stable-history hypotheses of the owner expansion. Suppose
\[
\beta_{j+1}=\beta_j+c_j,\qquad c_j=b_j+e_j,
\qquad |e_j|\le C_{\rm NP}u_j^{6/5},\quad |c_j|\le M_c,
\]
and $M_cu_j\le1/2$. Then
\begin{equation}\label{eq:physical-remainder}
u_{j+1}=u_j-b_ju_j^2+r_j,\qquad
|r_j|\le C_{\rm NP}u_j^{16/5}+2M_c^2u_j^3.
\end{equation}
If also $|b_j|\ge m_{\rm main}>0$, then at sufficiently small common $u_j$,
\begin{equation}\label{eq:sign-protection}
\frac{|r_j|}{|b_j|u_j^2}
\le\frac{C_{\rm NP}}{m_{\rm main}}u_j^{6/5}
+\frac{2M_c^2}{m_{\rm main}}u_j\le\frac12.
\end{equation}
Thus the exact real increment has the sign of $-b_ju_j^2$.
\end{criterion}
\begin{proof}
Use the exact identity
\[
\frac{u}{1+uc}=u-cu^2+\frac{c^2u^3}{1+uc}.
\]
Subtract the main increment $-bu^2$, bound $e$ and the denominator, and divide by $|b|u^2$. The common smallness threshold also retains every inherited owner and analytic-domain condition. \src{PAR}{V4}
\end{proof}

This is a protection theorem for a supplied principal margin. It does
not compute $b_j$, prove its sign, or identify its limiting value with
$b_{\rm YM}$. The principal coefficient must be that of the same calibrated
exact response. A margin measured with fixed-background Wilson banks
cannot be inserted as the physical principal margin before the background
and complete-determinant comparisons of \cref{sec:background-matching}
are established. Those discrepancies are not silently absorbed into the
$O(u^{6/5})$ nonlinear remainder. The principal-response calculation and the
remainder theorem are complementary, not substitutes.

\section{Zero theta, compact labels, and the correct thermodynamic target}
The exact compact construction retains its sector labels and the declared zero-theta section through the chosen transport. Reflection positivity is not used as a general argument that every possible theta parameter must vanish. Gauge quotient labels, commuting torons, and global charge labels require their own treatment.

A proposed volume-uniform tail for total sector complexity encounters a precise tensorization obstruction. A local nonzero charge costs finite action and finite one-component free energy, while its number of possible placements grows with volume. A global total-charge bound cannot be inferred from a local defect penalty.

\begin{theorem}[Tensorization obstruction to global sector tightness]\label{thm:sector-no-go}
Let $X$ be the bounded integer charge on a fixed finite component, and suppose $X$ is not almost surely zero. On $N$ disjoint copies, let $X_i$ be independent copies and set
\[
\Psi_N=\sum_{i=1}^N|X_i|,\qquad Q_N=\sum_{i=1}^NX_i.
\]
For every fixed $R<\infty$,
\begin{equation}\label{eq:sector-escape}
\Pr(\Psi_N\ge R)\longrightarrow1,\qquad
\Pr(|Q_N|\ge R)\longrightarrow1.
\end{equation}
Consequently no bound $Ce^{-cR}$ with positive $c$ and finite $C$ independent of $N$ controls either tail for all $N,R$.
\end{theorem}
\begin{proof}
If $p=\Pr(|X|\ge1)>0$, then $\Psi_N$ dominates a sum of independent Bernoulli variables with mean $Np$. If $\E X\ne0$, the law of large numbers gives escape of $Q_N$. If $\E X=0$, nontriviality gives positive variance and the central limit theorem implies that its probability in any fixed bounded interval tends to zero. A nonzero constant charge is covered by the nonzero-mean case. Choosing $R$ with $Ce^{-cR}<1/2$ contradicts the limiting probabilities. \src{PAR}{V5}
\end{proof}

The theorem applies exactly to the locality-compatible class allowing disjoint unions. It does not by itself prove anti-concentration on a sequence restricted to connected periodic tori; that would require a positive susceptibility or a suitable insertion/mixing estimate. Fixed physical volume can also change the target. The retained alternatives are local-window or charge-density bounds, explicit phase or sector decompositions, and exact compact transport. The global extensive tail is not an unpaid theorem that one should continue trying to obtain from the same insufficient energy--entropy mechanism.

\section{Refinement coherence is retained independently}
Let $P_j$ be the field block map, $q_j$ the label map, and $r_j$ the label refinement. Suppose $r_jq_{j+1}=q_jP_j$ outside an event of probability $\epsilon_j$, and
\[
d_j=\norm{(P_j)_*\mu_{j+1}-\mu_j}_{\rm TV}.
\]
Then data processing gives
\begin{equation}\label{eq:label-TV}
\norm{(r_j)_*(q_{j+1})_*\mu_{j+1}-(q_j)_*\mu_j}_{\rm TV}
\le\epsilon_j+d_j.
\end{equation}
Summable defects yield a projectively Cauchy family at every fixed level. This coherence estimate survives the tightness obstruction. Coherence and tightness are different: a coherent family can still place its mass on escaping labels. \src{PAR}{V5}
\sourcebasis{The exact owner partition, power bounds, stable-history control, and physical remainder are \src{PAR}{V1--V4}; the topology and refinement results are \src{PAR}{V5}.}

\chapter{Compact spectral geometry and local interface repair}
\label{ch:compact-repair}
The compact completion permits an auxiliary spectral analysis on the
actual fibre, without replacing it by a Gaussian or an independent
product reference. Two features must be treated together. A normally
coercive symmetry orbit supports a quantitative local Poincar\'e
inequality, but a second, higher-energy orbit creates metastability
for the full hidden diffusion. Local visible repair can nevertheless
be proved before integrating the hidden coordinates and transferred
to the original interface heat bath. This chapter develops that route
and its remaining local-to-global requirement. Every heat-bath rate
below is an auxiliary coordinate-update rate, not the physical clock
of \cref{ch:spectral}.

\section{The actual finite parent and its stationary orbits}
\label{sec:compact-orbit-data}
Use the $\SU(2)$ completion of
\cref{sec:su2-pivot-equation}, with $\rho=1/64$ and action convention
$s_p=2-\Tr Q_p$. On the fine torus of side six, fix the coarse field
$C=I$ and the parent of coarse cells $P=\{0,1\}^4$. Its independent
hidden carrier is
\[
 M=\SU(2)^{944},\qquad 944=272+672,\qquad \dim M=2832.
\]
The $272$ internal non-tree groups and $672$ nonpivot boundary groups
are independent; the $96$ pivots are smooth functions of them. Fix
the central exterior sheet specified by
\begin{equation}\label{eq:central-sheet}
 \epsilon_{(x,\mu)}=
 \begin{cases}
 -1,&\mu=2,\ x_2=1,\text{ and some }x_a\text{ is odd},\ a\in\{0,1,3\},\\
 +1,&\text{otherwise}.
 \end{cases}
\end{equation}
The parent law, including its inherited additive action normalization,
is
\begin{equation}\label{eq:actual-parent-law}
 \dd\mu_\beta=Z_\beta^{-1}a(Y)e^{-\beta S(Y)}\dd\operatorname{vol}_M(Y).
\end{equation}
Here $a$ is the full positive inverse-pivot amplitude, independent of
$\beta$, and the metric is product Hilbert--Schmidt.

\begin{certificate}[Two stationary orbits]
\label{def:two-orbit-input}
For this fixed parent and exterior, the finite certificates locate two
full stationary simultaneous-conjugation orbits
$\mathcal O_1,\mathcal O_2\simeq S^2$. Their Hessian kernels are
exactly the two orbit directions. On the $2830$-dimensional normal
spaces the respective lower bounds are $1/5000$ and $31/125000$ in
the Hilbert--Schmidt metric. Their energies satisfy
\[
 \frac{113936882369}{200000000}\le E_1
          \le\frac{569684411883}{10^9},\qquad
 569.68444636<E_2<569.68444638,
\]
and hence $E_2-E_1>34477/10^9>0$.
\end{certificate}
The first orbit and its full normal Hessian are reported in
\src{F12}{V89--V90}; the second, including the transition-source pivot
acceleration, is reported in \src{F13}{V8}. The associated interval
and integer factor certificates are not independently rerun here.
The analytic results below use these explicitly identified finite
inputs. No classification of the global ground set is assumed.
In particular, the second orbit is not globally minimizing, whereas
global minimality of the first is not established by these data.

\section{An equivariant tube inequality with the actual amplitude}
\label{sec:actual-orbit-tube}
Choose $p\in\mathcal O_1$ and let $K\simeq\mathrm U(1)$ be its
stabilizer. The normal bundle is $G\times_KN_p\mathcal O_1$, not an
assumed globally trivial product. For a sufficiently small fixed
normal ball $B_\varepsilon\subset N_p\mathcal O_1$, the map
\[
 F(g,z)=g\cdot\exp_pz
\]
parametrizes the orbit tube $\mathcal T_\varepsilon$ through that
associated bundle. Invariance of the action, amplitude, and metric
gives a smooth positive $K$-invariant function $b$ such that
\begin{equation}\label{eq:tube-law-pushforward}
 F_*(\dd g\otimes\dd\nu_\beta)
      =\mu_\beta(\,\cdot\mid\mathcal T_\varepsilon),\qquad
 \dd\nu_\beta(z)=
 \frac{b(z)e^{-\beta s(z)}\dd z}
      {\int_{B_\varepsilon}b(w)e^{-\beta s(w)}\dd w},
 \quad s(z)=S(\exp_pz).
\end{equation}
Indeed $b$ is the original amplitude times the normal-exponential
volume Jacobian. Haar measure pushes forward to normalized orbit
volume, while equivariance transports the same normal density around
the orbit. This retains the Jacobian on the whole tube; it does not
replace it by its value at $p$.

\begin{theorem}[Local hidden diffusion inequality]
\label{thm:actual-tube-gap}
For a fixed sufficiently small $\varepsilon>0$ and a finite threshold
$\beta_0$, every $\beta\ge\beta_0$ and every weighted $H^1$ function
on the tube satisfy
\begin{equation}\label{eq:actual-tube-gap}
 \Var_{\mu_\beta(\cdot\mid\mathcal T_\varepsilon)}f
 \le\left(3776+\frac{80000}{\beta}\right)
     \int_{\mathcal T_\varepsilon}|\nabla_Mf|^2
                              \dd\mu_\beta(\cdot\mid\mathcal T_\varepsilon).
\end{equation}
For conjugation-invariant $f$, the constant $3776$ can be omitted.
The local Neumann diffusion gap is therefore at least $1/83776$ for
$\beta\ge\max\{1,\beta_0\}$.
\end{theorem}
\begin{proof}
Normal positivity in \cref{def:two-orbit-input} and continuity permit
$D^2s\succeq I/10000$ and $\norm{D\exp_p|_N}\le2$ on the fixed
ball. With $K_b=\sup\norm{D^2\log b}$, take
$\beta_0\ge20000K_b$. Then
$D^2(\beta s-\log b)\succeq(\beta/20000)I$.
The weighted Bochner identity on a convex Euclidean ball, including
the nonnegative second-fundamental-form boundary term, gives the
Neumann Poincar\'e constant $20000/\beta$ for $\nu_\beta$.

Haar measure on $\SU(2)$ in this metric has a sufficient Poincar\'e
constant one: the group is the round three-sphere of radius $\sqrt2$
and its Ricci tensor equals the metric, so the integrated Bochner
identity yields this bound. For $h=f\circ F$,
\[
 |D_zh|^2\le4|\nabla_Mf|^2\circ F,\qquad
 |\nabla_Gh|^2\le4\cdot944|\nabla_Mf|^2\circ F.
\]
The latter uses $\norm{[Z,U]}_{\rm HS}\le2\norm Z_{\rm HS}$
in each of the $944$ coordinates; it is an operator-norm bound and
has no additional factor $\dim G$. Variance decomposition in the
product law of \cref{eq:tube-law-pushforward}, the two Poincar\'e
bounds, and Jensen's inequality for the differentiated Haar average
give \cref{eq:actual-tube-gap}. For invariant $f$ the group derivative
vanishes. Density of smooth functions in the weighted form domain
completes the proof. \src{F13}{V1}
\end{proof}
The tube radius, amplitude derivative bound, and onset $\beta_0$ are
existential constants, not numerically enclosed parameters. The
inequality is a local Neumann statement, with no boundary condition
imposed on its test functions. It is distinct from a killed escape
inequality and from the physical orbit criterion in \cref{thm:orbit}.

\section{A secondary well excludes a uniform hidden diffusion gap}
\label{sec:hidden-metastability}
Let $\lambda_\beta$ be the gap of the full hidden diffusion with
Dirichlet form
$\mathcal E_\beta(f)=\int_M|\nabla f|^2\dd\mu_\beta$.

\begin{theorem}[Metastability of the actual fixed-parent diffusion]
\label{thm:hidden-metastability}
For the law in \cref{eq:actual-parent-law} and the finite inputs in
\cref{def:two-orbit-input}, constants $C<\infty$, $\delta>0$, and
$\beta_1<\infty$ satisfy
\begin{equation}\label{eq:hidden-metastability}
 0<\lambda_\beta\le C\beta^{1415}e^{-\beta\delta}
                   \qquad(\beta\ge\beta_1).
\end{equation}
The same upper bound holds for the lowest Dirichlet eigenvalue on the
complement of any sufficiently small closed tube about $\mathcal O_1$
which is disjoint from a neighbourhood of $\mathcal O_2$.
\end{theorem}
\begin{proof}
Choose a small normal tube around $\mathcal O_2$, disjoint from
$\mathcal O_1$. Normal positivity gives
$E_2+c_-|z|^2\le S\le E_2+c_+|z|^2$ in finitely many tubular
charts. Let $f$ equal one on its radius-$r$ inner tube and vanish
outside the radius-$2r$ tube. Its gradient is supported where
$S\ge E_2+\delta$, with $\delta=c_-r^2>0$. Smooth positive amplitude
and volume bounds give
\[
 \mathcal E_\beta(f)\le Z_\beta^{-1}C e^{-\beta(E_2+\delta)},\qquad
 \int f^2\dd\mu_\beta\ge Z_\beta^{-1}c\beta^{-1415}e^{-\beta E_2}.
\]
The exponent is half the normal dimension $2830$. The lower-energy
orbit independently gives
$Z_\beta\ge c'\beta^{-1415}e^{-\beta E_1}$. Thus
$\mu_\beta(\supp f)\le C'\beta^{1415}e^{-\beta(E_2-E_1)}\to0$.
Cauchy--Schwarz implies
\[
 \Var_{\mu_\beta}f\ge(1-\mu_\beta(\supp f))\int f^2\dd\mu_\beta.
\]
Its Rayleigh quotient proves the upper bound. The density is positive
and smooth on a compact connected manifold at every finite $\beta$,
so its elliptic diffusion gap is positive. The same supported test
function belongs to the Dirichlet domain outside the first tube;
its uncentered Rayleigh quotient proves the killed bound.
\src{F13}{V8}
\end{proof}
The barrier $\delta$ is not the energy difference $E_2-E_1$.
Exponential smallness of the higher well's equilibrium mass does not
remove its small Rayleigh quotient. Covering both wells by separate
local tubes therefore cannot restore a temperature-uniform gap for
this same diffusion. This obstruction does not assert a corresponding
upper bound for a nonlocal heat bath and does not concern the physical
Yang--Mills Hamiltonian. It removes one sufficient auxiliary route,
not the continuum target.

\subsection{Localized marginal removal does not require hidden mixing}
\label{sec:well-removal}
There is a useful simultaneous positive conclusion. Let a fixed smooth
cutoff $0\le\chi\le1$ be supported in the second orbit tube and vanish
near the first. For coarse parameters $q$ near the declared background,
write $N_\beta(q)$ for the integral with factor $\chi$,
$p_\beta=N_\beta/Z_\beta$, and
$\widetilde Z_\beta=Z_\beta-N_\beta=Z_\beta(1-p_\beta)$.
This is a specified comparison integral, not a redefinition of the
original law.

\begin{proposition}[A local normalized well-comparison estimate]
\label{prop:well-removal}
On a sufficiently small fixed coarse neighbourhood $U$, and for every
fixed integer $k\ge0$,
\begin{align}
 0\le p_\beta(q)&\le C\beta^{1415}e^{-\sigma\beta},
                  &\sigma&=1/50000,\label{eq:well-mass}\\
 \norm{\log Z_\beta-\log\widetilde Z_\beta}_{C^k(\overline U)}
        &\le C_k\beta^{1415+k}e^{-\sigma\beta}.\label{eq:well-derivatives}
\end{align}
No hidden Poincar\'e inequality is required.
\end{proposition}
\begin{proof}
Shrink $U$ so $|S(q,x)-S(q_0,x)|\le\gamma=1/200000$ uniformly
in the compact hidden carrier. The numerator is at most
$Ce^{-\beta(E_2-\gamma)}$, while the first orbit supplies a denominator
at least $c\beta^{-1415}e^{-\beta(E_1+\gamma)}$.
The energy intervals give $E_2-E_1-2\gamma>\sigma$, proving
\cref{eq:well-mass}. Derivatives of the actual integrand divided by
that integrand are bounded by $C_\alpha(\beta+1)^{|\alpha|}$.
The fixed nonnegative cutoff contributes no parameter derivatives.
Induction in the differentiated identity $Z_\beta p_\beta=N_\beta$
gives
$|\partial^\alpha p_\beta|\le C_\alpha(\beta+1)^{|\alpha|}p_\beta$.
Differentiate $\log(1-p_\beta)$ and use $p_\beta\le1/2$ for large
$\beta$. This proves \cref{eq:well-derivatives}, including the
normalization derivatives. \src{F13}{V9}
\end{proof}
For normalized retained laws with the same direct retained weight and
partition factors $Z_\beta$ and $\widetilde Z_\beta$ on a compact
domain in $U$, the density-ratio oscillation is at most
$(1-\varepsilon_\beta)^{-1}$, where
$\varepsilon_\beta=\sup_U p_\beta$. Their same-metric diffusion gaps,
same-block heat-bath gaps, and killed floors on the same event therefore
obey the two-sided comparison factors
$1-\varepsilon_\beta$ and $(1-\varepsilon_\beta)^{-1}$.
For heat baths this follows directly from
\[
 \E_\nu\Var(f\mid q_{B^c})
    =\inf_{g\ \mathscr F_{B^c}\text{-measurable}}\int(f-g)^2\dd\nu,
\]
so the compared update blocks and rates are not changed. This
localized removal estimate does not classify the remaining wells or
supply an unrestricted retained gap.

\section{A same-law diffusion-to-heat-bath comparison}
\label{sec:global-resolvent}
Let $\mu$ be a positive smooth density proportional to $e^{-W}$
relative to product Haar in a finite volume. Write
$L=\sum_iL_i\succeq0$ for its reversible product-metric diffusion,
$\Gamma_i$ for the conditional expectation given all coordinates
except group $i$, and
\[
 \mathsf H_\mu=\sum_i(I-\Gamma_i),\qquad
 \mathcal E_L(u)=\ip u{Lu},\qquad
 \mathcal E_{\mathsf H_\mu}(f)=\sum_i\norm{(I-\Gamma_i)f}_2^2.
\]
The symbol $\mathsf H_\mu$ denotes the rate-one coordinate heat bath;
it is not the reconstructed physical Hamiltonian $H$.
Assume the intrinsic mixed Hessian row bound
\[
 \sup_i\sum_{j\ne i}\sup\norm{W_{ij}^{\rm Hess}}_{\op}\le K.
\]

\begin{theorem}[Global resolvent domination]
\label{thm:global-resolvent}
On the same $L^2(\mu)$ space,
\begin{equation}\label{eq:global-resolvent}
 \mathsf H_\mu\succeq L(L+K)^{-1}.
\end{equation}
The right side is defined by spectral calculus, with value zero on
$\Ker L$; if $K=0$ it is the projection onto $(\Ker L)^\perp$.
In particular a diffusion gap $g>0$ implies
$\gap\mathsf H_\mu\ge g/(g+K)$.
\end{theorem}
\begin{proof}
The product-manifold and one-coordinate weighted Bochner identities
give
\[
 \sum_i\norm{L_iu}_2^2\le\norm{Lu}_2^2+K\mathcal E_L(u).
\]
All diagonal curvature and Ricci terms cancel on subtraction; the
remaining mixed Hessian squares are nonnegative, and the mixed
potential terms are bounded by the row sum and
$ab\le(a^2+b^2)/2$. No commutation of the $L_i$ is assumed.
Since $\Gamma_iL_iu=0$,
\[
 \ip f{Lu}=\sum_i\ip{(I-\Gamma_i)f}{L_i u}.
\]
Completing a square in this direct sum yields, for every smooth $u$,
\[
 2\operatorname{Re}\ip f{Lu}-\norm{Lu}_2^2-K\mathcal E_L(u)
             \le\mathcal E_{\mathsf H_\mu}(f).
\]
Taking the supremum over $u$ by the spectral theorem gives
\cref{eq:global-resolvent}. Closure of the forms extends the estimate
to its full domain. The scalar gap consequence follows by restricting
to the centered subspace. \src{F12}{V60, V77--V80}
\end{proof}
For the unintegrated completed fibre,
$W_\beta=\beta S-\log j$ is a sum of fixed smooth finite-stencil
terms. Bounded incidence and compact input derivatives give
\begin{equation}\label{eq:preintegrated-row}
 K_\beta\le C_{\rm f}(\beta+1)
\end{equation}
uniformly in the ambient volume and compact coarse data. After hidden
integration, the score covariance in \cref{eq:compact-score} can
instead produce a quadratic-in-$\beta$ Hessian bound. Keeping
\cref{eq:global-resolvent} before integration avoids this unnecessary
loss; it does not posit a uniform hidden diffusion gap.

\begin{lemma}[Exact domination of marginal heat-bath forms]
\label{lem:marginal-form-domination}
Let $\mu(\dd x,\dd v)$ be the joint fine law and $\nu(\dd v)$ its
exact retained interface marginal, using the same independent group
partition. For $Jf(x,v)=f(v)$,
\[
 \norm{Jf}_{L^2(\mu)}=\norm f_{L^2(\nu)},\qquad
 \mathcal E_{\mathsf H_\nu}(f)\ge
                         \mathcal E_{\mathsf H_\mu}(Jf).
\]
\end{lemma}
\begin{proof}
Hidden-group updates leave $Jf$ unchanged. For a retained group,
conditioning additionally on the hidden variables reduces the expected
conditional variance, by variance decomposition. Summing the retained
coordinate inequalities proves the claim.
\end{proof}
To transfer a killed estimate by this lemma, the event must constrain
visible coordinates only. Its lift must contain every hidden field;
a hidden force-region restriction would introduce a further posterior
coverage problem.

\section{From visible negative curvature to an original-clock repair floor}
\label{sec:repair-transfer}
For an operator $A\succeq0$ and a visible event $D$, write
\[
 \kappa_D^A=\inf_{0\ne f=\one_Df}
                 \frac{\ip f{Af}}{\norm f_2^2}
\]
for its killed floor, with the infimum taken in the appropriate form
domain. It is not a reflected or Neumann gap on $D$.

\begin{lemma}[Spectral localization in a visible cylinder]
\label{lem:repair-localization}
Suppose the same global diffusion satisfies, for every form vector $u$,
\begin{equation}\label{eq:repair-localization}
 \norm{\one_Du}_2^2\le\frac2\lambda\mathcal E_L(u)
                           +\frac{2M_\chi}{\lambda}\norm u_2^2,
 \qquad \lambda\ge8M_\chi.
\end{equation}
Together with \cref{eq:global-resolvent}, this gives
\begin{equation}\label{eq:repair-floor-general}
 \kappa_D^{\mathsf H_\mu}\ge
                 \frac{\lambda}{2(\lambda+8K)}.
\end{equation}
The identical lower bound transfers to the exact marginal heat bath
when $D$ is a cylinder restricting only retained coordinates.
\end{lemma}
\begin{proof}
For the global diffusion spectral projection
$\Pi_t=\one_{[0,t]}(L)$, \cref{eq:repair-localization} gives
$\norm{\one_D\Pi_t}^2\le2(t+M_\chi)/\lambda$.
Taking adjoints gives the same estimate for $\Pi_t\one_D$; no
commutation is asserted. At $t=\lambda/8$ the bound is at most $1/2$,
so a vector supported in $D$ has at least half its squared norm above
$t$. The global resolvent bound gives
$\ip f{\mathsf H_\mu f}\ge\lambda\norm f^2/[2(\lambda+8K)]$.
Use \cref{lem:marginal-form-domination} for the last assertion.
\end{proof}
The resolvent in this proof is always that of the global diffusion for
the same law. A Dirichlet resolvent is not substituted for it.

Here is the local geometric input to \cref{eq:repair-localization}.
In product normal coordinates on a finite set of visible groups,
leave all other groups uncharted. If a constant unit chart direction
$a$ obeys
$\partial_a^2\Phi_\beta\le-\eta\beta$ for the actual measured
potential, uniformly in all uncharted inputs, and the inverse metric
is at least $c_gI$, then ground-state conjugation gives
\[
 \mathcal E_L(h)\ge\frac{c_g\eta\beta}{2}\norm h^2
\]
for functions supported in that chart cylinder. Indeed, with
$u=h e^{-\Phi_\beta/2}$, the one-direction potential term is
$|\partial_a\Phi_\beta|^2/4-\partial_a^2\Phi_\beta/2$.
A cutoff $\chi$ equal to one on a smaller cylinder yields
\cref{eq:repair-localization} with
$\lambda=c_g\eta\beta/2$ and $M_\chi=\sup|\nabla\chi|^2$.
For a fixed finite cover, normalized cutoffs with
$\sum_j\chi_j^2\le1$, equal to one on the event, give the same estimate
with $M_\chi=\sup\sum_j|\nabla\chi_j|^2$.

Thus any such fixed visible certificate with $\lambda=c_0\beta$ and
\cref{eq:preintegrated-row} gives, for sufficiently large $\beta$,
\begin{equation}\label{eq:uniform-repair-template}
 \kappa_D^{\mathsf H_\nu}\ge
 \frac{c_0\beta}{2\{c_0\beta+8C_{\rm f}(\beta+1)\}}
 \ge\frac{c_0}{2(c_0+16C_{\rm f})}>0.
\end{equation}
The constant is in the original rate-one single-group clock; no
collective rotation has been inserted into the generator.

\section{A non-wrapping all-hidden-field shell}
\label{sec:nonwrapping-repair}
The following geometry is defined on the integer lattice before
embedding into a torus. It therefore retains the outer boundary cost
which a small periodic shell would miss. Keep the parent $P$ and the
central signs \cref{eq:central-sheet}. For $q\ge3$ set
\[
 Q_q=\{-1,0,\ldots,q-2\}^3\setminus\{0,1\}^3
\]
and select the seven direction-$2$ free links
\[
 ((2a+u,2b+v,1,2c+w),2),\quad
 (a,b,c)\in Q_q,\quad
 (u,v,w)\in\{0,1\}^3\setminus\{0\}.
\]
Their set $T_q$ has size $7(q^3-8)$. A selected link occurs in two
nonroot paths over its own coarse edge and in no other pivot source
family. Call a plaquette uncontrolled when it contains a parent-hidden
physical link or a parent-dependent pivot; the other incident
plaquettes are determined by retained inputs.

The finite incidence enumeration gives the following data. The last
column is the upper second derivative under a common selected-link
rotation, with arbitrary uncontrolled fields.
\begin{center}\small
\begin{tabular}{rrrrrrr}
\toprule
$q$ & $|T_q|$ & Two selected & Fixed $-$ & Fixed $+$ & Uncontrolled & $S''$ upper\\
\midrule
3 &133&261&102&126&48&144\\
4 &392&828&324&288&84&96\\
5 &819&1845&690&450&84&$-312$\\
\bottomrule
\end{tabular}
\end{center}
Use $q=5$, write $T=T_5$, and let $F$ be the following sufficient
visible collar. Include $T$, all free inputs of its $117$ source
families, and all free inputs of the fixed incident plaquettes. For
this closure, discard tree links and replace a pivot by the free
links in its fourteen nonroot paths; those paths contain no pivots,
so no further iteration is required. The reported finite data are
\begin{equation}\label{eq:shell-collar-data}
 |T|=819,\quad |F|=10034,\quad F\cap I_P=\varnothing.
\end{equation}
The closure of all incident plaquettes, including the uncontrolled
ones, uses $10431$ free groups. Its required vertices lie within
$[-3,9]\times[-3,9]\times[0,3]\times[-3,9]$. Translation by
$(8,8,8,8)$ embeds the whole closure without wraparound into every
even fine torus of side $L\ge32$. These integer-lattice counts and
source exclusions are the geometric certificates of
\src{F13}{V14}, not an extrapolation of the periodic side-six shell.

\begin{theorem}[Non-wrapping negative curvature and visible repair]
\label{thm:nonwrapping-repair}
At $C=I$, prescribe $Y_l=\epsilon_lI$ on $F$ and leave every other
free group arbitrary. Under $Y_l(t)=-e^{tE}$ for all $l\in T$, with
$E^2=-I$ and $\norm E_{\rm HS}^2=2$, the completed action obeys
\[
 S''(0)\le-312,\qquad
 \operatorname{Hess}S(a,a)\le-4/21
\]
for the corresponding unit direction $a$. If $F$ is retained in the
interface, a fixed smaller visible collar cylinder $D_L$ has the
original-clock floor \cref{eq:uniform-repair-template} with $c_0=1/84$.
All constants can be chosen independently of $L\ge32$ and of the
compatible hidden index sets disjoint from $F$.
\end{theorem}
\begin{proof}
The selected source products are near $-I$, outside the logarithmic
cutoff. The associated $117$ completed pivots are exactly $I$ on a
fixed neighbourhood and have inverse Jacobian one there. Other pivots
are independent of the selected motion. Two-selected plaquettes have
opposite orientations and cancel the common rotation exactly. A fixed
negative one-selected plaquette contributes $-2$ to $S''$, a positive
one contributes $2$, and each uncontrolled plaquette contributes at
most $2$ by unitarity of its other quaternion factors. Hence
\[
 S''(0)\le2(-690+450+84)=-312.
\]
The incidence check is $2\cdot1845+690+450+84=6\cdot819$.
The independent-coordinate speed squared is $2\cdot819=1638$, so
normalization gives $-312/1638=-4/21$.

Only a fixed finite compact input closure enters these derivatives.
Uniform continuity gives a visible chart on which
$\partial_a^2S\le-2/21$ for all uncharted fields. The inverse-pivot
amplitude is independent of the selected groups there; the
coordinate-Haar second derivative is bounded. For large $\beta$,
$\partial_a^2\Phi_\beta\le-\beta/21$. Choose inverse-metric bound
$c_g=1/2$. The local diffusion coefficient is then
$c_0=c_g/(2\cdot21)=1/84$. All chart and cutoff constants are
volume independent. Apply \cref{lem:repair-localization} and
\cref{lem:marginal-form-domination}; the visibility condition on $F$
ensures that the lifted event permits all hidden variables.
\end{proof}
The positive upper bounds for $q=3,4$ only mean that this particular
worst-case sufficient test fails there. The $q=5$ theorem is a local
repair statement, not an assertion that these collar events cover all
defects or have large probability.

\subsection{Transport to locally curved coarse backgrounds}
\label{sec:coarse-gauge-repair}
The needed coarse input closure consists of $550$ edges inside a
contractible box with eight vertices in each coordinate direction.
Assume only that the elementary coarse plaquettes in this box satisfy
$d(I,Q_p(C))\le\varepsilon$. In a monotone path gauge, move the last
step of an edge past at most $3\cdot7=21$ later-direction steps.
Each swap changes the transport by a conjugate plaquette holonomy.
Bi-invariance and the triangle inequality give a gauge $h$ with
\begin{equation}\label{eq:local-box-gauge}
 d(I,C_e^h)\le21\varepsilon
\end{equation}
for every edge in the box. This uses no commuting-field approximation.

The cell-constant fine gauge
$Y_{x,\mu}^h=h_rY_{x,\mu}h_s^{-1}$, where $r,s$ are endpoint coarse
cells, preserves tree links and is a product Haar-preserving isometry
of the independent groups. The completed pivot maps are equivariant,
so their inverse Haar Jacobians and the full fixed-coarse law transform
exactly. The same transformation conjugates the single-group heat
baths and every marginal with the same index partition. Since $h$
depends only on the fixed coarse input, no additional field-dependent
Jacobian is introduced.

Finite-input continuity of the shell certificate therefore supplies
an $\varepsilon_*>0$ such that the repair floor with $c_0=1/84$
persists, uniformly in ambient volume, whenever
$\varepsilon\le\varepsilon_*$. The visible cylinder is the pullback
of the central collar chart by this gauge. It does not require that
all coarse links globally be near identity or that wrapping holonomies
vanish. The constant $\varepsilon_*$ and the chart radius are not
numerically enclosed. \src{F13}{V15}

\section{A noncentral visible matrix certificate}
\label{sec:matrix-repair}
The shell need not remain near its central pattern. Keep the fixed
$T,F$ and coarse input set $K$ above. Let $\mathcal P_1$ be its $1140$
visible one-selected plaquettes and $\mathcal P_2$ its $1845$ visible
two-selected plaquettes. Impose the strict source margin
\begin{equation}\label{eq:inactive-margin}
 d_{\rm ch}(I,C_e^{-1}P_{e,j}(Y_F))\ge\rho+\sigma,
 \qquad \sigma>0,
\end{equation}
for all fourteen paths of the $117$ selected edges. Since
$d_{\rm ch}\le d$, the corresponding cutoffs are inactive in an
open neighbourhood; their actual pivots equal $C_e$ and their inverse
Jacobians are one. All other coarse and fine inputs can be arbitrary.

Write a quaternion as $(q_0,\mathbf q)$, with $\Tr q=2q_0$.
Cyclically start a two-selected plaquette at its positive selected
link, writing its word as $UAV^{-1}B=XB$, where $X=(x_0,x)$ and
$B=(b_0,b)$. Define the symmetric matrix
\begin{equation}\label{eq:visible-matrix}
 M_{\rm vis}=2\sum_{p\in\mathcal P_1}q_{p,0}I_3
  +8\sum_{p\in\mathcal P_2}
       \left\{\frac{x_pb_p^T+b_px_p^T}{2}-(x_p\cdot b_p)I_3\right\}.
\end{equation}
It is a smooth function of the actual completed visible plaquettes
and their fixed compact coarse inputs. It is a criterion in the
declared tree coordinates; a common rotation axis is not asserted to
be invariant under independent cell gauges.

\begin{theorem}[Matrix-certified repair with arbitrary hidden fields]
\label{thm:matrix-repair}
For fixed $\sigma,\eta>0$, let $D_C^{\sigma,\eta}$ be the visible
event defined by \cref{eq:inactive-margin} and
\begin{equation}\label{eq:matrix-negative-test}
 \lambda_{\min}(M_{\rm vis})\le-168-\eta.
\end{equation}
If $F$ is retained, the exact interface heat bath has the floor
\cref{eq:uniform-repair-template} on this event with
$c_\eta=\eta/26208$. Its constants are uniform in even $L\ge32$,
in all compact coarse inputs, and in compatible hidden index sets.
They may depend on the fixed shell, $\sigma$, and $\eta$.
\end{theorem}
\begin{proof}
Left multiply every selected free link by
$R_n(t)=\exp(tE_n)$, where $n\in\R^3$ is unit. A one-selected
plaquette contributes $2q_{p,0}$ to the second derivative. A
two-selected plaquette has moving trace $\Tr(R_nXR_n^{-1}B)$.
The identity
\[
 [E_n,[E_n,X]]=(0,4\{n(n\cdot x)-x\})
\]
gives its action second derivative
$8\{(n\cdot x)(n\cdot b)-x\cdot b\}$. Each of the $84$
uncontrolled plaquettes contributes at most two. Consequently
\[
 S''(0)\le n^TM_{\rm vis}n+168.
\]
Choose a least-eigenvalue axis. The speed squared is $1638$, so the
unit-direction Hessian is at most $-\eta/1638$ for all uncontrolled
fields.

The closed inequalities defining the event give a compact subset of
the fixed finite carrier $G^K\times G^F$. At each point choose one
axis and freeze it in orthonormal normal coordinates. Continuity and
the strict source margin give a chart with action second derivative
at most $-\eta/(2\cdot1638)$. The inverse-pivot amplitude is constant
in the selected variables; bounded coordinate-Haar terms give measured
curvature at most $-\beta\eta/(4\cdot1638)$ at large $\beta$.
With inverse metric at least $1/2$, the local diffusion coefficient
is $\eta/(16\cdot1638)=\eta/26208$.

A finite cover of this fixed compact carrier admits normalized smooth
cutoffs whose squared sum is one on the certified set and at most one
elsewhere, with a bounded sum of squared fine-coordinate derivatives.
No derivative of an eigenvector through a crossing is needed. This
gives \cref{eq:repair-localization} uniformly in $C$; apply the global
resolvent and exact marginal domination. If the event is empty, the
form inequality is vacuous. \src{F13}{V16}
\end{proof}
At the central collar, $M_{\rm vis}=-480I_3$, so adding the worst-case
$168I_3$ recovers $-312I_3$. The event is nonempty for $0<\eta<312$
and sufficiently small $\sigma$, and includes noncentral perturbations.
A weaker scalar sufficient test is
\[
 \tfrac13\Tr M_{\rm vis}
 =2\sum_{\mathcal P_1}q_{p,0}
       -\tfrac{16}{3}\sum_{\mathcal P_2}x_p\cdot b_p
                         \le-168-\eta.
\]
This is not equivalent to the least-eigenvalue test. The finite
covering here is over configurations on one fixed patch, not over a
number of spatial patches growing with volume.

\section{Repair with genuinely active pivot sources}
\label{sec:active-repair}
The inactive-source margin excludes an important complementary family.
At $C=I$, reverse the central signs on a nonempty subset $R\subset T$,
leaving the other groups in $F$ at the central pattern. Let
$k_e\in\{0,\ldots,7\}$ count the reversed links over edge $e$.
Under the same common selected-link rotation, the fourteen source
products on that edge are $2k_e$ copies of $R_n(t)$ and
$14-2k_e$ copies of $-R_n(t)$. The first are in the inner logarithm
region and contribute to the pivot derivative.

\begin{lemma}[Exact active pivot motion]
\label{lem:active-pivot-motion}
On a common sufficiently small interval in $t$, the actual pivot is
\[
 B_e(t)=\exp(-a_etE_n),\qquad a_e=\frac{k_e}{8-k_e},
 \qquad B'_e(0)=-a_eE_n,\quad B''_e(0)=-a_e^2I.
\]
Every other pivot is independent of this selected motion.
\end{lemma}
\begin{proof}
At the candidate pivot the positive relative sources have logarithm
$(1+a_e)tE_n$ in the inner cutoff region, while the negative ones
remain outside the support. The completed equation reduces to
$-a_e+(2k_e/16)(1+a_e)=0$. Its solution is the displayed $a_e$.
Global uniqueness of the pivot identifies this candidate with the
actual solution. There are only eight occupancies, so the interval
can be chosen common. Source incidence excludes all other pivots.
\end{proof}
A zero logarithm at the centre does not make an inner active source
an inactive source. Its first derivative and the pivot acceleration
both enter the completed Hessian.

Define the rational curvature cost
\begin{equation}\label{eq:active-debit}
 d(a)=\max\{2a^2+4a,4a^2\},\quad
 \ell(k)=24k+6d\!\left(\frac{k}{8-k}\right),\quad
 \mathcal L(R)=\sum_e\ell(k_e).
\end{equation}
The values for $k=0,\ldots,7$ are
\[
 0,\quad1356/49,\quad172/3,\quad2268/25,\quad132,
                   \quad580/3,\quad360,\quad1344.
\]

\begin{theorem}[Active-source curvature and original-clock repair]
\label{thm:active-repair}
At every prescribed $R$-pattern, uniformly in all unprescribed fields,
\[
 S''(0)\le-312+\mathcal L(R).
\]
For fixed $\eta>0$, consider all nonempty subsets with
$\mathcal L(R)\le312-\eta$. A union $D_\eta^{\rm act}$ of sufficiently
small closed inner collar charts around those patterns has the floor
\cref{eq:uniform-repair-template} with $c_\eta=\eta/26208$, uniformly
in even $L\ge32$ and compatible hidden sets, provided $F$ is retained.
The charts can be chosen outside every all-sources-inactive event of
\cref{thm:matrix-repair}.
\end{theorem}
\begin{proof}
First hold pivots artificially fixed only to compare derivatives.
Two-selected central plaquettes still cancel the common rotation.
Reversing a selected sign changes a fixed one-selected second
derivative by at most four, at six incident plaquettes, so this
auxiliary derivative is at most $-312+24|R|$.

Restore every actual pivot term using \cref{lem:active-pivot-motion}.
A pivot of speed $a$ and acceleration magnitude $a^2$ contributes at
most $2a^2+4a$ when the other moving parallel link is a selected free
link, and at most $2a^2$ if it is fixed. For two moving pivots with
speeds $a,b$, all new terms are at most
\[
 2a^2+2b^2+4ab\le4a^2+4b^2\le d(a)+d(b).
\]
These product-rule bounds use only unitarity of the remaining links
and include their noncommuting configurations. Each pivot meets six
plaquettes; adding the costs gives \cref{eq:active-debit}. The
independent-coordinate speed squared remains $1638$: solved pivots
are not extra coordinates in the product metric.

There are finitely many admissible patterns on the fixed collar.
Their unit action Hessians are at most $-\eta/1638$. Uniform continuity
on the finite compact input closure gives nearby chart curvature at
most $-\eta/(2\cdot1638)$. Unlike the inactive case, the inverse-pivot
amplitude is not constant. Its logarithmic second derivatives are
bounded independently of volume, as are the coordinate-Haar terms.
Absorbing these actual order-one terms at large $\beta$ gives the same
measured local coefficient $c_\eta=\eta/26208$. The fixed finite
cover, global resolvent, and marginal domination prove the floor.

Every centre has a positive inner source with $C_e^{-1}P_{e,j}=I$.
Shrink its chart so this source remains inside chordal radius $\rho$;
then it violates \cref{eq:inactive-margin} for every $\sigma>0$.
This does not delete its nonconstant Jacobian. \src{F13}{V17}
\end{proof}
Eleven reversals on distinct coarse edges give
$S''\le-312+11(1356/49)=-372/49$, while five reversals on one edge
give $S''\le-312+580/3=-356/3$. The first margin yields
$c_\eta=31/107016$. Occupancies six and seven fail this particular
sufficient cost test; that failure does not exclude other negative
directions. General transition-cutoff configurations, high occupancy
loads, and other noncentral patterns remain outside this occupancy-based
certificate. The inverse-paired family in
\cref{sec:paired-transition-repair} gives a complementary mechanism
through the complete radial transition. Chart radii and thresholds
remain existential, not certified simulation parameters.

\section{Inverse-paired sources through the cutoff transition}
\label{sec:paired-transition-repair}
The occupancy cost in \cref{thm:active-repair} is sufficient but not
necessary for a repair direction. A complementary construction uses
inverse-paired source paths, for which the acceleration of the solved
pivot is constrained by symmetry rather than bounded independently.
It covers a noncommuting family through the complete radial cutoff
transition, without asserting a cover of all active configurations.

Use imaginary-quaternion coordinates $E_h$ with $|h|=1$, and the
smooth truncated logarithm of \cref{sec:global-pivots}. Its geodesic
cutoff radius is $\rho=1/64$. For
$Q=\cos u+\sin u\,E_n$, $|n|=1$, put
$s=\sqrt2u/\rho$ and define on the cutoff support
\begin{align}
 \lambda(Q)&=\chi(s)u\cot u,&
 d(Q)&=\chi(s)+s\chi'(s),&
 \gamma(Q)&=\chi(s)u,\notag\\
 D(Q)&=\lambda I+(d-\lambda)nn^{\mathsf T},&
 \Gamma(Q)&=\gamma J_n,& J_nw&=n\times w.
 \label{eq:paired-log-differential}
\end{align}
Set $D(I)=I$, $\Gamma(I)=0$, and both to zero off the support.
The left and right logarithmic differentials are $D\mp\Gamma$.
In particular $D\preceq I$, although $d$ can be negative in the
transition. The sign of $d$ is not discarded.

\begin{proposition}[Exact paired pivot response]
\label{prop:paired-pivot-response}
Consider four inverse pairs on a coarse edge, with weights $z_l>0$
and path variations
$Q_l\exp(tz_lE_h)$ and $\exp(tz_lE_h)Q_l^{-1}$.
The other source paths have inactive cutoff. The solved pivot obeys
$B(-t)=B(t)^{-1}$ and has
\begin{align}
 B(0)&=I,& B'(0)&=E_{a_e},& B''(0)&=-|a_e|^2I,\notag\\
 a_e&=-A_e^{-1}\sum_{l\sim e}z_l(D_l+\Gamma_l)h,&
 A_e&=8I-\sum_{l\sim e}D_l\succeq4I.
 \label{eq:paired-pivot-response}
\end{align}
Thus the Lie logarithm of the pivot has zero second derivative at the
center, but the group-valued acceleration remains present.
\end{proposition}
\begin{proof}
Inversion oddness and adjoint equivariance of the truncated logarithm
interchange the two paths when $t$ changes sign. Uniqueness of the
solved pivot gives $B(-t)=B(t)^{-1}$. Differentiating the actual pivot
equation gives the linear system in
\cref{eq:paired-pivot-response}; the four bounds $D_l\preceq I$
give its stated lower bound. Differentiating the exponential of its
odd Lie logarithm yields $B''(0)=E_{a_e}^2=-|a_e|^2I$.
\end{proof}

\begin{proposition}[All radial phases for coplanar axes]
\label{prop:paired-coplanar-cone}
Suppose one unit vector $h$ is perpendicular to every noncentral
source axis $n_l$. The axes can differ and need not commute. Writing
$a_e=-\alpha_e h+w_e$ with $w_e\perp h$, one has
\begin{equation}
 \alpha_e\ge0,\qquad |w_e|\le\tfrac3{64}\alpha_e.
 \label{eq:paired-coplanar-cone}
\end{equation}
This holds for every source phase, including the cutoff transition and
its endpoints, independently of the negative size of $d(Q)$.
\end{proposition}
\begin{proof}
Since $D_lh=\lambda_lh$, the matrix $A_e$ preserves $\R h$ and
$h^\perp$. If $L_e=\sum_l\lambda_l$ and
$P_e=\sum_lz_l\lambda_l$, then
\[
 \alpha_e=\frac{P_e}{8-L_e},\qquad
 w_e=-(A_e|_{h^\perp})^{-1}
                    \sum_lz_l\gamma_l(n_l\times h).
\]
On the support, $\gamma_l/\lambda_l=\tan u_l$ and
$u_l<\rho/\sqrt2$; zero terms are interpreted continuously.
Using $A_e|_{h^\perp}\succeq4I$ and $8-L_e\le8$ gives
$|w_e|\le2\tan(\rho/\sqrt2)\alpha_e
\le(3/64)\alpha_e$. At central or inactive sources the same result
follows by continuity. The radial derivative appears only in the
positive lower bound for $A_e$ and need not be bounded in absolute
value.
\end{proof}

\begin{certificate}[Reduced-collar curvature for the paired family]
\label{cert:paired-transition-curvature}
The fixed reduced collar of \src{F14}{V5--V11} has 133 selected links
on 19 coarse edges, 484 directly affected unrotated links, and 152
source-only companions, in 76 inverse pairs. The remaining groups
are unrestricted. For its specified positive integer rotation weights,
\[
 \sum_lz_l^2=21\,328\,389,\qquad
                \norm{\dot U}_{\rm HS}^2=42\,656\,778.
\]
The finite plaquette-incidence calculation gives, whenever every pivot
response satisfies $\alpha_e\ge0$ and $|w_e|\le\alpha_e$,
\begin{equation}
 \frac{\dd^2}{\dd t^2}S(U(t))\bigg|_{t=0}
       \le-\frac{981\,794\,718\,651}{131\,072},\qquad
            D^2S[\widehat{\dot U},\widehat{\dot U}]\le-\tfrac16.
 \label{eq:paired-transition-curvature}
\end{equation}
The unit direction uses the Hilbert--Schmidt metric. All affected
plaquettes, the solved-pivot velocities, and the group accelerations
are included. In particular \cref{prop:paired-coplanar-cone} populates
this certificate for arbitrary coplanar inverse-paired phases.
\end{certificate}

\begin{corollary}[Measured original-clock repair in a transition family]
\label{cor:paired-transition-repair}
For the certified collar, identity coarse data, and even side length
$L\ge32$, sufficiently small closed visible neighborhoods admit the
same-law killed heat-bath lower bound
\begin{equation}
 \kappa_D^{\mathsf H_\nu}
 \ge\frac{c_{\rm rep}\beta}
          {2\{c_{\rm rep}\beta+8C_f(\beta+1)\}}
 \ge\frac{c_{\rm rep}}{2(c_{\rm rep}+16C_f)},
 \qquad c_{\rm rep}=1/96,
 \label{eq:paired-transition-repair}
\end{equation}
for sufficiently large $\beta\ge1$. The radii and the threshold are
existential uniform choices on the declared compact source family.
All groups outside the visible collar remain unrestricted.
\end{corollary}
\begin{proof}
The strict unit curvature margin in
\cref{eq:paired-transition-curvature} survives a sufficiently small
closed neighborhood by compactness and smooth dependence of the
solved pivots. The actual inverse-Haar density has the same controlled
local derivatives there. The measured negative-curvature estimate
therefore supplies the diffusion coefficient $c_{\rm rep}$.
Apply the original-clock same-law resolvent comparison of
\cref{thm:global-resolvent}, with its complete conditional derivative
ceiling $C_f(\beta+1)$. This gives the first inequality; the second
uses $\beta+1\le2\beta$.
\end{proof}
The new family removes the blanket exclusion of cutoff-transition
sources. It does not remove the need to treat non-paired sources,
non-coplanar configurations outside the certified cone, or incompatible
large occupancy patterns. These are coverage questions for a complete
same-law assembly, not defects repaired by renaming a local floor.
\sourcebasis{The reduced-collar incidence and weighted direction are
specified in \src{F14}{V5--V10}; the full transition differential,
coplanar cone, and measured repair theorem are in \src{F14}{V11}.}


\section{A quantitative local-to-global interface criterion}
\label{sec:weighted-patch-criterion}
Local killed floors do not automatically imply the Poincar\'e
inequality of an unrestricted interface. A useful alternative
criterion keeps all interactions inside actual conditional patches
and makes their required strength explicit.

For the exact unrestricted interface law, let
$h_i=I-P_i$ and $\mathsf H=\sum_i h_i$ in the original rate-one clock.
Assume a commutation graph of degree at most $\Delta$, with owner
range $R_*$: nonadjacent $h_i,h_j$ commute, and adjacent owners differ
by at most $R_*$ in maximum norm. For each graph edge choose
$0\le c_{ij}\le1$ with
\[
 \{h_i,h_j\}\succeq-c_{ij}(h_i+h_j).
\]
The universal value one follows from $(h_i+h_j)^2\succeq0$ and
$h_i^2=h_i$. Smaller values require actual pair estimates.

For a nonnegative finitely supported owner weight $a$, put
\begin{align*}
 m_1&=\sum_y a(y),&m_2&=\sum_y a(y)^2,\\
 d_a(\delta)&=\tfrac12\sum_y|a(y)-a(y-\delta)|^2,
 & A_x&=\sum_i a(\operatorname{owner}(i)-x)h_i,\\
 D_a&=\max_i\sum_{j:\{i,j\}\in E}c_{ij}
               d_a(\operatorname{owner}(i)-\operatorname{owner}(j)).
\end{align*}
Periodize on a sufficiently large compatible torus, without a second
wraparound overlap. Only the deterministic weights are translated;
no translation invariance of the measure is assumed.

\begin{criterion}[Weighted conditional patch squares]
\label{crit:weighted-patches}
Suppose every $A_x$ has gap at least $\gamma_a>0$ in its actual
conditional patch law, uniformly over all admitted compact exterior
and retained coarse data. Then
\begin{equation}\label{eq:weighted-patches}
 \mathsf H^2\succeq\frac{m_1\gamma_a-D_a}{m_2}\mathsf H.
\end{equation}
If $m_1\gamma_a>D_a$, the coefficient is a volume-uniform lower bound
on $\gap\mathsf H$. For a rate-one patch gap $\gamma_\ell$ and
$\alpha_a=\min\{a(y):a(y)>0\}$, it suffices that
\[
 \gamma_\ell>\tau_a:=\frac{D_a}{\alpha_a m_1},\qquad
 \gap\mathsf H\ge\frac{\alpha_a m_1}{m_2}(\gamma_\ell-\tau_a).
\]
\end{criterion}
\begin{proof}
With $q_a(\delta)=\sum_y a(y)a(y-\delta)=m_2-d_a(\delta)$,
\[
 \sum_x A_x=m_1\mathsf H,\qquad
 \sum_x A_x^2=m_2\mathsf H+
       \sum_{i<j}q_a(\operatorname{owner}(i)-\operatorname{owner}(j))
                                                  \{h_i,h_j\}.
\]
Subtracting gives
$m_2\mathsf H^2-\sum_xA_x^2\succeq-D_a\mathsf H$;
nonedges contribute positively, and the edge row sum bounds the loss.
Positive compact conditional densities and positive rates on each
patch identify the kernel with its exterior-measurable functions.
The fibrewise gap and spectral theorem give
$A_x^2\succeq\gamma_a A_x$. Sum and divide by $m_2$.
The full generator has only constant functions in its kernel.
Finally $A_x\succeq\alpha_a\sum_{i\in B_x}h_i$ yields
$\gamma_a\ge\alpha_a\gamma_\ell$ by the conditional variational
principle. No operator inequality has been squared by a monotonicity
assumption. \src{F12}{V63}
\end{proof}

\begin{corollary}[An explicit four-dimensional finite-size threshold]
\label{cor:weighted-threshold}
Let $\ell$ be divisible by four, with
$\ell\ge4$, $\ell>2R_*$, $\ell\ge R_*^2$, and sufficiently large
ambient owner side at least $2\ell+2R_*$. Define
\[
 C_*=32\Delta(R_*^2+32R_*).
\]
If all the actual rate-one cube-patch gaps, with arbitrary compact
exterior, satisfy $\gamma_\ell>C_*\ell^{-3/2}$, then
\begin{equation}\label{eq:patch-explicit-floor}
 \gap\mathsf H\ge\ell^{-1/2}
                   (\gamma_\ell-C_*\ell^{-3/2})>0.
\end{equation}
\end{corollary}
\begin{proof}
On $\{1,\ldots,\ell\}^4$ use the floor-plus-pyramid weight
\[
 a_\ell(y)=\ell^{-1/2}
       +\ell^{-1}\min_{1\le b\le4}\{y_b,\ell+1-y_b\},
\]
and set it to zero elsewhere. Then $\alpha_a\ge\ell^{-1/2}$,
$m_2\le m_1$, and $m_1\ge\ell^4/64$. For
$\norm\delta_\infty\le R_*$, differences in the common interior
are at most $R_*/\ell$. The symmetric difference of the cubes has
at most $8R_*\ell^3$ sites, on which the nonzero weights are at most
$2\ell^{-1/2}$. Hence
\[
 d_a(\delta)\le(R_*^2/2+16R_*)\ell^2,
 \qquad D_a\le\Delta(R_*^2/2+16R_*)\ell^2.
\]
Insert these bounds in \cref{crit:weighted-patches} and use
$m_1/m_2\ge1$ after the positive difference is established.
\end{proof}
A bound $\gamma_\ell\ge c_\beta\ell^{-p}$ with $p<3/2$ would therefore
suffice at a sufficiently large fixed patch size and fixed $\beta$.
The available surface-order certificate
$\exp[-A(\beta+1)\ell^3]$ does not verify this threshold. It is a
lower certificate, not an upper estimate on the actual gap.

\subsection{The remaining coverage and physical-transport requirements}
The repair events of
\cref{thm:nonwrapping-repair,thm:matrix-repair,thm:active-repair,cor:paired-transition-repair}
are actual visible events with all hidden fields unrestricted. They
supply more than a bare rarity estimate, but not every all-exterior
conditional patch gap required by \cref{crit:weighted-patches}.
A finite configuration-space cover on one fixed patch cannot be
replaced by a union of an increasing number of spatial repair events
without a new localization or assembly estimate. Such union-avoidance
inferences are explicitly obstructed in \src{F12}{V79}.

The remaining auxiliary task is therefore to control the complementary
source and curvature cases, or to prove a different complete
same-law patch inequality which crosses the displayed threshold.
Neither the negative hidden-diffusion result nor the sufficient patch
criterion excludes another heat-bath mechanism. Even an unrestricted
volume-uniform auxiliary gap would require cofinal coupling control
and a proved comparison to the physical reflected transfer before it
could enter \cref{hyp:contraction}. This is the distinction maintained
by the native transfer analysis in \cref{ch:native-transfer} and the
physical endpoint in \cref{ch:spectral}.

\chapter{Native Wilson transfer and critical spectral limits}
\label{ch:native-transfer}
The static measured response and an auxiliary repair dynamics do not
identify the physical transfer. There is, however, a separate family
of results for the actual compact Wilson transfer on a spatial slice.
This chapter keeps its Hilbert space, magnetic factors, Perron
normalization, and limiting regime explicit. Fixed-regulator
asymptotics are followed by their controlled joint-regulator
extension; local fluctuation identities then isolate requirements
for further continuation. The spatial dimension is
three; the four-dimensional static integral of
\cref{ch:full-compact-sources} is not substituted for this operator.

\section{Exact spatial assembly and the normalized magnetic image}
\label{sec:native-assembly}
Let $G=\SU(2)$ and let the spatial torus have side $L=2m\ge8$.
Partition its vertices into $B=m^3$ disjoint cubes. There are $12B$
internal links and $12B$ bridge links, with $6B$ internal and $18B$
mixed plaquettes. A rooted tree in each cube gives exact Haar
coordinates consisting of seven gauge variables and five internal
chords per cube, together with the bridges. After nonroot gauge
reduction the full physical space is
\begin{equation}
 \mathcal H_{\rm phys}
       =L^2(G^{5B}\times G^{12B})^{G^B}.
 \label{eq:native-physical-space}
\end{equation}
The remaining root group conjugates the chords and acts at the bridge
endpoints. Define
\begin{equation}
 k_\gamma(U)=z_\gamma^{-1}e^{-\gamma s(U)},\qquad
 s(U)=1-\tfrac12\Tr U,\qquad
 z_\gamma=\int_Ge^{-\gamma s(U)}\dd H(U).
 \label{eq:native-Wilson-kernel}
\end{equation}
These are normalized Wilson convolution kernels, not exact heat
kernels. The native transfer factors as
\begin{equation}
 T_\gamma=M_\times
 \left[\left(\bigotimes_b T_{\Box,\gamma}^{(b)}\right)
                      \otimes\mathcal C_\gamma\right]M_\times,
 \qquad
 M_\times=e^{-\gamma S_\times/2},\qquad
 \mathcal C_\gamma=\bigotimes_{e\,\mathrm{bridge}}C_\gamma^{(e)}.
 \label{eq:native-assembly}
\end{equation}
Here $S_\times$ includes all mixed plaquettes. The regional and bridge
kinetics commute before the magnetic sandwich, but omission of that
sandwich changes the operator.

Let $\Phi_\gamma(X)$ be the normalized internal Gaussian carrier in
its actual Haar coordinates, including its exact chart normalization.
Choose the equivariant moving endpoint frame
$Z=\Theta_XY$, so at fixed $X$ the map $Y\mapsto Z$ preserves Haar.
For the isometry
$\widehat S_\gamma f(X,Y)=\Phi_\gamma(X)f(\Theta_XY)$, set
\begin{equation}
 F_{\gamma,t}(Z)=\int|\Phi_\gamma(X)|^2
            M_\times(X,\Theta_X^{-1}Z)^t\dd H(X).
 \label{eq:native-magnetic-Gram}
\end{equation}
Then
\[
 \widehat S_\gamma^*M_\times^t\widehat S_\gamma=M_{F_{\gamma,t}},
 \qquad
 \mathcal I_\gamma=M_\times\widehat S_\gamma
                                      M_{F_{\gamma,2}^{-1/2}},
 \qquad \mathcal I_\gamma^*\mathcal I_\gamma=I.
\]
The normalization is $F_{\gamma,2}$, not $F_{\gamma,1}^2$.
Its positivity at fixed regulator gives an exact normalized magnetic
image without deleting the exterior bridge configurations.

\begin{theorem}[Native bridge compression and conditional affinity]
\label{thm:native-affinity}
Define unit conditional vectors and their Gram kernel by
\begin{align*}
 \psi_{\gamma,Z}(X)&=
 \frac{\Phi_\gamma(X)M_\times(X,\Theta_X^{-1}Z)}
                    {\sqrt{F_{\gamma,2}(Z)}},\\
 \mathcal A_\gamma(Z,Z')&=
            \int\psi_{\gamma,Z}(X)\psi_{\gamma,Z'}(X)\dd H(X).
\end{align*}
The exact bridge-factor compression is
\begin{equation}
 (\mathcal I_\gamma^*\mathcal C_\gamma\mathcal I_\gamma)(Z,Z')
     =\left[\prod_{e\,\mathrm{bridge}}k_\gamma(Z_eZ_e'^{-1})\right]
                               \mathcal A_\gamma(Z,Z').
 \label{eq:native-affinity-kernel}
\end{equation}
The affinity is continuous, symmetric, positive semidefinite as a
Gram kernel, and satisfies $0<\mathcal A_\gamma\le1$ and
$\mathcal A_\gamma(Z,Z)=1$. The compressed operator is a positive
contraction. Root gauge invariance is simultaneous in $Z,Z'$.
\end{theorem}
\begin{proof}
The bridge factor leaves $X$ unchanged. Writing
$Y_e=h_sZ_eh_t^{-1}$ and $Y_e'=h_sZ_e'h_t^{-1}$ gives
$Y_eY_e'^{-1}=h_sZ_eZ_e'^{-1}h_s^{-1}$. Centrality of $k_\gamma$
and Haar invariance eliminate the frame factors. In these coordinates
$\mathcal I_\gamma f=\psi_{\gamma,Z}f(Z)$, giving
\cref{eq:native-affinity-kernel}. Cauchy--Schwarz gives the upper bound;
the Gram representation and compression give positivity. Strict
positivity and continuity follow from the finite compact integral.
\end{proof}
This is an identity for the bridge kinetic factor. The cube dynamics
and their excited-state returns are still present in
\cref{eq:native-assembly}.

In the fixed-$B$ Gaussian scaling, let $C_B$ be the internal squared
Gaussian covariance and let $A,D$ be the linearized mixed-plaquette
incidences on internal and bridge coordinates in the same moving
frame. The conditional Gaussian has
\[
 C_{\rm p}=(C_B^{-1}+A^{\mathsf T}A)^{-1},\qquad
 m_y=-C_{\rm p}A^{\mathsf T}Dy,
\]
with the colour tensor $I_3$ understood. Completing a square gives
\begin{equation}
 \langle\psi_y,\psi_z\rangle
       =\exp[-(y-z)^{\mathsf T}R(y-z)/8],\qquad
 R=D^{\mathsf T}AC_{\rm p}A^{\mathsf T}D\succeq0.
 \label{eq:native-Gaussian-affinity}
\end{equation}
The actual chart kernels converge to this form on fixed bounded
windows at fixed $B$. This is not global norm convergence uniform in
volume. Common parallel-bridge directions annihilated by
$A^{\mathsf T}D$ have $R=0$; a static fast conditional vector cannot
supply a restoring term in those directions.

\subsection{Why static normalization does not supply kinetic time}
For a Euclidean heat kernel with covariance metric $G>0$, multiplying
its kernel by the Gaussian affinity in
\cref{eq:native-Gaussian-affinity} gives exactly
\begin{align}
 H_t^G(y,z)e^{-(y-z)^{\mathsf T}R(y-z)/8}
       &=c_t H_t^{G_t}(y,z),\qquad G_t=(G^{-1}+tR/2)^{-1},\notag\\
 c_t&=\det(I+(t/2)G^{1/2}RG^{1/2})^{-1/2}.
 \label{eq:affinity-modified-heat}
\end{align}
The convention is
$H_t^G(y,z)=(4\pi t)^{-q/2}(\det G)^{-1/2}
\exp[-(y-z)^{\mathsf T}G^{-1}(y-z)/(4t)]$.
After division by $c_t$, the family generally fails the semigroup
law because $tG_t$ is not additive in $t$, unless $R=0$.
Thus even the exactly normalized static magnetic image does not
define a physical generator by itself. Full kinetic composition and
its return channels must still be controlled.
\sourcebasis{Exact spatial assembly and magnetic normalization are
proved in \src{F15}{V98--V100}. The conditional affinity and its
finite-time distinction are developed in \src{F16}{V1}.}

\subsection{The excitation budget of the magnetic image}
\label{sec:native-excitation-budget}
Static normalization also changes which internal oscillator levels are
occupied. This can be quantified before any claim of invariance under
the full transfer. Retain the actual cube matrices $A,D,C_B$ used in
\cref{eq:native-Gaussian-affinity}, and set
\[
 W=C_B^{1/2}A^{\mathsf T}AC_B^{1/2}.
\]
The specified spatial incidence satisfies
\begin{equation}
 0\preceq W,\qquad \norm W\le\rho_*=1/\sqrt2,\qquad
 \Tr W>B,\qquad \rank W\le5B.
 \label{eq:magnetic-incidence-stock}
\end{equation}
These are the finite cube-incidence bounds of \src{F15}{V99--V100}.
The matrices here act on path indices; all three colour directions
are retained by $\mathsf W=W\otimes I_3$. In whitened internal
coordinates $u\in\R^{n}$, $n=15B$, define
\[
 Q=I+\mathsf W,\qquad
 b=(C_B^{1/2}A^{\mathsf T}D\otimes I_3)y,\qquad \mu_y=-Q^{-1}b.
\]
The original product ground vector and the normalized magnetic vector
are
\begin{align}
 \Omega_0(u)&=(2\pi)^{-n/4}e^{-|u|^2/4},\notag\\
 \psi_y(u)&=(2\pi)^{-n/4}\det(Q)^{1/4}
                   e^{-(u-\mu_y)^{\mathsf T}Q(u-\mu_y)/4}.
 \label{eq:magnetic-whitened-state}
\end{align}
The total old oscillator degree is
$\mathcal N=-\Delta+|u|^2/4-n/2$, with spectral projection
$P_{\le k}=\one_{[0,k]}(\mathcal N)$. It is not a physical
Hamiltonian. Put
\[
 S=\mathsf W(2I+\mathsf W)^{-1},\qquad
 d=-(2I+\mathsf W)^{-1}b,\qquad
 s_*=\frac{\rho_*}{2+\rho_*}<1/3.
\]

\begin{theorem}[Exact occupation law of the Gaussian magnetic image]
\label{thm:magnetic-excitation-law}
For $0\le z<s_*^{-1}$,
\begin{align}
 \langle\psi_y,z^{\mathcal N}\psi_y\rangle
  ={}&\left(\frac{\det(I-S^2)}{\det(I-z^2S^2)}\right)^{1/2}\notag\\
 &\times\exp\left\{(z-1)d^{\mathsf T}
           (I+zS)^{-1}(I+S)^{-1}d\right\}.
 \label{eq:magnetic-excitation-pgf}
\end{align}
At $z=0$ this is the ground-projection weight. In particular, with
\[
 c_-=\frac{9}{40(2+\rho_*)^2},\qquad
 c_+=\frac{15}{2}\log\frac{1-s_*^2}{1-4s_*^2},
\]
one has
\begin{align}
 \norm{P_{\le k}\psi_0}^2
       &\le\exp\{k\log2-c_-B\},\notag\\
 \norm{(I-P_{\le k})\psi_y}^2
       &\le\exp\{c_+B+(\rho_*/4)\norm{Dy}^2-k\log2\}.
 \label{eq:magnetic-excitation-tails}
\end{align}
The colour tensor in $Dy$ is understood. Every subextensive budget
$k=o(B)$ loses all asymptotic weight at $y=0$. An extensive budget
with the displayed displacement cost gives an exponentially controlled
upper tail.
\end{theorem}
\begin{proof}
An orthogonal change of variables diagonalizes $\mathsf W$ and
preserves $\mathcal N$. In one coordinate, the Hermite coefficients
of the displaced Gaussian have generating function
$k_0\exp(-sw^2/2+dw)$. Squaring the coefficients with weight $z^j$
and integrating in the complex Gaussian representation of the
Hermite basis gives
$k_0^2(1-z^2s^2)^{-1/2}\exp(zd^2/(1+zs))$. Divide by its
value at $z=1$ and multiply the $n$ factors. The identity
\[
 \frac{z}{1+zs}-\frac{1}{1+s}
       =\frac{z-1}{(1+zs)(1+s)}
\]
gives \cref{eq:magnetic-excitation-pgf} with its exact normalization.

For $0<z<1$,
$\log[(1-s^2)/(1-z^2s^2)]\le-(1-z^2)s^2$. Also
\[
 \Tr S^2=3\Tr\bigl[W^2(2I+W)^{-2}\bigr]
    \ge\frac{3B}{5(2+\rho_*)^2},
\]
using $\Tr W^2\ge(\Tr W)^2/(5B)$. At $y=0$ and $z=1/2$
this bounds the generating function by $e^{-c_-B}$. On degrees
at most $k$, $2^{-\mathcal N}\ge2^{-k}$, proving the lower-tail
bound. At $z=2$, the determinant term is at most $e^{c_+B}$.
The two positive resolvents in the exponent commute and are
contractions, while
$\norm d^2\le\norm b^2/4\le(\rho_*/4)\norm{Dy}^2$.
The upper Markov bound proves the second inequality.
\src{F15}{V100}
\end{proof}

On a bridge window $\max_e|y_e|\le R$, the literal mixed-plaquette
incidence gives $\norm{Dy}^2\le144BR^2$: there are $12B$
two-term rows and $6B$ four-term rows. Hence
\[
 k=O\bigl(B(1+R^2)+\log\delta^{-1}\bigr)
\]
suffices to lose at most $\delta$. The old total-degree space has
rank $\binom{15B+k}{k}$. This is a limitation of a truncation in the
old product oscillator basis, not a lower bound on the rank of an
adaptive carrier: for fixed $y$, $\psi_y$ itself spans one line.

\begin{proposition}[The genuine compact projection at fixed volume]
\label{prop:magnetic-compact-projection}
For fixed $B,k$, restrict the Hermite vectors of degree at most $k$
to the rescaled compact chart $E_\gamma\subset\R^{15B}$, and let
$V_{\gamma,k}$ be their synthesis map. Then
\begin{equation}
 G_{\gamma,k}=V_{\gamma,k}^*V_{\gamma,k},\qquad
 \Pi_{\gamma,k}=V_{\gamma,k}G_{\gamma,k}^{-1}V_{\gamma,k}^*
 \label{eq:magnetic-compact-projection}
\end{equation}
defines an orthogonal projection. Transporting it through the actual
Haar chart and moving frame gives a compact internal projection
respecting the remaining root-gauge action. For the actual normalized
conditional vectors at $Z_e=\exp(y_e/\sqrt\gamma)$,
\begin{equation}
 \sup_{\max_e|y_e|\le R}
 \left|\norm{\Pi_{\gamma,k}\psi_{\gamma,y}}^2
                       -\norm{P_{\le k}\psi_y}^2\right|
                       \longrightarrow0
 \label{eq:magnetic-compact-band-limit}
\end{equation}
for fixed $B,k,R$ as $\gamma\to\infty$.
\end{proposition}
\begin{proof}
A nonzero polynomial times a Gaussian cannot vanish on the nonempty
open chart, so $G_{\gamma,k}$ is positive definite. Formula
\cref{eq:magnetic-compact-projection} is therefore the projection onto
the restricted span, rather than a mere chart restriction of the
Euclidean projection. As $E_\gamma$ exhausts Euclidean space, each
of the finitely many restricted vectors converges in $L^2$ and
$G_{\gamma,k}\to I$. The projections, extended by zero, converge
in operator norm. The normalized actual conditional densities converge
in $L^1$ on each fixed bridge window at fixed volume; positivity and
$(\sqrt a-\sqrt b)^2\le|a-b|$ give $L^2$ convergence of their
amplitudes. This proves \cref{eq:magnetic-compact-band-limit}.
The chart, Gaussian, and complete total-degree space are invariant
under root colour rotations; the Gram inverse and equivariant moving
frame preserve this invariance. \src{F15}{V99--V100}
\end{proof}

The extensive choice of $k$ is applied separately for each fixed $B$.
No rate uniform in $B$ is supplied by this compact projection limit.
In particular, neither a local internal oscillator separation nor
normalization of the magnetic image bounds the complement of the
full transfer in \cref{eq:native-assembly}. Conditional affinity,
internal excitation return, and physical kinetic composition remain
separate parts of the same operator problem.

\section{Resistance-normalized gauge geometry on growing slices}
\label{sec:native-resistance-geometry}
Uniform gauge geometry must be distinguished from uniform native dynamics.
On the three-dimensional torus of volume $V=L^3$, fix a root $o=0$ and
let $d_o$ be the based vertex-to-link difference. Put $A_L=d_o^*d_o$.
For centered momenta define
\[
 \nu_L(k)=4\sum_{\mu=1}^3\sin^2(\pi k_\mu/L),\qquad
 G_L(x)=V^{-1}\sum_{k\ne0}\frac{e^{2\pi\ii k\cdot x/L}}{\nu_L(k)}.
\]

\begin{proposition}[Grounded inverse and bounded point evaluation]
\label{prop:native-resistance}
The grounded inverse has kernel
\begin{equation}
 A_L^{-1}(x,y)=K_L(x,y)
 =G_L(x-y)-G_L(x)-G_L(y)+G_L(0),\qquad x,y\ne0.
 \label{eq:native-grounded-kernel}
\end{equation}
With constants independent of $L\ge8$,
\begin{equation}
 0\le K_L(x,x)\le C,\quad
 \|A_L^{-1}\|\asymp L^3,\quad
 \|\eta\|_\infty\le C\|d_o\eta\|_2\quad(\eta_o=0).
 \label{eq:native-resistance-evaluation}
\end{equation}
Moreover $\Tr K_L=2VG_L(0)$ and
$|\|K_L\|-VG_L(0)|\le CL^2$.
\end{proposition}
\begin{proof}
The Fourier identity $\Delta G_L=\delta_0-V^{-1}$ verifies
\cref{eq:native-grounded-kernel}, including its zero root row.
Three-dimensional shell summation gives $c\le G_L(0)\le C$ and
$\sum_x|G_L(x)|^2\le CL$. Positivity gives
$K_L(x,x)=2(G_L(0)-G_L(x))\le4G_L(0)$.
Extending by a zero root row and column, write
$K_L=G-g\one^*-\one g^*+G_L(0)\one\one^*$, where
$g_x=G_L(x)$. The first three terms have norm $O(L^2)$ and the last
has norm $VG_L(0)$. Direct summation gives the trace identity.
Finally
$|\eta_x|\le\sqrt{K_L(x,x)}\|A_L^{1/2}\eta\|_2$ proves the
point-evaluation bound, componentwise for Lie-algebra-valued fields.
\end{proof}
Thus the large norm of a grounded inverse is not the cost of every
local gauge evaluation. The distinction is specific to the stated
three-dimensional slice; it is not a uniform assertion on arbitrary graphs.

Let $\mathcal N_L=\ker(d_o^*\otimes I_3)$ be the horizontal link space.
For $s\in\mathcal N_L$, use based gauge coordinates normalized by
$A_L^{1/2}$, and retain the actual exponential-slice quotient metric
$\mathsf g_s$ and coarea density $J_L(s)$. These are coordinate
representations of the original metric and Haar law.

\begin{theorem}[Uniform based tube, quotient metric, and relative density]
\label{thm:native-uniform-geometry}
There is $r_0>0$ independent of $L$ such that the based gauge map
$(g,s)\mapsto g\cdot\operatorname{Exp}s$, $\|s\|_2<r_0$, is a
smooth one-to-one parametrization of its image. Its eight center
translates are disjoint; global colour conjugation remains unfixed.
On this tube,
\begin{align}
 &\tfrac12I\preceq\mathsf g_s\preceq I,\qquad
 \|I-\mathsf g_s\|_1+\|\mathsf g_s^{-1}-I\|_1\le C\|s\|_2^2,\notag\\
 &\left|\log\frac{J_L(s)}{J_L(0)}\right|\le C\|s\|_2^2,\qquad
 J_L(0)=(2\pi^2)^{-(2V+1)}(\det A_L)^{3/2},\notag\\
 &\det A_L=V^{-1}\prod_{k\ne0}\nu_L(k).
 \label{eq:native-uniform-geometry}
\end{align}
The first and second directional derivatives of the metric in trace
norm are bounded by $C\|s\|_2\|t\|_2$ and
$C\|t\|_2\|v\|_2$, respectively. The same bounds hold for the
corresponding derivatives of $\log J_L$.
\end{theorem}
\begin{proof}
Write $P_L=A_L^{-1/2}\otimes I_3$ and normalize the orbit differential
as $Z_s=Y_sP_L$, with
$(Y_s\eta)_{xy}=\eta_x-\operatorname{Ad}_{e^{s_{xy}}}\eta_y$.
Then $Z_0^*Z_0=I$, and \cref{prop:native-resistance} gives
$\|Z_s-Z_0\|_{\rm HS}\le C\|s\|_2$, including bounded derivatives.
Together with the horizontal differential $X_s$, this proves uniform
local invertibility. If two horizontal representatives in this ball
are based-gauge equivalent, the same point-evaluation bound puts the
intervening gauge transformation in one logarithm chart. Pairing the
interpolated link logarithm with $d_o\eta$ gives
$0\ge(1-Cr_0)\|d_o\eta\|_2^2$, hence uniqueness.
Distinct central-flat orbits are separated by at least $\pi\sqrt L$
by summing the holonomy distance over disjoint fundamental loops.

Set $\mathsf V_s=Z_s^*Z_s$ and $\mathsf C_s=X_s^*Z_s$.
The exact quotient metric is
$\mathsf g_s=X_s^*X_s-\mathsf C_s\mathsf V_s^{-1}\mathsf C_s^*$.
The single-link exponential metric and
$\|\mathsf C_s\|_{\rm HS}\le C\|s\|_2$ give its trace bounds.
The coarea identity is
$J_L(s)/J_L(0)=\sqrt{\det\mathsf V_s}\sqrt{\det\mathsf g_s}$.
For its vertical factor the linear trace vanishes: explicitly
$\Tr(\mathsf V_s-I)=8\sum_{xy}K_L(x,y)\sin^2|s_{xy}|$.
Combining this identity with Hilbert--Schmidt bounds and
$\log\det(I+D)-\Tr D=O(\|D\|_{\rm HS}^2)$ proves the relative
density and differentiated bounds. At zero the orbit Gram is
$A_L\otimes I_3$; probability Haar and the Laplacian cofactor formula
give the displayed absolute scalar. \src{F16B}{V50}
\end{proof}

For the hard Gaussian field $u$ and a bounded critical coordinate $q$,
$s=\gamma^{-1/2}u+h\mathsf Kq$, where $\mathsf K^*\mathsf K=VI$,
obeys
\begin{equation}
 \|s\|_2^2=\gamma^{-1}|u|^2+L\gamma^{-2/3}|q|^2,
 \qquad \E|u|^2\asymp V.
 \label{eq:native-tube-occupation}
\end{equation}
Hence high Gaussian occupation of this fixed total-$\ell^2$ tube
requires a volume budget, for example $L^3/\gamma\to0$ for bounded
$q$. The theorem removes an artificial dimension factor from gauge
geometry; it does not prove that the native vacuum occupies the tube,
or that a much larger sup-norm chart has a unique representative.
\sourcebasis{Resistance evaluation, the uniform based tube, trace-ideal
metric control, and the exact relative Haar normalization are established
in \src{F16B}{V50}.}


\section{All-label kinetic comparison with actual Perron normalization}
\label{sec:native-subordination}
A comparison reference can preserve the large-representation shape of the
Wilson kernel without being identified with the physical dynamics. Let
$\mathcal L_G$ be the round group Laplacian, whose spin-$\ell/2$
eigenvalue is $s_\ell=\ell(\ell+2)$. The one-link Wilson convolution has
multiplier $I_{\ell+1}(\gamma)/I_1(\gamma)$; write its negative logarithm
as $\Phi_{\gamma,\ell}$. Define
\begin{equation}
 F_\gamma(s)=\int_1^{\sqrt{1+s}}\operatorname{arsinh}(x/\gamma)\dd x,
 \qquad R_\gamma=e^{-F_\gamma(\mathcal L_G)}.
 \label{eq:native-Bernstein-dispersion}
\end{equation}

\begin{proposition}[Positive subordinated reference and all-label envelope]
\label{prop:native-Bernstein}
The dispersion has the representation
\begin{equation}
 F_\gamma(s)=\int_0^\infty(1-e^{-st})
       \frac{e^{-t}}{4\sqrt\pi\,t^{3/2}}
       \mathrm E_1(\gamma^2t)\dd t,
 \qquad \mathrm E_1(x)=\int_x^\infty e^{-v}\frac{\dd v}{v}.
 \label{eq:native-Bernstein-density}
\end{equation}
Consequently $R_\gamma$ is a normalized positive mixture of group heat
kernels. For $\gamma\ge2$, with $p_\gamma=1+2/\gamma$,
\begin{equation}
 F_\gamma(s_\ell)\le\Phi_{\gamma,\ell}
                   \le p_\gamma F_\gamma(s_\ell),\qquad \ell\ge0.
 \label{eq:native-all-label-envelope}
\end{equation}
There is no representation cutoff in this comparison.
\end{proposition}
\begin{proof}
Differentiation gives
$F_\gamma'(s)=\tfrac12\int_0^1[\gamma^2+(1+s)u^2]^{-1/2}\dd u$.
The gamma integral and integration in $s$ give
\cref{eq:native-Bernstein-density}. Its nonnegative density satisfies
$\int(1\wedge t)n_\gamma(t)\dd t<\infty$. Compound-Poisson
approximations therefore construct a probability mixture with Laplace
transform $e^{-F_\gamma(s)}$. Since $F_\gamma(s)\to\infty$, the
mixture has no atom at zero. The fast decay of its group multipliers
gives a smooth strictly positive kernel.

For the scalar envelope, put $f(x)=\operatorname{arsinh}(x/\gamma)$.
The Bessel-ratio barriers used in \src{F16B}{V52--V53} imply
\[
 \sum_{\nu=1}^{\ell}f(\nu+1/2)
 \le\Phi_{\gamma,\ell}
 \le(1+4/(3\gamma))\sum_{\nu=1}^{\ell}f(\nu+1/2).
\]
Concavity and the midpoint remainder, using
$-f''\le\gamma^{-2}f$, bound this sum between
$F_\gamma(s_\ell)$ and
$(1+1/(8\gamma^2))F_\gamma(s_\ell)$.
Their product is at most $1+2/\gamma$ when $\gamma\ge2$.
\end{proof}

Retain the complete magnetic multiplication
$b_\gamma=e^{-\gamma S_{\rm spatial}/2}$ and put
\[
 T^{\rm B}_{\gamma,L}=b_\gamma R_\gamma^{\otimes3V}b_\gamma,
 \qquad T^{\rm W}_{\gamma,L}=b_\gamma C_\gamma^{\otimes3V}b_\gamma.
\]
Both act on the same compact gauge-invariant space. The superscript
$\rm B$ denotes this Bernstein reference, not an ordinary heat step.

\begin{theorem}[Full magnetic envelope and normalized transfer comparison]
\label{thm:native-Perron-comparison}
Let $\Lambda^{\rm B},\Lambda^{\rm W}$ be the respective actual
Perron eigenvalues and $E_0^{\rm B}=-\log\Lambda^{\rm B}$.
For every admitted finite lattice and $\gamma\ge2$,
\begin{align}
 &(T^{\rm B})^{p_\gamma}\preceq T^{\rm W}\preceq T^{\rm B},
       \qquad 0\le E_0^{\rm B}\le C_0V,
       \quad C_0=30+3\log48,\notag\\
 &\left\|\frac{T^{\rm W}}{\Lambda^{\rm W}}
                  -\frac{T^{\rm B}}{\Lambda^{\rm B}}\right\|
 \le\min\left\{1,\frac2\gamma(2E_0^{\rm B}+e^{-1})\right\}.
 \label{eq:native-Perron-comparison}
\end{align}
For $N$ steps the right side may be replaced by
$\min\{1,2N(2E_0^{\rm B}+e^{-1})/\gamma\}$.
The envelope holds on each center sector and every common orthogonal
compression, with powers taken on the compressed space.
\end{theorem}
\begin{proof}
Adding \cref{eq:native-all-label-envelope} over representation labels
gives $(R_\gamma^{\otimes3V})^{p_\gamma}
\preceq C_\gamma^{\otimes3V}\preceq R_\gamma^{\otimes3V}$.
Contractive Jensen for $x^p$, $1\le p\le2$, gives
$(BKB)^p\preceq BK^pB$ for positive contractions $B,K$.
Apply it to the complete magnetic sandwich. This argument does not
commute the magnetic and kinetic factors. The same inequality proves
the assertion for a common compression.

Write $A=T^{\rm B}$, $W=T^{\rm W}$, $a=\|A\|$, $w=\|W\|$,
and $\delta=2/\gamma$. The envelope gives
\[
 0\preceq A-W\preceq A-A^{1+\delta}
                  \preceq\delta A(-\log A),\qquad
 0\le a-w\le\delta aE_0^{\rm B}.
\]
Since $x(-\log x)\le a(E_0^{\rm B}+e^{-1})$ on $[0,a]$, use
\[
 \frac Ww-\frac Aa=\frac{W-A}{a}
                    +\left(\frac1w-\frac1a\right)W.
\]
The second term has norm $(a-w)/a$, so no reciprocal lower Perron
bound is lost. This proves the normalized estimate. The compact trial
bound of \src{F16B}{V51} is
$\Lambda^{\rm W}\ge e^{-C_0V}$ and proves the stated energy bound.
Telescoping the normalized contraction powers proves the $N$-step
estimate. Both powers are positive contractions, giving the unit cap.
\end{proof}

For completeness let $J^{\rm B}=-\log T^{\rm B}$ and similarly
$J^{\rm W}$. Operator monotonicity of the logarithm, with positive
regularization at zero, gives equality of their closed form domains and
\begin{equation}
 J^{\rm B}\preceq J^{\rm W}\preceq(1+2/\gamma)J^{\rm B}.
 \label{eq:native-log-form-envelope}
\end{equation}
After subtraction of their own bottom energies, the actual generators
satisfy
\begin{equation}
 H^{\rm B}-\frac{2E_0^{\rm B}}\gamma I
 \preceq H^{\rm W}
 \preceq(1+2/\gamma)H^{\rm B}+\frac{2E_0^{\rm B}}\gamma I.
 \label{eq:native-subtracted-envelope}
\end{equation}
The two vacuum vectors are not identified. The extensive vacuum term
is an explicit comparison debit, not a vanished normalization constant.

On the critical scale $h=\gamma^{-1/3}/L$,
\cref{thm:native-Perron-comparison} gives
\begin{equation}
 h^{-1}\left\|\widehat T^{\rm W}-\widehat T^{\rm B}\right\|
       \le C L^4\gamma^{-2/3},
 \label{eq:native-critical-comparison}
\end{equation}
where each hat uses its own Perron value. Thus if
$L^4\gamma^{-2/3}\to0$, their powers through bounded critical times
and their fixed near-top logarithmic levels agree asymptotically.
This compares two operators; an individual limiting spectrum still
requires recovery and no-escape estimates. The mixture variable in
\cref{eq:native-Bernstein-density} is not a physical time calibration.
\sourcebasis{The all-label subordination, magnetic power envelope,
closed-form comparison, and actual-Perron estimate are established in
\src{F16B}{V53}.}


\section{The actual fixed-volume critical spectrum}
\label{sec:native-critical-spectrum}
Fix $L=2m\ge8$, set $V=L^3$, and let $\gamma\to\infty$ in the
unchanged native Wilson transfer. Denote its Perron eigenvalue by
$\Lambda_\gamma$ and the top of its fixed-regulator hard Gaussian
limit, in the same normalization, by $\lambda_*$. Put
\begin{equation}
 h=(\gamma V)^{-1/3}.
 \label{eq:native-critical-scale}
\end{equation}
The soft critical variables are three imaginary-quaternion vectors
$q_\mu\in\R^3$. On their simultaneous colour-rotation invariant space,
the comparison Hamiltonian is
\begin{equation}
 H_{\rm mat}=-\tfrac12\Delta_q
              +2\sum_{\mu<\nu}|q_\mu\times q_\nu|^2,
 \qquad
 \mathcal H_{\rm mat}=L^2(\R^9)^{\mathrm{SO}(3)}.
 \label{eq:native-quartic-model}
\end{equation}
Its commuting valleys do not imply continuous spectrum. The transverse
oscillator comparison gives, in quadratic-form sense,
\[
 H_{\rm mat}\succeq2\sum_\mu|q_\mu|,
 \qquad
 H_{\rm mat}\succeq-\tfrac14\Delta_q+\sum_\mu|q_\mu|.
\]
Consequently its Friedrichs realization has compact resolvent and a
simple positive ground function, also after the invariant restriction.
Write its ordered invariant eigenvalues as
$e_0<e_1\le e_2\le\cdots$ and $g_*=e_1-e_0>0$.
This is the model confinement input of \src{F16}{V44}; it is not alone
an identification of native eigenvalues.

The physical transfer has eight center-character sectors, indexed by
$a\in(\Z/2\Z)^3$. Order their eigenvalues decreasingly as
$\lambda_{a,1},\lambda_{a,2},\ldots$, with
$\lambda_{0,1}=\Lambda_\gamma$ in the neutral sector. The based-gauge
tube and the extracted critical coordinates include the actual
half-density. The return metric of the regional Feshbach comparison
is kept; it is not replaced by the identity on the whole region.

\begin{theorem}[Ordered native spectrum at fixed spatial regulator]
\label{thm:native-critical-spectrum}
For every fixed $L$ and every fixed $j\ge1$,
\begin{equation}
 \frac{1-\lambda_{a,j}(\gamma)/\lambda_*}{h}
                \longrightarrow e_{j-1}
                   \qquad\text{in every center sector}.
 \label{eq:native-ordered-spectrum}
\end{equation}
In particular,
\begin{align}
 \Lambda_\gamma&=\lambda_*[1-he_0+o_L(h)],\notag\\
 -\log\frac{\lambda_{0,2}}{\Lambda_\gamma}
                      &=h g_*+o_L(h).
 \label{eq:native-neutral-gap}
\end{align}
For any fixed threshold outside the invariant model spectrum, the
number of scaled native deficits below that threshold is eventually
exactly the corresponding model count in each center sector.
\end{theorem}
\begin{proof}
The physical recovery maps give the ordered min--max upper bound
using finitely many invariant model tests. They preserve the exact
based-gauge half-density and reconstruct the same native top.
For the reverse bound, the native band estimates give capture by the
hard Gaussian ground, local compactness of the soft profiles, and a
corrected profile equation. The linear hard variation has zero ground
expectation but is not the zero operator: its off-diagonal return is
retained in that equation.

The remaining issue is loss of mass at large soft position. The
native annular and localization estimates give, for each fixed band
parameter $K$, a bound of the form
\[
 \limsup_{\gamma\to\infty}
 \sup_{\substack{\psi\text{ in the native }K\text{-band}\\\norm\psi=1}}
 \int_{|q|>4R}|\mathfrak E_\gamma\psi(u,q)|^2\dd u\dd q
                    \le C_L(K/R+R^{-3}).
\]
Together with the native semiboundedness estimate, this supplies
global strong compactness without an annulus-counting loss.
Extract a finite ordered orthonormal family. Asymptotic isometry
preserves its Gram matrix, and the corrected equation makes every
limit a model eigenfunction. Model min--max gives the ordered lower
bound. Sector and one-region comparison errors are exponentially
small at fixed $L$, hence $o(h)$, proving
\cref{eq:native-ordered-spectrum} in every sector. Subtract the first
two neutral levels and expand the logarithm to obtain
\cref{eq:native-neutral-gap}. The ordered limits on both sides of a
fixed spectral threshold give the exact eventual count.
The native recovery, hard capture, profile equation, and radial
estimate are established successively in \src{F16}{V41--V48}.
\end{proof}

\begin{corollary}[Critical vacuum and two finite-volume spectral scales]
\label{cor:native-vacuum-two-scales}
After the specified center unfolding and critical extraction,
\begin{equation}
 \mathfrak E_{\gamma,0}\Omega_\gamma
       \longrightarrow\varphi_{\rm hard}\otimes\phi_0
                       \quad\text{in }L^2,
 \label{eq:native-critical-vacuum}
\end{equation}
where $\phi_0$ is the positive unit ground function of
\cref{eq:native-quartic-model}. The extracted squared densities converge
in $L^1$. For sufficiently large $\gamma$, exactly eight native
physical eigenvalues lie above
$\Lambda_\gamma-\tfrac12\lambda_*hg_*$, one in each center sector.
Their charged top splittings obey fixed-$L$ bounds
$C_Le^{-c_L\sqrt\gamma}$, while the next neutral level is separated
by $\lambda_*hg_*+o_L(h)$.
\end{corollary}
\begin{proof}
Strong compactness and the profile equation identify all ground
subsequential limits with the simple positive model ground. Positivity
removes the phase ambiguity, so the whole sequence converges.
The $L^1$ statement follows from
$\norm{|f|^2-|g|^2}_1\le(\norm f_2+\norm g_2)\norm{f-g}_2$.
The first two ordered limits in every sector, together with the
fixed-regulator center-splitting estimate, give the count and the
two scales. No matching tunneling lower bound is used.
\end{proof}

\begin{certificate}[A rigorous bracket for the model gap]
\label{cert:quartic-model-gap}
The sector-resolved finite calculations and spectral comparison
bounds recorded in \src{F16}{V51, V57, V60} give
\[
 6.5322\le e_0\le6.5350,\qquad e_1\ge6.936064,
 \qquad 0.4011\le g_*\le3.0156.
\]
The lower and upper endpoints use different certified inputs. A
floating approximation near $3.01$ is not a certified identification
of the first excited sector or a replacement for this interval.
\end{certificate}

\begin{scope}[Fixed regulator, center sectors, and physical units]
The limit in \cref{thm:native-critical-spectrum} fixes $L$ and $j$
before taking $\gamma\to\infty$. Its error constants are not uniform
on a growing spatial lattice. The neutral critical gap is not the
smallest gap on the union of the eight center sectors, whose tops are
exponentially close. The extracted vacuum statement is not an unscaled
$L^2$ limit on the original compact coordinates. Finally a lattice
logarithmic deficit $hg_*+o_L(h)$ becomes a physical energy only after
the temporal spacing is independently calibrated. These results are
actual native spectral asymptotics, but not an interacting continuum
mass-gap theorem.
\end{scope}

\section{The joint-regulator critical branch}
\label{sec:native-joint-regulator}
The preceding comparison can be combined with measured recovery and
localization estimates on a growing slice. Retain the quartic model
$H_{\rm mat}$, its invariant eigenvalues $e_j$, and the center sectors
of \cref{sec:native-critical-spectrum}. Define the explicit joint budget
\begin{equation}
 \Theta_{\gamma,L}
    =L^2\gamma^{-1/3}+L^4\gamma^{-2/3}+L^3\gamma^{-1}.
 \label{eq:native-joint-budget}
\end{equation}
The three terms record soft--hard return, comparison and geometric
errors on the critical scale, and extensive hard nonlinear errors,
respectively. They are not replaced by fixed-$L$ remainder notation.

A useful hard-mode estimate retains the full Fourier weights. Set
\[
 \omega_L(k)=2\operatorname{arsinh}(\sqrt{\nu_L(k)}/2),
 \qquad d_L(k)=1-e^{-\omega_L(k)}.
\]
Then
\begin{equation}
 \sum_{k\ne0}d_L(k)^{-1}\le CL^3,
 \qquad \sum_{k\ne0}d_L(k)^{-2}\le CL^3.
 \label{eq:native-hard-susceptibilities}
\end{equation}
Indeed $d_L(k)$ is comparable to
$\min\{1,|\widetilde k|/L\}$, and three-dimensional shell summation
bounds both sums. Application to a particular source also requires its
Fourier coefficient bounds; the smallest hard gap cannot simply be
charged once to every mode.

\begin{theorem}[Ordered critical spectrum on a growing regulator]
\label{thm:native-joint-spectrum}
On the measured critical branch with
$\gamma\to\infty$, even $L=L(\gamma)\ge8$, and
$\Theta_{\gamma,L}\to0$, the original Wilson transfer satisfies, for
every fixed center character $a$ and fixed index $j\ge1$,
\begin{equation}
 -\frac1h\log\frac{\lambda^{\rm W}_{a,j}(\gamma,L)}
                         {\Lambda^{\rm W}_{\gamma,L}}
       \longrightarrow e_{j-1}-e_0,
 \qquad h=\gamma^{-1/3}/L.
 \label{eq:native-joint-spectrum}
\end{equation}
In particular the neutral first excited level has limiting coefficient
$e_1-e_0>0$. A sufficient family is
$L=O(\gamma^b)$ with fixed $0\le b<1/6$.
The statement uses the measured recovery and all-vector no-escape
estimates on this branch, not a frozen Gaussian trial calculation.
\end{theorem}
\begin{proof}
We specify the analytic inputs to the spectral passage. In the exact
based tube of \cref{thm:native-uniform-geometry}, the hard-return
corrected recovery maps of \src{F16B}{V54} satisfy on each finite
invariant model graph core
\begin{equation}
 \left\|\widehat T^{\rm B}\Psi_{\gamma,L}f
       -\Psi_{\gamma,L}\{f-h(H_{\rm mat}-e_0)f\}\right\|
       \le C_fh\Theta_{\gamma,L},
 \label{eq:native-joint-recovery}
\end{equation}
with Gram error $O_f(\Theta_{\gamma,L})$. Their volume accounting uses
the relative density and trace-ideal metric estimates, together with
\cref{eq:native-hard-susceptibilities} applied to the actual finite-incidence
sources. The accompanying no-escape estimate controls every unit vector
in a fixed actual-top band $[e^{-Kh},1]$: the mass outside the measured
central tubes and the hard-ground complement tends to zero, while the
extracted soft profiles have uniformly vanishing position and Fourier
tails. These are the separate recovery and compactness inputs proved
for the subordinated reference in \src{F16B}{V54}.

They yield collective compactness and the model resolvent limit after
the measured extraction. Testing against the corrected graph core
identifies every finite limiting cluster with $H_{\rm mat}-e_0$;
min--max and compactness give the ordered reference spectrum, including
multiplicities. The reference uses its actual Perron value throughout.
By \cref{eq:native-critical-comparison}, the original Wilson operator
has normalized error $o(h)$ on this joint domain. Min--max in each
common center sector transfers the fixed near-top eigenvalues.
The logarithm is Lipschitz on a fixed neighborhood of one, proving
\cref{eq:native-joint-spectrum}. The power condition follows directly
from \cref{eq:native-joint-budget}.
\end{proof}

The theorem advances the fixed-regulator statement without changing its
physical meaning. It concerns a growing number of lattice sites, not
an already calibrated fixed physical box. In particular, a possible
smaller splitting between center-sector tops is not a positive local
mass. The physical scale comparison in
\cref{sec:native-bare-calibration} and the source-invisibility test in
\cref{sec:local-toron-Grams} must be imposed separately.
\sourcebasis{The volume-uniform geometry and winding estimates are
in \src{F16B}{V49--V50}; the corrected joint recovery, hard capture,
soft no-escape estimate, and spectral passage are in
\src{F16B}{V54}. The passage from the subordinated reference to Wilson
uses \cref{thm:native-Perron-comparison}.}


\section{Local kinetic control without a volume-dependent chart radius}
\label{sec:native-supnorm}
The local algebra suggests a route beyond a deterministic global
$\ell^2$ tube, but it does not complete that route. Use imaginary-
quaternion coordinates $U=\operatorname{Exp}x$, with round metric
$|x|$. For $|x|,|y|\le r\le1$,
\begin{equation}
 \frac{\sin r}{r}|x-y|
 \le d(\operatorname{Exp}x,\operatorname{Exp}y)\le|x-y|.
 \label{eq:native-chart-distance}
\end{equation}
The ball is geodesically convex; the differential of the exponential
has radial singular value one and tangential singular value
$\sin|x|/|x|$. Applying these bounds to straight paths and inverse
images of minimizing geodesics proves
\cref{eq:native-chart-distance}.

\begin{proposition}[Per-link Gaussian domination and normalized Haar]
\label{prop:native-supnorm-kernel}
Define
\[
 1-\delta(r)=\left(\frac{\sin r}{r}\right)^2(1-r^2/3).
\]
On a product chart with $\max_e(|x_e|,|x'_e|)\le r\le1$,
\begin{equation}
 z_\gamma^{-E}e^{-\frac\gamma2\sum_e|x_e-x'_e|^2}
 \le\prod_{e=1}^E k_\gamma(U_eU_e'^{-1})
 \le z_\gamma^{-E}
            e^{-\frac\gamma2(1-\delta(r))\sum_e|x_e-x'_e|^2}.
 \label{eq:native-supnorm-sandwich}
\end{equation}
The error coefficient is independent of $E$ and
$0\le\delta(r)\le\tfrac23r^2+\tfrac15r^4$. Normalized Haar is
\begin{equation}
 \dd H(\operatorname{Exp}x)=c_HJ(x)\dd x,\quad
 c_H=\frac1{2\pi^2},\quad
 J(x)=\left(\frac{\sin|x|}{|x|}\right)^2,
 \label{eq:native-normalized-Haar}
\end{equation}
and, for $|x|\le1$,
\begin{equation}
 -|x|^2/3-|x|^4/45\le\log J(x)\le-|x|^2/3.
 \label{eq:native-Haar-remainder}
\end{equation}
\end{proposition}
\begin{proof}
For $d\le2r$, Taylor's inequality gives
$(d^2/2)(1-r^2/3)\le1-\cos d\le d^2/2$.
Combine it with \cref{eq:native-chart-distance} and multiply the exact
one-link kernels. The bound on $\delta$ follows from the sine series
on $[0,1]$. The round volume of the unit three-sphere is $2\pi^2$,
which fixes $c_H$ for normalized Haar. The convergent expansion
\[
 \log J(x)=-|x|^2/3-|x|^4/90-2|x|^6/2835-\cdots
\]
and its tail bound give \cref{eq:native-Haar-remainder}.
\end{proof}

For clarity, in flat product coordinates the full fine-link kernel
therefore has the exact form
\begin{align}
 T_\gamma^{\rm flat}(x,x')
 ={}&(c_H/z_\gamma)^E b_\gamma(x)b_\gamma(x')\rho(x,x')\notag\\
 &\quad\times\exp\left[-\frac\gamma2\sum_e(1-\delta_e)|x_e-x'_e|^2
                     -\frac16\sum_e(|x_e|^2+|x'_e|^2)\right],\notag\\
 &0\le\delta_e\le\delta(r),\qquad
 \exp\left[-\frac1{90}\sum_e(|x_e|^4+|x'_e|^4)\right]
                      \le\rho(x,x')\le1,
 \label{eq:native-flat-factorization}
\end{align}
where $b_\gamma=e^{-\gamma S_{\rm spatial}/2}$. Both factors
$c_H$ and $z_\gamma$ belong to the normalization. They are
source-independent scalar factors and cancel in spectral ratios,
whereas the sums in the exponent can be extensive. A dimension-free
per-link coefficient is not a volume-uniform relative operator error;
control of the fluctuating and field-dependent sums is a separate
analytic obligation.

\subsection{Gaussian maxima and the scope of a sup-norm region}
\label{sec:native-Gaussian-maxima}
Let $y=\gamma^{-1/2}u$ have the frozen transverse hard Gaussian law.
Assume each color component has variance at most $\sigma^2/\gamma$,
with $\sigma$ independent of $L$. There are $9V$ such components, so
no independence assumption is needed for the union bound
\begin{equation}
 \mathbb P\left\{\max_{e,a}|y_e^a|>
 \sigma\sqrt{\frac{2(\log(18V)+\lambda)}\gamma}\right\}
 \le e^{-\lambda},\qquad \lambda\ge0.
 \label{eq:native-Gaussian-maximum}
\end{equation}
The corresponding expectation is bounded by
$\sigma\sqrt{2\log(18V)/\gamma}$. The component maximum in this
formula is not the round link radius: a valid Euclidean link radius
is $\sqrt3$ times the displayed threshold. In coordinates
$x_e=y_e+hq_{\mu(e)}$, it must also include
$h\max_\mu|q_\mu|$.

Thus a frozen Gaussian configuration has small individual link angles
whenever $\log V=o(\gamma)$, even when its total $\ell^2$ norm need
not be small. This identifies a useful probabilistic region, not a
new gauge slice or a native-vacuum concentration theorem. The uniform
slice in \cref{thm:native-uniform-geometry} is proved on a total
$\ell^2$ tube; uniqueness on the larger sup-norm region is a separate
geometric requirement. Likewise a reference-Gaussian tail cannot be
substituted for the Doob-chain bounds in
\cref{eq:native-uniform-packet}.

\subsection{An exact mean-one kinetic factor}
\label{sec:native-mean-one}
A second improvement removes a deterministic local normalization
before any extensive estimate. Write $\theta(U)=d(U,I)$ and define,
with respect to normalized Haar,
\begin{equation}
 g_\gamma(U)=Z_\gamma^{-1}e^{-\gamma\theta(U)^2/2},\qquad
 k_\gamma(U)=z_\gamma^{-1}e^{-\gamma(1-\cos\theta(U))}.
 \label{eq:native-geodesic-reference}
\end{equation}
Their exact ratio is
\begin{equation}
 \rho_\gamma(U)=\frac{Z_\gamma}{z_\gamma}
 \exp\left\{\gamma\left(\frac{\theta(U)^2}2-1+
                                      \cos\theta(U)\right)\right\},
 \qquad \int g_\gamma\rho_\gamma\,\dd H=1.
 \label{eq:native-mean-one-ratio}
\end{equation}
In particular the full magnetic Wilson kernel is the full magnetic
$g_\gamma$-kernel multiplied by
$\prod_e\rho_\gamma(U_eU_e'^{-1})$. This is an exact kernel identity,
not an identification of the geodesic reference with a heat semigroup.

\begin{proposition}[One-link centering and fluctuation scale]
\label{prop:native-mean-one}
For $c_\gamma=(2\pi/\gamma)^{3/2}/(2\pi^2)$,
\begin{align}
 z_\gamma/c_\gamma
   &=1-\frac{3}{8\gamma}-\frac{15}{128\gamma^2}
                                      +O(\gamma^{-3}),\notag\\
 Z_\gamma/c_\gamma
   &=1-\frac1\gamma+\frac{2}{3\gamma^2}
                                      +O(\gamma^{-3}),\notag\\
 \mathbb E_{g_\gamma}\log\rho_\gamma
   &=-\frac{5}{8\gamma^2}+O(\gamma^{-3}),\qquad
 \operatorname{Var}_{g_\gamma}(\log\rho_\gamma)
     =\frac{5}{4\gamma^2}+O(\gamma^{-3}).
 \label{eq:native-mean-one-moments}
\end{align}
On $\theta\le2r$ the uncentered exponent satisfies
\begin{equation}
 0\le\gamma(\theta^2/2-1+\cos\theta)
       \le\gamma\theta^4/24\le\tfrac23\gamma r^4.
 \label{eq:native-ratio-local-remainder}
\end{equation}
\end{proposition}
\begin{proof}
The radial Haar density is $(2/\pi)\sin^2\theta\,\dd\theta$.
Rescale $\theta=\gamma^{-1/2}t$, expand the sine and cosine on a fixed
small ball, and bound the exterior by its exponential tail. Gaussian
radial moments give the first two lines. Expanding
$\log(Z_\gamma/z_\gamma)+\gamma(\theta^2/2-1+\cos\theta)$ under
the normalized $g_\gamma$ law gives the last line; the order
$\gamma^{-1}$ mean cancels. The global scalar inequalities
$0\le\theta^2/2-1+\cos\theta\le\theta^4/24$ give the local bound.
\end{proof}

The centering in \cref{eq:native-mean-one-ratio} is with respect to
one isolated normalized kinetic increment. Complete magnetic weights,
a transverse constraint, or a Perron tilt change that law. Therefore
\cref{eq:native-mean-one-moments} is not, without a connected estimate
under the new law, a volume-uniform estimate for a product of factors
or for the actual normalized transfer. Removing a scalar mean and
controlling its field-dependent fluctuations are different operations.

\subsection{The cubic magnetic term and its Wiener-chaos structure}
\label{sec:native-cubic-chaos}
The leading odd magnetic term must be retained in any such connected
estimate. For a plaquette with oriented links
$e^{y_1}e^{y_2}e^{-y_3}e^{-y_4}$, put
\begin{align}
 F(y)&=y_1+y_2-y_3-y_4,\notag\\
 B(y)&=y_1\mathbin\times y_2-y_1\mathbin\times y_3
       -y_1\mathbin\times y_4-y_2\mathbin\times y_3
       -y_2\mathbin\times y_4+y_3\mathbin\times y_4.
 \label{eq:native-plaquette-cubic-polynomial}
\end{align}
The quaternion convention of this chapter gives the exact local
expansion
\begin{equation}
 s(U_p)=\tfrac12|F(y)|^2+F(y)\cdot B(y)+R_{4,p}(y),\qquad
 |R_{4,p}(y)|\le C\max_i|y_i|^4
 \label{eq:native-plaquette-cubic-expansion}
\end{equation}
on a fixed small link ball. It follows either by the degree-two
Baker--Campbell--Hausdorff expansion, for which
$[u,v]=2u\times v$, or by direct quaternion multiplication. Thus the
half-magnetic cubic factor
$\xi_p=(\gamma/2)F(y)\cdot B(y)$ is bounded by $C\gamma r^3$,
and its fourth-order action remainder by $C\gamma r^4$. At a
Gaussian maximum radius these are local activities of orders
$(\log V)^{3/2}/\sqrt\gamma$ and $(\log V)^2/\gamma$; their
sums over plaquettes are not dimension-free bounds.

\begin{proposition}[Exact covariance accounting for the cubic term]
\label{prop:native-cubic-Wick}
Suppose the frozen Gaussian covariance is color-isotropic,
$\mathbb E[y_e^a y_{e'}^b]=\gamma^{-1}\delta_{ab}C_{ee'}$.
Then each $\xi_p$ is a centered homogeneous third Wiener-chaos
variable. With $Y=\sqrt\gamma\,y$, write
$W_{x,\alpha}=F(Y)\cdot B(Y)$ for the plaquette of spatial
orientation $\alpha$ based at $x$. Translation invariance gives
\begin{equation}
 \operatorname{Var}\left(\sum_p\xi_p\right)
   =\frac{V}{4\gamma}
      \sum_{\alpha,\beta}\sum_z
       \mathbb E[W_{0,\alpha}W_{z,\beta}].
 \label{eq:native-cubic-complete-covariance}
\end{equation}
All orientation pairs occur. If the frozen link covariance obeys
$|C_{ee'}|\le K_L(1+\operatorname{dist}(e,e'))^{-2}$, then
\begin{equation}
 \left|\mathbb E[W_{0,\alpha}W_{z,\beta}]\right|
       \le C K_L^3(1+\operatorname{dist}(0,z))^{-6},
 \qquad
 \operatorname{Var}\left(\sum_p\xi_p\right)
       \le C\frac{VK_L^3}{\gamma}.
 \label{eq:native-cubic-Wick-bound}
\end{equation}
\end{proposition}
\begin{proof}
Every monomial in the cubic polynomial contracts its color indices
with the alternating tensor. Each internal Gaussian contraction
therefore contains $\epsilon_{abc}\delta_{ab}$, or its permutation,
and vanishes. The polynomial is its own Wick-ordered cubic part.
The product of two such parts contains only complete cross-pairings;
each is a product of three link covariances. Finite plaquette incidence
gives the first bound in \cref{eq:native-cubic-Wick-bound}, and
three-dimensional shell summation gives the second. Expanding the
variance and translating one plaquette to the origin proves
\cref{eq:native-cubic-complete-covariance}.
\end{proof}

\begin{scope}[What the cubic calculation does not establish]
A cubic polynomial of a Gaussian field is not itself Gaussian.
Consequently a variance computation does not justify an exact
lognormal identity for its exponential or a lower bound for the full
transfer norm. An uncut exponential of the isolated cubic term also
cannot replace the stabilizing higher magnetic orders. Moreover a
bound with $K_L=O(1+\log L)$ does not by itself produce a
volume-independent limiting covariance density. Such a density
requires a sharper summable bound, with cross-orientation terms
retained. The valid conclusions here are the local expansion,
its third-chaos structure, and the stated covariance estimate.
An all-order connected estimate under the actual magnetic and
Perron-weighted law remains a separate obligation.
\end{scope}
\sourcebasis{The exact Perron identity and normalized kinetic
factor are developed in \src{F16}{V65--V66}; the local cubic
expansion and Wick structure are in \src{F16}{V67}. The covariance
sum above retains all plaquette orientations and distinguishes
third Wiener chaos from a Gaussian random variable.}


\section{A quantitative near-top reduction and its missing inputs}
\label{sec:native-Feshbach-reduction}
Let $P=T_\gamma/\Lambda_\gamma$ on the physical Haar space. It is a
positive self-adjoint contraction. Multiplication by the positive
Perron vector gives its unitary realization as the reversible Doob
chain on $\Omega_\gamma^2\dd H$. Region projections commute with this
unitary, so the following norms can be computed in either realization.

\begin{theorem}[Same-top regional reduction]
\label{thm:native-regional-Feshbach}
Let $A$ be a measurable physical region, $Q=\one_{A^c}$, and assume
$\norm{QPQ}\le1-\kappa$ with $0<\kappa\le1$. For
$1-\kappa/2<z\le1$, $z$ is an eigenvalue of $P$ if and only if it
is an eigenvalue of
\begin{equation}
 F_A(z)=\one_AP\one_A+
         \one_APQ(z-QPQ)^{-1}QP\one_A
                    \quad\text{on }L^2(A),
 \label{eq:native-regional-Feshbach}
\end{equation}
with equal multiplicity. Moreover
\begin{equation}
 \norm{F_A(z)-\one_AP\one_A}
                         \le\frac{2\norm{\one_APQ}^2}{\kappa}.
 \label{eq:native-regional-return-bound}
\end{equation}
For the reversible Doob kernel,
\begin{align}
 \norm{\one_APQ}^2
 &\le\sup_{U\in A}P(U,A^c)\,
                   \sup_{U\in A^c}P(U,A),\notag\\
 \norm{QPQ}&\le\sup_{U\in A^c}P(U,A^c).
 \label{eq:native-Doob-Schur}
\end{align}
Essential suprema are sufficient.
\end{theorem}
\begin{proof}
The complement resolvent exists with norm at most $2/\kappa$.
Solving the complement component of the eigenvalue equation gives
\cref{eq:native-regional-Feshbach}, including the reconstruction map
and multiplicities. The resolvent bound gives
\cref{eq:native-regional-return-bound}. Reversibility identifies the
column and row integrals needed for the Schur test, proving
\cref{eq:native-Doob-Schur}.
\end{proof}

For an approximation on the scale $h$, a sufficient quantitative
packet is
\begin{equation}
 \kappa\gg h,\qquad \norm{\one_APQ}^2=o(\kappa h),
 \qquad\text{and an }o(h)\text{ comparison of the retained pencil}.
 \label{eq:native-uniform-packet}
\end{equation}
A small unconditional probability of $A^c$ proves neither of the
first two estimates. Nor does a Gaussian jump estimate from a strictly
interior subregion prove the boundary supremum over all $A$ in
\cref{eq:native-Doob-Schur}. Buffered regions, weighted cutoffs, or a
direct operator estimate require their own argument. The complement
can contain long-lived states even when local displacements are small.

\begin{proposition}[A native but insufficient Perron-ratio bound]
\label{prop:native-crude-Harnack}
If $U'$ differs from $U$ in one spatial link by round angle $\theta$,
then the positive native Perron vector satisfies
\begin{equation}
 e^{-3\gamma\theta}\le
     \frac{\Omega_\gamma(U')}{\Omega_\gamma(U)}
                              \le e^{3\gamma\theta}.
 \label{eq:native-crude-Harnack}
\end{equation}
\end{proposition}
\begin{proof}
Use the positive Perron integral equation. Changing one endpoint link
changes its kinetic action by at most $\theta$. Four spatial
plaquettes contain that link; the half magnetic action therefore
changes by at most $2\theta$. The ratio of every positive integrand
lies between the stated bounds, as does the ratio of their integrals.
\end{proof}
A quadratic-scale estimate such as
\[
 \left|\log\frac{\Omega_\gamma(U')}{\Omega_\gamma(U)}\right|
          \le c_1\sqrt\gamma\,\theta+c_2\gamma\theta^2
\]
on a suitable common region would improve the local jump comparison.
It has not been established here. Even such a local estimate would
have to be combined with complement control and boundary transport
to verify \cref{eq:native-uniform-packet}; it is not by itself a
proof of those global hypotheses. The joint theorem
\cref{thm:native-joint-spectrum} already controls one restricted
growing-volume domain; the regional reduction concerns continuation
beyond it. Neither that extension nor an independent physical clock
is supplied by a local Perron bound.
\sourcebasis{The sup-norm kinetic algebra and native Perron estimate
are developed in \src{F16}{V64}. The regional implication above keeps
its complement and boundary hypotheses explicit.}


\subsection{An exact transport equation for the Perron ratio}
\label{sec:native-Perron-transport}
The desired improvement of \cref{eq:native-crude-Harnack} has an
exact starting identity. Write the full transfer kernel as
$T(U,U')=b(U)c(U,U')b(U')$, and let $\tau$ preserve Haar measure
and the simultaneous kinetic kernel:
$c(\tau U,\tau U')=c(U,U')$. For the positive Perron function
$\Omega$ define
\[
 w(U)=\frac{\Omega(\tau U)}{\Omega(U)},\qquad
 \rho(U)=\frac{b(\tau U)}{b(U)}.
\]
Here $\tau$ may be a simultaneous fixed link rotation on the full
configuration space; it need not preserve a chosen gauge chart.

\begin{proposition}[Multiplicative and infinitesimal Perron transport]
\label{prop:native-Perron-transport}
For the actual Perron Doob kernel $P$,
\begin{equation}
 w(U)=\rho(U)\int P(U,\dd U')\rho(U')w(U').
 \label{eq:native-Perron-transport}
\end{equation}
Iteration gives, for $n\ge1$,
\begin{equation}
 w(U)=\mathbb E_U\left[
  \rho(X_0)\left\{\prod_{r=1}^{n-1}\rho(X_r)^2\right\}
  \rho(X_n)w(X_n)\right].
 \label{eq:native-Perron-path-transport}
\end{equation}
For a differentiable one-parameter family with
$\rho_\theta=e^{-\theta s+O(\theta^2)}$ and
$g=\left.\partial_\theta\log w_\theta\right|_{\theta=0}$,
\begin{equation}
 (I-P)g=-(I+P)s.
 \label{eq:native-Perron-Poisson}
\end{equation}
For the Wilson magnetic factor,
$s=(\gamma/2)\partial_\theta S_{\rm spatial}(\tau_\theta U)|_0$.
\end{proposition}
\begin{proof}
Apply the positive Perron integral equation at $\tau U$ and change
variables $U'=\tau V$. Kinetic invariance leaves the two magnetic
ratios and the Perron ratio at $V$. Division by
$\Lambda\Omega(U)$ proves \cref{eq:native-Perron-transport}.
The Markov property gives the iterated identity, and differentiation
at zero gives $g=-s+P(g-s)$.
\end{proof}

This formula keeps the magnetic forcing and the actual reversible law
in the same equation. It does not yet bound the inverse of $I-P$ on
this forcing. A local derivative estimate must be obtained without
assuming the very global soft spectral gap that it is intended to
establish. In particular, replacing the native law by an independent
Gaussian law or invoking an unspecified uniform Poisson bound would
not pay the Perron-ratio requirement of the regional reduction.
\sourcebasis{The exact multiplicative identity and its
infinitesimal equation are established in \src{F16}{V65}.}



\section{Electric-dominated native spectral control}
\label{sec:native-electric-control}
The critical-spectrum results above organize a weak-magnetic scaling of
the compact transfer.  A complementary calculation works directly with
the finite spatial Hamiltonian
\begin{equation}\label{eq:native-electric-H}
 H=T+b|\mathcal P_N|+V,\qquad
 T=\kappa\sum_e L_e,\qquad
 V=-b\sum_{p\in\mathcal P_N}w_p,
 \qquad w_p=\tfrac12\Tr U_p,
\end{equation}
where $L_e$ is the positive sphere Casimir.  The additive term
$b|\mathcal P_N|$ has no effect on spectral gaps.  This formulation is
useful because the electric spectrum supplies an exact local energy
grade while the Wilson interaction remains the original bounded
plaquette multiplication operator.  It is not a new physical clock and
it does not identify the continuum trajectory.

\subsection{An actual volume-uniform lattice gap}
Let $\Delta=3\kappa$ be the first nonzero one-link electric energy and
let $a\ge0$ be the connected-support weight used in the creation norm.
The nonlinear creation equation of \src{F18}{V5} keeps the original
Wilson interaction inside a finite-degree commutator algebra rather than
iterating an ever growing effective action.

\begin{theorem}[Electric-dominated native ground state and gap]
\label{thm:native-electric-gap}
If
\begin{equation}\label{eq:native-electric-smallness}
 \frac{be^a}{\Delta}\le\frac1{1024},
\end{equation}
then the native finite-lattice Hamiltonian \cref{eq:native-electric-H}
has a simple gauge-invariant ground state and
\begin{equation}\label{eq:native-electric-gap}
 \gap(H)\ge\frac{25}{32}\Delta=\frac{75}{32}\kappa .
\end{equation}
The lower bound is uniform in the side length of the spatial torus and
remains valid after restriction to the gauge-invariant subspace.
\end{theorem}
\begin{proof}[Proof mechanism]
The creation map is a contraction with constant at most $7/32$ on the
closed weighted ball of radius $32be^a$.  Its fixed point $C_*$ gives an
actual eigenvector $e^{C_*}\Omega_{\rm el}$.  The derivative of the
same nonlinear map descends to the quotient of exact excitation
sectors and has norm at most $7/32$, while $T^{-1}$ costs at most
$\Delta^{-1}$.  An independent eigenvalue at displacement $\delta$
would therefore force
$1\le7/32+|\delta|/\Delta$.  Continuation from $b=0$ identifies the
constructed simple eigenvalue with the ground state and yields
\cref{eq:native-electric-gap}.  \src{F18}{V5}
\end{proof}

The regime is deliberately explicit: at $a=0$ it requires
$b/\kappa\le3/1024$.  Thus \cref{thm:native-electric-gap} is a genuine
volume-uniform lattice gap, but in an electric-dominated domain.  No
claim is made that the asymptotically free continuum trajectory remains
inside this domain after physical calibration.

\subsection{The first physical plaquette band}
The same construction admits a sharper, source-resolved spectral
statement.  Put
\begin{equation}\label{eq:native-band-radius}
 \rho_a=\frac{\kappa e^{-a}}{2688}.
\end{equation}
For $|b|<\rho_a$, the complete physical band continuing the electric
plaquette eigenspace is represented by an analytic matrix $F_N(b)$ in
the physical plaquette basis.  If $d(p,q)$ is plaquette adjacency
distance, its matrix columns are exponentially summable uniformly in
volume.

\begin{theorem}[Localized analytic plaquette band]
\label{thm:native-plaquette-band}
For $|b|<\rho_a$ the continued plaquette band is holomorphic and, for
real $0\le b<\rho_a$, consists of the first $|\mathcal P_N|$ nonzero
physical eigenvalues.  Its expansion is
\begin{equation}\label{eq:native-band-expansion}
 F_N(b)=12\kappa I+
 \frac{b^2}{\kappa}\left(-\frac{71}{1680}I+
                  \frac{\mathcal L_N^{\rm pl}}{336}\right)
 +\mathcal R_N(b),
\end{equation}
where $\mathcal L_N^{\rm pl}$ is the plaquette graph Laplacian.  With
$\tau=|b|/\rho_a$,
\begin{equation}\label{eq:native-band-remainder}
 \sup_p\sum_q e^{a d(p,q)}|\mathcal R_N(b)_{qp}|
 \le\varepsilon_a(b),\qquad
 \norm{\mathcal R_N(b)}_{2\to2}\le\varepsilon_a(b),
 \quad
 \varepsilon_a(b)=\frac{2\kappa e^a}{3}
                     \frac{\tau^3}{1-\tau}.
\end{equation}
Consequently the physical finite-lattice gap $\Delta_N$ satisfies
\begin{equation}\label{eq:native-band-gap}
 \left|\Delta_N-
 \left(12\kappa-\frac{71b^2}{1680\kappa}\right)\right|
 \le\varepsilon_a(b),
\end{equation}
with constants independent of $N$.
\end{theorem}
\begin{proof}[Proof mechanism]
A uniformly convergent creation-family resolvent construction gives a
Riesz projection and a spatially weighted matrix norm on the entire
continued band.  Exact second-order degenerate perturbation theory
retains both the disconnected vacuum return and the shared-link
channel, yielding the coefficient in
\cref{eq:native-band-expansion}.  Banach-valued Cauchy estimates sum the
higher coefficients geometrically in the same rooted norm, which gives
\cref{eq:native-band-remainder}.  Translation and cubic symmetry turn
the weighted row and column bounds into the $\ell^2$ estimate; the
actual spectral identification then gives
\cref{eq:native-band-gap}.  \src{F18}{V12}
\end{proof}

This band theorem is stronger than a formal strong-coupling series:
its remainder and spatial tail are uniform in the torus size.  It is
still a lattice statement in the small ratio $b/\kappa$ and therefore
serves as a controlled anchor, not as the sought weak-coupling
continuum result.

\subsection{High-harmonic Schur control for a whole spectral window}
Small $b/\kappa$ is not needed to control sufficiently high one-link
harmonics.  Let $P_\ell$ project one selected link onto sphere harmonic
degree $\ell$, $F_L=\sum_{\ell=0}^LP_\ell$, $Q_L=I-F_L$, and set
\begin{equation}\label{eq:native-high-dL}
 \mu_L=(L+1)(L+3),\qquad d_L=\kappa\mu_L-4b,
 \qquad K=H-E_0.
\end{equation}
The exterior field is never frozen in the following statements.

\begin{theorem}[Native high block, Schur return, and factorial tails]
\label{thm:native-high-Schur}
If $d_L>0$, then
\begin{equation}\label{eq:native-high-block}
 C_L:=Q_LKQ_L\succeq d_LI,
 \qquad B_L:=Q_LKF_L=P_{L+1}U_eP_L,
 \qquad \norm{B_L}\le2b,
\end{equation}
where $U_e$ is the complete bounded interaction incident on the selected
link.  For real $z<d_L$,
\begin{equation}\label{eq:native-high-return}
 \Sigma_L(z)=B_L^*(C_L-z)^{-1}B_L,
 \qquad
 0\preceq\Sigma_L(z)\preceq\frac{4b^2}{d_L-z}I.
\end{equation}
The recursive truncation of the high return at degree $R$ has error
\begin{equation}\label{eq:native-high-return-tail}
 \norm{\Sigma_L(z)-\Sigma_L^{[R]}(z)}
 \le \frac{4b^2}{d_R-z}
       \prod_{\ell=L}^{R-1}
          \frac{4b^2}{(d_\ell-z)^2},
\end{equation}
uniformly in lattice volume.

If $\mathcal P_*=\one_{[0,E_*]}(K)$ and
$T_n=\sum_{\ell\ge n}P_\ell$, then whenever
$\kappa n(n+2)>4b+E_*$,
\begin{equation}\label{eq:native-spectral-window-tail}
 \norm{T_n\mathcal P_*}
 \le \frac{2b}{\kappa n(n+2)-4b-E_*}
       \norm{T_{n-1}\mathcal P_*}.
\end{equation}
Thus, for any $n_0$ satisfying
$\kappa n_0(n_0+2)\ge2(4b+E_*)$,
\begin{equation}\label{eq:native-spectral-window-factorial}
 \norm{T_n\mathcal P_*}
 \le\left(\frac{4b}{\kappa}\right)^{n-n_0+1}
       \left(\frac{(n_0-1)!}{n!}\right)^2,
 \qquad n\ge n_0.
\end{equation}
All statements remain valid on the gauge-invariant subspace.
\end{theorem}
\begin{proof}[Proof mechanism]
After separating one link, the exterior contribution to $K$ is
nonnegative.  The omitted one-link Casimir contributes
$\kappa\mu_L$, whereas the complete incident Wilson multiplication is
bounded below by $-4b$, proving the high-block floor.  Multiplication by
a coordinate on $S^3$ changes harmonic degree only by one and its
raising norm is exactly $1/2$, which gives the boundary-source identity
and the bound $\norm{B_L}\le2b$.  Schur complementation gives
\cref{eq:native-high-return}; iterating the degree-by-degree resolvent
identity gives \cref{eq:native-high-return-tail}.  For the spectral
window, the exact Sylvester equation between $T_n\mathcal P_*$ and its
complement is solved by a convergent semigroup integral using the
energy separation $\kappa n(n+2)-4b-E_*$.  Iteration gives the
factorial tail.  \src{F17}{V99}
\end{proof}

The theorem controls an entire low-energy spectral projection, not a
chosen list of eigenvectors.  It does not prove a uniform gap for the
retained Schur head: the latter still contains the complete collective
exterior dynamics.

\subsection{Resonances forbid whole-head elimination}
The same exact electric grading also identifies a structural limit of
high/low block-diagonalization.  For every link cutoff $\ell$ there are
physical gauge-invariant states on an embedded theta graph with equal
electric energy
\begin{equation}\label{eq:native-resonant-energy}
 E_\ell=\kappa(4\ell^2+12\ell+18)
\end{equation}
which lie on opposite sides of the selected-link degree-$\ell$ cutoff
and satisfy
\begin{equation}\label{eq:native-resonant-matrix-element}
 \langle\Psi'_\ell,V\Psi_\ell\rangle
 =-\frac{b}{2(\ell+1)(\ell+2)}\ne0\qquad(b>0).
\end{equation}
Consequently no whole-head solution of the homological equation
$[T,X]=-QVF$ exists on the complete physical retained space.  This does
not contradict \cref{thm:native-high-Schur}: the resonant states occur
at high total energy, whereas the Schur/Sylvester estimate separates a
fixed low-energy window from the high block.  \src{F17}{V100}

The correct finite-lattice operation is therefore to preserve the
resonant part.  Since the electric dynamics has period
$\tau=2\pi/\kappa$, define for a bounded local operator $A$
\begin{equation}\label{eq:native-resonant-average}
 \mathcal R(A)=\frac1\tau\int_0^\tau e^{itT}Ae^{-itT}\dd t,
 \qquad
 \mathcal J(A)=\frac{i}{\tau}\int_0^\tau
        (t-\tau/2)e^{itT}Ae^{-itT}\dd t.
\end{equation}
For a plaquette term $V_p$, $R_p=\mathcal R(V_p)$ commutes with $T$,
$S_p=\mathcal J(V_p)$ obeys $[T,S_p]=V_p-R_p$, and
\begin{equation}\label{eq:native-resonant-local-bounds}
 \norm{R_p}\le b,\qquad \norm{S_p}\le\frac{\pi b}{2\kappa}.
\end{equation}
Both remain supported on the same plaquette and preserve gauge
invariance.

\begin{theorem}[One resonance-preserving local normal-form step]
\label{thm:native-resonant-normal-form}
Let $S=\sum_pS_p$, $R=\sum_pR_p$, and assume
\begin{equation}\label{eq:native-resonant-q}
 q=13\pi b/\kappa<1.
\end{equation}
Then at every finite volume
\begin{equation}\label{eq:native-resonant-normal-form}
 e^S(H-b|\mathcal P_N|)e^{-S}=T+R+W
\end{equation}
exactly.  The bounded remainder has a connected plaquette expansion
$W=\sum_ZW_Z$ with the volume-uniform local strength
\begin{equation}\label{eq:native-resonant-local-remainder}
 \sup_e\sum_{Z\ni e}\norm{W_Z}
 \le b\sum_{n\ge1}3(4n+5)^3q^n<\infty.
\end{equation}
The right side is $O(b^2/\kappa)$ as $b/\kappa\to0$, and the tail after
$n$ commutator steps is bounded by the corresponding tail of the same
convergent series.
\end{theorem}
\begin{proof}[Proof mechanism]
The first commutator with the unbounded $T$ is replaced exactly by the
bounded operator $-(V-R)$.  The remaining Baker--Campbell--Hausdorff
series contains only bounded local commutators.  A new plaquette can
contribute only if it overlaps the accumulated support; the rooted
order-$n$ choice count is at most $13^nn!$.  The $1/n!$ coefficient and
$2\norm{S_p}$ per commutator give the geometric factor $q^n$.
Counting roots within $n$ overlap steps of a fixed link yields the
polynomial factor $3(4n+5)^3$ and proves the spatially summable bound.
\src{F17}{V100}
\end{proof}

The comparison operator $T+R$ itself has the constant ground and gap at
least $3\kappa-4b$ when $4b<3\kappa$, but
\cref{thm:native-resonant-normal-form} retains the nonzero $W$ rather
than converting that comparison gap into a claim about $H$.  The
normal-form domain \cref{eq:native-resonant-q} is again strong-electric;
it is not the asymptotically free physical scaling regime.

\section{Finite-time conditional deformation and the failure of monotonicity}
\label{sec:native-finite-time-deformation}
A separate continuation studies the actual nonlinear conditional fibre
under a small native flow parameter.  Its strongest uniform local
statement is useful as an interface estimate: after cancellation of the
complete second normal source, the corrected block map has a uniform
actual conditional gap and a uniformly small centered source inverse
on its certified domain.  In the notation of that construction, for
all spatial side lengths $n\ge5$, every coarse field $g$, and
$|\beta|\le10^{-13}$,
\begin{equation}\label{eq:native-conditional-gap-late}
 \gap L_{g,\beta}
 \ge e^{-2\Lambda_\beta}(1-H_\beta)>\frac12,
 \qquad
 \Lambda_\beta=\frac83|\beta|+
       \frac{157525571}{19274112}\beta^2,
\end{equation}
and the full centered normal-source inverse is $O(|\beta|^6)$ in
operator norm, uniformly in volume and coarse configuration.
\src{F18}{V48--V49}  This is a gap of the actual conditional fibre
generator for a very small deformation parameter; it is not the
physical Hamiltonian gap and it does not supply the long block-time
transport by itself.

The subsequent exact finite-jet calculation also closes an important
negative shortcut.  Let $\mathcal C(t,g)$ denote the native conditional
cost in its normal-chart convention on a cubic spatial torus of side
$n$, and let
\begin{equation}\label{eq:native-coarse-W}
 \mathcal W_{\rm c}(g)=\sum_y\Re H_y
\end{equation}
be the sum of coarse plaquette holonomies.  The complete quadratic
coefficient, including the previously unpaid cubic-density contraction,
is
\begin{equation}\label{eq:native-cost-C2}
 \mathcal C(t,g)=3n^3+t^2\mathcal C_2(g)+O_{n,g}(t^4),
 \qquad
 \mathcal C_2(g)=\frac{1079217}{30590}n^3+
                         \frac{333}{8}\mathcal W_{\rm c}(g).
\end{equation}
The coefficient is positive at the flat coarse field but negative at a
legitimate all-negative coarse-plaquette field:
\begin{equation}\label{eq:native-cost-signs}
 \mathcal C_2(g_{\rm flat})=
       \frac{19596573}{122360}n^3>0,
 \qquad
 \mathcal C_2(g_{\rm minus})=
      -\frac{10962837}{122360}n^3<0
\end{equation}
for the stated even tori.  \src{F18}{V59}
Hence neither global increase nor global decrease of this native cost
at small positive time can be used as a uniform coarse-field principle.
The result is finite-regulator and its $O_{n,g}(t^4)$ remainder is not
uniform in volume; its role in the present monograph is to eliminate a
monotonicity shortcut and to identify the source-specific finite-time
remainder and Poisson-Hessian estimates that remain necessary.

\sourcebasis{The electric-dominated ground-state construction and
localized first plaquette band are in \src{F18}{V5,V12}. The
volume-independent native high block, Schur return, spectral-window
tails, resonance obstruction, and resonance-preserving local normal
form are in \src{F17}{V99--V100}. The uniform corrected conditional
fibre estimate and the completed small-time cost are in
\src{F18}{V48--V49,V59}.}

\chapter{The physical spectral endpoint}
\label{ch:spectral}
The auxiliary results of \cref{ch:compact-repair} remain upstream of
this chapter. The full hidden diffusion can be metastable even though
specific visible interface events have uniform killed heat-bath floors.
Neither generator is identified with the physical reflected transfer by
its name or its invariant density. A physical application requires the
same-law source identification and rate-calibrated form comparison
specified below. No such comparison follows solely from the pivot
convexity of \cref{prop:pivot-potential}.

The actual native spectral theorems in
\cref{thm:native-critical-spectrum,thm:native-joint-spectrum} give
lattice-time deficits on specified critical domains. The continuum
criterion below instead requires a compatible family at a fixed
calibrated physical delay. The exact regional reduction in
\cref{thm:native-regional-Feshbach} provides one possible interface,
with its complement and boundary hypotheses retained. A local route
can also separate a reducing soft band from the source algebra, but
only through a vanishing source-compression estimate. The loop and
smeared-probe results below make the two possible outcomes explicit.

\section{The local reflected Hilbert space}
Let $\mathfrak A_+^{\rm loc}$ be the bounded gauge-invariant positive-time cylinder algebra with bounded physical support. For a candidate Euclidean state $\omega$ and reflection $\Theta$, set
\[
(F,G)_{\OS}=\omega((\Theta F)G).
\]
Reflection positivity permits quotienting by the null space and completion. The remaining Euclidean hypotheses must supply a strongly continuous positive-time semigroup $e^{-tH}$, $H\ge0$, with vacuum $\Omega$. Denote the centered local vector by
\[
v(F)=[F]-\omega(F)\Omega.
\]
This construction is conditional on the candidate correlations and axioms. A finite positive measure or a finite positive kernel is not itself the continuum Euclidean theory.

For a finite family $\mathcal F=(F_1,\ldots,F_m)$, define $V_{\mathcal F}c=\sum_ic_iv(F_i)$ and
\begin{equation}\label{eq:localGrams}
G_{\mathcal F}^{0}=V_{\mathcal F}^*V_{\mathcal F},\qquad
G_{\mathcal F}^{\tau}=V_{\mathcal F}^*e^{-2\tau H}V_{\mathcal F}.
\end{equation}
The factor $2\tau$ is deliberate: the delayed Gram is the squared norm after evolution by time $\tau$.

\begin{theorem}[Complete local Gram characterization]\label{thm:local-Gram}
Suppose centered local vectors span a dense subspace of $\Omega^\perp$. For a fixed physical $\tau>0$ and $0<q<1$, the following are equivalent:
\begin{enumerate}[label=\textup{(\roman*)},leftmargin=2.5em]
\item $G_{\mathcal F}^{\tau}\preceq q^2G_{\mathcal F}^{0}$ for every finite local family;
\item $\norm{e^{-\tau H}|_{\Omega^\perp}}\le q$;
\item $H|_{\Omega^\perp}\succeq-\tau^{-1}\log q$.
\end{enumerate}
In particular, this packet implies uniqueness of the vacuum in the reconstructed local Hilbert space.
\end{theorem}
\begin{proof}
For every coefficient vector $c$,
\[
c^*G_{\mathcal F}^{\tau}c
=\norm{e^{-\tau H}V_{\mathcal F}c}^2,
\qquad c^*G_{\mathcal F}^{0}c=\norm{V_{\mathcal F}c}^2.
\]
Thus (i) is the contraction inequality on the dense centered local span and extends by continuity to (ii). The spectral theorem makes (ii) equivalent to (iii). A second zero-energy vector orthogonal to $\Omega$ would contradict the strict lower bound. The reverse implications are immediate. \src{F1}{V1}
\end{proof}

All mixed linear combinations are part of the theorem. Separate diagonal tests on named loops do not suffice: a matrix with diagonal entries $0.6$ and off-diagonal entries $0.3$ has largest eigenvalue $0.9$, despite both diagonal entries being below $0.7$. The complete Gram order is the relevant finite datum.

\section{Passing through changing cutoff null spaces}
At cutoff $j$, use the corresponding source family and write $G_j^0,G_j^{\tau_j}$. Suppose for each fixed family
\begin{equation}\label{eq:Gram-convergence}
G_j^0\to G^0,\qquad G_j^{\tau_j}\to G^\tau,
\qquad \tau_j\to\tau>0,
\end{equation}
and
\begin{equation}\label{eq:Gram-error}
G_j^{\tau_j}\preceq q^2G_j^0+E_j,\qquad \norm{E_j}\to0.
\end{equation}
Then $G^\tau\preceq q^2G^0$. This is simply passage to the limit in every quadratic form, with no inverse Gram matrix. It remains valid when the limiting null space grows.

A lower-frame estimate has a different purpose. A bound $G_{j,m}^0\succeq c_mG_m^{\rm phys}$ with $c_m>0$ prevents a predeclared physical coefficient direction from disappearing and stabilizes generalized eigenvalue comparisons. It is not needed merely to pass a common Loewner inequality to the surviving quotient. Some nonzero limiting local states and a suitable dense algebra are nevertheless indispensable for a nontrivial theory. An absolute error estimate can conceal normalized soft states whose entrance norm vanishes. \src{F1}{V1}; \cite{GT}

\section{Local toron invisibility and quotient-stable Gram excision}
\label{sec:local-toron-Grams}
A low toron level obstructs a global torus spectral gap but need not
be visible to a limiting local source. The distinction can be tested
at the level of the actual source compression, rather than by
removing states by definition.

\subsection{Physically bounded Wilson loops}
Let $C_L$ be a contractible spatial loop whose lift closes in
$\mathbb Z^3$, with successive oriented directions
$(\sigma_r,\mu_r)$ and length $m_L\le\kappa L$. Thus
$\sum_r\sigma_r e_{\mu_r}=0$. In the constant soft coordinate put
\[
 W^{\rm tor}_{C_L}(q)=\tfrac12\operatorname{Tr}
       \prod_{r=1}^{m_L}\operatorname{Exp}(h\sigma_rq_{\mu_r}),
 \qquad h=\gamma^{-1/3}/L.
\]
The length condition includes lattice approximations of fixed
rectifiable loops in a box of fixed physical size.

\begin{proposition}[Uniform local-loop invisibility of the soft coordinate]
\label{prop:local-toron-loops}
For every such loop,
\begin{equation}
 0\le1-W^{\rm tor}_{C_L}(q)
 \le\frac{\kappa^4}{8}\gamma^{-4/3}|q|^4
             e^{2\kappa\gamma^{-1/3}|q|}.
 \label{eq:local-toron-loop-blindness}
\end{equation}
The bound is uniform in $L$, including exponential growth in
$\gamma$. For a background configuration in a fixed exponential
ball, the weaker but background-uniform estimate is
\begin{equation}
 |W_{C_L}(s;hq)-W_{C_L}(s;0)|
                    \le C\kappa\gamma^{-1/3}|q|.
 \label{eq:local-toron-background}
\end{equation}
The latter estimate does not require contractibility.
\end{proposition}
\begin{proof}
If matrices $X_r$ satisfy $\sum_rX_r=0$ and
$A=\sum_r\|X_r\|$, the ordered product series gives
$\|\prod_re^{X_r}-I\|\le e^A-1-A\le A^2e^A/2$.
For $U\in\SU(2)$,
$1-\tfrac12\operatorname{Tr}U\le\|U-I\|^2/2$.
Apply these facts with $X_r=h\sigma_rq_{\mu_r}$ and
$A\le\kappa\gamma^{-1/3}|q|$. For a nonzero background,
telescope the product and use the uniform Lipschitz bound for the
exponential on the declared ball. The total change is at most
$Cm_Lh|q|$, proving \cref{eq:local-toron-background}.
\end{proof}

Let $\Phi_L$ be a fixed polynomial of finitely many normalized
traces of such loops and put $\Phi_*=\Phi(1,\ldots,1)$.
These multipliers are uniformly bounded because all normalized
traces lie in $[-1,1]$. \Cref{eq:local-toron-loop-blindness} therefore gives
$M_{\Phi_L^{\rm tor}}\to\Phi_*I$ strongly on $L^2(\mathbb R^9)$,
uniformly in $L$. On any moving family of soft subspaces whose unit
balls have uniformly vanishing position tails, the convergence is
uniform in operator norm after restriction to that family. The proof
splits at a fixed $|q|\le R$, uses uniform convergence on that ball,
and then sends $R\to\infty$.

This conclusion also interfaces with the hard variables on the
controlled joint branch of \cref{thm:native-joint-spectrum}. For
its normalized hard ground function $\varphi_L$, define the compressed
symbol, with the same measured tube cutoff,
\[
 m_{\gamma,L}(q)=\int |\varphi_L(u)|^2
                         \Phi_L(\gamma^{-1/2}u,hq)\,\dd u.
\]
\Cref{eq:local-toron-background} gives
$\sup_{|q|\le R}|m_{\gamma,L}(q)-m_{\gamma,L}(0)|
 \le C_{\Phi,R}\gamma^{-1/3}+o(1)$.
The existing hard-tail estimate pays the omitted tube mass. Soft
tightness and convergence of the corresponding isolated spectral
projections then make the local source asymptotically scalar on each
fixed finite toron band, for $\Theta_{\gamma,L}\to0$.
This last spectral assertion is not extended to an exponentially
growing physical trajectory by the geometric loop bound alone.

\subsection{A reducing soft band and the local entrance form}
Let $H_j\ge0$ have normalized vacuum $\Omega_j$, and let $\Pi_j$
be a reducing orthogonal projection with $\Pi_j\Omega_j=\Omega_j$.
Write
\[
 \Pi_j^\circ=\Pi_j-|\Omega_j\rangle\langle\Omega_j|,
 \qquad Q_j=I-\Pi_j^\circ,
 \qquad
 V_jc=\sum_{\alpha=1}^m c_\alpha
  (F_{\alpha,j}-\langle\Omega_j,F_{\alpha,j}\Omega_j\rangle)
                                                        \Omega_j.
\]
Here $F_{\alpha,j}$ are the bounded local sources whose limiting
Gram is being tested. The projection $Q_j$ retains the vacuum;
$V_j$ is centered independently of that convention.

\begin{theorem}[Soft-band invisibility and exact Gram excision]
\label{thm:local-toron-Gram-excision}
Suppose there are scalars $c_{\alpha,j}$ such that
\begin{equation}
 \epsilon_j=\max_{\alpha\le m}
  \|\Pi_j(F_{\alpha,j}-c_{\alpha,j}I)\Pi_j\|\longrightarrow0.
 \label{eq:local-toron-scalar-hypothesis}
\end{equation}
Then
\begin{equation}
 \|\Pi_j^\circ V_j\|\le2\sqrt m\,\epsilon_j,
 \qquad
 0\preceq V_j^*\Pi_j^\circ e^{-2tH_j}\Pi_j^\circ V_j
                 \preceq4m\epsilon_j^2I_m\quad(t\ge0).
 \label{eq:local-toron-Gram-invisibility}
\end{equation}
If, for one delay $t_j\ge0$ and $0\le q_j<1$,
\begin{equation}
 V_j^*Q_je^{-2t_jH_j}Q_jV_j
              \preceq q_j^2V_j^*Q_jV_j+E_j,
 \label{eq:local-toron-complement-contraction}
\end{equation}
then the full entrance and delayed Grams satisfy
\begin{equation}
 G_j^{t_j}\preceq q_j^2G_j^0+E_j
                         +(1-q_j^2)4m\epsilon_j^2I_m,
 \qquad
 G_j^0=V_j^*V_j,\quad G_j^t=V_j^*e^{-2tH_j}V_j.
 \label{eq:local-toron-full-contraction}
\end{equation}
No inverse entrance Gram is required.
\end{theorem}
\begin{proof}
Since the vacuum belongs to $\operatorname{Ran}\Pi_j$, its
expectation of $F_{\alpha,j}$ differs from $c_{\alpha,j}$ by at
most $\epsilon_j$. Each centered projected source thus has norm at
most $2\epsilon_j$. Cauchy--Schwarz in $\mathbb C^m$ gives the
synthesis bound. Positivity, contractivity, and the reducing property
give the delayed bound. Split both Grams using the orthogonal reducing
decomposition $I=Q_j+\Pi_j^\circ$. On the second component use
$e^{-2t_jH_j}\preceq I$, while on the first use
\cref{eq:local-toron-complement-contraction}. The excess over
$q_j^2G_j^0+E_j$ is bounded by
$(1-q_j^2)V_j^*\Pi_j^\circ V_j$, proving the assertion.
\end{proof}

If $t_j\to\tau>0$, $q_j\to q<1$, $E_j\to0$, and the local
Grams converge, this theorem supplies the inequality in
\cref{thm:local-Gram} even when their limiting entrance matrix is
singular. It does not prove the complement contraction, source
noncollapse, or density of the limiting source family. Those must
hold on the same calibrated physical trajectory. In addition, a
diverging source-renormalization factor can undo the raw loop's
invisibility rate. The hypothesis
\cref{eq:local-toron-scalar-hypothesis} must therefore be verified
for the actual normalized sources, not merely their unrenormalized
loop representatives.
\sourcebasis{Physically bounded loop estimates, hard-ground
compression on the joint critical branch, and the exact Gram
excision argument are established in \src{F16B}{V56}.}


\section{The physical clock}
Let $T_j$ be a one-step positive transfer at lattice time spacing $a_j\downarrow0$.

\begin{proposition}[Fixed one-step contraction can collapse the limit]\label{prop:step-collapse}
Under local kernel convergence to a nonzero strongly continuous continuum sector, a regulator-independent bound $\norm{T_j}\le q<1$ on a reducing sector containing a noncollapsed local family is impossible. A sufficient rate-calibrated condition is instead
\begin{equation}\label{eq:physicalrate}
\liminf_{j\to\infty}\frac{-\log\rho_j}{a_j}\ge m>0,
\qquad \rho_j=\norm{T_j|_{\Omega_j^\perp}}.
\end{equation}
Under the same convergence and density hypotheses it yields $H_\infty|_{\Omega_\infty^\perp}\succeq mI$.
\end{proposition}
\begin{proof}
A fixed one-step $q$ gives $\norm{T_j^{\lfloor t/a_j\rfloor}}\le q^{\lfloor t/a_j\rfloor}\to0$ for every $t>0$. The limiting kernels vanish at every positive time, contradicting continuity at zero unless the entrance family collapses. Under \cref{eq:physicalrate}, for each $m'<m$, the powers are eventually bounded by $e^{-m'a_j\lfloor t/a_j\rfloor}$. Pass the corresponding finite Gram inequalities to the limit and let $m'\uparrow m$. \src{F1}{V1}
\end{proof}

For $\rho_j=1-\gamma_j\to1$, the physical rate is $(\gamma_j/a_j)(1+o(1))$, not the raw deficit. The fixed-delay criterion of \cref{thm:local-Gram} avoids imposing invariance on each finite local synthesis range.

\subsection{Bare isotropic calibration and the domain of the toron limit}
\label{sec:native-bare-calibration}
The native transfer provides a concrete example of the clock
requirement. In the round convention
$\mathcal L_G=4\mathbf E^2$, the leading Wilson logarithm has
kinetic coefficient $1/(2\gamma)$ in front of
$\sum_e\mathcal L_G$ and magnetic coefficient $\gamma$ in front
of $\sum_ps(U_p)$. Define the bare Hamiltonian convention by
\begin{equation}
 H_{\rm KS}
   =\frac{g_0^2}{2a_s}\sum_e\mathbf E_e^2
                  +\frac4{g_0^2a_s}\sum_ps(U_p).
 \label{eq:native-bare-Hamiltonian-convention}
\end{equation}
Matching both coefficients in $a_tH_{\rm KS}$, rather than only
the kinetic term, gives
\begin{equation}
 \gamma=\frac{4a_s}{g_0^2a_t}
        =\frac{4a_t}{g_0^2a_s},
 \qquad a_t=a_s=:a,\qquad \gamma=4/g_0^2.
 \label{eq:native-bare-isotropic-calibration}
\end{equation}
This is the bare isotropic convention of the native action, not an
independent construction of a renormalized physical observable.

\begin{proposition}[The joint critical domain is not a fixed-box
asymptotically free trajectory]
\label{prop:native-physical-domain-separation}
Suppose a physical scale is selected with the asymptotic relation
\begin{equation}
 \log\frac1{a\Lambda_{\rm lat}}
               =\frac{3\pi^2}{11}\gamma+O(\log\gamma).
 \label{eq:native-AF-scale-premise}
\end{equation}
For a fixed physical box length $\ell=aL$, this implies
$\log L=(3\pi^2/11)\gamma+O(\log\gamma)$. Such a trajectory is
outside the domain $\Theta_{\gamma,L}\to0$ of
\cref{thm:native-joint-spectrum}.
Moreover, on any branch for which a native logarithmic level has
$-\log(\lambda/\Lambda)=h\mu_{\gamma,L}$ with
$\mu_{\gamma,L}$ bounded, its physical energy at fixed $\ell$ is
\begin{equation}
 \frac{-\log(\lambda/\Lambda)}a
          =\frac{\gamma^{-1/3}}\ell\mu_{\gamma,L}
                                                     \longrightarrow0.
 \label{eq:native-toron-physical-energy}
\end{equation}
A nonzero limiting toron energy would instead require
$\mu_{\gamma,L}$ to grow at least on the scale $\gamma^{1/3}$.
\end{proposition}
\begin{proof}
Insert $a=\ell/L$ into \cref{eq:native-AF-scale-premise}.
Exponential growth in $L$ makes the positive terms of
\cref{eq:native-joint-budget} nonvanishing. Independently,
$h/a=(\gamma^{-1/3}/L)/(\ell/L)$ proves
\cref{eq:native-toron-physical-energy}.
\end{proof}

\Cref{eq:native-AF-scale-premise} is a stated physical
scale premise here; it is not deduced from the monograph's measured
local covariance coefficient. The second part is conditional on the
continued existence of a bounded toron coefficient and does not
extrapolate the joint spectral theorem to that physical trajectory.
Together these statements prevent a critical matrix-model level from
being relabeled as a positive continuum mass. They leave two distinct
tasks: control the actual running physical generator, or establish a
local-vacuum gap through sources that become insensitive to the soft
band, as in \cref{thm:local-toron-Gram-excision}.
\sourcebasis{The bare convention, joint-domain exclusion, and
conditional physical-energy calculation are developed in
\src{F16B}{V55}.}


\section{A nonvanishing local soft-spectrum test}
\label{sec:smeared-vacuum-probe}
Soft-band excision requires vanishing local spectral weight, not
merely a lattice gap tending to zero. A complementary test is
available for a canonically smeared magnetic-energy probe under the
actual infinite-time Wilson vacuum. It also separates the selected
weak-field mesoscopic branch from the desired interacting physical
theory.

\subsection{The actual vacuum covariance input}
Consider the cylinder
$\mathbb Z\times(\mathbb Z/M\mathbb Z)^3$, with the selected
dyadic side lengths and couplings $\gamma_M$ of
\src{F15}{V60--V64}. This is not the finite periodic four-dimensional
law of \cref{ch:full-compact-sources}. Let
\begin{align*}
 P_M&=T_M/\Lambda_M,\qquad H_M=-\log P_M,\\
 X_{M,x}&=\gamma_Ms(U_{(0,x;12)}),\qquad
 \overline X_{M,x}=X_{M,x}-\mathbb E_{\rm vac}X_{M,x}.
\end{align*}
The actual stationary slice law is $\Omega_M^2\dd H$. Its inherited
mean-energy excess $\vartheta_M$ and local moments obey
\begin{equation}
 \vartheta_M\longrightarrow0,\qquad
 \mathbb E_{\rm vac}T_0\le\tfrac92+\vartheta_M,
 \qquad
 \mathbb E_{\rm vac}X_p^q\le e^3 16^q q!,
 \qquad 2\le\gamma_M\le2\log M,
 \label{eq:smeared-vacuum-inputs}
\end{equation}
with $\gamma_M\to\infty$. Here $T_0$ is the symmetrically assigned
cell energy, not the transfer operator. Define
\begin{equation}
 \delta_M=\sqrt{\vartheta_M}+\gamma_M^{-1/16}+M^{-2}
                                                        \longrightarrow0.
 \label{eq:smeared-vacuum-error}
\end{equation}
The local Haar identities with linear and cubic tests, the moment
bounds, and the finite-range curl filter give the uniform-in-time
covariance estimate
\begin{equation}
 \left|\operatorname{Cov}_{\rm vac}(X_{M,x}(0),X_{M,y}(n))
                    -\tfrac32 A(n,x-y)^2\right|\le C\delta_M
 \label{eq:smeared-vacuum-covariance}
\end{equation}
whenever each component of the lifted spatial separation has
absolute value at most $M/4$. This is the actual-vacuum input of
\src{F15}{V64}, not a differentiation of a characteristic-function
error or a conditional Gaussian assumption.

For $n\ge1$ the comparison curvature kernel is
\begin{equation}
 A(n,z)=\int_{[-\pi,\pi]^3}e^{ip\cdot z}
              \frac{K_{12}(p)}{2\sinh E(p)}e^{-nE(p)}
                                  \frac{\dd p}{(2\pi)^3},
 \label{eq:smeared-lattice-curvature-kernel}
\end{equation}
where
$K(p)=4\sum_{\mu=1}^3\sin^2(p_\mu/2)$,
$K_{12}(p)=4\{\sin^2(p_1/2)+\sin^2(p_2/2)\}$, and
$E(p)=2\operatorname{arsinh}(\sqrt{K(p)}/2)$.
Write
\begin{equation}
 a(\tau,z)=\int_{\mathbb R^3}e^{ik\cdot z}
                    w(k)e^{-\tau|k|}\frac{\dd k}{(2\pi)^3},
 \qquad w(k)=\frac{k_1^2+k_2^2}{2|k|},\quad w(0)=0.
 \label{eq:smeared-continuum-curvature-kernel}
\end{equation}
Rescaling $p=k/b$ and applying dominated convergence gives
$b^4A(n,z)-a(n/b,z/b)\to0$ uniformly in $z$ when $n/b$ stays
in a compact subset of $(0,\infty)$. The exponential damping is
essential to this positive-time assertion.

\subsection{Smearing with its full error budget}
Fix real $f\in C_c^\infty(\mathbb R^3)$ with
$m_f=\int f\ne0$, and integers $b_M$ satisfying
\begin{equation}
 b_M\to\infty,\qquad b_M/M\to0,
                       \qquad \delta_M b_M^8\to0.
 \label{eq:smeared-probe-window}
\end{equation}
This window is nonempty: eventually
$b_M=\lfloor\delta_M^{-1/16}\rfloor$ works because
$\delta_M\ge M^{-2}$ implies $b_M\le M^{1/8}$.
Using a centered spatial lift of the support, put
\begin{equation}
 F_M^f=b_M\sum_x f(x/b_M)\overline X_{M,x},
 \qquad v_M^f=F_M^f\Omega_M,
 \qquad
 K_f(\tau)=\tfrac32\int\!\!\int
                   f(x)f(y)a(\tau,x-y)^2\,\dd x\,\dd y.
 \label{eq:smeared-probe-and-covariance}
\end{equation}
The coefficient $b_M=b_M^{-3}b_M^4$ declares spatial integration
of a dimension-four thermal composite. It is not an established
renormalization of a physical interacting observable.

\begin{theorem}[Positive-time covariance of the canonically smeared probe]
\label{thm:smeared-probe-covariance}
If $n_M/b_M\to\tau>0$, then
\begin{equation}
 \langle v_M^f,P_M^{n_M}v_M^f\rangle\longrightarrow K_f(\tau).
 \label{eq:smeared-probe-covariance-limit}
\end{equation}
Convergence is uniform for $n/b_M$ in compact positive time
intervals. No assertion about the undamped norm at time zero is
required.
\end{theorem}
\begin{proof}
For a centered slice multiplication operator, the exact vacuum
transfer identity identifies the left side with its covariance at
time separation $n$. There is no slab-endpoint shift in this exponent.
The support condition makes \cref{eq:smeared-vacuum-covariance}
applicable to every term of the double sum. Since
$\sum_x|f(x/b_M)|\le C_fb_M^3$, the complete comparison error is
at most $C_f\delta_Mb_M^8$. The remaining sum is
\[
 \tfrac32b_M^{-6}\sum_{x,y}f(x/b_M)f(y/b_M)
                                  [b_M^4A(n,x-y)]^2.
\]
The uniform positive-time kernel limit and a compactly supported
Riemann sum give \cref{eq:smeared-probe-covariance-limit}.
\end{proof}

\subsection{Surviving damped vectors and the soft spectral endpoint}
For $s>0$, define the actual damped vector and rescaled generator
\[
 w_{M,s}^f=P_M^{\lfloor sb_M\rfloor}v_M^f,
 \qquad \mathscr H_M=b_MH_M,
 \qquad
 \mu_{M,f,s}(B)=\langle w_{M,s}^f,
                     \one_B(\mathscr H_M)w_{M,s}^f\rangle.
\]
If a zero eigenspace of $P_M$ is present, integer damping removes
it; the logarithm is taken on the positive spectral subspace.
Use $\widehat f(k)=\int e^{ik\cdot x}f(x)\dd x$ and put
\begin{equation}
 \mu_{f,s}(B)=\frac{3}{2(2\pi)^6}
   \int_{\{|k|+|\ell|\in B\}}
    |\widehat f(k+\ell)|^2 w(k)w(\ell)
                      e^{-2s(|k|+|\ell|)}\,\dd k\,\dd\ell.
 \label{eq:smeared-probe-spectral-measure}
\end{equation}

\begin{theorem}[Nonvanishing local soft spectral weight]
\label{thm:smeared-probe-spectrum}
Under \cref{eq:smeared-probe-window},
\begin{equation}
 \|w_{M,s}^f\|^2\longrightarrow K_f(2s)>0,
 \qquad \mu_{M,f,s}\Longrightarrow\mu_{f,s}
 \quad\hbox{as finite measures on }[0,\infty).
 \label{eq:smeared-probe-spectral-convergence}
\end{equation}
The limiting measure is finite, nonzero, and atomless. It has
positive mass in every interval $(0,E)$, with
\begin{equation}
 \mu_{f,s}((0,E))\sim\frac{|m_f|^2}{26880\pi^4}E^8
       \quad(E\downarrow0),
 \qquad
 K_f(\tau)\sim\frac{3|m_f|^2}{2\pi^4}\tau^{-8}
       \quad(\tau\to\infty).
 \label{eq:smeared-probe-soft-asymptotics}
\end{equation}
Consequently the unit damped vectors have a positive limiting
spectral weight in $(0,E)$ for every $E>0$. If $g_M$ is the
bottom centered spectral energy of $H_M$, then $b_Mg_M\to0$.
\end{theorem}
\begin{proof}
The norm assertion follows from
\cref{thm:smeared-probe-covariance} at time $2s$.
Fubini in \cref{eq:smeared-continuum-curvature-kernel} gives
\[
 K_f(\tau)=\frac3{2(2\pi)^6}\int
  |\widehat f(k+\ell)|^2w(k)w(\ell)
                         e^{-\tau(|k|+|\ell|)}\,\dd k\,\dd\ell.
\]
Thus $\mu_{f,s}$ has Laplace transform $K_f(2s+\tau)$.
The integrand is integrable for $s>0$, and each level set of
$|k|+|\ell|$ has six-dimensional Lebesgue measure zero, proving
atomlessness. Since $\widehat f(0)=m_f\ne0$, every sufficiently
small energy ball contains positive mass.

For the actual measures, integer powers with exponents
$2\lfloor sb_M\rfloor+\lfloor b_M\tau\rfloor$ and
$2\lfloor sb_M\rfloor+\lceil b_M\tau\rceil$ squeeze the
Laplace transform at every $\tau>0$. Their limits are both
$K_f(2s+\tau)$. Together with the norm limit this also gives
energy tightness, since
\[
 \mu_{M,f,s}([R,\infty))\le
 \frac{\|w_{M,s}^f\|^2-
            \int e^{-\tau u}\,\dd\mu_{M,f,s}(u)}
      {1-e^{-\tau R}}.
\]
First take the upper limit in $M$, then use continuity at $2s$
and choose $\tau$ small and $R$ large. Uniqueness of finite
Laplace transforms proves weak convergence. Atomlessness makes
$(0,E)$ a continuity set. Its positive limiting mass divided by
$K_f(2s)$ gives the unit-vector statement, which forces
$b_Mg_M\to0$.

Finally scale $k=Ep$, $\ell=Eq$ in
\cref{eq:smeared-probe-spectral-measure}. The two homogeneous
weights and six-dimensional volume give $E^8$.
The angular integral is $\int_{S^2}w(r\omega)\dd\omega=4\pi r/3$,
and $\int_{r+u<1}r^3u^3\dd r\dd u=1/1120$.
They give the first coefficient in
\cref{eq:smeared-probe-soft-asymptotics}. Rescaling by $1/\tau$
in the formula for $K_f$ gives the second, using
$\int w(p)e^{-|p|}\dd p/(2\pi)^3=1/\pi^2$.
\end{proof}

\begin{scope}[A branch-selection test, not a physical gap conclusion]
Unlike a disappearing toron contribution, the soft spectrum here
has nonzero weight in vectors whose norms survive the limit. It
cannot be discarded by the vanishing-Gram mechanism of
\cref{thm:local-toron-Gram-excision} while this probe and scaling
are retained. The theorem rules out relabeling this selected
mesoscopic window as a massive theory with these observables.
It constructs neither the full interacting continuum law nor a
counterexample to the physical Yang--Mills mass-gap target.
Only the two-point spectral law of a quadratic curvature composite
is identified; the composite is not asserted to be Gaussian, and
no higher joint law follows from the covariance calculation.
\end{scope}
\sourcebasis{The uniform actual-vacuum covariance is established
in \src{F15}{V64}; canonical smearing, positive-time damping,
and the nonvanishing soft spectral measure are in
\src{F15}{V65}. Their scale and source normalization remain
explicit in \cref{eq:smeared-probe-window,eq:smeared-probe-and-covariance}.}


\section{Separating the retained and conditional temporal forms}
\label{sec:two-level-transfer}
The conditional covariance identity \cref{eq:totalcov} has an exact
operator counterpart for a physical transfer. Let $r:X\to Z$ be a
retained field map for the physical law $\pi$, with
$\overline\pi=r_*\pi$. Write
\[
\mathsf Jm=m\circ r,\qquad
\mathsf Kf=\E_\pi[f\mid r],\qquad \mathsf K=\mathsf J^*.
\]
Then $\mathsf J$ is an isometry and $\mathsf K\mathsf J=I$. Every
$f\in L^2(\pi)$ decomposes uniquely as
\begin{equation}\label{eq:two-level-decomposition}
f=\mathsf Jm+g,\qquad m=\mathsf Kf,\quad \mathsf Kg=0,
\qquad \norm f^2=\norm m^2+\norm g^2.
\end{equation}
For centered $f$, the retained field $m$ is centered as well.

\begin{proposition}[Exact two-level comparison transfer]
\label{prop:two-level-transfer}
Let $Q$ be a positive self-adjoint Markov contraction on
$L^2(\overline\pi)$, and define $R=\mathsf JQ\mathsf K$.
Then $R$ is a positive self-adjoint Markov contraction and
\begin{equation}\label{eq:two-level-form}
\ip{f}{(I-R)f}=\norm g^2+\ip{m}{(I-Q)m}.
\end{equation}
Its centered deficit is exactly the retained deficit:
\begin{equation}\label{eq:two-level-gap}
1-\norm{R|_{\one^\perp}}
=1-\norm{Q|_{\one^\perp}}.
\end{equation}
\end{proposition}
\begin{proof}
Positivity, self-adjointness, and preservation of constants follow
from $\mathsf K=\mathsf J^*$ and the corresponding properties of $Q$.
Since $R(\mathsf Jm+g)=\mathsf JQm$, orthogonality proves
\cref{eq:two-level-form}. The restriction of $R$ to $\Ker\mathsf K$
is zero, while its restriction to $\Ran\mathsf J$ is unitarily
equivalent to $Q$. This proves \cref{eq:two-level-gap}.
\src{F6}{V56}
\end{proof}

The comparison transfer completely refreshes the conditional detail.
Nevertheless, its slowest nonconstant mode is retained. Thus even
perfect conditional refresh does not remove the need for a coarse
physical estimate.

\begin{criterion}[Two-level physical coercivity]
\label{crit:two-level-physical}
Let $P$ be the actual positive self-adjoint physical transfer on
$L^2(\pi)$. Suppose, for every centered decomposition in
\cref{eq:two-level-decomposition},
\begin{equation}\label{eq:two-level-domination}
\ip{f}{(I-P)f}\ge
 a\norm g^2+b\ip{m}{(I-Q)m},
\qquad a,b\ge0,
\end{equation}
and $\ip{m}{(I-Q)m}\ge\gamma_Q\norm m^2$ for centered $m$.
Then
\begin{equation}\label{eq:two-level-physical-gap}
\ip{f}{(I-P)f}\ge\gamma_P\norm f^2,
\qquad \gamma_P=\min\{a,b\gamma_Q\}.
\end{equation}
For a physical step $\delta$ and $0<\gamma_P<1$, the corresponding
spectral lower bound is
\begin{equation}\label{eq:two-level-physical-mass}
-\delta^{-1}\log\norm{P|_{\one^\perp}}
\ge-\delta^{-1}\log(1-\gamma_P)
\ge\gamma_P/\delta.
\end{equation}
\end{criterion}
\begin{proof}
Insert the coarse inequality in \cref{eq:two-level-domination} and
use the orthogonal norm identity \cref{eq:two-level-decomposition}.
The spectral statement follows from positivity of $P$ and
$-\log(1-x)\ge x$ for $0\le x<1$. \src{F6}{V56}
\end{proof}

This supplies a physical mass along a cofinal family only with the
same convergence, noncollapse, and time identification as in
\cref{prop:step-collapse}. It is enough that both $a_j$ and
$b_j\gamma_{Q,j}$ retain the required order-$\delta_j$ margin.
The actual-law form domination \cref{eq:two-level-domination} is an
input: the definition of $R$ does not identify it with $P$.

\begin{lemma}[Comparison accuracy at the physical deficit scale]
\label{lem:transfer-comparison-debit}
For positive self-adjoint transfers $P,R$ on the same centered carrier,
\begin{equation}\label{eq:transfer-comparison-debit}
1-\norm{P|_{\one^\perp}}
\ge1-\norm{R|_{\one^\perp}}
  -\norm{(P-R)|_{\one^\perp}}.
\end{equation}
Hence a reference deficit $m\delta_j+o(\delta_j)$ is preserved at that
coefficient by an $o(\delta_j)$ transfer error; an unscaled $o(1)$
error alone is insufficient.
\end{lemma}
\begin{proof}
Apply the triangle inequality to the centered operator norms. The
last assertion follows by division by $\delta_j$.
\src{F6}{V56}
\end{proof}

The exact half-slab form in \cref{prop:half-variance} provides one
same-law route to \cref{eq:two-level-domination}. The core--bad-set
criterion \cref{crit:core-smoothing} provides another possible upstream
functional inequality, but only after its energy form has been
identified with the physical transfer or generator. Neither route
licenses a change of law or of physical clock.

\section{Orbit coercivity provides a legitimate static-to-dynamic incidence}
For a finite synthesis $V$ into the form domain of the actual $H$, define
\[
G(s)=V^*e^{-2sH}V,\qquad
K(s)=V^*e^{-sH}He^{-sH}V.
\]
Spectral calculus gives
\begin{equation}\label{eq:orbitidentity}
G'(s)=-2K(s),\qquad
G(\tau)=G(0)-2\int_0^\tau K(s)\dd s.
\end{equation}

\begin{criterion}[Orbit coercivity with a residual]\label{thm:orbit}
Suppose $K(s)\succeq mG(s)-r(s)R$ for almost every $s\in[0,\tau]$, where $m>0$, $R\succeq0$, and $r\ge0$ is integrable. Then
\begin{equation}\label{eq:orbitbound}
G(\tau)\preceq e^{-2m\tau}G(0)
+2\left(\int_0^\tau e^{-2m(\tau-s)}r(s)\dd s\right)R.
\end{equation}
If $R\preceq CG(0)$ and the resulting scalar coefficient is below one, this supplies the delayed Gram contraction.
\end{criterion}
\begin{proof}
Test on a coefficient vector, use \cref{eq:orbitidentity}, multiply the scalar differential inequality by $e^{2ms}$, and integrate. \src{F1}{V2}
\end{proof}

An entrance energy bound alone is not orbit coercivity: the evolving normalized state can concentrate in a slow component. Alternatively, a static source form $\mathbb N_j$ can contribute only after a proved relation
\begin{equation}\label{eq:static-incidence}
G_j^0-G_j^\tau\succeq c_\tau R_j^*\mathbb N_jR_j-E_j
\end{equation}
for the same history law and physical delay. Combining $R_j^*\mathbb N_jR_j\succeq\delta G_j^0$ and $E_j\preceq\varepsilon_jG_j^0$ gives the contraction factor $1-c_\tau\delta+\varepsilon_j$. Whitening a static matrix does not construct \cref{eq:static-incidence}.

\section{Vacuum-corrected Feshbach coercivity}
Let
\[
H=\begin{pmatrix}A&B\\B^*&C\end{pmatrix},\qquad C\succeq\gamma I>0,
\]
and set
\begin{equation}\label{eq:Feshbach}
S=A-BC^{-1}B^*,\quad \mathcal R=C^{-1}B^*,\quad
Jx=(x,-\mathcal Rx),\quad M=J^*J=I+\mathcal R^*\mathcal R.
\end{equation}
Then $\Ker H=J\Ker S$. If $\Omega=Jx_0$ is a normalized vacuum, the correct head complement is $x_0^{\perp_M}$, not the complement of a guessed bare constant.

\begin{criterion}[Vacuum-corrected effective head]\label{thm:Feshbach}
If $\Ker S=\C x_0$ and
\[
\ip{x}{Sx}\ge\delta\ip{x}{Mx}\qquad(x\perp_Mx_0),
\qquad\delta>0,
\]
then
\begin{equation}\label{eq:Feshbachgap}
H|_{\Omega^\perp}\succeq\frac{\delta\gamma}{\delta+\gamma}I.
\end{equation}
The same estimate applies on an orbit set when the head inequality is required only for the actual centered head coordinates in that set.
\end{criterion}
\begin{proof}
Write a centered vector as a graph part plus its innovation, and remove the graph-vacuum component in the $M$ metric. If $a$ is the norm of the centered graph part and $b$ the innovation norm, Schur completion gives energy at least $\delta a^2+\gamma b^2$, while the centered norm is at most $a+b$. Weighted Cauchy--Schwarz yields
\[
(a+b)^2\le(\delta a^2+\gamma b^2)(\delta^{-1}+\gamma^{-1}),
\]
which proves \cref{eq:Feshbachgap}. \src{F1}{V3--V4}
\end{proof}

\subsection{A retained infrared obstruction survives fine coercivity}
The Feshbach criterion has a converse obstruction at the level of a
nonnegative transfer form. It explains why a large detail floor cannot
replace the head condition.

\begin{proposition}[Massless Schur edge]
\label{prop:massless-Schur-edge}
On a declared centered carrier, let the bounded nonnegative form
$L=I-P$ have blocks
\[
L=\begin{pmatrix}A&B\\B^*&C\end{pmatrix},\qquad C\succeq cI>0,
\qquad S_L=A-BC^{-1}B^*.
\]
If there are unit retained vectors $x_n$ with
$\ip{x_n}{S_Lx_n}\to0$, then $L$ has no positive lower spectral bound.
For a family, such a sequence rules out a common positive
dimensionless lower bound on its respective centered carriers.
\end{proposition}
\begin{proof}
Lift $x_n$ by the minimizing detail:
$\xi_n=(x_n,-C^{-1}B^*x_n)$. Schur completion gives
$\ip{\xi_n}{L\xi_n}=\ip{x_n}{S_Lx_n}$, while
$\norm{\xi_n}\ge1$. The Rayleigh quotients of the normalized lifts
therefore tend to zero. \src{F6}{V57}
\end{proof}

For a cofinal physical application, the vanishing Rayleigh quotient
must be compared with the step duration using
\cref{eq:two-level-physical-mass,eq:physicalrate}. A dimensionless
deficit tending to zero can still be of order $\delta_j$ and therefore
be compatible with a positive physical mass. The obstruction is not
used without this rate comparison.

If $C\succeq0$ has a null space, positivity first forces
$B n=0$ for every $n\in\Ker C$: test the quadratic form on
$(x,tn)$ for both signs of real $t$, and then on $(\ii x,tn)$.
Only when $C$ has closed range and a positive lower bound on its
orthogonal complement may the same completion use $C^+$.
Gauge directions must therefore be removed or retained in the declared
physical carrier before inversion. This restriction is the transfer-form
analogue of the transverse inverse convention in \cref{ch:tower}.

The Casimir analysis gives finite ultraviolet floors and exploits the finite representation bandwidth of a Wilson plaquette. Nested cuts shield an inner Casimir region from a remote router, leaving mixed terms in a finite collar. These are useful exact reductions. A total magnetic norm extensive in volume is not thereby converted into a cofinal infrared invariant; growing-head transport needs its collective tail and mixed-stock bounds. \src{F1}{V3--V6}

\section{Common-time source transport and signed return}
Let $J$ be an isometry from an old centered source range into a refined range, and let $T_X,T_Y$ be transfers at the same physical time. Set $P=JJ^*$ and
\[
M_1=J^*T_YJ,\quad M_2=J^*T_Y^2J,\quad
C=M_1-T_X,\quad L=(I-P)T_YJ.
\]
Then the exact source--detail identity is
\begin{equation}\label{eq:transport-defect}
T_YJ-JT_X=L+JC,\qquad
(T_YJ-JT_X)^*(T_YJ-JT_X)=\mathbb V+C^2,
\end{equation}
where $\mathbb V=L^*L=M_2-M_1^2\succeq0$.

If the refined tail block $D$ obeys $D\preceq dI$ with $d<q$, the Schur complement gives the sufficient threshold condition
\begin{equation}\label{eq:transport-threshold}
qI-M_1-\frac{\mathbb V}{q-d}\succeq0
\quad\Longrightarrow\quad T_Y\preceq qI.
\end{equation}
The compression mismatch $C$ is signed. A negative direction can increase the threshold reserve, and replacing it everywhere by $\norm C I$ loses that information. This is the precise common-time alternative to a generic norm-only cutoff transport. \cite{GT}

For a simpler positive block transfer $T=\left(\begin{smallmatrix}A&B\\B^*&D\end{smallmatrix}\right)$, $\norm A=a<1$, $\norm B=b$, and $\norm D\le1-\mu$, the condition $b^2<(1-a)\mu$ gives
\begin{equation}\label{eq:2x2-transfer}
\norm T\le\frac{a+1-\mu+\sqrt{(a-1+\mu)^2+4b^2}}2<1.
\end{equation}
Thus a source-head contraction needs an explicit complement reserve and mixing bound. A finite bank cannot represent the full gap by itself.

\section{A complete effective-action terminal criterion}
For a compact effective Gibbs interaction $\Phi=\{\Phi_X\}$, define the mixed rectangular oscillation $\partial_{ef}\Phi_X$ as the supremum of
\[
|\Phi_X(u_e,u_f)-\Phi_X(u'_e,u_f)-\Phi_X(u_e,u'_f)+\Phi_X(u'_e,u'_f)|
\]
over the displayed variables and the remaining exterior. Put
\[
\kappa_{ef}(\Phi)=\sum_{X\ni e,f}\partial_{ef}\Phi_X,
\qquad
\mathfrak D_\mu(\Phi)=\sup_e\sum_{f\ne e}e^{\mu d_E(e,f)}
\frac{e^{\kappa_{ef}(\Phi)}-1}{2}.
\]

\begin{criterion}[Weighted Dobrushin terminal ball]\label{thm:Dobrushin}
If $\mathfrak D_\mu(\Phi)=\alpha_\mu<1$ for some $\mu>0$, then the complete Gibbs specification has a unique infinite-volume state and
\begin{equation}\label{eq:Dobrushin-cov}
|\Cov(F,G)|\le
\frac{4|S_F|\norm F_\infty\norm G_\infty}{1-\alpha_\mu}
 e^{-\mu\dist(S_F,S_G)}
\end{equation}
for bounded cylinder observables. If that state is reflection positive at terminal physical spacing $a_*$ and the gauge-invariant cylinders generate its physical OS space, its Hamiltonian gap is at least $\mu/a_*$.
\end{criterion}
\begin{proof}[Proof structure]
A change in exterior link $f$ tilts the conditional density at $e$ by a function with oscillation at most $\kappa_{ef}$. The conditional influence is bounded by $(e^{\kappa_{ef}}-1)/2$. Its weighted row norm is below one, so the comparison matrix has convergent series $\sum_{n\ge0}C^n$ with weighted norm at most $(1-\alpha_\mu)^{-1}$. The coupling comparison gives \cref{eq:Dobrushin-cov}. Exponential decay under physical time translations, reflection positivity, and the dense local core imply the spectral bound. \src{F7}{V82--V83}
\end{proof}

The effective interaction here is complete, including its generated quasi-local terms and retained labels. A regular conditional fibre, an unevaluated reset reference, or a truncation without a tail estimate is not the $\Phi$ in this theorem. The rare-defect and bridge-resolvent developments give alternative stability criteria around already gapped regular kernels; they still require their selected-path, response, reflection, and physical-clock hypotheses. \src{F7}{V84--V91}

\section{An explicit strong-electric region}
For the finite-volume $\SU(3)$ heat-kernel/Wilson slab, let $\epsilon_\kappa=\norm{K_\kappa-1}_\infty$. The source proves that
\[
2\epsilon_\kappa+12\tanh\eta<1
\]
implies a volume-uniform neutral-sector contraction
\begin{equation}\label{eq:strong-electric}
q\le\frac{1-12\tanh\eta-
\sqrt{(1-12\tanh\eta)^2-4\epsilon_\kappa^2}}
{2\epsilon_\kappa}<1,
\end{equation}
with its continuous interpretation at $\epsilon_\kappa=0$. In particular, $\kappa\ge4$ and $0\le\eta\le1/18$ give $q\le9/13$. The temporal influence and the twelve same-slice magnetic neighbors provide the comparison row, and the Doob transform identifies the resulting boundary contraction with the physical slab transfer. \cite{GT}

This supplies a concrete terminal domain. It does not establish that the weak-coupling continuum trajectory enters that domain after exact renormalization, nor that all generated interactions satisfy the same bound. A proof using this terminal route must establish that bridge.
\sourcebasis{Local Gram, quotient passage, physical time, and orbit coercivity: \src{F1}{V1--V2}. Vacuum-corrected Feshbach and ultraviolet shielding: \src{F1}{V3--V6}. Common-time signed transport and strong-electric criterion: \cite{GT}. Complete-action terminal clustering: \src{F7}{V82--V91}.}

\chapter{The integrated conditional closure theorem}
\label{ch:closure}
\section{The sufficient packet}
The preceding results isolate a short logical endpoint. It is useful to state it without requiring any one historical sufficient mechanism for the final contraction.

\begin{hypothesis}[Common compact construction]\label{hyp:compact}
For the chosen gauge group and declared local-vacuum branch, there is a cofinal regulator family of actual finite renewal-loaded compact gauge laws with compatible source writers and exact blocking. All generated interactions, Haar factors, quotient data, calibration marks, and required sector labels are retained. The comparison objects used in estimates are connected to these laws by explicit proved identities or controlled defects.
\end{hypothesis}

\begin{hypothesis}[Renormalized local continuum theory]\label{hyp:continuum}
The family has a physically calibrated renormalized trajectory and convergent local gauge-invariant Euclidean correlations. The limiting correlations satisfy the required Euclidean symmetry, regularity, reflection positivity, and continuity properties for the stated Yang--Mills reconstruction. The coupling is fixed by a nonzero physical normalization at a declared scale, rather than by an uncontrolled over-weak bare limit. The resulting theory is nontrivial, and bounded local gauge-invariant observables generate its physical Hilbert space.
\end{hypothesis}

\begin{hypothesis}[Noncollapsed local probes]\label{hyp:noncollapse}
The centered local source banks have well-defined limiting entrance and delayed Gram matrices. At least the specified nontrivial local states have nonzero limiting norm; lower frames are supplied where physical coefficient directions are required to survive. The source identification and required non-Gaussian response packet establish the intended interacting Yang--Mills branch rather than a zero or purely auxiliary theory.
\end{hypothesis}

\begin{hypothesis}[Uniform physical contraction]\label{hyp:contraction}
There exist one physical $\tau>0$ and one $0<q<1$ such that for every finite centered bounded-support local source family,
\[
G_{j,\mathcal F}^{\tau_j}\preceq q^2G_{j,\mathcal F}^0+E_{j,\mathcal F},
\qquad \tau_j\to\tau,\qquad \norm{E_{j,\mathcal F}}\to0.
\]
The same $q$ applies to every family, although convergence rates may depend on the fixed family.
\end{hypothesis}

\begin{theorem}[Conditional local-vacuum Yang--Mills closure]\label{thm:master}
Under \cref{hyp:compact,hyp:continuum,hyp:noncollapse,hyp:contraction}, the reconstructed local Yang--Mills theory has a unique vacuum and
\begin{equation}\label{eq:master-gap}
\boxed{\displaystyle
H_\infty|_{\Omega_\infty^\perp}\succeq m_*I,
\qquad m_*=-\frac{\log q}{\tau}>0.}
\end{equation}
\end{theorem}
\begin{proof}
The first three hypotheses give the stated nontrivial local OS theory and its dense centered local source span. For each fixed family, pass the Loewner inequality of \cref{hyp:contraction} to the limiting Grams using \cref{eq:Gram-convergence,eq:Gram-error}. It holds with the same $q$ for every finite family. Apply \cref{thm:local-Gram}. No inverse entrance Gram, fixed source dimension, or independent auxiliary generator is needed. The upstream construction and renormalization assumptions identify the Yang--Mills application of this local closure implication. \src{F1}{V1}
\end{proof}

\begin{corollary}[Conditional closure after a locally invisible soft band]
\label{cor:closure-local-soft-band}
Retain \cref{hyp:compact,hyp:continuum,hyp:noncollapse}. Suppose
that for every fixed finite local source family the scalar-compression
hypothesis and complement contraction of
\cref{thm:local-toron-Gram-excision} hold on the same physical
trajectory, with $t_j\to\tau>0$, $q_j\to q<1$, and $E_j\to0$.
Assume the same $q$ and $\tau$ apply to all families. Then the
conclusion of \cref{thm:master} holds, even if the discarded
finite-regulator reducing bands contain arbitrarily small positive
energies.
\end{corollary}
\begin{proof}
For each fixed family the additional error in
\cref{eq:local-toron-full-contraction} tends to zero. Boundedness
of its convergent entrance Gram absorbs $q_j^2-q^2$. Thus
\cref{hyp:contraction} holds and \cref{thm:master} applies.
Only a vanishing source contribution has been removed; no state of
the limiting noncollapsed local Hilbert space is deleted.
\end{proof}

\section{Established inputs and their precise limits}
The explicit finite inputs do not have a single common limiting
regime. Their contribution to \cref{thm:master} can nevertheless be
stated without repeating their proofs. Exact compact completion and
measured redundancy fix the finite action and its density. Mixed-shell
Schur ownership fixes determinant accounting. The local branch has a
moving minimum, an explicit measured two-loop coefficient, and an
$O(\beta^{-2})$ remainder through four source derivatives under the
stated higher-order hypotheses. The finite coefficient map composes
with the complete incoming density and closes
$(J^8S,J^6Q,J^4B;J^7F)$; its algebraic closure does not extend those
analytic hypotheses to a new scale automatically.

The balanced root has a depth-independent small counting-$\ell^4$
domain. Finite-resolution approximation makes its continuum source
weakly continuous, and exact-source recovery with constrained
convexity gives strong convergence of the full local classical
minima. This supplies the common branch rather than merely a
homogeneous trial or a formal source jet. On that branch the cofinal
analysis supplies a complete measured determinant, a convergent
subtracted trace tail, an exact flat-holonomy limit, and an evaluated
noncommuting quartic logarithm with a finite part. The low-order Born terms, source matrix, and endpoint series
are retained. This resolves the corresponding measured coefficient
calculation. The fourth-source finite part independently controls a
second-chaos counterterm and a summable remainder on the free
Gaussian carrier. It retains the complete matrix filter, not just
the measured scalar trace. Neither result proves joint nonlinear
cutoff control.

For the unrestricted bare Wilson law, the selected source window
controls arbitrary cell-energy arrays without a support-size loss.
Its entropy inequality gives a quantitative comparison mechanism
whenever the actual changed law has a controlled information budget.
For a fixed source-independent block, the exact likelihood and entropy
chain rule identify both the generated law and the conditional cost.
Conditional covariance bounds and the regression test distinguish
upper control of retained information from its survival. The measured
linearization estimate is a probability-law approximation, not a
pointwise local-action bound. Neither the selected window nor these
finite pushforward identities construct a physically calibrated
interacting trajectory.

The compact spectral certificates imply both local tube coercivity
and hidden metastability. Same-law original-clock repair remains
valid, including the inverse-paired coplanar family through the full
cutoff transition. The all-exterior weighted patch criterion still
requires a complete conditional gap, not merely these local events.

The native Wilson transfer has exact magnetic and affinity identities.
The Gaussian image's exact occupation law requires an extensive total
old-oscillator degree budget, while its genuine Gram-normalized
compact projections converge at fixed volume. Neither conclusion
bounds the full-transfer complement. Separately, volume-uniform based
geometry and an all-label comparison retain the two actual Perron
values. Measured recovery and no-escape
estimates extend the ordered critical spectrum to
$\Theta_{\gamma,L}\to0$. A neutral lattice gap and charged center
splittings remain distinct, and bare calibration excludes identifying
this joint domain with the usual fixed-box physical trajectory.
The exact Perron equation and local kinetic and cubic identities
specify further analytic requirements without supplying a global
native fluctuation theorem.

The local endpoint is strengthened in two complementary directions.
Physically bounded Wilson loops give a scalar-compression mechanism
on controlled toron bands; the Gram-excision theorem states exactly
how a separately proved complement contraction would imply the full
local criterion. The selected vacuum's canonically smeared probe,
on the other hand, has surviving soft spectral weight. It therefore
rules out a massive interpretation of that particular mesoscopic
window with that probe retained. These results preserve rather than
replace the separate continuum and physical-contraction hypotheses.

\section{A conditional construction route through the Gaussian tower}
The explicit Gaussian construction contributes to
\cref{hyp:continuum} through the following sufficient packet:
\begin{enumerate}[label=\textup{(R\arabic*)},leftmargin=3em]
\item A compatible exact compact blocking realizes the tower as its
principal Gaussian family, with a common analytic domain, declared
quotient and sector transport, and control of the complete transported
vertices. The one-step smooth construction of \cref{cor:compact-product}
provides a concrete candidate carrier, but its common germ does not by
itself prove this all-level compatibility.
\item The fluctuation covariance is placed in the representative of the
vertex banks, with common bounds on the growing alias trace and its
external derivatives, including the treatment of the toron region.
\item The full-torus complete-principal response, including Schur-owned
mixed returns and its endpoint, has the consistently normalized upward
coefficient $b^*=b_{\rm YM}$. It is evaluated either in
the harmonic background of the exact blocked action or in the fixed
background together with \cref{eq:background-Cesaro}. An equivalent
SCIC/ADMC comparison is sufficient; equality of corner cards is not.
The measured $\SU(2)$ quartic result of
\cref{thm:measured-quartic-log} supplies one evaluated input, with its
own explicit normalization, rather than automatically verifying this
entire physical comparison.
\item Nonlinear and stable-coordinate estimates remain uniform along the
same family. The finite packet of
\cref{thm:closed-loop-jet-packet} is transported with its complete
incoming density, common domains, and marked norms. The degree-faithful
forest growth condition, complete four-mark source activities, and
local current jets control the measured Jacobian response. Owner bounds, local sector control, and reset or
boundary terms satisfy their common regulator hypotheses. The principal margin used in remainder protection
belongs to that same physically matched response.
\item The resulting local correlations have the compactness, regularity,
reflection, continuity, and nontriviality required by
\cref{hyp:continuum,hyp:noncollapse}.
\end{enumerate}
Under (R1)--(R4), the fixed-point response criterion, the background
comparison where used, the owner remainder, the stable cone, and the
Ces\`aro recursion provide the intended weak-coupling marginal asymptotics
and error bounds. Condition (R5) remains the continuum-existence handoff.
The exact finite representative map supplies an algebraic part of (R2);
it does not supply its uniform analytic estimates or (R3). None of these
five requirements replaces the fixed-physical-time contraction in
\cref{hyp:contraction}.

A spectral route can then use the complete-action terminal ball, an
actual-law response/capacity bound with same-generator incidence, the
two-level form criterion, the orbit-Feshbach criterion, or direct local
Gram estimates. A finite source
head requires its complement and mixed return. An effective terminal
action requires a demonstrated flow into its contraction domain. These
are alternative sufficient mechanisms for the same physical endpoint.

\subsection{The marked local input to the compact continuation}
The analytic part of (R4) has a more explicit internal structure than a
single unspecified remainder estimate. Exact blocking supplies the
conditional differential and its retained unit sector
(\cref{thm:exact-block-differential}). Measured field-redefinition
tangents pass to the retained quotient by
\cref{thm:measured-redundancy}, through the conditional current of the
same law. The Jacobian Hessian spends the third jet of that current
(\cref{prop:Jacobian-third-jet}), and the normalized score hierarchy
identifies the connected cumulants which must be controlled.

Finite measured composition now specifies which primitive coefficients
must be retained: \cref{thm:closed-loop-jet-packet} fixes the
$(8,6,4;7)$ field-degree budget, including the density terms generated
at preceding stages. Independently,
\cref{thm:generated-likelihood,prop:retained-source-information}
identify the actual output law and its fine-source derivatives. Their
information bounds neither control retained-coordinate current jets
nor establish locality of the generated interaction.


Once the actual compact source representation, local jet bounds,
analytic radii, and contour activities in
\cref{crit:nonlinear-current-gate} are supplied uniformly, the
four-mark theorem provides those current and Jacobian estimates. The
Gaussian covariance bounds and degree-faithful forest census are
upstream inputs to this marked representation, not substitutes for
it. In particular, this chain does not infer a uniform conditional
four-point theorem from a quadratic vacuum Hessian or from pair decay.

On the spectral side, \cref{crit:two-level-physical} separates the
conditional detail estimate from the retained physical deficit, while
\cref{prop:massless-Schur-edge} records the obstruction when the retained
form remains infrared soft. The core--bad-set compiler can supply an
upstream energy bound only with its physical-law, tensorization, and
incidence assumptions. These refinements leave the terminal closure
packet unchanged: construction, noncollapse, and a common physical
contraction must ultimately hold for the same theory.

\section{Obstructions that determine the admissible route}
\subsection{The massless balanced baseline}
For the Gaussian bridge analyzed in the Grand Tensor, both the reflected head and its balanced reference can equal $2\beta/L$ despite the model being massless. A positive raw balanced ratio is therefore not a gap test. The meaningful quantity is the reference-subtracted normalized excess. In the scalar Mehler calibration,
\[
\delta=\sinh^2(x/2)=\frac{(1-e^{-x})^2}{4e^{-x}},
\qquad
e^{-x}=(\sqrt{1+\delta}-\sqrt\delta)^2.
\]
A strictly positive excess then has a calibrated meaning in that model. The identity is a baseline test, not a universal substitution for the Yang--Mills delayed Gram. \cite{GT}

\subsection{The Gaussian infrared channel}
For the transverse quadratic colour theory,
\[
\omega_a(k)=a^{-1}\operatorname{arcosh}(1+q(k)/2),\qquad
q(k)=4\sum_{\ell\ {\rm spatial}}\sin^2(k_\ell/2).
\]
At fixed physical torus size $L$, the lowest nonzero frequency is $2\pi/L+O(a^2L^{-3})$. As $L\to\infty$, its half-slab correlation approaches one. A norm-small perturbation of that infrared Gaussian channel cannot open an order-one uniform contraction below one. The ultraviolet Gaussian construction must therefore be connected to a genuinely interacting infrared mechanism; its Gaussian positivity is not enough. \cite{GT}

\subsection{The over-weak toron trajectory}
The Grand Tensor constructs a center-neutral wrapping soft sequence on a specified over-weak bare trajectory. In its parameters, normalized vectors satisfy an energy asymptotic of order
\[
\ip{v_j}{H_j^{\rm phys}v_j}
=\frac{\mu_8+o(1)}{LK_j^{1/3}},\qquad K_j\to\infty,\quad\mu_8>0.
\]
At every fixed physical delay, their squared evolved norms tend to one. This rules out a uniform full finite-torus gap on that trajectory. It does not rule out a correctly renormalized local infinite-volume theory. A replacement trajectory must enforce a nonzero physical coupling condition and its step scaling; a replacement algebra must be declared openly. Wrapping vectors cannot simply be deleted from a theorem that originally quantified over them. \cite{GT}; \src{F1}{V1}

\subsection{Invisible toron modes and surviving local soft spectrum}
The geometric loop estimate of \cref{prop:local-toron-loops} holds
uniformly for lengths of order $L$, but its transfer to actual finite
toron bands is proved only on the joint critical domain. Even there,
a physical source renormalization must preserve scalar compression.
The exact Gram excision theorem is therefore a conditional local
route, not permission to delete wrapping sectors from a global
finite-torus assertion.

Conversely, \cref{thm:smeared-probe-spectrum} produces nonzero
limiting damped source norms and positive weight in every interval
$(0,E)$ of the mesoscopically rescaled generator. This sharpens a
bare lattice-gap-collapse test: the soft modes are visible to the
declared probe. The selected branch and normalization cannot supply
\cref{hyp:contraction}; a genuinely interacting physical scale and
its source identification require separate construction.

\subsection{The actual compact hidden-well obstruction}
For the specified $\SU(2)$ parent, \cref{thm:hidden-metastability}
excludes a temperature-uniform full hidden diffusion gap. Consequently
that gap cannot be retained as an already acquired premise in a
low-temperature source-transport argument. The obstruction does not
exclude a same-law heat-bath route, a localized marginal comparison,
or a different complete physical criterion. It does exclude inferring
a global hidden gap from normal positivity at a single orbit or from
the small probability of the second well.

\subsection{Finite banks and continuum visibility}
A fixed bank can converge while normalized near-null directions outside it remain soft. A source-complete argument therefore requires either all fixed finite local families with one contraction constant, as in \cref{hyp:contraction}, or a collective tail mechanism in addition to a growing finite head. For example, the source graph-Rellich condition
\[
\norm{(I-\Pi_{j,m})v}^2
\le\eta_m\bigl(\ip{v}{H_jv}+\norm v^2\bigr),\qquad\eta_m\downarrow0,
\]
on the stated branch supplies a quantitative high-energy complement. It is a sufficient separate input, not a consequence of finite-dimensional diagonalization. One should not impose global compactness unnecessarily on a local infinite-volume route when the all-local Gram theorem already applies. \cite{GT}; \src{F1}{V1--V5}

\section{The exact remaining obligations}
\subsection{Physical matching of the evaluated measured response}
The nested covariance tower, representative map, mixed-shell ownership,
and complete one-cell quartic logarithm are explicit. The unresolved
point is no longer evaluation of the leading measured matrix in the
balanced-root $\SU(2)$ branch. It is a compatible comparison to the
physical coupling and background, with uniform momentum and source
control, growing alias traces, toron treatment, and endpoint accounting.
Reuse of fixed-background banks still requires their averaged defect
relative to the harmonic background to be controlled. Neither the
finite part nor the convergent Schatten tail deletes the low-order
measured terms which enter that comparison. The finite incoming
fourth-source counterterm must also remain in its subtraction:
convergence of a counterterm without a rate does not permit its
replacement by the limit before multiplication by the logarithmic
scale.

\subsection{Uniform nonlinear transport and the compact exterior}
The local two-loop remainder becomes summable along a linearly growing
reciprocal coupling only when its hypotheses and transport constants
remain uniform. Its $\beta^{-1}\mathcal B_1$ term generally remains in
the running action. One must control the genuine complete density,
its normalized source ratios, Haar factors, and large-field part on
the same trajectory; a truncated polynomial exponent or a small
exceptional-set probability is not such control. The finite incoming
jet packet is already closed by \cref{thm:closed-loop-jet-packet};
what remains is its uniform transport in the required domains and
marked norms, together with the degree-faithful forests, conditional
current bounds, and owner decomposition. Likewise, the strong local
classical limit of \cref{thm:classical-minimum-limit} does not supply
these quantum or compact-exterior estimates.

The unrestricted bare-Wilson source window is an additional tool,
not a universal law replacement. \Cref{thm:generated-likelihood}
identifies its exact transport under a fixed block, with separate
conditional and output entropy costs. Using that comparison for a
different blocked or renewal-loaded measure still requires a proved
budget, a compatible source map, and the correct coupling regime.
Smallness of a conditional error alone does not control the output
normalizer, and no small entropy budget for a general renormalization
step is assumed.

\subsection{Complete same-law interface control}
The certified repair families now include inverse-paired noncommuting
sources through the entire cutoff transition. Arbitrary non-paired
sources, configurations outside the response cone, and other large
active-source patterns remain to be covered or controlled by another
mechanism. A proof must assemble a genuine same-law cover with
controlled losses or establish all-exterior conditional patch gaps
strong enough for \cref{crit:weighted-patches}. Its explicit sufficient
threshold is $\gamma_\ell>C_*\ell^{-3/2}$; the available surface-order
certificate does not verify it. A growing spatial union, a rare bad
set, or a temperature-uniform hidden diffusion shortcut cannot replace
that form estimate.

\subsection{Beyond the joint critical domain}
The ordered native spectrum is established both at fixed $L$ and
under $\Theta_{\gamma,L}\to0$. The latter condition is a genuine
growing-regulator result, but excludes the exponential lattice growth
of the fixed-box scale premise in
\cref{prop:native-physical-domain-separation}. Continuation toward
that regime requires a new uniform recovery and localization
argument or a different local physical criterion. The sufficient
packet \cref{eq:native-uniform-packet} keeps complement deficit and
boundary return distinct; neither a per-link Gaussian estimate nor
an interior Perron bound proves both.

The Gaussian occupation law of \cref{thm:magnetic-excitation-law}
also excludes a subextensive cutoff in the original total oscillator
degree as a uniform carrier for the magnetic image. Extensive bands
have a controlled Gaussian tail, but
\cref{prop:magnetic-compact-projection} takes the compact limit at
fixed volume and fixed chosen degree. A joint-regulator construction
must control that passage uniformly as well as the complement and
mixed return of the actual transfer.

The later electric-Hamiltonian analysis improves a different part of
this interface.  \Cref{thm:native-high-Schur} gives a genuinely
volume-independent high-harmonic complement and factorial tails for an
entire fixed low-energy window, while
\cref{thm:native-electric-gap,thm:native-plaquette-band} give an actual
uniform lattice gap and analytic low band in a strong-electric domain.
These results remove finite-rank truncation as an obstruction in that
regime, but do not transport the gap to the weak-coupling physical
trajectory.  Moreover \cref{eq:native-resonant-matrix-element} shows
that a whole-head homological inverse fails exactly; any continuation
must retain the resonant dynamics or use a different block variable.
The local normal form of \cref{thm:native-resonant-normal-form} provides
one controlled finite-lattice organization, not the missing continuum
transport.


A probabilistically enlarged sup-norm region further requires native
concentration and a valid gauge representation. Mean-one kinetic
centering must be transported through the complete magnetic and
Perron-weighted law. The centered cubic is a third-chaos observable,
so its covariance alone gives neither a lognormal identity nor a
relative operator-norm theorem. A connected estimate must retain
all plaquette orientations and stabilizing higher magnetic orders.

\subsection{Physical local contraction and source survival}
The interacting local correlations must satisfy the regularity,
reflection, nontriviality, and continuity required for reconstruction.
A mass bound needs one contraction constant on every finite centered
local reflected family at one fixed physical delay, on that same
calibrated branch. The soft-band route additionally requires actual
scalar compression after source normalization, a reducing projection,
and a complement contraction; the smallness of raw loop variations
alone does not verify this packet.

The surviving soft spectral measure of
\cref{thm:smeared-probe-spectrum} is a test that an admissible
physical replacement must pass rather than ignore. Neither its
selected mesoscopic normalization nor a bounded critical toron
coefficient supplies the required physical mass. A proved
physical-time and source transport remains essential.

The exact finite-time conditional coefficient
\cref{eq:native-cost-C2,eq:native-cost-signs} further shows that this
transport cannot be obtained from a coarse-field-independent
monotonicity of the native conditional cost.  A valid argument must
control the source-specific Poisson Hessian, coarse return, and
finite-time remainder on the actual law rather than infer them from
the sign of the zero-time jet at one background.

\section{Interpretation}
The construction separates the sources of physical content. The common
finite law determines which source responses and conditional operations
belong to the theory. Exact blocking determines the action and its
normalization. The Gaussian tower supplies a controlled multiscale
principal object, while the representative map specifies how its shells
enter the finite Wilson response. The background comparison determines
whether that response is the marginal derivative of the intended
blocked action. The explicit nonlinear fibre then fixes the actual
constraint, moving minimum, and measure derivatives. Its critical-norm
source limit and exact-source recovery provide a genuine local
classical continuum branch. The local two-loop theorem states which
uniform source norms are controlled, while finite measured composition
specifies the density and field jets to be carried to the next step.
Complete measured determinant assembly evaluates the noncommuting
quartic logarithm without discarding its endpoints. Full compact
pressure and entropy transport add source bounds on a separately
declared law, with exact generated likelihoods keeping its output
normalization and retained information explicit.
The compact spectral analysis retains the higher well rather than
confusing its small probability with fast mixing; the visible repair
mechanism works through the original marginal conditionals. Measured field changes remain redundant after blocking,
but their local analytic control is carried by the complete normalized
current, not by an omitted Jacobian. The native transfer analysis then retains the magnetic image,
conditional affinity, the occupation budget of its Gaussian image,
actual Perron normalization, and controlled fixed and joint critical
spectra. Its resistance geometry avoids
charging a global grounded inverse to each local gauge evaluation;
its full magnetic comparison retains the extensive vacuum debit.
The regional reduction and exact Perron equation keep further
complement, boundary, and fluctuation control explicit. Local Gram
excision distinguishes a source-invisible soft band from the
surviving soft spectrum of a normalized smeared probe. Finally,
the two-level form retains the infrared channel, and the complete
local Gram criterion determines what physical spectral decay must mean.

This separation turns finite positivity, analyticity, and computational
identities into usable components without promoting them beyond their
scope. The conditional conclusion is quantitative: once the interacting
continuum law and a noncollapsed fixed-time local contraction are
established on the same calibrated branch, the Hamiltonian has the
strictly positive mass bound of \cref{thm:master}. The remaining work is
localized to the displayed analytic and physical interfaces, rather than
hidden in changes of law, representative, background, or clock.

\appendix
\chapter{Notation and interface ledger}
\label{app:notation}
\section{Principal notation}
\begin{longtable}{L{0.23\textwidth}L{0.70\textwidth}}
\toprule Symbol & Meaning and restriction \\ \midrule
\endfirsthead
\toprule Symbol & Meaning and restriction \\ \midrule
\endhead
$\mu,\pi,\omega$ & Respectively a declared finite law, the actual Perron ground law when so identified, and a candidate limiting Euclidean state. They are not interchangeable. \\
$a_j,\tau$ & Lattice time spacing and a fixed physical delay. The continuum gap uses $\tau$ or the rate $-a_j^{-1}\log\rho_j$. \\
$u_j,\beta_j$ & Physical coupling coordinate and its reciprocal, $u_j=\beta_j^{-1}$, in the declared calibration. A polymer activity denoted by $u$ elsewhere must be distinguished by context. \\
$G_{\mathcal F}^0,G_{\mathcal F}^\tau$ & Entrance and delayed local reflected Grams, with the delay matrix containing $e^{-2\tau H}$. \\
$\mathcal A,\mathcal V,\mathcal K,\mathcal M$ & Positive one-form comparison, defect synthesis, relative return channel, and physical landing metric in \cref{eq:woodbury-objects}. \\
$S,M$ & Feshbach head and its graph metric $M=I+\mathcal R^*\mathcal R$. Vacuum shortening is in the $M$ metric. \\
$\theta_{ij}$ & Normalized coordinate-mark overlap of \cref{eq:marks}; not a raw Fourier coefficient or a representation-dimension probability. \\
$z,K_{\rm HK}$ & Primitive compact analytic activity and the complete weighted forest contraction stock of \cref{eq:KHK}. \\
$\kappa_a,\Pi(I)$ & Normalized joint cumulants and labelled set partitions in \cref{eq:connected-score-hierarchy}; this $\Pi(I)$ is not the transverse momentum projection. \\
$\mathcal O_V,\mathcal J_V,\overline V$ & Measured field-redefinition tangent, pushed vector density, and normalized conditional current in \cref{thm:measured-redundancy}. \\
$\mathsf K_S,\mathsf J_P$ & Conditional expectation and retained pullback for the exact action differential, \cref{thm:exact-block-differential}. \\
$F_a,L_a,A(a),\mathcal R_a$ & Conditional cumulant, normalized generated likelihood, retained score, and nonlinear source correction in \cref{sec:generated-source-laws}; all use the actual baseline law. \\

$Q_{a,\mu}(u),\Delta$ & Rooted contour-activity majorant and actual contour degree; filling neighbourhoods require $\Delta=80$. \\
$\gamma_P,\gamma_Q,\delta$ & Full and retained dimensionless transfer deficits and physical step duration in \cref{crit:two-level-physical}. \\

$\mathfrak b$ & The calibrated transverse $F^2$ response extractor with its action, colour, and time-orientation convention fixed. \\
ADMC & Averaged determinant matching card, \cref{def:ADMC}. \\
SCIC & Sublinear compact--continuum intertwiner card, \cref{def:SCIC}. \\
$\Pi,\lambda,G^{(j)},S^{(j)}$ & Free transverse projection, lattice dispersion, level-$j$ tower covariance, and transverse inverse kernel. \\
$\Gamma_\mu^*,\Gamma_\ell$ & A diagonal component of the fixed-point lattice sum and, separately, a positive lifted covariance shell. The index and context distinguish them. \\
$\kappa_{\rm W}$ & Finite Wilson relative derivative ceiling below $4/3$, not a stable-flow contraction constant below one. \\
$T_{{\rm g},j},\Xi_j$ & Nested tree-gauge representative and the covariance transported into its full fine-edge vertex banks, \cref{eq:nested-tree-cov}. \\
$b^{\rm fix},b_j^{\rm harm}$ & Fixed fine router background and the harmonic extension of the level-$(j+1)$ coarse background. \\
$\delta_j^{\rm bg}$ & Difference of the consistently normalized harmonic- and fixed-background principal shell coefficients; ADMC requires its Ces\`aro mean to vanish when the fixed route is used. \\
$\widetilde K,B(Y,C),j$ & Smooth compact root-path block, its actual solved pivots, and the product inverse-Haar amplitude; \cref{cor:compact-product}. \\
$G(c),H_c,\widehat Q$ & Minimized local action, eliminated fibre Hessian, and normalized measured one-loop coefficient; the retained Schur form need not have a gap. \\
$d,d_{\rm ch}$ & Hilbert--Schmidt geodesic and chordal distances in the explicit $\SU(2)$ fibre; the cutoff uses $d$, while $d_{\rm ch}\le d$ gives sufficient inactive margins. \\
$\theta,\mathcal I_0$ & Actual one-step transverse $p^2$ coefficient and flat constant response integrand, \cref{thm:theta-integral,thm:flat-bloch}; neither is a physical mass. \\
$L,\mathsf H_\mu,\kappa_D$ & Auxiliary product diffusion, rate-one coordinate heat bath, and killed floor on a visible event; distinct from the physical Hamiltonian $H$. \\
$T,F,M_{\rm vis}$ & Fixed 819-link repair shell, its 10034-group visible collar, and the actual noncentral $3\times3$ curvature certificate. \\
$\mathcal L(R),c_\eta$ & Active-source curvature cost including pivot acceleration, and the measured local diffusion coefficient $c_\eta=\eta/26208$. \\
$\gamma_\ell,\tau_a$ & All-exterior conditional patch gap and sufficient weighted patch threshold, \cref{crit:weighted-patches}; local killed floors do not supply $\gamma_\ell$. \\

$\mathcal B_1,\mathcal R_\beta^{(2)}$ & Complete local measured two-loop coefficient and four-mark $O(\beta^{-2})$ remainder, \cref{thm:measured-two-loop}. \\
$J^8S,J^6Q,$\newline $J^4B;J^7F$ & Finite field-degree packet for measured two-loop composition, \cref{thm:closed-loop-jet-packet}; these are not momentum orders. \\
$\mathcal F_{L,N},\mathcal O_{\infty,N}$ & Actual balanced root composite and its weakly continuous continuum source on the small critical ball. \\
$\mathcal G_L,\mathcal G_\infty$ & One-cell local classical minima in \cref{thm:classical-minimum-limit}, distinct from measured or full compact free energies. \\

$M_L,K_L,G_L^{\rm gh}$ & Normal source differential, relative constrained precision, and based-gauge Jacobian of the one-cell branch; their domains differ. \\
$E_4,\mathsf C_*$ & Homogeneous noncommuting classical quartic and the complete measured source--precision matrix, \cref{sec:measured-quartic-log}. \\
$\beta_M^\diamond,T_x,k_M,R_M$ & Selected bare-Wilson coupling, shared-face cell energy, pressure constant, and source radius; not a prescribed physical trajectory. \\
$\mathcal I_\gamma,\mathcal A_\gamma$ & Exactly normalized magnetic isometry and conditional affinity for the native bridge compression. \\
$\mathcal N,P_{\le k},\Pi_{\gamma,k}$ & Original oscillator degree, Euclidean degree band, and genuine compact Gram-normalized projection; not physical spectral projections. \\

$\Lambda_\gamma,\lambda_*,h,g_*$ & Actual native Perron value, hard Gaussian limiting top, critical scale $(\gamma L^3)^{-1/3}$, and invariant quartic model gap. \\
$\widehat{\mathcal V}_{L,M},\mathcal H^{\rm loc}_{L,M}$ & Fourth-source finite part and its finite incoming second-chaos counterterm; distinct from the measured scalar quartic coefficient. \\
$A_L,K_L,\mathsf g_s$ & Grounded based Laplacian, its inverse, and the actual quotient metric in \cref{sec:native-resistance-geometry}; $K_L$ here is not the one-cell precision matrix. \\
$F_\gamma,T^{\rm B},T^{\rm W}$ & Bernstein dispersion, full magnetic subordinated reference, and original Wilson transfer; each normalized operator uses its own Perron value. \\
$\Theta_{\gamma,L}$ & Joint critical error budget of \cref{eq:native-joint-budget}; not a physical running-coupling condition. \\
$\Pi_j^\circ,\epsilon_j$ & Reducing soft-band projection with its vacuum removed, and scalar-compression error for actual normalized local sources. \\
$\delta_M,b_M,\mu_{f,s}$ & Actual-vacuum covariance error, mesoscopic smearing scale, and limiting damped probe spectral measure; distinct from a constructed physical continuum law. \\
\bottomrule
\end{longtable}

\section{Which conclusions require which inputs}
\begin{longtable}{L{0.27\textwidth}L{0.34\textwidth}L{0.29\textwidth}}
\toprule Desired conclusion & Required interface & What does not replace it \\ \midrule
\endfirsthead
\toprule Desired conclusion & Required interface & What does not replace it \\ \midrule
\endhead
Physical Poincar\'e or Witten gap & Actual ground law, full signed sources, physical form coefficient, controlled response/capacity & An unweighted reduced fibre or an arbitrary spatial sampler \\
Strong all-order source analyticity & Exponential fixed-tree grade and a valid complete spatial compiler & A factorial-graded inequality at every arity \\
Compact forest convergence & One analytic radius budget, degree-faithful sum, $8zK_{\rm HK}<1$ & A degree-free estimate that discards a star factorial \\
Measured redundant quotient & Exact pushed current, its divergence, and the same pushed measure & Strict verticality or action transport without its Jacobian \\
Local pushed Jacobian Hessian & Three current jets, four-entry connected bounds, source radii, and corrected contour activity & Pair decay, a determinant expansion alone, or two current jets \\
Core-based physical coercivity & Conditional core gap, restricted smoothing, tensorization, and same-form incidence & A bare core floor or a small unconditional bad-set probability \\
Gaussian ADMC & Exact nested determinant, representative-matched shell response, physical scalar, and sublinear endpoints & A corner mean or a frozen-chart polynomial coefficient \\
Use of fixed-background banks & Complete normalization and a vanishing averaged defect relative to the harmonic prescription & Gauge equivalence of fluctuations or equality at finitely many corners \\
Stable ultraviolet trajectory & Calibrated marginal recursion and invariant stable cone & Strict contraction of the full retained action \\
Interacting continuum theory & Local correlation compactness, regularity, reflection, physical normalization, nontriviality & Convergence of Gaussian kernels alone \\
Two-level temporal gap & Physical form domination plus both conditional and retained deficits at the correct time scale & Complete detail refresh with no retained infrared estimate \\
Physical mass gap & Complete local reflected contraction at a fixed time, or a populated equivalent route & A finite eliminated Hessian floor or positivity on a selected bank \\
Volume-compatible sector control & A specified local, density, fixed-volume, or phase-resolved target & A uniform tail for an extensive total charge without entropy control \\
Smooth compact fibre & Actual globally solved pivots, inverse-Haar factors, and a declared completion & Independent pivot caps or equality of local Taylor germs \\
Local measured one-loop limit & Moving minimum, full density jets, fixed-domain wall comparison, and localized derivative norms & A frozen-background determinant or a bound on the compact complement \\
Local visible repair & Uniform all-hidden curvature, original-clock resolvent, and a visible-only support condition & A force bound restricted to a hidden posterior event \\
Unrestricted auxiliary interface gap & Complete same-law assembly or all-exterior weighted patch gaps above their threshold & One good tube, rare bad fields, or an uncontrolled spatial union \\

Finite measured composition & Complete incoming density and the $(8,6,4;7)$ field-degree packet & Quartic truncation of every input or a reset to product Haar \\
Local classical continuum limit & Small critical domain, weak source continuity, exact-source recovery, and constrained convexity & A homogeneous trial or source convergence without energy control \\

Measured two-loop remainder & Higher local jets, rooted connected bounds, parity, and normalized four-mark wall return & Hessian positivity alone or a full-compact assertion \\
Cofinal measured logarithm & Complete source, constrained, gauge and Haar factors, low-Born matching, retained endpoints & Schatten-tail convergence alone \\
Entropy source comparison & A declared full law, source window and actual relative-entropy budget & An assumed small budget for general blocking \\
Generated source law & Normalized conditional likelihood, output entropy, and a separate retained-information test & A conditional estimate with its output normalizer omitted \\
Compact magnetic degree band & Genuine Gram normalization and a controlled order of volume, degree, and temperature limits & Raw restriction of a Euclidean projection or subextensive old degree \\

Fixed-regulator native spectrum & Physical recovery, same-top return metric, corrected profile equation and radial tightness & A frozen Gaussian spectrum or a regional compression alone \\
Joint critical native spectrum & Uniform measured geometry, actual-Perron comparison, recovery and no-escape with $\Theta_{\gamma,L}\to0$ & Fixed-volume remainder notation or a physical-trajectory extrapolation \\
Beyond-domain native comparison & Complement deficit, boundary return, uniform retained pencil and clock & Per-link exponent control or an interior jump estimate alone \\
Fourth-source finite part & Full matrix filter, finite incoming quadratic-chaos subtraction, and summable local shell errors & The measured scalar quartic coefficient alone \\
Soft-band local Gram excision & Reducing projection, scalar compression of normalized sources, and physical complement contraction & Low band energy or invisibility of unrenormalized loops alone \\
Surviving local soft-spectrum test & Actual-vacuum covariance, full smearing error, and nonzero positive-time norm & A shrinking lattice energy window whose probe weight may vanish \\
\bottomrule
\end{longtable}

\chapter{Corrections retained in the mathematical chain}
\label{app:corrections}
The following restrictions are part of the mathematical hypotheses and conclusions. Each records a distinction needed by the corresponding implication.

\begin{longtable}{L{0.24\textwidth}L{0.49\textwidth}L{0.17\textwidth}}
\toprule Earlier shortcut & Retained replacement & Main location \\ \midrule
\endfirsthead
\toprule Earlier shortcut & Retained replacement & Main location \\ \midrule
\endhead
Raw quartic numerator as normalized MTA & Correct first-connectivity and product-state coefficient-capacity argument, including globally summed rescued suffixes & \cref{ch:trees} \\
All-order factorial-graded MTA as strong MTA & Keep the valid $m!K^m$ bound, withdraw its uniform analytic closure inference; restore exact-Haar assembly & \cref{prop:factorial,thm:Haar} \\
Forest fixed-tree bound without degree factorials & Keep $\prod d_i!$ and use the exact Pr\"ufer sum & \cref{lem:degree-census,thm:forest} \\
A matching Wilson exponent as the Perron law & Prove the incidence equation or repair the actual-law density defect & \cref{prop:law-no-go} \\
Rare bad fields as a uniform conditional estimate & Establish complete-source incidence, restricted smoothing, and capacity & \cref{ch:source,ch:geometry} \\
Finite action-closed update block & Retain exterior plaquettes and control boundary response & \cref{prop:no-finite-block} \\
Strict verticality as the definition of redundancy & Push the measured tangent through its conditional current & \cref{thm:measured-redundancy} \\
Two current jets for the Haar-Jacobian Hessian & Retain the third jet and four-entry normalized cumulants & \cref{prop:Jacobian-third-jet,thm:four-mark-clustering} \\
Degree-eight entropy for filling contours & Use the actual degree-eighty contour graph and the $6401$ threshold & \cref{lem:polymer-entropy-gate} \\
Pointwise centered remainder as a local analytic norm & Supply source-uniform polymer control; the fibre range may be extensive & \cref{thm:exact-block-differential} \\

Conditional fast mixing as full mixing & Retain the slow covariance and the exact two-level temporal form & \cref{eq:totalcov,prop:two-level-transfer} \\
Unscaled small transfer error as physical gap transport & Compare the error to the physical step deficit & \cref{lem:transfer-comparison-debit} \\
Fine floor as retained infrared coercivity & Control the actual Schur head; a soft head has soft minimizing lifts & Prop.~\ref{prop:massless-Schur-edge} \\
Finite grid signs as global certificates & Supply an analytic quadrature remainder or an actual all-torus proof & \cref{eq:quadrature,ch:oracle} \\
Direct transverse covariance in fixed tree-gauge banks & Push the covariance through the explicit representative map before contraction & \cref{sec:tree-gauge} \\
Fixed background as the harmonic blocked action & Retain the distinct prescriptions and prove the averaged background comparison where required & \cref{crit:background-ADMC} \\
Corner agreement as a tower coefficient & Separate finite algebraic checks, torus integration, and the uniform limiting response & \cref{sec:fixed-background-cards} \\
Frozen endpoint mean as the anomaly & Use an actual nested tower; a nonzero fixed reset mean obstructs that comparison & \cref{ch:anomaly,ch:tower} \\
Covariance nondegeneracy as a physical mass & Retain the massless axis behavior and distinguish inverse-kernel control & \cref{thm:tower-floor,eq:axis} \\
Finite one-step contraction as a continuum gap & Calibrate the physical rate or prove a fixed-delay Gram inequality & \cref{prop:step-collapse} \\
Bare head constant as the vacuum & Short in the exact Feshbach graph metric & \cref{thm:Feshbach} \\
Diagonal probe decay as a full gap & Control every finite mixed local Gram with one constant & \cref{thm:local-Gram} \\
Global sector tail from local action cost & Keep the tensorization no-go and use an explicitly different target & \cref{thm:sector-no-go} \\
Over-weak wrapping trajectory as the physical limit & Impose nonzero physical normalization or declare a different local branch & \cref{ch:closure} \\
Laurent representative at the toron & Use the rational punctured-torus map and its integrable weak derivative bounds & \cref{prop:toron-regularity} \\
Linear nonabelian averaging as a homomorphism & Retain endpoint path dressing and the actual constrained Hessian & \cref{prop:nonabelian-average,eq:actual-constrained-Hessian} \\
Identical local germs as identical compact laws & Distinguish completion choices and compare normalized pushforwards & \cref{scope:completion-germ} \\
Dropping pivot acceleration or Haar terms & Use the complete measured jets, including active-source response & \cref{eq:complete-fourth-jet,thm:active-repair} \\
Nonzero finite flat mean as a bulk constant & Identify the exact Fourier aliases after Haar cancellation & \cref{thm:flat-bloch} \\
Small higher-well mass as fast hidden mixing & Keep its Rayleigh quotient and distinguish marginal removal & \cref{thm:hidden-metastability,prop:well-removal} \\
Periodic shell transplanted into larger volume & Count the true outer plaquettes and embed the complete input closure & \cref{thm:nonwrapping-repair} \\
Fixed-patch cover as a growing spatial union & Require a proved global assembly or all-exterior patch criterion & \cref{sec:weighted-patch-criterion} \\
Scalar measured trace as the source filter & Retain the cross-oriented matrix and the finite incoming source counterterm & \cref{sec:fourth-source-renormalization} \\
Grounded inverse norm charged to every local gauge variation & Use bounded resistance evaluation and trace-ideal relative geometry & \cref{sec:native-resistance-geometry} \\
Reference top substituted for the actual Perron value & Normalize each transfer at its own top and keep the vacuum-energy debit & \cref{thm:native-Perron-comparison} \\
Gaussian component maximum as a link radius & Include the Euclidean color factor and any soft displacement; keep the reference law explicit & \cref{sec:native-Gaussian-maxima} \\
Cubic Gaussian-field polynomial treated as Gaussian & Retain third-chaos Wick structure; no lognormal or norm obstruction follows from variance alone & \cref{sec:native-cubic-chaos} \\
Joint critical limit as fixed physical volume & Keep the error domain and independent bare/physical scale comparison & Prop.~\ref{prop:native-physical-domain-separation} \\
Soft toron band deleted without a source test & Verify scalar compression after normalization, then use exact Gram excision & \cref{thm:local-toron-Gram-excision} \\
Lattice-gap collapse with vanishing probe weight as a local spectral limit & Prove nonzero damped norm and convergence of actual spectral measures & \cref{thm:smeared-probe-spectrum} \\
Electric-dominated lattice gap identified with the continuum mass & Keep the explicit $b/\kappa$ domain and perform physical scale transport & \cref{thm:native-electric-gap,thm:native-plaquette-band} \\
Whole-head Schrieffer--Wolff elimination across an electric cutoff & Retain the exact equal-energy resonant block and eliminate only the separated complement & \cref{eq:native-resonant-matrix-element,thm:native-high-Schur} \\
Small-time native cost assumed monotone in every coarse background & Retain the holonomy-dependent quadratic coefficient and control finite-time transport directly & \cref{eq:native-cost-C2,eq:native-cost-signs} \\
\bottomrule
\end{longtable}

\chapter{Result concordance and computational scope}
\label{app:concordance}
\section{Mathematical dependencies}
The bibliographic keys F1--F13 identify the successive Yang--Mills
series. F11 contains the nested Gaussian construction, its representative
and background analysis, and the nonlinear local fibre and measured
coefficient. F12 develops the compact completion, local one-loop
comparison, and conditional spectral machinery. F13 develops the
actual orbit geometry, metastability, and fixed-patch visible repair. PAR contains the independent owner-remainder
and sector arguments. The Grand Tensor reference supplies the inherited
finite construction and physical-transfer results. F14 supplies the
measured quartic and source finite parts, repair extension, and
full-law pressure bounds; F15 includes entropy transport and the
actual-vacuum smeared spectral test. F16 and F16B are parallel
continuations of series 16 after their common V48 material. F16 is
the V67 branch containing the model certificates and local kinetic
and Perron developments; F16B is the V56 branch containing uniform
based geometry, subordination, the joint critical limit, physical
calibration, and local toron excision. Their version numbers are not
interchanged. F17 is the later consolidated continuation through V100;
its terminal native results provide the high-harmonic Schur bounds,
spectral-window tails, resonance obstruction, and local resonant normal
form used here. F18 is the subsequent continuation through V59; only
its results with independent analytic closure relevant to the present
architecture are imported, notably the electric-dominated actual gap,
the localized plaquette band, the corrected conditional-fibre estimate,
and the completed finite-time cost. These are unpublished programme
manuscripts, not external literature used to discharge an unproved
input.

The following table records the mathematical result locations. Original
labels, version locations, and, where retained, one-based manuscript
line numbers identify the statements;
the exposition above standardizes notation and displays the implications
used in the conditional chain. Reported computational results remain
distinct from the algebraic consequences proved here.

\section{Principal result map}
\small
\begin{longtable}{L{0.32\textwidth}L{0.12\textwidth}L{0.477\textwidth}}
\toprule Result in this monograph & Source & Original label and line \\ \midrule
\endfirsthead
\toprule Result in this monograph & Source & Original label and line \\ \midrule
\endhead
Half-slab and physical Doob channel (\cref{thm:half}) & \cite{GT} & \nolinkurl{thm:YM-half-slab-Doob}\newline Line 62,008 \\
Actual Perron law incidence (\cref{prop:law-no-go}) & \cite{F1} & \nolinkurl{thm:V21-incidence-equation}\newline Line 13,157 \\
Sequential physical response matrix (\cref{thm:response-matrix}) & \cite{F1} & \nolinkurl{thm:V14-response-matrix}\newline Line 8,721 \\
Physical landing metric (\cref{thm:landing}) & \cite{F1} & \nolinkurl{thm:V17-landing-ledger}\newline Line 10,363 \\
Complete-source capacity factorization (\cref{eq:bad-capacity}) & \cite{F1} & \nolinkurl{thm:V20-incidence-capacity}\newline Line 12,803 \\
Corrected quartic weak-cut argument (\cref{ch:trees}) & \cite{F2} & \nolinkurl{thm:V1-quartic-weak-cut}\newline Line 292 \\
Completed normalized quartic theorem (\cref{thm:quartic}) & \cite{F3} & \nolinkurl{thm:V96-global-quartic-MTA}\newline Line 46,293 \\
Projected-chain degree filtration (\cref{ch:trees}) & \cite{F4} & \nolinkurl{lem:V100-filtered-chain}\newline Line 44,273 \\
Withdrawal of the overstrong all-order grade (\cref{prop:factorial}) & \cite{F5} & \nolinkurl{thm:newV12-correction}\newline Line 5,072 \\
Exact assembled Haar derivatives (\cref{thm:Haar}) & \cite{F5} & \nolinkurl{thm:newV12-exact-Haar-card}\newline Line 5,248 \\
Strong forest projection, not source-occurrence closure (\cref{ch:trees}) & \cite{F6} & \nolinkurl{thm:s6v7-VAFPC}\newline Line 3,135 \\
Connected-carrier activity criterion (\cref{eq:carrier-row}) & \cite{F6} & \nolinkurl{thm:s6v16-carrier}\newline Line 6,944 \\
Exact half-slab conditional-variance form (\cref{prop:half-variance}) & \cite{F6} & \nolinkurl{thm:s6v58-conditional-variance}\newline Line 25,626 \\
Same-carrier half-slab density comparison (\cref{lem:half-density-comparison}) & \cite{F6} & \nolinkurl{lem:s6v58-density-comparison}\newline Line 25,896 \\
Complete normalized score hierarchy (\cref{thm:connected-score-hierarchy}) & \cite{F8} & \nolinkurl{thm:e8v60-all-score}\newline Line 25,780 \\
Conditional core-mean gluing (\cref{lem:core-mean-gluing}) & \cite{F6} & \nolinkurl{lem:s6v97-mean}\newline Line 45,631 \\
Core gap plus restricted smoothing (\cref{crit:core-smoothing}) & \cite{F6} & \nolinkurl{thm:s6v97-bootstrap}\newline Line 45,675 \\
Exact retained-pullback sector (\cref{thm:exact-block-differential}) & \cite{F7} & \nolinkurl{thm:s7v99-pullback}\newline Line 44,998 \\
Centered exact nonlinear remainder (\cref{eq:exact-centered-remainder}) & \cite{F7} & \nolinkurl{thm:s7v99-exact-split}\newline Line 45,085 \\
Measured redundancy through pushforward (\cref{thm:measured-redundancy}) & \cite{F8} & \nolinkurl{thm:e8v58-pushforward-redundancy}\newline Line 25,050 \\
Gaussian measured EOM has no marginal feed (\cref{cor:Gaussian-EOM-no-feed}) & \cite{F8} & \nolinkurl{cor:e8v58-no-Gaussian-feed}\newline Line 25,187 \\
Third current jet in the Haar-Jacobian Hessian (\cref{prop:Jacobian-third-jet}) & \cite{F8} & \nolinkurl{prop:e8v60-third-jet-debit}\newline Line 25,736 \\
Volume-uniform four-mark clustering (\cref{thm:four-mark-clustering}) & \cite{F8} & \nolinkurl{thm:e8v61-four-mark-cluster}\newline Line 26,291 \\
Corrected nonlinear-current activity criterion (\cref{crit:nonlinear-current-gate}) & \cite{F8} & \nolinkurl{thm:e8v63-revised-gate}\newline Line 27,112 \\
Exact two-level comparison transfer (\cref{prop:two-level-transfer}) & \cite{F6} & \nolinkurl{thm:s6v56-two-level}\newline Line 24,778 \\
Two-level physical temporal coercivity (\cref{crit:two-level-physical}) & \cite{F6} & \nolinkurl{thm:s6v56-compiler}\newline Line 24,837 \\
Physical-scale transfer comparison (\cref{lem:transfer-comparison-debit}) & \cite{F6} & \nolinkurl{lem:s6v56-additive-debit}\newline Line 24,898 \\
Massless retained Schur obstruction (\cref{prop:massless-Schur-edge}) & \cite{F6} & \nolinkurl{thm:s6v57-massless-Schur}\newline Line 25,301 \\
Corrected filling-contour entropy gate (\cref{lem:polymer-entropy-gate}) & \cite{F8} & \nolinkurl{prop:e8v63-contour-entropy}\newline Line 27,057 \\
Single analytic radius budget (\cref{ch:forests}) & \cite{F9} & \nolinkurl{lem:n1v1-one-radius}\newline Line 478 \\
Correct regulator-dependent forest power (\cref{eq:forest-beta}) & \cite{F10} & \nolinkurl{thm:newv100-degree-scale}\newline Line 21,200 \\
Averaged determinant matching definition (\cref{def:ADMC}) & \cite{F9} & \nolinkurl{def:n1v14-ADMC}\newline Line 6,600 \\
ADMC marginal implication (\cref{ch:anomaly}) & \cite{F9} & \nolinkurl{thm:n1v14-ADMC}\newline Line 6,617 \\
Sublinear compact--continuum intertwiner (\cref{def:SCIC}) & \cite{F9} & \nolinkurl{def:n1v15-SCIC}\newline Line 7,097 \\
One-dimensional analytic alias cohomology (\cref{thm:alias}) & \cite{F9} & \nolinkurl{thm:n1v17-cohomology}\newline Line 7,755 \\
Exact remaining mean criterion (\cref{eq:edge-scalar}) & \cite{F9} & \nolinkurl{thm:n1v17-edge-SCIC}\newline Line 7,854 \\
Reported all-torus relative margin (\cref{eq:relative-floor}) & \cite{F10} & \nolinkurl{thm:newv97-full-torus}\newline Line 20,399 \\
Normalized covariance-derivative consequence (\cref{thm:covariance-margin}) & \cite{F10} & \nolinkurl{thm:newv100-normalized-margin}\newline Line 21,318 \\
Fixed-point Gaussian matching criterion (\cref{thm:fixedpoint-anomaly}) & \cite{F11} & \nolinkurl{thm:v34-card}\newline Line 12,168 \\
Explicit finite-level block-average covariance (\cref{thm:tower-explicit}) & \cite{F11} & \nolinkurl{thm:v37-explicit}\newline Line 13,356 \\
Positive lattice-sum fixed point (\cref{thm:tower-fixed}) & \cite{F11} & \nolinkurl{thm:v37-fixed-point}\newline Line 13,378 \\
Global transverse floor (\cref{thm:tower-floor}) & \cite{F11} & \nolinkurl{thm:v37-nondegenerate}\newline Line 13,419 \\
Composed harmonic extension (\cref{eq:tower-minimizer}) & \cite{F11} & \nolinkurl{thm:v38-harmonic}\newline Line 13,634 \\
Positive conditional covariance shells (\cref{thm:shells}) & \cite{F11} & \nolinkurl{thm:v38-shells}\newline Line 13,684 \\
Harmonic tree-level shell functional and fixed prescription (\cref{sec:background-prescription}) & \cite{F11} & \nolinkurl{def:v38-shell-functional}\newline Line 13,742; \nolinkurl{dec:v41-background}\newline Line 14,334 \\
Fluctuation-representative dependence (\cref{prop:representative-dependence}) & \cite{F11} & \nolinkurl{prop:v41-gauge-dependence}\newline Line 14,271 \\
Tree-gauge map and finite route agreement (\cref{sec:tree-gauge}) & \cite{F11} & \nolinkurl{def:v41-map}\newline Line 14,296; \nolinkurl{thm:v41-map}\newline Line 14,309 \\
Fixed-background level-zero cards (\cref{eq:corners-fixed}) & \cite{F11} & \nolinkurl{thm:v41-cards}\newline Line 14,357 \\
Averaged background compatibility (\cref{crit:background-ADMC}) & \cite{F9,F11} & SCIC specialization of \cref{def:SCIC}; the unproved background comparison is identified in \nolinkurl{rem:v41-why}. \\
Non-double-counted response identity (\cref{eq:owner-ledger}) & \cite{PAR} & \nolinkurl{thm:exact-ledger}\newline Line 179 \\
Physical principal-margin protection (\cref{thm:remainder}) & \cite{PAR} & \nolinkurl{thm:margin-protection}\newline Line 670 \\
Global sector tensorization obstruction (\cref{thm:sector-no-go}) & \cite{PAR} & \nolinkurl{thm:global-sector-no-go}\newline Line 1,167 \\
Label refinement data processing (\cref{eq:label-TV}) & \cite{PAR} & \nolinkurl{prop:refinement-defect}\newline Line 1,262 \\
Complete local OS Gram criterion (\cref{thm:local-Gram}) & \cite{F1} & \nolinkurl{thm:local-Gram-criterion}\newline Line 231 \\
Passage through changing null spaces (\cref{eq:Gram-error}) & \cite{F1} & \nolinkurl{thm:quotient-stable-passage}\newline Line 382 \\
Physical-time collapse obstruction (\cref{prop:step-collapse}) & \cite{F1} & \nolinkurl{thm:fixed-step-collapse}\newline Line 558 \\
Orbit-tube coercivity (\cref{thm:orbit}) & \cite{F1} & \nolinkurl{thm:V2-orbit-coercivity}\newline Line 988 \\
Exact vacuum metric and Feshbach gap (\cref{thm:Feshbach}) & \cite{F1} & \nolinkurl{thm:V4-vacuum-corrected-gap}\newline Line 2,344 \\
Common-time signed return packet (\cref{eq:transport-threshold}) & \cite{GT} & \nolinkurl{thm:YM-signed-threshold-packet}\newline Line 66,113 \\
Head--tail transfer compiler (\cref{eq:2x2-transfer}) & \cite{GT} & \nolinkurl{thm:YM-regulated-criterion}\newline Line 66,078 \\
Complete effective-action terminal ball (\cref{thm:Dobrushin}) & \cite{F7} & \nolinkurl{thm:s7v83-ball}\newline Line 36,610 \\
Physical mass from terminal clustering (\cref{thm:Dobrushin}) & \cite{F7} & \nolinkurl{cor:s7v83-gap}\newline Line 36,651 \\
Explicit volume-uniform terminal region (\cref{eq:strong-electric}) & \cite{GT} & \nolinkurl{thm:YM-strong-electric-gap}\newline Line 66,036 \\
Local-vacuum conditional closure (\cref{thm:master}) & \cite{F1} & \nolinkurl{thm:local-vacuum-closure}\newline Line 726 \\
Toron regularity (\cref{prop:toron-regularity}) & \cite{F11} & \nolinkurl{prop:v42-not-laurent}; \nolinkurl{lem:v42-derivatives}\newline V42 \\
Nonabelian average obstruction (\cref{prop:nonabelian-average}) & \cite{F11} & V56 \\
Global compact pivot diffeomorphism (\cref{thm:global-pivot}) & \cite{F12} & \nolinkurl{thm:f12v18-pivot}\newline V18 \\
Global product carrier and smooth density (\cref{cor:compact-product}) & \cite{F12} & \nolinkurl{thm:f12v18-global}; \nolinkurl{cor:f12v18-density}\newline V18 \\
Actual pivot residual and strong convexity (\cref{prop:pivot-potential}) & \cite{F12} & \nolinkurl{thm:f12v98-potential}\newline V98 \\
Uniform local one-loop remainder (\cref{thm:uniform-local-loop}) & \cite{F12} & \nolinkurl{thm:f12v12-local}\newline V12 \\
Differentiated transverse coefficient (\cref{thm:theta-integral}) & \cite{F12} & \nolinkurl{thm:f12v5-integral}; \nolinkurl{thm:f12v5-quadrature}\newline V5 \\
Flat identity and Fourier aliases (\cref{thm:flat-bloch}) & \cite{F12} & \nolinkurl{thm:f12v9-flat}; \nolinkurl{prop:f12v9-alias}\newline V9 \\
Actual local orbit-tube inequality (\cref{thm:actual-tube-gap}) & \cite{F13} & \nolinkurl{thm:tube}\newline V1 \\
Second well and hidden metastability (\cref{thm:hidden-metastability}) & \cite{F13} & \nolinkurl{thm:f13v8-minimum}; \nolinkurl{thm:f13v8-metastability}\newline V8 \\
Localized effective-action comparison (\cref{prop:well-removal}) & \cite{F13} & \nolinkurl{thm:f13v9-potential}\newline V9 \\
Same-law diffusion and heat-bath comparison (\cref{thm:global-resolvent}) & \cite{F12} & \nolinkurl{thm:f12v60-convert}\newline V60; resolvent development V77--V80 \\
Non-wrapping original-clock repair (\cref{thm:nonwrapping-repair}) & \cite{F13} & \nolinkurl{thm:f13v14-curvature}; \nolinkurl{thm:f13v14-floor}\newline V14 \\
Locally curved-background transport (\cref{sec:coarse-gauge-repair}) & \cite{F13} & \nolinkurl{thm:f13v15-floor}\newline V15 \\
Noncentral visible matrix criterion (\cref{thm:matrix-repair}) & \cite{F13} & \nolinkurl{thm:f13v16-repair}\newline V16 \\
Active-source pivot and repair (\cref{thm:active-repair}) & \cite{F13} & \nolinkurl{thm:f13v17-curvature}; \nolinkurl{thm:f13v17-repair}\newline V17 \\
Weighted conditional patch criterion (\cref{crit:weighted-patches}) & \cite{F12} & \nolinkurl{thm:f12v63-weighted}\newline V63 \\
Explicit finite-size threshold (\cref{cor:weighted-threshold}) & \cite{F12} & \nolinkurl{thm:f12v63-threshold}\newline V63 \\

Mixed-shell response and Schur ownership (\cref{prop:mixed-response,prop:Schur-owned-response}) & \cite{SM} & V1--V4; exact quadratic determinant allocation \\
Measured two-loop $C^4$ remainder (\cref{thm:measured-two-loop}) & \cite{SM} & V7--V8; centered expansion and wall return \\
Finite measured composition and closed field packet (\cref{thm:measured-loop-composition,thm:closed-loop-jet-packet}) & \cite{SM} & V8; measured step, associative composition, and necessary field degrees \\
Uniform nonlinear root domain (\cref{thm:classical-root-domain}) & \cite{SM} & V10--V12; balanced orbit and counting-norm induction \\
Weak root source and classical continuum minimum (\cref{thm:classical-weak-source,thm:classical-minimum-limit}) & \cite{SM} & V13; finite-resolution source, full Wilson lower limit, and exact-source recovery \\

Complete one-cell determinant and trace tail (\cref{prop:cofinal-measured-determinant,thm:cofinal-trace-tail}) & \cite{SM} & V15--V16; normal, constrained, gauge and Haar factors \\
Exact flat cofinal coefficient (\cref{thm:cofinal-flat}) & \cite{SM} & V16; complete flat formula and $L^{-2}$ rate \\
Measured quartic logarithm and finite part (\cref{thm:measured-quartic-log}) & \cite{F14} & V19--V22, V28--V30; full measured bank reduction \\
Full compact source window (\cref{thm:full-Wilson-source-window}) & \cite{F14} & V95--V100; unrestricted Wilson pressure and selected coupling \\
Entropy source transport (\cref{thm:entropy-source-transport}) & \cite{F15} & V1--V2; finite information budget and plaquette tilts \\
Generated likelihood, retained information, and linearization (\cref{thm:generated-likelihood,prop:retained-source-information,prop:generated-source-linearization}) & \cite{F15} & V2; normalized pushforward and information split; V4, V7 for retained-score tests \\

Inverse-paired cutoff-transition repair (\cref{cor:paired-transition-repair}) & \cite{F14} & V5--V11; reduced collar and coplanar source cone \\
Native magnetic assembly and affinity (\cref{thm:native-affinity}) & \cite{F15,F16} & F15 V98--V100; F16 V1 \\
Magnetic occupation and compact projection (\cref{thm:magnetic-excitation-law,prop:magnetic-compact-projection}) & \cite{F15} & V99--V100; exact degree distribution and fixed-volume compact band limit \\

Ordered fixed-regulator native spectrum (\cref{thm:native-critical-spectrum}) & \cite{F16} & V41--V48; recovery, profile equation and radial tightness \\
Model gap bracket (\cref{cert:quartic-model-gap}) & \cite{F16} & V51, V57, V60; sector-resolved certificates \\
Sup-norm kinetic and Haar control (\cref{prop:native-supnorm-kernel}) & \cite{F16} & V64; exact per-link exponent bounds \\
Same-top regional reduction (\cref{thm:native-regional-Feshbach}) & \cite{F16} & V64; complement resolvent and Doob Schur test \\
Fourth-source subtraction and cofinal finite part (\cref{prop:fourth-source-subtraction,thm:fourth-source-finite-part}) & \cite{F14} & V31--V32; full matrix source filter, second-chaos logarithm, locally summed finite part \\
Resistance-normalized based geometry (\cref{thm:native-uniform-geometry}) & \cite{F16B} & V50; uniform tube, relative Haar density and quotient metric \\
All-label actual-Perron comparison (\cref{thm:native-Perron-comparison}) & \cite{F16B} & V51--V53; positive subordination and complete magnetic envelope \\
Ordered joint-regulator critical spectrum (\cref{thm:native-joint-spectrum}) & \cite{F16B} & V54; measured recovery and all-vector no-escape on the explicit joint domain \\
Mean-one kinetic factor and cubic-chaos accounting (\cref{prop:native-mean-one,prop:native-cubic-Wick}) & \cite{F16} & V66--V67; local coefficients and Wick structure, without the Gaussian/lognormal inference \\
Exact native Perron transport (\cref{prop:native-Perron-transport}) & \cite{F16} & V65; multiplicative path identity and Poisson equation \\
Bare scale and toron physical energy (\cref{prop:native-physical-domain-separation}) & \cite{F16B} & V55; separate clock and fixed-box scale premises \\
Local toron loop and Gram excision (\cref{prop:local-toron-loops,thm:local-toron-Gram-excision}) & \cite{F16B} & V56; physically bounded loops, actual source compression and quotient-stable contraction \\
Smeared actual-vacuum soft spectrum (\cref{thm:smeared-probe-covariance,thm:smeared-probe-spectrum}) & \cite{F15} & V64--V65; uniform covariance, dimension-four smearing and surviving damped spectral weight \\
Electric-dominated actual lattice gap (\cref{thm:native-electric-gap}) & \cite{F18} & V5; nonlinear creation fixed point and excitation exclusion, uniform in spatial volume \\
Localized analytic plaquette band (\cref{thm:native-plaquette-band}) & \cite{F18} & V12; complete first physical band, exponential spatial columns and cubic uniform remainder \\
Native high block and spectral-window tails (\cref{thm:native-high-Schur}) & \cite{F17} & V99; exact one-link source shell, Schur recursion and factorial low-energy tails \\
Resonance obstruction and local normal form (\cref{thm:native-resonant-normal-form}) & \cite{F17} & V100; exact equal-energy matrix element and summable connected remainder \\
Uniform corrected conditional fibre and completed small-time cost (\cref{eq:native-conditional-gap-late,eq:native-cost-C2}) & \cite{F18} & V48--V49, V59; actual conditional gap/source inverse and holonomy-dependent quadratic cost \\
\bottomrule
\end{longtable}
\normalsize
\section{Reproducibility and scope}
The normalized coefficient and connected-source arguments depend on
\cite{F1,F2,F3,F4,F5,F6}; the conditional geometry and exact compact
blocking on \cite{F6,F7,F8,F9}; and the finite response oracle on
\cite{F9,F10}. The explicit tower, full-edge vertices, and representative
and background analysis, together with the actual local nonlinear
fibre, are in \cite{F11}. The smooth compact completion, local one-loop
error, global resolvent, and weighted patch criterion are in \cite{F12};
the two-orbit geometry, normalized well comparison, and visible repair
families are in \cite{F13}. The mixed-shell algebra, local two-loop
expansion, finite measured composition, critical root-source domain,
full local classical continuum minimum, and cofinal measured
determinant are developed in \cite{SM}. The noncommuting quartic
evaluation, inverse-paired repair, and full compact source window are
in \cite{F14}; finite-entropy transport, generated likelihoods and
retained information, the actual-vacuum spectral probe, native spatial
assembly, and the magnetic excitation budget in \cite{F15}; native affinity, fixed-regulator spectral
completeness, and local kinetic, cubic, and Perron identities in
\cite{F16}; uniform geometry, subordination, the joint critical
spectrum, and local toron excision in \cite{F16B}. The consolidated
native high-block, resonance, and local-normal-form results are in
\cite{F17}; the subsequent electric-dominated gap, localized band,
conditional-fibre, and finite-time coefficient results used here are in
\cite{F18}. The nonlinear remainder and sector conclusions are in
\cite{PAR}. The local physical endpoint is developed in \cite{F1} and
in the Grand Tensor \cite{GT}.

The collar topology, all-torus finite Hessian margin, response banks,
corner identities, compact stationary-orbit Hessians, and non-wrapping
shell incidence data, the complete measured Laurent-bank reductions,
paired-collar weights, and sector-resolved quartic-model enclosures
refer to finite computational certificates documented
in these references. Their auxiliary executables are not reproduced in
this monograph, and no independent rerun of the complete certificate
chain is asserted. In particular, the finite alias--slice agreement and
corner means do not constitute a uniform full-torus or higher-shell
calculation. The covariance-transport propositions are algebraic
consequences of the displayed representative constraints;
\cref{crit:background-ADMC} is the explicit SCIC specialization, not a
proof of background independence.

The higher score, measured pushforward, core-gluing, and two-level
transfer statements are proved on their displayed finite or
same-carrier domains. The four-mark result is an analytic implication
from an exact source-dependent activity representation; it is not
claimed as an independently verified bound for every retained Wilson
background. The existing Gaussian tower and response certificates are
retained with their original scope. The actual orbit implications use
the specified finite certificates, not an independently reconstructed
numerical critical point. The local repair proofs use the stated
finite geometric incidence data, together with the displayed analytic
localization and same-law operator comparisons. Their chart radii and
large-$\beta$ thresholds are existential constants unless explicitly
stated otherwise.

The cubic analysis in \cref{sec:native-cubic-chaos} retains the
local expansion and complete orientation covariance sum, but not an
inference from third Wiener chaos to a Gaussian variable. Neither an
all-order exponential estimate nor a transfer-norm obstruction is
claimed from those low-order moments. The actual-vacuum spectral
limit remains a statement on its displayed mesoscopic window.

No unproved interacting construction, continuum measure, infrared
contraction, or anomaly identification is imported from another sector.
The conditional statements retain those requirements explicitly.

\backmatter

\begin{thebibliography}{F18}

\bibitem[GT]{GT}

A.~P\'elissier, \emph{The Renewal Grand Tensor: Relative Primitive Minimality, Operational Spectralization, and Domain-Specific SM--Spacetime and Arithmetic Loadings}, version V075. Unpublished manuscript, 2026.

\bibitem[F1]{F1}

A.~P\'elissier, \emph{Renewal Yang--Mills mass-gap programme}, series 1, version V135\_f1. Unpublished manuscript, 2026.

\bibitem[F2]{F2}

A.~P\'elissier, \emph{Renewal Yang--Mills mass-gap programme}, series 2, version V100\_f2. Unpublished manuscript, 2026.

\bibitem[F3]{F3}

A.~P\'elissier, \emph{Renewal Yang--Mills mass-gap programme}, series 3, version V100\_f3. Unpublished manuscript, 2026.

\bibitem[F4]{F4}

A.~P\'elissier, \emph{Renewal Yang--Mills mass-gap programme}, series 4, version V100\_f4. Unpublished manuscript, 2026.

\bibitem[F5]{F5}

A.~P\'elissier, \emph{Renewal Yang--Mills mass-gap programme}, series 5, version V100\_f5. Unpublished manuscript, 2026.

\bibitem[F6]{F6}

A.~P\'elissier, \emph{Renewal Yang--Mills mass-gap programme}, series 6, version V100\_f6. Unpublished manuscript, 2026.

\bibitem[F7]{F7}

A.~P\'elissier, \emph{Renewal Yang--Mills mass-gap programme}, series 7, version V100\_f7. Unpublished manuscript, 2026.

\bibitem[F8]{F8}

A.~P\'elissier, \emph{Renewal Yang--Mills mass-gap programme}, series 8, version V65\_f8. Unpublished manuscript, 2026.

\bibitem[F9]{F9}

A.~P\'elissier, \emph{Renewal Yang--Mills mass-gap programme}, series 9, version V100\_f9. Unpublished manuscript, 2026.

\bibitem[F10]{F10}

A.~P\'elissier, \emph{Renewal Yang--Mills mass-gap programme}, series 10, version V100\_f10. Unpublished manuscript, 2026.

\bibitem[F11]{F11}

A.~P\'elissier, \emph{Renewal Yang--Mills mass-gap programme}, series 11, version V100\_f11. Unpublished manuscript, 2026.

\bibitem[F12]{F12}
A.~P\'elissier, \emph{Renewal Yang--Mills mass-gap programme}, series 12,
version V100\_f12. Unpublished manuscript, 2026.

\bibitem[F13]{F13}
A.~P\'elissier, \emph{Renewal Yang--Mills mass-gap programme}, series 13,
version V17\_f13. Unpublished manuscript, 2026.

\bibitem[SM]{SM}
A.~P\'elissier, \emph{Renewal Yang--Mills shell mixing and measured
multiscale response}, version V16. Unpublished manuscript, 2026.

\bibitem[F14]{F14}
A.~P\'elissier, \emph{Renewal Yang--Mills mass-gap programme}, series 14,
version V100\_f14. Unpublished manuscript, 2026.

\bibitem[F15]{F15}
A.~P\'elissier, \emph{Renewal Yang--Mills mass-gap programme}, series 15,
version V100\_f15. Unpublished manuscript, 2026.

\bibitem[F16]{F16}
A.~P\'elissier, \emph{Renewal Yang--Mills mass-gap programme}, series 16,
version V67, spectral-certificate and local-transfer branch.
Unpublished manuscript, 2026.

\bibitem[F16B]{F16B}
A.~P\'elissier, \emph{Renewal Yang--Mills mass-gap programme}, series 16,
version V56, parallel joint-regulator and local-vacuum branch.
Unpublished manuscript, 2026.

\bibitem[F17]{F17}
A.~P\'elissier, \emph{Renewal Yang--Mills mass-gap programme}, series 17,
version V100, consolidated native-resonance and local-averaging line.
Unpublished manuscript, 2026.

\bibitem[F18]{F18}
A.~P\'elissier, \emph{Renewal Yang--Mills mass-gap programme}, series 18,
version V59, native conditional-deformation and electric spectral line.
Unpublished manuscript, 2026.

\bibitem[PAR]{PAR}

A.~P\'elissier, \emph{Renewal Yang--Mills theory: Nonperturbative margin protection, sector tightness, and topological transport}, version V5. Unpublished manuscript, 2026.

\end{thebibliography}

\end{document}

